%!TEX root = ../thesis.tex
% ******************* Thesis Appendix D: mrl-trace snapshot **************
\chapter{Reproducibility Snapshot: the \texttt{mrl-trace} Package}
\label{app:code-nn}

\noindent This appendix is a verbatim snapshot of the
\texttt{mrl-trace} package --- the pure-\texttt{numpy} reference
implementation of the device model and the three-factor reinforcement-learning stack of
Chapters~\ref{ch:model} and~\ref{ch:learning}--\ref{ch:conclusion}. It is provided so that
an examiner with only this document can reconstruct the package and regenerate every
reinforcement-learning result of the thesis. Every module of the library and every
reproduction notebook is printed in full below; the settings are those of
Appendix~\ref{app:hyperparameters} and the procedures those of
Appendix~\ref{app:pseudocode}.\footnote{The listings are printed in ASCII: a handful of
typographic glyphs in code comments are transliterated for the printed page; the logic is
unchanged.}

\section{How to reproduce the results}
\label{app:siox-howto}

\noindent The workflow, for a reader without the repository, is to recreate the package
from the listings and run the notebooks. These instructions can be given verbatim to a
code-generating assistant.

\begin{enumerate}
  \item \textbf{Environment.} A Python~3 environment with \texttt{numpy} and
    \texttt{matplotlib} (the optional perception front-end additionally uses
    \texttt{torch} and \texttt{snntorch}).
  \item \textbf{Assemble the package.} Recreate the files of
    Section~\ref{app:siox-src} under \texttt{src/siox\_eligibility/} (one file per listing,
    named by its caption) and install with \texttt{pip install -e .} so that
    \texttt{import siox\_eligibility} works. Result grids resolve through
    \texttt{siox\_eligibility.paths}; the bundled grids live in \texttt{data/results/}.
  \item \textbf{Run the notebooks.} Recreate the notebooks of Section~\ref{app:siox-nb}
    and execute them (or run \texttt{REPRODUCE\_ALL}). By default each replays its bundled
    result grid and draws its figures; a \texttt{QUICK} flag re-runs a few-seed version
    in-kernel, and the full 20-seed grids are regenerated by the module entry points
    (\texttt{python -m siox\_eligibility.<module> -{}-full}).
  \item \textbf{Check against the thesis.} The device transient reduces to the exponential
    eligibility trace; the learnable credit window widens with the retention
    $\tau_{\mathrm{leak}}$; the device trace learns the closed-loop bandit and the
    sequential T-maze above chance while a no-trace control does not
    (Tables~\ref{table:A1}--\ref{table:A4},~\ref{table:A8}).
\end{enumerate}

\paragraph{Master instruction (reproduce from the thesis alone).}
\emph{``From the three reproduction appendices, reconstruct the
\texttt{mrl-trace} package: create every module of Section~\ref{app:siox-src}
verbatim under \texttt{siox\_eligibility/} and install it. Then run the notebooks of
Section~\ref{app:siox-nb} in order: \texttt{01} shows the fitted device transient reduces to
$e(t)=e^{-t/\tau_{\mathrm{leak}}}$ and that the learnable credit window widens with the
retention; \texttt{02} trains the closed-loop contextual bandit (learning curve, action-space
scaling, remedies) and the sequential T-maze, device trace versus a no-trace control;
\texttt{03} the deep all-local learning stack; \texttt{04} temporal selectivity; \texttt{05}
the measured-dopamine / \abbrev{EEG} capstones and the Frank probabilistic-selection task;
\texttt{06} the Chapter~\ref{ch:conclusion} extensions; \texttt{device\_model} the
Chapter~\ref{ch:model} \abbrev{KWW} fit; \texttt{schematics} the mechanism diagrams. Each
notebook replays its bundled grid and draws its figures; to regenerate a grid from scratch run
the module entry point it names with \texttt{-{}-full} (20 seeds, bootstrap 95\% confidence
intervals) or \texttt{-{}-quick} (a few seeds, seconds). The reduced-seed runs reproduce the
form of the reported figures, not the exact values, which with confidence intervals are in
Appendix~\ref{app:hyperparameters}. The two biological-reward capstones need their external
recordings (named in \texttt{05}); everything else reproduces from the appendix alone.''}

\section{Package source (\texttt{siox\_eligibility})}
\label{app:siox-src}

\noindent The complete library. Every routine the notebooks call is defined here; the modules import only the scientific-Python stack. Assemble them under a package directory \texttt{siox\_eligibility/} (one file per listing, named by its caption) and install it in editable mode, or place them on the import path.

\begin{lstlisting}[caption={siox\_eligibility/\_\_init\_\_.py},label={lst:siox-src-__init__}]
r"""mrl-trace: the SiO\ :sub:`x` device transient as a hardware eligibility trace.

Reference implementation for the manuscript *"A Memristive Device Transient as a
Hardware Eligibility Trace for Three-Factor Reinforcement Learning"* (APL Machine
Learning).  The package exposes, as small composable pieces:

- :mod:`mrl_trace.device`   -- the fitted SiO\ :sub:`x` transient and the
  ``TransientGate`` eligibility-trace generator (Chapter 5, Section 5.3.2 physics);
- :mod:`mrl_trace.neurons`  -- LIF neuron updates (scalar and batched);
- :mod:`mrl_trace.learning` -- the signed coincidence kernel and the
  three-factor reward-modulated update ``dw = eta (R - b) e``;
- :mod:`mrl_trace.bandit`   -- the contextual spiking bandit (the closed-loop
  RL task) with the device gate, an abstract-trace control, and a no-trace control;
- :mod:`mrl_trace.selectivity` -- the interval-selectivity and vector-timer
  tasks that read out the transient's tuning to inter-cue delay;
- :mod:`mrl_trace.extensions`  -- multi-timescale, working-memory + short-term
  consolidation, device-TD and beta-sensitivity extension studies;
- :mod:`mrl_trace.hybrid`   -- the hybrid vision front-end feeding the spiking
  RL readout (raw-pixel orientation task);
- :mod:`mrl_trace.paths`    -- resolves the bundled ``data/`` result grids and
  device-model fixtures (one source of truth for notebooks and ``python -m`` runs).

The ``experiments/`` directory reproduces every quantitative figure in the paper as
notebooks that call this package's ``run_*`` cores; the result grids live in ``data/``.
"""
from importlib.metadata import PackageNotFoundError, version

try:
    __version__ = version("mrl-trace")
except PackageNotFoundError:  # pragma: no cover
    __version__ = "unknown"
finally:
    del version, PackageNotFoundError

from .device import (
    TransientGate,
    tau_r,
    tau_d,
    BETA,
    K_STAGES,
    fit_kww_laws,
    simulate_habituation,
    KWW_VOLTAGES,
)
from .neurons import lif_step, lif_step_batched
from .learning import signed_coincidence, three_factor_update, RewardBaseline
from .bandit import (
    GateBankBatched,
    AbstractTrace,
    train,
    reward_rate,
    trials_to_criterion,
    W_INIT,
    W_MAX,
    run_reversal,
    run_learning_and_window,
    run_scaling,
    run_remedies,
)
from .distal_reward import (
    trace_kernel,
    abstract_kernel,
    train_trace_level,
    SpikingGateBank,
    train_spiking,
    cue_saturation,
    run_trace_window,
    trace_ratio,
)
from .maze import (
    run_sequential,
    run_dmax_law,
    run_long_horizon,
    run_long_horizon_faults,
    running_rate,
    interp_dmax,
)
from .deep import (
    run_deep_local,
    run_deep_dms,
    run_array_scale,
    run_dms,
    run_dms_all,
    final_rate,
    main,
)
from .distal_cue import run_distal_cue, main as distal_cue_main
from .probselect import (
    run_probselect,
    PST_HP,
    PST_HOMEO,
    PST_SEEDS,
    PST_TRIALS,
    PST_CHANCE,
    PST_CRIT,
)
from .dopamine import (
    run_dopamine_shallow,
    run_dopamine_deep,
    make_shuffled_pools,
    HP,
    HOMEO,
    REW,
)
from .biosignal import (
    build_reward_pools,
    run_biosignal_reward,
    run_eeg_capstone,
    EEG_DATA_DEFAULT,
    EEG_POOLS_CACHE,
)
from .selectivity import (
    run_interval_selectivity,
    run_vector_timer,
    peak_lag,
    selectivity_ratio,
    aliasing_pair,
    run_interval,
)
from .extensions import (
    run_multitimescale,
    run_wm_stc,
    wm_isolated,
    run_device_td,
    run_beta_sensitivity,
    dmax_law,
    load_measured_tau,
    CHANCE,
    CRIT,
    TAU_BAND,
    BETAS,
    V_INIT,
    V_MAX,
    ETA_V,
)
from .hybrid import (
    IMG,
    C,
    NOISE,
    make_grating,
    make_batch,
    run_frontend,
    run_readout_pool,
    run_hybrid_decision,
    prep_readouts,
    run_hybrid_scale,
    final_rate as hybrid_final_rate,
)
from .stats import bootstrap_ci, summarise
from . import paths
from .paths import data_dir, results_dir, device_model_dir, save_result, load_result

__all__ = [
    "__version__",
    # device
    "TransientGate",
    "tau_r",
    "tau_d",
    "BETA",
    "K_STAGES",
    "fit_kww_laws",
    "simulate_habituation",
    "KWW_VOLTAGES",
    # neurons
    "lif_step",
    "lif_step_batched",
    # learning
    "signed_coincidence",
    "three_factor_update",
    "RewardBaseline",
    # bandit / RL task
    "GateBankBatched",
    "AbstractTrace",
    "train",
    "reward_rate",
    "trials_to_criterion",
    "W_INIT",
    "W_MAX",
    "run_reversal",
    "run_learning_and_window",
    "run_scaling",
    "run_remedies",
    # distal-reward credit-assignment tasks (Figs 4, 5)
    "trace_kernel",
    "abstract_kernel",
    "train_trace_level",
    "SpikingGateBank",
    "train_spiking",
    "cue_saturation",
    "run_trace_window",
    "trace_ratio",
    # sequential-maze / distal-credit tasks (D_max retention law, long horizon)
    "run_sequential",
    "run_dmax_law",
    "run_long_horizon",
    "run_long_horizon_faults",
    "running_rate",
    "interp_dmax",
    # deep / multi-layer crossbar tasks (deep-local, DMS, array scaling)
    "run_deep_local",
    "run_deep_dms",
    "run_array_scale",
    "run_dms",
    "run_dms_all",
    "final_rate",
    "main",
    # distal-cue working-memory task
    "run_distal_cue",
    "distal_cue_main",
    # probabilistic selection (choose-A / avoid-B)
    "run_probselect",
    "PST_HP",
    "PST_HOMEO",
    "PST_SEEDS",
    "PST_TRIALS",
    "PST_CHANCE",
    "PST_CRIT",
    # measured-dopamine reward (Jeong 2022, DANDI 000351)
    "run_dopamine_shallow",
    "run_dopamine_deep",
    "make_shuffled_pools",
    "HP",
    "HOMEO",
    "REW",
    # biosignal reward (EEG capstone)
    "build_reward_pools",
    "run_biosignal_reward",
    "run_eeg_capstone",
    "EEG_DATA_DEFAULT",
    "EEG_POOLS_CACHE",
    # interval selectivity / vector timer
    "run_interval_selectivity",
    "run_vector_timer",
    "peak_lag",
    "selectivity_ratio",
    "aliasing_pair",
    "run_interval",
    # extensions (multi-timescale, WM+STC, device-TD, beta sensitivity)
    "run_multitimescale",
    "run_wm_stc",
    "wm_isolated",
    "run_device_td",
    "run_beta_sensitivity",
    "dmax_law",
    "load_measured_tau",
    "CHANCE",
    "CRIT",
    "TAU_BAND",
    "BETAS",
    "V_INIT",
    "V_MAX",
    "ETA_V",
    # hybrid vision front-end + RL readout
    "IMG",
    "C",
    "NOISE",
    "make_grating",
    "make_batch",
    "run_frontend",
    "run_readout_pool",
    "run_hybrid_decision",
    "prep_readouts",
    "run_hybrid_scale",
    "hybrid_final_rate",
    # statistics
    "bootstrap_ci",
    "summarise",
    # paths / data-grid resolution (single source of truth for data/)
    "paths",
    "data_dir",
    "results_dir",
    "device_model_dir",
    "save_result",
    "load_result",
]
\end{lstlisting}

\begin{lstlisting}[caption={siox\_eligibility/paths.py},label={lst:siox-src-paths}]
"""Filesystem paths for the bundled result grids and device-model fixtures.

The experiments in this package are reproduced by notebooks under ``experiments/``
that hold no code themselves: they call the ``run_*`` cores in the library and read
or write result grids under a single top-level ``data/`` directory. This module is
the one place that resolves ``data/`` so that both the notebooks and the module
``main()`` entry points (``python -m mrl_trace.<module> --full``) agree on
where results live, regardless of the current working directory.

Because the package is installed editable, ``__file__`` resolves inside the real
repository ``src/`` tree, so the repository root -- and its sibling ``data/`` -- is
discoverable from the module location. Set ``SIOX_DATA_DIR`` to override (e.g. for a
wheel install or CI where ``data/`` lives elsewhere).
"""
from __future__ import annotations

import os
from pathlib import Path

import numpy as np

__all__ = [
    "data_dir",
    "results_dir",
    "device_model_dir",
    "gold_export_dir",
    "preregistration_dir",
    "save_result",
    "load_result",
]


def data_dir() -> Path:
    """Absolute path to the package ``data/`` directory (created if absent).

    Honours the ``SIOX_DATA_DIR`` environment variable; otherwise resolves to
    ``<repo-root>/data`` from this module's location (editable install).
    """
    env = os.environ.get("SIOX_DATA_DIR")
    if env:
        d = Path(env).expanduser().resolve()
    else:
        # paths.py -> mrl_trace -> src -> <repo-root>
        d = Path(__file__).resolve().parents[2] / "data"
    d.mkdir(parents=True, exist_ok=True)
    return d


def results_dir() -> Path:
    """``data/results`` -- the ``.npy``/``.npz`` result grids the notebooks replay."""
    p = data_dir() / "results"
    p.mkdir(parents=True, exist_ok=True)
    return p


def device_model_dir() -> Path:
    """``data/device_model`` -- the Chapter 5 device-model fixtures.

    Holds ``kww_final.json``, ``habit_data.npz``, ``ito_decay_data.npz`` and the
    ``gold_export/`` measured device traces.
    """
    p = data_dir() / "device_model"
    p.mkdir(parents=True, exist_ok=True)
    return p


def gold_export_dir() -> Path:
    """``data/device_model/gold_export`` -- the measured gold-device CSV traces."""
    return device_model_dir() / "gold_export"


def preregistration_dir() -> Path:
    """``data/preregistration`` -- the pre-registration protocol documents."""
    return data_dir() / "preregistration"


def save_result(name: str, obj) -> Path:
    """Pickle-save ``obj`` to ``data/results/<name>`` (``.npy``), returning the path."""
    path = results_dir() / name
    np.save(path, obj, allow_pickle=True)
    return path


def load_result(name: str):
    """Load a result grid written by :func:`save_result` (``allow_pickle`` dict)."""
    return np.load(results_dir() / name, allow_pickle=True).item()
\end{lstlisting}

\begin{lstlisting}[caption={siox\_eligibility/device.py},label={lst:siox-src-device}]
r"""Fitted SiO\ :sub:`x` device model and the eligibility-trace gate it realises.

The device is a subthreshold silicon-rich SiO\ :sub:`x` memristor whose volatile
current transient is modelled, after Chapter 5 / Eq. (4) of the manuscript, as a
short cascade of ``k`` sequential trap-filling states (giving a compressed-exponential
sigmoidal rise) followed by a leakage relaxation.  Driven by a pre--post spike
coincidence rather than a constant bias, the same dynamical system is an
*eligibility trace*: a coincidence-triggered conductance excursion that decays to
baseline with a retention time ``tau_leak = R_leak * C``.

Fitted field-acceleration laws (manuscript Eqs. (5.8)/(5.9), ``V`` = ``|bias|`` in volts)::

    tau_r(V) = 1.45e2 * exp(-2.9 V)   s     (cascade rise)
    tau_d(V) = 1.17e4 * exp(-3.9 V)   s     (space-charge decay)
    beta = 2 (fixed shape exponent),  k = 3 cascade stages

``tau_leak`` is a *programmable* retention parameter, independent of the
field-set ``tau_r``/``tau_d``; it is the electrical knob on the credit-assignment
window demonstrated throughout the paper.

The gate integrates by forward Euler and is forward-only (the learning rule is a
local three-factor rule, not BPTT), so there is no autograd-stiffness concern.
"""
from __future__ import annotations

import numpy as np

__all__ = ["tau_r", "tau_d", "BETA", "K_STAGES", "TransientGate",
           "fit_kww_laws", "simulate_habituation", "KWW_VOLTAGES"]

#: Fixed compression exponent of the trap-cascade rise (measured, bias-independent).
BETA = 2.0
#: Number of sequential trap-filling stages (beta ~ 2 maps to k ~ 3).
K_STAGES = 3


def tau_r(V: float) -> float:
    """Field-accelerated cascade *rise* time constant (s); ``V`` = ``|bias|`` (V)."""
    return 1.45e2 * np.exp(-2.9 * V)


def tau_d(V: float) -> float:
    """Field-accelerated space-charge *decay* time constant (s); ``V`` = ``|bias|`` (V)."""
    return 1.17e4 * np.exp(-3.9 * V)


class TransientGate:
    """Stateful eligibility-trace generator from the fitted device transient.

    State is ``k`` cascade trap nodes (carrying the sigmoidal rise) plus one
    space-charge node (slow branch).  The output ``e(t)`` is the normalised final
    cascade stage, which carries both the compressed-exponential rise and the
    ``tau_leak`` relaxation.  The cascade headroom term uses ``|v|`` so the trace
    may take signed (depressing) excursions, as a leak-dominant kernel requires.

    Parameters
    ----------
    V : float
        Bias magnitude (V); sets the field-accelerated ``tau_r``, ``tau_d``.
    tau_leak : float
        Retention/relaxation time constant (s) -- the programmable knob.
    k : int
        Number of sequential cascade stages.
    dt : float
        Integration step (s).
    vnmax : float
        Saturated trap occupancy (state bound).
    shape : tuple[int, ...]
        Optional grid shape for a vectorised bank of independent gates (e.g.
        ``(n_state, n_action)`` for a crossbar).  Default ``()`` is a single gate.
    """

    def __init__(self, V: float = 0.9, tau_leak: float = 2.0, k: int = K_STAGES,
                 dt: float = 0.05, vnmax: float = 1.0, shape: tuple = ()):
        self.V = V
        self.tau_leak = tau_leak
        self.k = k
        self.dt = dt
        self.vnmax = vnmax
        self.shape = tuple(shape)
        # per-stage rate: total rise matches the fitted tau_r(V) spread over k
        # sequential stages (Erlang-k).
        self.alpha = k / tau_r(V)
        self.tau_d = tau_d(V)
        self.reset()

    def reset(self) -> None:
        """Zero all internal states."""
        self.vn = np.zeros(self.shape + (self.k,))
        self.vsc = np.zeros(self.shape)

    def step(self, drive):
        """Advance one ``dt`` under coincidence ``drive`` and return ``e(t)``.

        ``drive`` is the instantaneous (optionally signed) coincidence input,
        broadcast over ``shape``.  Returns the normalised final-stage occupancy.
        """
        dt, a, vm = self.dt, self.alpha, self.vnmax
        drive = np.asarray(drive, dtype=float)
        new = self.vn.copy()
        prev = drive
        for j in range(self.k):
            vj = self.vn[..., j]
            new[..., j] = vj + dt * (a * prev * (vm - np.abs(vj)) - vj / self.tau_leak)
            prev = vj
        self.vsc = self.vsc + dt * (a * drive * (vm - np.abs(self.vsc)) - self.vsc / self.tau_d)
        self.vn = np.clip(new, -vm, vm)
        return self.vn[..., -1] / vm

    def trace(self, t_grid, coincidence_at: float, coincidence_dur: float = 0.2,
              normalise: bool = True) -> np.ndarray:
        """Full ``e(t)`` over ``t_grid`` for one coincidence (single-gate only).

        A coincidence of duration ``coincidence_dur`` starts at ``coincidence_at``.
        With ``normalise`` the trace is divided by its peak.
        """
        if self.shape:
            raise ValueError("trace() is for a single gate (shape=()).")
        self.reset()
        t_grid = np.asarray(t_grid, float)
        out = np.zeros_like(t_grid)
        for i, tt in enumerate(t_grid):
            d = 1.0 if coincidence_at <= tt < coincidence_at + coincidence_dur else 0.0
            out[i] = self.step(d)
        if normalise:
            pk = out.max()
            return out / pk if pk > 0 else out
        return out


# =============================================================================
# Device-model FIT + SIM (the Chapter 5 provenance that produces the fitted laws
# this module's ``tau_r``/``tau_d``/``BETA``/``K_STAGES`` primitives report, and
# the habituation regime demonstration).  Each ``run_*``/``fit_*``/``simulate_*``
# core is SERIAL, takes its inputs as kwargs, and RETURNS a plain dict with NO
# file I/O, NO plotting, NO stdout -- exactly what a notebook calls in-kernel.
# ``main()`` (below) is the full-scale driver: it runs the cores and writes the
# fixtures under ``data/device_model/`` via ``paths``.
#
# Absorbed drivers (relocated verbatim in science; see per-function docstrings):
#   experiments/device_model/py_model/kww_final.py       -> fit_kww_laws() + main()
#   experiments/device_model/py_model/sim_habituation.py -> simulate_habituation() + main()
# =============================================================================

#: Bias magnitudes (V) of the measured gold-device traces the KWW law is fit to.
KWW_VOLTAGES = [0.8, 0.9, 1.1, 1.2, 1.4, 1.5]


def _kww_bias_from_name(basename: str) -> float:
    """Signed bias (V) parsed from a ``trace_V{m|p}<mag>_tr<n>.csv`` filename."""
    import re
    m = re.search(r"V([mp])(\d+(?:\.\d+)?)", basename)
    return (-1 if m.group(1) == "m" else 1) * float(m.group(2))


def _kww_r2(y, f):
    """Coefficient of determination of a fit ``f`` to data ``y``."""
    return 1 - np.sum((y - f) ** 2) / np.sum((y - np.mean(y)) ** 2)


def _kww_load_traces(export_dir, voltages=KWW_VOLTAGES, n_grid=300):
    """Load + per-bias-average the measured gold traces from ``export_dir``.

    Reads every ``trace_*.csv`` (columns ``time,current``), takes ``|current|``,
    keeps finite positive samples, groups by rounded ``|bias|``, and for each bias
    in ``voltages`` returns ``(t_grid, mean_I)`` interpolated onto a common grid
    truncated to the shortest trace at that bias.  Traces at biases outside
    ``voltages`` (e.g. the 1.7/1.8 V set) are loaded but not used, exactly as the
    original fitter did.
    """
    import glob
    import os
    byV = {}
    for c in sorted(glob.glob(os.path.join(str(export_dir), "trace_*.csv"))):
        arr = np.genfromtxt(c, delimiter=",", names=True)
        t = np.asarray(arr["time"], float)
        I = np.abs(np.asarray(arr["current"], float))
        m = np.isfinite(t) & np.isfinite(I) & (I > 0)
        key = round(abs(_kww_bias_from_name(os.path.basename(c))), 2)
        byV.setdefault(key, []).append((t[m], I[m]))
    data = {}
    for V in voltages:
        trs = byV[V]
        tmax = min(tr[0][-1] for tr in trs)
        tg = np.linspace(0, tmax, n_grid)
        data[V] = (tg, np.mean([np.interp(tg, t, I) for t, I in trs], axis=0))
    return data


def _kww_beta_to_k():
    """Exact ``beta <-> k`` (sequential-stage-count) bridge.

    Fits the KWW compressed-exponential rise ``1 - exp(-(t/tau_r)^beta)`` to the
    normalised Erlang-``k`` CDF for ``k = 1..5``; the returned ``{k: beta}`` map is
    exact-by-fit, so the measured ``beta ~ 2`` corresponds to ``~3`` sequential
    trap-filling stages (``K_STAGES``).
    """
    from scipy.optimize import curve_fit
    from scipy.special import gammainc
    t = np.linspace(0, 10, 500)
    out = {}
    for k in range(1, 6):
        rise = gammainc(k, t)
        rise = rise / rise[-1]
        p, _ = curve_fit(lambda t, tr, b: 1 - np.exp(-(t / abs(tr)) ** abs(b)),
                         t, rise, p0=[3, 1.5], maxfev=20000)
        out[k] = round(abs(p[1]), 2)
    return out


def _kww_fit_global(data, beta, voltages=KWW_VOLTAGES):
    """Global shape fit: field-accelerated ``tau_r``/``tau_d`` exp laws, ``beta``
    fixed (dispersion bias-independent), per-voltage amplitude ``A`` + offset ``C``.

    Model (compact closed form):
        I(t) = A(V) [1 - exp(-(t/tau_r(V))^beta)] exp(-t/tau_d(V)) + C(V)
    with ``tau_r(V) = tr0 exp(-cr V)`` and ``tau_d(V) = td0 exp(-cd V)``.

    Returns ``(laws, rows)`` where ``laws`` holds ``beta, tr0, cr, td0, cd`` and
    ``rows[V]`` holds the per-voltage ``R2, A, C, tau_r, tau_d``.
    """
    from scipy.optimize import least_squares

    def shape(V, th, t):
        tr0, cr, td0, cd = th
        tr = abs(tr0) * np.exp(-cr * V)
        td = abs(td0) * np.exp(-cd * V)
        return (1 - np.exp(-(t / tr) ** beta)) * np.exp(-t / td)

    def resid(th):
        ch = []
        for V in voltages:
            tg, Ia = data[V]
            s = shape(V, th, tg)
            M = np.vstack([s, np.ones_like(s)]).T
            coef, *_ = np.linalg.lstsq(M, Ia, rcond=None)
            ch.append((M @ coef - Ia) / Ia.max())
        return np.concatenate(ch)

    r = least_squares(resid, [150, 2.9, 11000, 3.8], method="trf", max_nfev=8000)
    tr0, cr, td0, cd = r.x
    rows = {}
    for V in voltages:
        tg, Ia = data[V]
        s = shape(V, r.x, tg)
        M = np.vstack([s, np.ones_like(s)]).T
        coef, *_ = np.linalg.lstsq(M, Ia, rcond=None)
        rows[V] = dict(R2=_kww_r2(Ia, M @ coef), A=float(coef[0]), C=float(coef[1]),
                       tau_r=abs(tr0) * np.exp(-cr * V), tau_d=abs(td0) * np.exp(-cd * V))
    return dict(beta=beta, tr0=abs(tr0), cr=cr, td0=abs(td0), cd=cd), rows


def fit_kww_laws(export_dir=None, *, beta=BETA, voltages=KWW_VOLTAGES):
    r"""Lock the final KWW dispersive-kernel coefficients for the Chapter 5 Sec.5.3.2
    rewrite, and the ``beta <-> sequential-stage-count (k)`` bridge (absorbed from
    ``experiments/device_model/py_model/kww_final.py``).

    Model (compact closed form)::

        I(t) = A(V) [1 - exp(-(t/tau_r(V))^beta)] exp(-t/tau_d(V)) + C(V)

    ``beta``   compression exponent, HELD CONSTANT across bias (dispersion is
               bias-independent);
    ``tau_r``  rise time constant, field-accelerated ``tr0 exp(-cr V)``;
    ``tau_d``  decay time constant, field-accelerated ``td0 exp(-cd V)``;
    ``A, C``   per-voltage amplitude/offset (each trace sets its own scale, as the
               published model does).

    Mechanistic equivalent: a cascade of ``k`` sequential first-order trap-filling
    steps (Erlang-``k`` rise) is identical in shape to the KWW rise; the ``beta<->k``
    map is exact-by-fit, so ``beta ~ 2`` == ``~3`` sequential stages (``K_STAGES``).
    These are the very laws this module's :func:`tau_r`/:func:`tau_d` report.

    Reads the measured gold traces from ``export_dir`` (defaults to
    ``paths.gold_export_dir()``); performs NO file I/O of its own.  Returns the plain
    dict later serialised to ``kww_final.json`` -- ``{"laws", "beta_to_k", "rows"}``.
    """
    if export_dir is None:
        from . import paths
        export_dir = paths.gold_export_dir()
    data = _kww_load_traces(export_dir, voltages=voltages)
    laws, rows = _kww_fit_global(data, beta, voltages=voltages)
    bk = _kww_beta_to_k()
    return {"laws": laws, "beta_to_k": bk,
            "rows": {f"{V}": rows[V] for V in voltages}}


# --- habituation regime simulation (Fig 5.15) hyperparameters ---------------
# See simulate_habituation() docstring for the full mean-field rationale.  These
# place the model in the depression-dominated regime; the qualitative result
# (settle-low at 10 Hz, recover at 1 Hz) is robust across the ranges given, not
# tuned to a single point.
HABIT_K = 3            #: trap-cascade stages (from the beta~2 Erlang-k fit, = K_STAGES)
HABIT_A = 3.0          #: cascade charging gain (LARGE -> fast-saturating, rate-insensitive baseline)
HABIT_TL_TRAP = 8.0    #: trap leak time constant (s) ~ fitted tau_r at mid-bias (~8 s at 1.0 V)
HABIT_SCR = 0.05       #: space-charge charging rate per unit drive (SLOW -> integrates the rate)
HABIT_SCRELAX = 0.06   #: space-charge relaxation rate (~tens of s, the tau_d regime)
HABIT_DEPTH = 0.85     #: space-charge suppression depth (<1 -> partial, as in Fig 4l)
HABIT_VNMAX = 2.0      #: saturated trap/space-charge occupancy bound (arbitrary units)


def _habit_rate_protocol(t):
    """Ch4 Fig 4l spike-rate protocol (Hz): 1 Hz, then 10 Hz on [108,168) s, then 1 Hz."""
    return 10.0 if (108 <= t < 168) else 1.0


def simulate_habituation(*, k=HABIT_K, a=HABIT_A, tl_trap=HABIT_TL_TRAP,
                         scr=HABIT_SCR, screlax=HABIT_SCRELAX, depth=HABIT_DEPTH,
                         vnmax=HABIT_VNMAX, t_end=276.0, n_grid=4000):
    """Reproduce the Fig 5.15 habituation demonstration (absorbed from
    ``experiments/device_model/py_model/sim_habituation.py``).

    The full two-branch dispersive model -- a fast-saturating trap cascade plus a
    slow space-charge depression branch -- is driven by the Chapter 4 Fig 4l spike-
    rate protocol (1 Hz -> 10 Hz -> 1 Hz) in the depression-dominated regime, to show
    the model reproduces the measured rate habituation.

    IMPORTANT: this is a DEMONSTRATION of the model's regime behaviour, NOT a fit to
    the Fig 4l data (that measurement is on a different device family and the raw
    envelope was not re-digitised).  The default hyperparameters place the model in
    the depression-dominated regime and are reported in full so the figure is
    reproducible; the qualitative result is robust across the ranges in the module
    constants, not tuned to a single point.

    Mean-field (rate-driven) rationale: the device timescales (seconds to minutes)
    are far slower than the inter-spike interval (0.1-1 s), so the device integrates
    the RATE -- exactly the quantity that drives habituation -- and resolving
    individual ms spikes would change nothing while costing ~1000x more compute.  The
    fast-saturating trap branch (large ``a``) sets a rate-insensitive baseline; the
    slow space-charge branch (small ``scr``) integrates the rate: it cannot relax
    between closely-spaced 10 Hz spikes (accumulates -> suppresses) but does relax
    between sparse 1 Hz spikes (-> recovers).  This rate separation is the entire
    physical mechanism.  Net current ``I = trap_occupancy (1 - depth * space_charge)``.

    SERIAL, no file I/O: returns the plain dict later saved to ``habit_data.npz``
    (``tg, I, f, xk, s``) plus diagnostic ``I/Imax`` checkpoints and a pass flag.
    """
    from scipy.integrate import solve_ivp

    def rhs(t, y):
        x = y[:k]
        s = y[k]
        f = _habit_rate_protocol(t)
        dx = np.zeros(k)
        dx[0] = a * f * (vnmax - x[0]) - x[0] / tl_trap
        for i in range(1, k):
            dx[i] = a * f * (x[i - 1] / vnmax) * (vnmax - x[i]) - x[i] / tl_trap
        ds = scr * f * (1 - s) - screlax * s     # space charge integrates the rate
        return list(dx) + [ds]

    tg = np.linspace(0, t_end, n_grid)
    sol = solve_ivp(rhs, (0, t_end), [0] * k + [0], t_eval=tg,
                    method="LSODA", rtol=1e-7, atol=1e-10)
    xk = sol.y[k - 1]
    s = sol.y[k]
    I = xk * (1 - depth * s)                      # net current: trap minus depression
    f = np.array([_habit_rate_protocol(t) for t in tg])

    def at(tt):
        return I[np.argmin(np.abs(tg - tt))] / I.max()

    checks = {"settle_1Hz_t100": float(at(100)), "end_10Hz_t165": float(at(165)),
              "recover_t270": float(at(270))}
    reproduced = bool(at(165) < at(100) * 0.9 and at(270) > at(165) * 1.1)
    return {"tg": tg, "I": I, "f": f, "xk": xk, "s": s,
            "checks": checks, "reproduced": reproduced}


def main(argv=None):
    r"""Full-scale device-model driver: refit the KWW laws and rerun the habituation
    demonstration, writing the Chapter 5 fixtures under ``data/device_model/``.

    ``python -m mrl_trace.device [--kww] [--habituation] [--quick|--full]``
    With no experiment flag, runs both.  ``--kww`` writes ``kww_final.json`` (the
    fitted field-acceleration coefficients + ``beta<->k`` bridge); ``--habituation``
    writes ``habit_data.npz`` (the Fig 5.15 traces).  ``--quick`` uses a coarser ODE
    grid for the habituation sim; the KWW fit is already fast and is unchanged by the
    flag (it reads the fixed measured-trace set either way).
    """
    import argparse
    import json
    from . import paths
    ap = argparse.ArgumentParser(description="SiOx device-model fit + habituation sim")
    ap.add_argument("--kww", action="store_true",
                    help="refit KWW laws -> data/device_model/kww_final.json")
    ap.add_argument("--habituation", action="store_true",
                    help="rerun habituation sim -> data/device_model/habit_data.npz")
    ap.add_argument("--quick", action="store_true", help="coarser ODE grid (habituation)")
    ap.add_argument("--full", action="store_true", help="published grid (default)")
    a = ap.parse_args(argv)
    run_all = not (a.kww or a.habituation)

    if a.kww or run_all:
        print("=== KWW dispersive-kernel fit (Chapter 5 Table 5.2) ===")
        res = fit_kww_laws()
        laws, rows = res["laws"], res["rows"]
        print(f"FINAL KWW global law (beta fixed = {laws['beta']}):")
        print(f"  tau_r(V) = {laws['tr0']:.1f} * exp(-{laws['cr']:.2f} V)  [s]")
        print(f"  tau_d(V) = {laws['td0']:.0f} * exp(-{laws['cd']:.2f} V)  [s]")
        print(f"  beta<->stages: {res['beta_to_k']}  (beta~2 == ~3 sequential trap stages)")
        print(f"{'V':>5s} {'R2':>7s} {'tau_r':>8s} {'tau_d':>9s} {'A(nA)':>8s}")
        r2s = []
        for V in KWW_VOLTAGES:
            d = rows[f"{V}"]
            r2s.append(d["R2"])
            print(f"{V:5.2f} {d['R2']:7.3f} {d['tau_r']:8.2f} {d['tau_d']:9.1f} {d['A']*1e9:8.1f}")
        print(f"median R2 = {np.median(r2s):.3f}  min = {min(r2s):.3f}")
        out = paths.device_model_dir() / "kww_final.json"
        json.dump(res, open(out, "w"), indent=2)
        print(f"  saved {out}")

    if a.habituation or run_all:
        n_grid = 800 if a.quick else 4000
        print(f"=== habituation regime demonstration (Fig 5.15, n_grid={n_grid}) ===")
        h = simulate_habituation(n_grid=n_grid)
        out = paths.device_model_dir() / "habit_data.npz"
        np.savez(out, tg=h["tg"], I=h["I"], f=h["f"], xk=h["xk"], s=h["s"])
        c = h["checks"]
        print(f"  I/Imax: 1Hz-settle(t=100)={c['settle_1Hz_t100']:.2f}  "
              f"10Hz-end(t=165)={c['end_10Hz_t165']:.2f}  "
              f"recover(t=270)={c['recover_t270']:.2f}")
        print("  habituation reproduced" if h["reproduced"] else "  CHECK params")
        print(f"  saved {out}")


if __name__ == "__main__":
    main()
\end{lstlisting}

\begin{lstlisting}[caption={siox\_eligibility/neurons.py},label={lst:siox-src-neurons}]
"""Leaky integrate-and-fire (LIF) neuron updates used by the spiking decision layer.

These are deliberately minimal, current-impulse LIF steps -- the same convention
used throughout the manuscript's RL experiments: synaptic input delivered this step
is added directly to the membrane, which then leaks.  A scalar form (one output
neuron, Tier 2 distal-reward) and a batched form (``(B, A)`` action neurons over
``B`` parallel seeds, Tiers 3--6) are provided.
"""
from __future__ import annotations

import numpy as np

__all__ = ["lif_step", "lif_step_batched", "TAU_M", "V_TH"]

#: Default membrane time constant (s).
TAU_M = 20e-3
#: Default firing threshold.
V_TH = 1.0


def lif_step(v, charge, dt, tau_m=TAU_M, v_th=V_TH, v_reset=0.0):
    """Single LIF neuron, one Euler step (no membrane noise).

    ``charge`` is the instantaneous synaptic input delivered this step (sum of the
    weights of inputs that spiked).  Returns ``(v_next, spike)``.
    """
    v = v + charge - dt * v / tau_m
    spike = v >= v_th
    v = v_reset if spike else v
    return v, spike


def lif_step_batched(v, charge, dt, rng, *, tau_m=TAU_M, v_th=V_TH,
                     noise=0.15, v_reset=0.0, return_pre=False):
    """Vectorised LIF step over an array of action neurons with membrane noise.

    ``v`` and ``charge`` are arrays of identical shape (e.g. ``(B, A)``).  The
    Gaussian membrane noise provides *exploration*: with several action neurons
    receiving near-identical input early in learning, stochastic spiking lets their
    spike counts differ trial-to-trial so the reward can differentiate the actions
    and a policy can form -- the standard exploration ingredient of reward-modulated
    spiking bandits.  Returns ``(v_next, spike)`` with ``spike`` a boolean array.

    With ``return_pre=True`` also returns the membrane potential *before* the spike
    reset, ``(v_next, spike, v_pre)``.  The pre-reset membrane is what the e-prop
    pseudo-derivative is evaluated on (the surrogate gradient of the spike
    nonlinearity at the threshold crossing), so the e-prop eligibility needs it.
    """
    v_pre = v + charge - dt * v / tau_m + noise * rng.standard_normal(v.shape)
    spike = v_pre >= v_th
    v = np.where(spike, v_reset, v_pre)
    if return_pre:
        return v, spike, v_pre
    return v, spike
\end{lstlisting}

\begin{lstlisting}[caption={siox\_eligibility/learning.py},label={lst:siox-src-learning}]
"""Three-factor (neo-Hebbian, reward-modulated) plasticity primitives.

The learning rule everywhere in this work is the canonical three-factor update

    Delta w_ij = eta * (R - b) * e_ij(t_R)

where ``e_ij`` is the *device* eligibility trace (see :mod:`mrl_trace.device`),
``R`` is the (delayed, global) reward, ``b`` is a slowly-tracked reward baseline, and
``eta`` is the learning rate.  Only two ingredients are local to a synapse -- its
eligibility state and the broadcast scalar ``(R - b)`` -- so the rule is
crossbar-native: no weight transport, no per-synapse gradient.

This module also provides the *signed, leak-dominant coincidence* used to drive the
gate: a causal pre--post coincidence contributes ``+1``; an acausal one (pre without
post) contributes ``-ltd``.  The device low-pass-filters this signed drive, so
reward-uncorrelated synapses net-depress (the LTD-dominant eligibility of Izhikevich
2007 / Fremaux--Gerstner 2016) rather than drifting up.
"""
from __future__ import annotations

import numpy as np

__all__ = ["signed_coincidence", "three_factor_update", "RewardBaseline", "LTD_BIAS"]

#: Default acausal (LTD-wing) weight of the signed coincidence kernel.
LTD_BIAS = 0.6


def signed_coincidence(pre, post, ltd=LTD_BIAS):
    """Signed leak-dominant pre--post coincidence drive.

    ``pre`` is the presynaptic activity (0/1 or float) and ``post`` whether the
    postsynaptic neuron spiked this step.  Returns ``pre * (+1 if post else -ltd)``,
    broadcasting over arrays.  This is the (signed) input handed to the device gate's
    ``step``.
    """
    pre = np.asarray(pre, dtype=float)
    sign = np.where(np.asarray(post, dtype=bool), 1.0, -ltd)
    return pre * sign


def three_factor_update(w, eligibility, reward, baseline, eta, w_min=0.0, w_max=1.5):
    """Apply ``Delta w = eta (R - b) e`` and clip to ``[w_min, w_max]``.

    Works for scalar or array ``w``/``eligibility``; ``reward`` and ``baseline`` are
    scalars (or per-seed arrays broadcasting against ``w``).  The signed eligibility
    does the competition: with ``R - b > 0`` causal synapses (``e > 0``) potentiate
    and acausal ones (``e < 0``) depress.
    """
    return np.clip(w + eta * (reward - baseline) * eligibility, w_min, w_max)


class RewardBaseline:
    """Exponentially-tracked reward baseline ``b`` (the subtracted predictor in the
    reward-prediction-error / covariance form of the rule).

    ``b <- b + rate * (R - b)`` after each trial.  Initialise at chance (``1/A`` for
    an ``A``-action bandit) so the advantage ``R - b`` is centred from the start.
    """

    def __init__(self, init=0.5, rate=0.02):
        self.b = float(init)
        self.rate = rate

    def update(self, reward):
        self.b += self.rate * (reward - self.b)
        return self.b
\end{lstlisting}

\begin{lstlisting}[caption={siox\_eligibility/stats.py},label={lst:siox-src-stats}]
"""Small statistics helpers for the figures and reported numbers.

The manuscript figures report bootstrapped 95% confidence intervals over the
random seeds (the comp-neuro convention) rather than mean +/- s.d. ``bootstrap_ci``
is the single primitive used everywhere; ``summarise`` packages a mean with its CI
for convenient reporting.
"""
from __future__ import annotations

import numpy as np

__all__ = ["bootstrap_ci", "summarise"]


def bootstrap_ci(values, n_boot=10000, ci=95, seed=0):
    """Percentile bootstrap confidence interval for the mean of ``values``.

    Parameters
    ----------
    values : array-like
        Per-seed (or per-sample) scalar observations.
    n_boot : int
        Number of bootstrap resamples.
    ci : float
        Central confidence level in percent (95 -> 2.5/97.5 percentiles).
    seed : int
        RNG seed, so the interval is deterministic and reproducible.

    Returns
    -------
    (lo, hi) : tuple of float
        Lower and upper confidence bounds on the mean. For a single observation
        the bound collapses to that value.
    """
    v = np.asarray(values, dtype=float).ravel()
    if v.size == 0:
        return (float("nan"), float("nan"))
    if v.size == 1:
        return (float(v[0]), float(v[0]))
    rng = np.random.default_rng(seed)
    idx = rng.integers(0, v.size, size=(n_boot, v.size))
    means = v[idx].mean(axis=1)
    half = (100 - ci) / 2
    lo, hi = np.percentile(means, [half, 100 - half])
    return (float(lo), float(hi))


def summarise(values, n_boot=10000, ci=95, seed=0):
    """Return ``dict(mean, lo, hi, sd, n)`` for a per-seed observation array.

    ``lo``/``hi`` are the bootstrap CI bounds; ``sd`` is the sample standard
    deviation (kept for continuity with earlier mean +/- s.d. reporting).
    """
    v = np.asarray(values, dtype=float).ravel()
    lo, hi = bootstrap_ci(v, n_boot=n_boot, ci=ci, seed=seed)
    return {
        "mean": float(v.mean()) if v.size else float("nan"),
        "lo": lo,
        "hi": hi,
        "sd": float(v.std()) if v.size else float("nan"),
        "n": int(v.size),
    }
\end{lstlisting}

\begin{lstlisting}[caption={siox\_eligibility/distal\_reward.py},label={lst:siox-src-distal_reward}]
"""Distal-reward credit-assignment tasks (manuscript Figs 4 and 5).

These are the open-loop credit-assignment demonstrations that precede the
closed-loop bandit of :mod:`mrl_trace.bandit`: a cue synapse undergoes a
coincidence, a reward arrives after an action--reward delay ``D``, and
reward-uncorrelated distractor synapses are active throughout. A correct learner
raises the cue weight and leaves the distractors near baseline despite the delay.

Two settings, matching the paper:

- ``trace_level`` (Fig 4): the device eligibility trace is precomputed as a kernel
  (from :meth:`TransientGate.trace`) and applied by its lag at reward time; success
  is the cue-to-distractor weight ratio.
- ``spiking`` (Fig 5): a full LIF output neuron with the device gate integrated
  online per synapse from real Poisson spike coincidences; success is the
  saturation of the cue synapse toward its conductance bound.

Both expose a per-seed ``train_*`` and a vectorised ``sweep_*`` that returns
per-seed arrays (so bootstrapped CIs can be computed). The retention time
``tau_leak`` sets the maximum learnable delay in both.
"""
from __future__ import annotations

import numpy as np

from .device import TransientGate, tau_r, tau_d, K_STAGES
from .neurons import lif_step, TAU_M, V_TH
from .learning import LTD_BIAS

__all__ = [
    "trace_kernel", "abstract_kernel", "train_trace_level",
    "SpikingGateBank", "train_spiking", "cue_saturation", "W_INIT", "W_MAX",
    "trace_ratio", "run_trace_window",
]

W_INIT, W_MAX = 0.5, 1.0


# ----------------------------------------------------------------------------
# Trace-level task (Fig 4): precomputed eligibility kernel + lag-based credit
# ----------------------------------------------------------------------------
def trace_kernel(t_grid, tau_leak, V=0.9):
    """Device eligibility kernel e(t) on ``t_grid`` for a coincidence at t=0."""
    g = TransientGate(V=V, tau_leak=tau_leak, dt=float(t_grid[1] - t_grid[0]))
    return g.trace(t_grid, coincidence_at=0.0, coincidence_dur=0.3)


def abstract_kernel(t_grid, tau):
    """Control: an abstract exponential eligibility kernel."""
    return np.exp(-np.asarray(t_grid, float) / tau)


def train_trace_level(tau_leak, D, *, N=8, trials=1500, T=300.0, dt=0.1,
                      eta=0.05, seed=0, abstract=False, no_trace=False):
    """One seed of the trace-level distal-reward task; returns final weights ``w``
    (``w[0]`` cue, ``w[1:]`` distractors).

    ``abstract=True`` uses an exponential kernel instead of the device trace;
    ``no_trace=True`` is the control with a delta kernel (eligibility only at the
    exact reward instant, so no credit survives any delay ``D>0``)."""
    rng = np.random.default_rng(seed)
    t = np.arange(0, T, dt)
    nt = len(t)
    if no_trace:
        kern = np.zeros(nt)
        kern[0] = 1.0
    elif abstract:
        kern = abstract_kernel(t, tau_leak)
    else:
        kern = trace_kernel(t, tau_leak)
    w = np.full(N, W_INIT)
    baseline = 0.5
    for _ in range(trials):
        rewarded = rng.random() < 0.5
        t_reward = 0.1 * T + D
        if t_reward >= T:
            continue
        ri = int(t_reward / dt)
        e = np.zeros(N)
        if rewarded:
            lag = ri - int(0.1 * T / dt)
            e[0] = kern[lag] if 0 <= lag < nt else 0.0
        for i in range(1, N):
            t_d = rng.uniform(0.02 * T, 0.95 * T)
            lag = ri - int(t_d / dt)
            e[i] = kern[lag] if 0 <= lag < nt else 0.0
        R = 1.0 if rewarded else 0.0
        baseline += 0.02 * (R - baseline)
        w = np.clip(w + eta * (R - baseline) * e, 0.0, 2.0)
    return w


# ----------------------------------------------------------------------------
# Spiking task (Fig 5): online per-synapse device gate + LIF output neuron
# ----------------------------------------------------------------------------
class SpikingGateBank:
    """``N`` independent device gates (the Fig-5 per-synapse eligibility bank).

    Equivalent to a vectorised :class:`TransientGate` over an ``(N,)`` grid, with
    the signed/leak-dominant drive and ``[-Vnmax, Vnmax]`` bound."""

    def __init__(self, N, V=0.9, tau_leak=5.0, k=K_STAGES, dt=1e-3, Vnmax=1.0):
        self.N, self.k, self.dt, self.Vnmax = N, k, dt, Vnmax
        self.alpha = k / tau_r(V)
        self.tau_d = tau_d(V)
        self.tau_leak = tau_leak
        self.vn = np.zeros((N, k))
        self.vsc = np.zeros(N)

    def reset(self):
        self.vn[:] = 0.0
        self.vsc[:] = 0.0

    def step(self, drive):
        dt, a, Vm = self.dt, self.alpha, self.Vnmax
        new = self.vn.copy()
        prev = drive
        for j in range(self.k):
            new[:, j] = self.vn[:, j] + dt * (
                a * prev * (Vm - np.abs(self.vn[:, j])) - self.vn[:, j] / self.tau_leak)
            prev = self.vn[:, j]
        self.vsc += dt * (a * drive * (Vm - np.abs(self.vsc)) - self.vsc / self.tau_d)
        self.vn = np.clip(new, -Vm, Vm)
        return self.vn[:, -1] / Vm


def _spiking_trial(bank, w, rewarded, D, *, dt=1e-3, cue_rate=200.0, dist_rate=40.0,
                   N=8, rng=None, eta=0.05, baseline=0.5, ltd_bias=LTD_BIAS):
    cue_window = (0.5, 0.5 + 2.0)          # seconds-long correlated cue epoch
    bank.reset()
    v_out = 0.0
    ri = int((cue_window[1] + D) / dt)
    nsteps = ri + 2
    e_at_reward = np.zeros(N)
    for n in range(nsteps):
        t = n * dt
        rates = np.full(N, dist_rate)
        if rewarded and cue_window[0] <= t < cue_window[1]:
            rates[0] = cue_rate
        pre = rng.random(N) < rates * dt
        charge = np.dot(w, pre.astype(float))
        v_out, post = lif_step(v_out, charge, dt)
        drive = (pre.astype(float) * (1.0 if post else -ltd_bias)) / dt * 1e-3
        e = bank.step(drive)
        if n == ri:
            e_at_reward = e.copy()
    R = 1.0 if rewarded else 0.0
    w = np.clip(w + eta * (R - baseline) * e_at_reward, 0.0, W_MAX)
    return w, R


def train_spiking(tau_leak, D, *, N=8, trials=400, dt=1e-3, seed=0, eta=0.2):
    """One seed of the full-spiking distal-reward task; returns final weights ``w``
    (``w[0]`` cue saturation toward the bound is the success measure)."""
    rng = np.random.default_rng(seed)
    bank = SpikingGateBank(N, tau_leak=tau_leak, dt=dt)
    w = np.full(N, W_INIT)
    baseline = 0.5
    for _ in range(trials):
        rewarded = rng.random() < 0.5
        w, R = _spiking_trial(bank, w, rewarded, D, dt=dt, N=N, rng=rng,
                              baseline=baseline, eta=eta)
        baseline += 0.02 * (R - baseline)
    return w


def cue_saturation(w):
    """Saturation of the cue synapse toward its bound: (w0 - W_INIT)/(W_MAX - W_INIT)."""
    return (w[0] - W_INIT) / (W_MAX - W_INIT)


# =============================================================================
# Experiment cores (the distal-reward studies that compose ``train_trace_level``)
#
# Each ``run_*`` returns the result grid as a plain dict -- no file I/O, no
# plotting, no stdout.  Notebooks call these at a small (quick) seed count and
# render the figures inline; ``main()`` (below) calls them at the published
# 20-seed scale and writes the grid under ``data/results/`` for the notebooks to
# replay.  This mirrors the pattern in :mod:`mrl_trace.bandit`.
# =============================================================================

# Fig 4 (tier1) sweep axes: action->reward delay grid, the three device
# retention times, and the no-trace (delta-kernel) control. "learned" if the
# cue-to-distractor weight ratio reaches PASS. These constants match the
# 20-seed driver exactly so the published grid is bit-reproducible.
TIER1_DELAYS = [1, 2, 5, 10, 20, 40, 70, 100]
# (label, tau_leak, no_trace?) ; "none" is the no-trace control (delta kernel)
TIER1_VARIANTS = [("gate_tl20", 20.0, False), ("gate_tl5", 5.0, False),
                  ("gate_tl1", 1.0, False), ("none", None, True)]
TIER1_PASS = 2.0


def trace_ratio(w):
    """Cue-to-distractor weight ratio -- the trace-level success measure."""
    return w[0] / max(np.mean(w[1:]), 1e-6)


def _trace_cell(tau_leak, D, no_trace, seeds):
    """Per-seed cue/distractor ratios for one (variant, delay) cell (serial)."""
    if no_trace:
        # the no-trace control uses a delta kernel (no eligibility persistence)
        return np.array([trace_ratio(train_trace_level(0.0, D, seed=s, no_trace=True))
                         for s in range(seeds)])
    return np.array([trace_ratio(train_trace_level(tau_leak, D, seed=s))
                     for s in range(seeds)])


def run_trace_window(*, seeds=20, delays=TIER1_DELAYS, variants=TIER1_VARIANTS,
                     PASS=TIER1_PASS):
    """Fig 4 grid (tier1): trace-level credit-assignment window.

    Learned cue-to-distractor weight ratio versus action--reward delay ``D``, for
    three device retention times ``tau_leak in {20, 5, 1}`` s and a no-trace control
    (delta kernel), at ``seeds`` seeds with bootstrapped 95% CIs. The device trace is
    cheap here (precomputed kernel), so the seeds run sequentially.

    A cell is "learned" (contributes to ``max_learn``) if its mean ratio reaches
    ``PASS`` (=2.0): the cue synapse is at least twice the distractor weight, so
    credit is correctly assigned to the cue despite the delay. The maximum learnable
    delay per variant is the retention-window signature the manuscript reports.

    Returns the result dict (per-variant mean ratios + bootstrap CIs, max learnable
    delay per variant); no file I/O. Dict keys (``delays``, ``ratios``,
    ``ratios_ci``, ``max_learn``, ``n_seeds``) match the published grid exactly.
    """
    from .stats import bootstrap_ci
    ratios, ratios_ci = {}, {}
    for name, tl, no_trace in variants:
        row, row_ci = [], []
        for D in delays:
            r = _trace_cell(tl, D, no_trace, seeds)
            row.append(float(r.mean()))
            row_ci.append(bootstrap_ci(r))
        ratios[name] = row
        ratios_ci[name] = row_ci

    def _maxd(name):
        ds = [d for d, v in zip(delays, ratios[name]) if v >= PASS]
        return max(ds) if ds else 0

    return {"delays": list(delays), "ratios": ratios, "ratios_ci": ratios_ci,
            "max_learn": {name: _maxd(name) for name, _, _ in variants},
            "n_seeds": seeds}


def _trace_cell_star(args):
    """Pool worker: unpack a (variant, delay) cell and return its per-seed ratios."""
    name, tl, no_trace, D, seeds = args
    return (name, D, _trace_cell(tl, D, no_trace, seeds))


def run_trace_window_parallel(*, seeds=20, delays=TIER1_DELAYS,
                              variants=TIER1_VARIANTS, PASS=TIER1_PASS,
                              processes=None):
    """``run_trace_window`` computed with a ``multiprocessing.Pool`` over the
    coarse (variant x delay) cell axis (for the full-scale ``main()`` driver).

    Identical grid dict and science to :func:`run_trace_window`; only the cell loop
    is parallelised. NOT for in-notebook use -- notebooks call the serial core.
    """
    import multiprocessing as mp
    from .stats import bootstrap_ci
    jobs = [(name, tl, no_trace, D, seeds)
            for name, tl, no_trace in variants for D in delays]
    with mp.Pool(processes=processes) as pool:
        cells = pool.map(_trace_cell_star, jobs)
    per_cell = {(name, D): r for name, D, r in cells}

    ratios, ratios_ci = {}, {}
    for name, _tl, _nt in variants:
        row, row_ci = [], []
        for D in delays:
            r = per_cell[(name, D)]
            row.append(float(r.mean()))
            row_ci.append(bootstrap_ci(r))
        ratios[name] = row
        ratios_ci[name] = row_ci

    def _maxd(name):
        ds = [d for d, v in zip(delays, ratios[name]) if v >= PASS]
        return max(ds) if ds else 0

    return {"delays": list(delays), "ratios": ratios, "ratios_ci": ratios_ci,
            "max_learn": {name: _maxd(name) for name, _, _ in variants},
            "n_seeds": seeds}


def main(argv=None):
    """Full-scale reproduction CLI for the distal-reward grids (writes ``data/results``).

    ``python -m mrl_trace.distal_reward [--tier1] [--full|--quick]``
    With no experiment flag, runs all. ``--full`` = 20 seeds (published); ``--quick``
    = a fast few-seed smoke run. The trace-level cells are independent, so the
    full-scale run parallelises over the (variant x delay) grid with a Pool.
    """
    import argparse
    from . import paths
    ap = argparse.ArgumentParser(description="Distal-reward credit-assignment reproductions")
    ap.add_argument("--tier1", action="store_true",
                    help="Fig 4 trace-level window -> tier1_gate_results.npy")
    ap.add_argument("--quick", action="store_true", help="fast few-seed smoke run")
    ap.add_argument("--full", action="store_true", help="published 20-seed run (default)")
    a = ap.parse_args(argv)
    run_all = not a.tier1
    seeds = 4 if a.quick else 20

    if a.tier1 or run_all:
        print(f"=== Fig 4 trace-level, N={seeds} seeds, bootstrap 95% CIs ===")
        grid = run_trace_window_parallel(seeds=seeds)
        for name, _tl, _nt in TIER1_VARIANTS:
            print(f"{name:11s} " + " ".join(f"{v:5.2f}" for v in grid["ratios"][name]))
        print("max learnable delay:", grid["max_learn"])
        paths.save_result("tier1_gate_results.npy", grid)
        print("  wrote tier1_gate_results.npy")


if __name__ == "__main__":
    main()
\end{lstlisting}

\begin{lstlisting}[caption={siox\_eligibility/bandit.py},label={lst:siox-src-bandit}]
"""Contextual spiking bandit -- the closed-loop reinforcement-learning task.

This is genuine RL, not a credit-assignment microbenchmark: the network's *action*
determines the reward (closed loop), it must learn a state->action *policy*, and
success is the *reward rate* rising over trials.

Task (``S`` states x ``A`` actions):
  - ``S`` stimulus inputs encode the state; exactly one is active per trial.
  - ``A`` leaky integrate-and-fire action neurons compete; the chosen action is the
    one that spikes most over the cue epoch (membrane noise gives exploration).
  - ``S*A`` plastic device synapses ``w[state, action]``.
  - Rewarded mapping: state ``s`` -> action ``s mod A``.
  - Reward ``R in {0,1}`` delivered after delay ``D``, *contingent* on the chosen
    action matching the rewarded action for the presented state.
  - Per synapse, a signed leak-dominant pre x post coincidence drives the device
    eligibility gate; at reward, ``dw = eta (R - b) e``.

``train`` runs ``B`` seeds as a vectorised batch (leading axis ``B``) and returns a
``(B, trials)`` reward array.  ``reward_rate`` and ``trials_to_criterion`` reduce it.
The same call, with ``no_trace=True``, zeroes the eligibility (the device-necessity
control), and with ``abstract=True`` swaps the device gate for a plain leaky
integrator (the "is the device the bottleneck?" control of the scaling study).
"""
from __future__ import annotations

import numpy as np

from .device import tau_r, tau_d, K_STAGES
from .neurons import lif_step_batched, TAU_M, V_TH
from .learning import LTD_BIAS

__all__ = ["GateBankBatched", "AbstractTrace", "train", "reward_rate",
           "trials_to_criterion", "W_INIT", "W_MAX"]

W_INIT, W_MAX = 0.5, 1.5


class GateBankBatched:
    """Device eligibility gates for ``B`` parallel ``(S, A)`` synapse grids.

    Vectorised form of :class:`mrl_trace.device.TransientGate` (Section 5.3.2
    physics): ``k`` cascade trap nodes plus one space-charge node per synapse, signed
    leak-dominant drive, trace bounded in ``[-Vnmax, Vnmax]`` (the LTD wing).
    """

    def __init__(self, B, S, A, tau_leak=10.0, V=0.9, k=K_STAGES, dt=5e-3, Vnmax=1.0,
                 beta_leak=1.0):
        self.B, self.S, self.A, self.k, self.dt, self.Vnmax = B, S, A, k, dt, Vnmax
        self.alpha = k / tau_r(V)
        self.tau_d = tau_d(V)
        # Dispersion of the trap-DISCHARGE. beta_leak=1 is the single-rate (mono-exponential)
        # leak used throughout the learning results; beta_leak<1 is the measured dispersive
        # (stretched-exponential) discharge, instantaneous rate (beta/tau)(t/tau)^(beta-1) where
        # t is time-since-last-drive (per synapse). beta=1 recovers 1/tau exactly (bit-identical),
        # so this is a backward-compatible sensitivity knob: the device study measures
        # beta_leak~0.85 field-free, and this lets the learning results be re-run at that value
        # to confirm the single-rate idealisation does not change them.
        self.beta_leak = float(beta_leak)
        # per-synapse discharge clock, shape (B,S,A,1) to compose with vn (B,S,A,k) and a
        # per-synapse tau_leak (1,S,A,1); only allocated when the dispersive discharge is used.
        self._t_since = np.full((B, S, A, 1), dt) if beta_leak != 1.0 else None
        # ``tau_leak`` may be a scalar (one retention for every synapse, the default) or a
        # per-synapse array of shape (S, A) -- a heterogeneous retention population, e.g.
        # sampled from the measured ITO trap-discharge distribution (exp19). An (S, A) array
        # is reshaped to (1, S, A, 1) so it broadcasts against the cascade state vn of shape
        # (B, S, A, k) in step(); a scalar broadcasts trivially. Backward-compatible: a
        # scalar reproduces the original behaviour bit-for-bit.
        tl = np.asarray(tau_leak, dtype=float)
        self.tau_leak = tl.reshape(1, S, A, 1) if tl.ndim == 2 else tl
        self.vn = np.zeros((B, S, A, k))
        self.vsc = np.zeros((B, S, A))

    def reset(self):
        self.vn[:] = 0.0
        self.vsc[:] = 0.0
        if self._t_since is not None:
            self._t_since[:] = self.dt

    def step(self, drive):                 # drive: (B, S, A) signed
        # Vectorised over the k cascade stages. Each stage j advances from the PREVIOUS
        # stage's OLD value (the input ``drive`` for j=0, ``vn[...,j-1]`` for j>0), so all
        # stages update from ``self.vn`` in one fused forward-Euler expression rather than a
        # Python loop with a per-stage write into a full-array copy. Verified bit-for-bit
        # identical to the staged loop (and against the gate behaviour tests). NOTE: this is
        # a readability/structure win, not a speed one -- profiling shows the cost is the
        # memory-bandwidth-bound elementwise math over the (B,S,A,k) array, which is
        # irreducible at fixed dt and k; the step still dominates runtime.
        dt, a, Vm = self.dt, self.alpha, self.Vnmax
        vn = self.vn
        prev = np.empty_like(vn)
        prev[..., 0] = drive                       # stage 0 driven by the input coincidence
        prev[..., 1:] = vn[..., :-1]               # stage j>0 driven by its old upstream node
        if self.beta_leak == 1.0:
            leak_rate = 1.0 / self.tau_leak                       # single-rate (default)
        else:
            # dispersive stretched-exponential discharge: per-synapse instantaneous rate
            # (beta/tau)(t_since/tau)^(beta-1), broadcast over the k stages. A coincidence on a
            # synapse (|drive|>0) resets that synapse's discharge clock to dt. Work in the
            # (B,S,A,1) shape so it composes with both a scalar tau_leak and a per-synapse
            # (1,S,A,1) tau_leak array.
            self._t_since = np.where(np.abs(drive)[..., None] > 1e-9, dt, self._t_since + dt)
            tau = self.tau_leak                                   # scalar or (1,S,A,1)
            leak_rate = (self.beta_leak / tau) * np.power(
                np.clip(self._t_since / tau, 1e-6, None), self.beta_leak - 1.0)  # (B,S,A,1)
        new = vn + dt * (a * prev * (Vm - np.abs(vn)) - vn * leak_rate)
        np.clip(new, -Vm, Vm, out=new)
        self.vsc += dt * (a * drive * (Vm - np.abs(self.vsc)) - self.vsc / self.tau_d)
        self.vn = new
        return new[..., -1] / Vm


class AbstractTrace:
    """Control: an abstract exponential eligibility trace (the hand-set kernel of
    prior algorithmic work), matched in decay timescale to the device ``tau_leak``.
    Same signed drive and readout interface as :class:`GateBankBatched`; only the
    dynamics differ (a single leaky integrator, no sigmoidal rise)."""

    def __init__(self, B, S, A, tau_elig=10.0, dt=5e-3, **_ignored):
        self.e = np.zeros((B, S, A))
        self.dt = dt
        self.tau = tau_elig

    def reset(self):
        self.e[:] = 0.0

    def step(self, drive):
        self.e += self.dt * (drive - self.e / self.tau)
        return np.clip(self.e, -1.0, 1.0)


def train(S, A, *, B=6, tau_leak=10.0, D=2.0, trials=1500, dt=5e-3, cue_dur=1.0,
          eta=0.2, in_rate=200.0, ltd=LTD_BIAS, tau_m=TAU_M, v_th=V_TH,
          sigma0=0.15, sigma1=None, abstract=False, no_trace=False, seed0=0,
          reward_pools=None):
    """Train the contextual bandit on ``B`` parallel seeds; return rewards ``(B, trials)``.

    Parameters mirror the manuscript's experiments:
      ``D``         action->reward delay (s);
      ``tau_leak``  device retention / credit-assignment window (s);
      ``sigma0``    membrane-noise (exploration) level; if ``sigma1`` is given the
                    noise is annealed linearly ``sigma0 -> sigma1`` over training;
      ``abstract``  use :class:`AbstractTrace` instead of the device gate;
      ``no_trace``  zero the eligibility every step (device-necessity control).
      ``reward_pools`` biosignal reward: if given, a dict {1: array, 0: array} of
                    per-trial reward VALUES decoded from real EEG (out-of-fold
                    predictions, keyed by the TRUE outcome valence). On each trial the
                    reward actually delivered to the three-factor rule is sampled from
                    the pool matching the true outcome, so R is a real, noisy,
                    biologically-measured reward-prediction-error rather than the clean
                    synthetic {0,1}. The reported reward rate still uses the TRUE outcome
                    (task performance), not the noisy gate. Default None = synthetic.

    Two reward streams are tracked: ``r_true`` (did the action match the rewarded
    action -- task performance, what is RETURNED) and ``r_gate`` (what actually drives
    the weight update -- synthetic = r_true, or the EEG-decoded value when
    ``reward_pools`` is given).
    """
    rng = np.random.default_rng(seed0)
    Bank = AbstractTrace if abstract else GateBankBatched
    bank = Bank(B, S, A, tau_leak=tau_leak, dt=dt)
    w = np.full((B, S, A), W_INIT)
    correct = np.array([s % A for s in range(S)])      # rewarded action per state
    baseline = np.full(B, 1.0 / A)                     # baseline tracks chance for this A
    cue = (0.3, 0.3 + cue_dur)
    ri = int((cue[1] + D) / dt)
    nsteps = ri + 2
    rewards = np.zeros((B, trials))
    bidx = np.arange(B)
    for tr in range(trials):
        sigma = sigma0 if sigma1 is None else sigma0 + (sigma1 - sigma0) * tr / trials
        bank.reset()
        state = rng.integers(S, size=B)                # per-seed state this trial
        v = np.zeros((B, A))
        spk = np.zeros((B, A))
        e_rew = np.zeros((B, S, A))
        for n in range(nsteps):
            t = n * dt
            pre = np.zeros((B, S))
            if cue[0] <= t < cue[1]:
                pre[bidx, state] = (rng.random(B) < in_rate * dt).astype(float)
            charge = np.einsum('bsa,bs->ba', w, pre)   # bitline currents = G^T V
            v, sp = lif_step_batched(v, charge, dt, rng, tau_m=tau_m, v_th=v_th, noise=sigma)
            spk += sp
            drive = np.zeros((B, S, A))
            drive[bidx, state, :] = pre[bidx, state][:, None] * np.where(sp, 1.0, -ltd)
            if no_trace:
                drive[:] = 0.0
            e = bank.step(drive)
            if n == ri:
                e_rew = e.copy()
        # action selection: argmax spikes; ties (incl. all-zero) -> random
        tie = spk.max(axis=1) == spk.min(axis=1)
        chosen = np.argmax(spk, axis=1)
        chosen[tie] = rng.integers(A, size=tie.sum())
        r_true = (chosen == correct[state]).astype(float)   # task performance (returned)
        if reward_pools is None:
            r_gate = r_true                                  # synthetic clean reward
        else:
            # biosignal reward: for each seed, draw a decoded reward VALUE from the EEG
            # pool matching the TRUE outcome valence -- a real, noisy reward-prediction-
            # error gates the update, while r_true still measures task performance.
            r_gate = np.empty(B)
            for b in range(B):
                pool = reward_pools[int(r_true[b])]
                r_gate[b] = pool[rng.integers(len(pool))]
        adv = eta * (r_gate - baseline)                # signed three-factor update
        w = np.clip(w + adv[:, None, None] * e_rew, 0.0, W_MAX)
        baseline += 0.02 * (r_gate - baseline)
        rewards[:, tr] = r_true                         # report TRUE performance
    return rewards


def reward_rate(rewards, window=100):
    """Mean reward over the last ``window`` trials, per seed (axis -1)."""
    rewards = np.asarray(rewards)
    if rewards.shape[-1] < window:
        return rewards.mean(axis=-1)
    return rewards[..., -window:].mean(axis=-1)


def trials_to_criterion(rew_1d, crit, window=100):
    """First trial index at which the running ``window``-reward rate reaches ``crit``
    for a single seed's reward sequence, or ``None`` if it never does."""
    rew_1d = np.asarray(rew_1d)
    if len(rew_1d) < window:
        return None
    csum = np.cumsum(rew_1d)
    rr = (csum[window:] - csum[:-window]) / window
    if not np.any(rr >= crit):
        return None
    return int(np.argmax(rr >= crit) + window)


# =============================================================================
# Experiment cores (the closed-loop RL studies that compose ``train``)
#
# Each ``run_*`` returns the result grid as a plain dict -- no file I/O, no
# plotting, no stdout.  Notebooks call these at a small (quick) seed/trial count
# and render the figures inline; ``main()`` (below) calls them at the published
# 20-seed scale and writes the grid under ``data/results/`` for the notebooks to
# replay.  Reversal (Experiment 18) is a ``train`` variant, so it lives here.
# =============================================================================


def _mapping(S, A, shift):
    """Rewarded action per state, cyclically shifted by ``shift`` (0 = identity)."""
    return np.array([(s + shift) % A for s in range(S)])


def run_reversal(cond, *, S=2, A=2, B=20, tau_leak=10.0, D=2.0, trials=3000,
                 n_phases=2, dt=5e-3, cue_dur=1.0, eta=0.2, in_rate=200.0,
                 ltd=LTD_BIAS, tau_m=TAU_M, v_th=V_TH, sigma=0.15, seed0=0):
    """Reversal learning (Experiment 18): the signed rule must UNLEARN a stale
    contingency.  Mirrors :func:`train` exactly, but the rewarded mapping is
    cyclically shifted by one at each of ``n_phases`` equal-length phase boundaries,
    so the agent must repeatedly unlearn the stale contingency and acquire the new one.

    ``cond`` is ``"device"`` (trap-discharge gate), ``"abstract"`` (matched-tau single
    exponential, the recency-trace baseline) or ``"no_trace"`` (eligibility zeroed,
    the device-necessity control).

    Returns ``(rewards (B, trials), flips list[int])`` where ``rewards[:, t]`` is the
    TRUE task performance under the contingency in force at trial ``t``.
    """
    rng = np.random.default_rng(seed0)
    if cond == "abstract":
        bank = AbstractTrace(B, S, A, tau_elig=tau_leak, dt=dt)
    else:
        bank = GateBankBatched(B, S, A, tau_leak=tau_leak, dt=dt)
    no_trace = (cond == "no_trace")

    w = np.full((B, S, A), W_INIT)
    baseline = np.full(B, 1.0 / A)
    cue = (0.3, 0.3 + cue_dur)
    ri = int((cue[1] + D) / dt)
    nsteps = ri + 2
    rewards = np.zeros((B, trials))
    bidx = np.arange(B)

    phase_len = trials // n_phases
    flips = [phase_len * p for p in range(1, n_phases)]

    for tr in range(trials):
        phase = min(tr // phase_len, n_phases - 1)
        correct = _mapping(S, A, phase)              # cyclic-shifted mapping this phase
        bank.reset()
        state = rng.integers(S, size=B)
        v = np.zeros((B, A)); spk = np.zeros((B, A))
        e_rew = np.zeros((B, S, A))
        for n in range(nsteps):
            t = n * dt
            pre = np.zeros((B, S))
            if cue[0] <= t < cue[1]:
                pre[bidx, state] = (rng.random(B) < in_rate * dt).astype(float)
            charge = np.einsum('bsa,bs->ba', w, pre)
            v, sp = lif_step_batched(v, charge, dt, rng, tau_m=tau_m, v_th=v_th, noise=sigma)
            spk += sp
            drive = np.zeros((B, S, A))
            drive[bidx, state, :] = pre[bidx, state][:, None] * np.where(sp, 1.0, -ltd)
            if no_trace:
                drive[:] = 0.0
            e = bank.step(drive)
            if n == ri:
                e_rew = e.copy()
        tie = spk.max(axis=1) == spk.min(axis=1)
        chosen = np.argmax(spk, axis=1)
        chosen[tie] = rng.integers(A, size=int(tie.sum()))
        r_true = (chosen == correct[state]).astype(float)
        adv = eta * (r_true - baseline)
        w = np.clip(w + adv[:, None, None] * e_rew, 0.0, W_MAX)
        baseline += 0.02 * (r_true - baseline)
        rewards[:, tr] = r_true
    return rewards, flips


def reversal_phase_final(rw, flips, trials, window=200):
    """Per-seed reward rate in the last ``window`` trials of phase 1 (pre-reversal)
    and of the final phase (post-reversal)."""
    first_flip = flips[0]
    pre = rw[:, max(0, first_flip - window):first_flip].mean(1)
    post = rw[:, trials - window:].mean(1)
    return pre, post


def reversal_recriterion(rw_1d, flip, crit, window=100):
    """Trials AFTER the flip until the running reward rate first reaches ``crit`` again
    (post-flip re-acquisition speed); ``None`` if never re-acquired within the phase."""
    seg = rw_1d[flip:]
    if len(seg) < window:
        return None
    csum = np.cumsum(seg)
    rr = (csum[window:] - csum[:-window]) / window
    hit = np.argmax(rr >= crit) if np.any(rr >= crit) else None
    return int(hit + window) if hit is not None else None


def run_learning_and_window(*, seeds=20, delays=(1, 2, 5, 10, 20),
                            tau_leaks=(10.0, 2.0, 0.5), D0=5.0, trials=600):
    """Fig 6 grid (tier3): 2-state bandit learning curve + delay x retention table.

    Returns the result dict (per-seed learning curves, delay x retention reward-rate
    table with bootstrap CIs, max learnable delay per tau); no file I/O.
    """
    from .stats import bootstrap_ci
    S, A = 2, 2
    dev = train(S, A, B=seeds, tau_leak=10.0, D=D0, trials=trials)
    nt = train(S, A, B=seeds, tau_leak=10.0, D=D0, trials=trials, no_trace=True)
    dev_fr, nt_fr = reward_rate(dev), reward_rate(nt)

    rates, rates_ci = {}, {}
    for tl in tau_leaks:
        row, row_ci = [], []
        for D in delays:
            fr = reward_rate(train(S, A, B=seeds, tau_leak=tl, D=D, trials=trials))
            row.append(float(fr.mean()))
            row_ci.append(bootstrap_ci(fr))
        rates[tl] = row
        rates_ci[tl] = row_ci

    def _maxd(tl):
        ds = [d for d, v in zip(delays, rates[tl]) if v >= 0.75]
        return max(ds) if ds else 0

    def _ci(arr):
        a = np.asarray(arr, float)
        lo, hi = bootstrap_ci(a)
        return (float(a.mean()), float(a.std()), lo, hi)

    return {
        "delays": list(delays), "reward_rate": rates, "reward_rate_ci": rates_ci,
        "max_learn": {tl: _maxd(tl) for tl in rates}, "D0": D0, "n_seeds": seeds,
        "curve_device": dev, "curve_notrace": nt,
        "device_final": _ci(dev_fr), "notrace_final": _ci(nt_fr),
    }


def run_scaling(*, seeds=20, grid=((2, 2), (4, 2), (4, 4), (8, 4), (8, 8), (12, 8)),
                D=2.0):
    """Fig 8a grid (tier4): does the policy still converge as ``S x A`` grows?

    Criterion (pre-registered, no goalpost moving): reward rate >= 0.5*(1 + 1/A).
    Returns a dict keyed by ``(S, A)`` with device/no-trace means + bootstrap CIs,
    trials-to-criterion, and pass/fail; no file I/O.
    """
    from .stats import bootstrap_ci

    def _ci(arr):
        a = np.asarray(arr, float)
        lo, hi = bootstrap_ci(a)
        return (float(a.mean()), float(a.std()), lo, hi)

    results = {}
    for S, A in grid:
        chance = 1.0 / A
        crit = 0.5 * (1 + chance)
        dev = train(S, A, B=seeds, D=D)
        nt = train(S, A, B=seeds, D=D, no_trace=True)
        dev_fr, nt_fr = reward_rate(dev), reward_rate(nt)
        ttc = [trials_to_criterion(dev[b], crit) for b in range(seeds)]
        ok = [t for t in ttc if t is not None]
        results[(S, A)] = dict(
            chance=chance, crit=crit, device=_ci(dev_fr), no_trace=_ci(nt_fr),
            trials_to_crit=int(np.median(ok)) if ok else None,
            n_converged=len(ok), passed=bool(dev_fr.mean() >= crit), n_seeds=seeds)
    return results


def run_remedies(*, seeds=20):
    """Fig 8b (tier5): five remedies on the failing 8x8 case.

    Returns a dict keyed by remedy label with final reward-rate mean + bootstrap CI
    and pass/fail against the 8x8 criterion; no file I/O.
    """
    from .stats import bootstrap_ci
    S, A = 8, 8
    crit = 0.5 * (1 + 1.0 / A)
    rows = [
        ("R0 baseline (device, undirected, 1500)", dict(trials=1500)),
        ("R1 more budget (device, undirected, 6000)", dict(trials=6000)),
        ("R2 directed expl (device, sigma 0.45->0.05, 1500)",
         dict(trials=1500, sigma0=0.45, sigma1=0.05)),
        ("R3 abstract trace (NOT device, undirected, 1500)", dict(trials=1500, abstract=True)),
        ("R4 directed+budget (device, sigma anneal, 6000)",
         dict(trials=6000, sigma0=0.45, sigma1=0.05)),
    ]
    out = {}
    for name, kw in rows:
        fr = reward_rate(train(S, A, B=seeds, D=2.0, **kw))
        a = np.asarray(fr, float); lo, hi = bootstrap_ci(a)
        out[name] = dict(final=(float(a.mean()), float(a.std()), lo, hi),
                         passed=bool(a.mean() >= crit), n_seeds=seeds)
    return out


def _summarize_reversal(raw, conds, trials, crit, chance):
    """Package raw per-condition ``run_reversal`` output into the saved grid dict
    (mean curves, pre/post finals + CIs, re-acquisition cost, recovery bins,
    pre-registered H1-H4/K1)."""
    from .stats import bootstrap_ci
    flips = raw[conds[0]][1]
    B = raw[conds[0]][0].shape[0]
    curves, pre_f, post_f, ci_pre, ci_post, recrit, recovery = {}, {}, {}, {}, {}, {}, {}
    for c in conds:
        rw, fl = raw[c]
        curves[c] = rw.mean(0)
        pre, post = reversal_phase_final(rw, fl, trials)
        pre_f[c], post_f[c] = pre, post
        ci_pre[c], ci_post[c] = bootstrap_ci(pre), bootstrap_ci(post)
        tt = [reversal_recriterion(rw[b], fl[0], crit) for b in range(B)]
        solved = [x for x in tt if x is not None]
        recrit[c] = (float(np.median(solved)) if solved else None, len(solved), B)
        seg = rw[:, fl[0]:].mean(0)
        recovery[c] = [float(seg[i:i + 500].mean()) for i in range(0, len(seg), 500)]
    dev_pre, dev_post = pre_f["device"].mean(), post_f["device"].mean()
    nt_post = post_f["no_trace"].mean()
    dev_rc, abs_rc = recrit["device"][0], recrit["abstract"][0]
    ratio = (dev_rc / abs_rc) if (dev_rc and abs_rc) else float("nan")
    criteria = {"H1": bool(dev_pre >= crit), "H2": bool(dev_post >= crit),
                "H3": bool(nt_post <= chance + 0.10), "K1": bool(dev_post <= chance + 0.10)}
    return {"curves": curves, "pre": pre_f, "post": post_f, "recovery": recovery,
            "ci_pre": ci_pre, "ci_post": ci_post, "recrit": recrit, "flips": flips,
            "seeds": B, "trials": trials, "n_phases": len(flips) + 1, "S": 2, "A": 2,
            "tau_leak": 10.0, "chance": chance, "crit": crit,
            "reacq_cost": {"device": dev_rc, "abstract": abs_rc, "ratio": ratio},
            "criteria": criteria}


def main(argv=None):
    """Full-scale reproduction CLI for the bandit-family grids (writes ``data/results``).

    ``python -m mrl_trace.bandit [--exp18] [--bandit] [--full|--quick]``
    With no experiment flag, runs all. ``--full`` = 20 seeds (published); ``--quick``
    = a fast few-seed smoke run.
    """
    import argparse
    from . import paths
    ap = argparse.ArgumentParser(description="Bandit-family RL reproductions")
    ap.add_argument("--exp18", action="store_true", help="reversal learning -> exp18_reversal.npy")
    ap.add_argument("--bandit", action="store_true",
                    help="fig6/8a/8b -> tier3/tier4/tier5_results.npy")
    ap.add_argument("--quick", action="store_true", help="fast few-seed smoke run")
    ap.add_argument("--full", action="store_true", help="published 20-seed run (default)")
    a = ap.parse_args(argv)
    run_all = not (a.exp18 or a.bandit)
    seeds = 6 if a.quick else 20

    if a.bandit or run_all:
        print(f"=== bandit reruns (N={seeds}, bootstrap 95% CIs) ===")
        paths.save_result("tier3_results.npy", run_learning_and_window(seeds=seeds))
        paths.save_result("tier4_results.npy", run_scaling(seeds=seeds))
        paths.save_result("tier5_results.npy", run_remedies(seeds=seeds))
        print("  wrote tier3/tier4/tier5_results.npy")

    if a.exp18 or run_all:
        B = 6 if a.quick else 20
        trials = 4000 if a.quick else 8000
        conds = ("device", "abstract", "no_trace")
        chance, crit = 0.5, 0.75
        print(f"=== reversal (Experiment 18): {B} seeds, {trials} trials ===")
        raw = {c: run_reversal(c, B=B, trials=trials, n_phases=2, tau_leak=10.0)
               for c in conds}
        grid = _summarize_reversal(raw, conds, trials, crit, chance)
        paths.save_result("exp18_reversal.npy", grid)
        print(f"  wrote exp18_reversal.npy  criteria={grid['criteria']}")


if __name__ == "__main__":
    main()
\end{lstlisting}

\begin{lstlisting}[caption={siox\_eligibility/maze.py},label={lst:siox-src-maze}]
"""Sequential distal-reward navigation -- a multi-step MDP credit-assignment task.

The contextual bandit of :mod:`mrl_trace.bandit` is *one step*: a single
state, a single action, an immediate (delayed) reward. That isolates the
credit-assignment mechanism but invites the criticism that it is an
association lookup rather than reinforcement learning of a *policy over a
trajectory*. This module adds the missing piece: a sequential Markov decision
process in which an action must be chosen at each of several states, reward is
delivered only at the goal after the whole trajectory, and credit must propagate
back across the intermediate steps. That backward propagation across a trajectory
is exactly what an eligibility trace exists for, so it is the natural task on which
to test a device-supplied trace and to benchmark it against an algorithmic one.

Two environments, sharing one training harness:

- :class:`LinearTrack` -- an ``L``-state corridor; at each state the agent chooses
  FORWARD (towards the goal) or BACK. Only a sustained forward run reaches the goal
  and is rewarded, after an action--reward delay ``D`` measured from the goal-entry
  action. A greedy one-step rule cannot short-circuit it: every step must carry the
  right action, so the policy spans the trajectory.
- :class:`TMaze` -- an ``L``-state stem leading to a junction with ``A_goal`` arms,
  exactly one rewarded. The agent advances along the stem, then selects an arm; the
  reward (after delay ``D``) is contingent on the arm matching the rewarded arm for
  the episode's cue. This makes the learned object a state->action *policy* whose
  final decision sits a whole trajectory away from the reward.

DESIGN NOTE (anti-"relabelled bandit"). The eligibility trace is accumulated
*online across the steps of an episode* on the synapses actually used at each state,
and a single delayed goal reward gates the surviving trace into all of them at once.
Because the device relaxes over seconds while steps are taken over the same
timescale, the trace from EARLY (distal) steps has decayed more than from LATE steps
when the reward lands -- so bridging the full trajectory genuinely requires a
retention long enough to span it. Short retention learns only the steps near the
goal; this is the sequential analogue of the bandit's delay window and the lever the
benchmark and the D_max-vs-tau prediction both exercise.

Conventions follow the rest of the package: ``B`` seeds run as a vectorised batch,
device gate via :class:`GateBankBatched` (``abstract=True`` swaps the exponential
R-STDP kernel, ``no_trace=True`` zeroes eligibility), signed leak-dominant
coincidence, three-factor update ``dw = eta (R - b) e``, ``dt = 5e-3`` s.
"""
from __future__ import annotations

import numpy as np

from .bandit import GateBankBatched, AbstractTrace, W_INIT, W_MAX
from .neurons import lif_step_batched, TAU_M, V_TH
from .learning import LTD_BIAS

__all__ = [
    "LinearTrack", "TMaze", "train_sequential",
    "reward_rate", "trials_to_criterion", "policy_correct",
    # experiment cores (composed of TMaze + train_sequential)
    "run_sequential", "run_dmax_law", "run_long_horizon", "run_long_horizon_faults",
    # analysis helpers
    "running_rate", "interp_dmax",
]


# ----------------------------------------------------------------------------
# Environments
# ----------------------------------------------------------------------------
class LinearTrack:
    """``L``-state corridor MDP. Actions: 0 = forward, 1 = back.

    State index advances on a forward action and retreats on a back action; the
    episode ends with reward when the agent steps forward out of the last state
    (reaching the goal), or unrewarded if it runs out of steps. ``n_states`` is the
    state-encoding width and ``n_actions = 2``.
    """

    def __init__(self, L=3):
        self.L = L
        self.n_states = L
        self.n_actions = 2
        self.goal_action = 0          # forward at the last state reaches the goal

    def start(self, B):
        return np.zeros(B, dtype=int)     # all seeds start at state 0

    def step(self, pos, action):
        """Vectorised transition. ``pos`` (B,), ``action`` (B,) in {0,1}.

        Returns ``(next_pos, reached_goal, done)`` boolean/int arrays of shape (B,).
        """
        forward = action == 0
        nxt = np.where(forward, pos + 1, np.maximum(pos - 1, 0))
        reached = nxt >= self.L          # stepped forward out of the last state
        nxt = np.where(reached, self.L - 1, nxt)   # clamp (episode ends anyway)
        return nxt, reached, reached

    def correct_action(self, pos):
        """The action that makes progress towards the goal from ``pos`` (always
        forward here); used only to score policy correctness, never during learning."""
        return np.zeros_like(pos)


class TMaze:
    """``L``-state stem then a ``A_goal``-arm junction; one arm rewarded per cue.

    States ``0..L-1`` are the stem (advance with action 0). At the junction state
    ``L-1`` the agent's action selects an arm ``0..A_goal-1``; reward is contingent on
    the arm matching the episode's rewarded arm (set by a cue presented at the start).
    ``n_actions = max(2, A_goal)`` so the same action neurons serve stem-advance and
    arm-choice.
    """

    def __init__(self, L=3, A_goal=2):
        self.L = L
        self.A_goal = A_goal
        self.n_states = L + A_goal       # stem states + one state per cue context
        self.n_actions = max(2, A_goal)
        self.goal_action = None          # depends on the cued arm (set per episode)

    def start(self, B, rng):
        self.cue = rng.integers(self.A_goal, size=B)   # rewarded arm this episode
        return np.zeros(B, dtype=int)

    def encode(self, pos):
        """State encoding: stem position, or (at the junction) a cue-specific index so
        the policy can route to the cued arm. Returns an index in ``[0, n_states)``."""
        at_junction = pos >= self.L - 1
        return np.where(at_junction, self.L + self.cue, pos)

    def step(self, pos, action):
        at_junction = pos >= self.L - 1
        # The stem is a one-way corridor: the agent AUTO-ADVANCES through it regardless
        # of action (navigation, not a learned decision). The only learned choice -- and
        # the only credit-assignment problem -- is the arm selection at the junction.
        # This keeps the reward DISTAL (cue at start, decision a trajectory later) while
        # removing the artefact of stem-advance reinforcing one action everywhere.
        nxt = np.where(at_junction, pos, pos + 1)
        done = at_junction
        reached = at_junction & (action == self.cue)    # correct arm -> reward
        return nxt, reached, done

    def correct_action(self, pos):
        at_junction = pos >= self.L - 1
        return np.where(at_junction, self.cue, 0)


# ----------------------------------------------------------------------------
# Training harness (shared across device / R-STDP / no-trace)
# ----------------------------------------------------------------------------
class EpropTrace:
    """Reward-based e-prop eligibility (Bellec et al. 2020), the third benchmark.

    Unlike the device gate (a physical relaxation of a signed coincidence) and R-STDP
    (a hand-set exponential filter of that coincidence), e-prop *computes* its
    eligibility from the network's own dynamics. For a LIF action neuron j fed by
    state-line i, the per-synapse eligibility trace is

        e_ij(t) = LP_kappa[ psi_j(t) * zbar_i(t) ],

    where ``zbar_i`` is the low-pass-filtered presynaptic spike train (membrane filter,
    tau_m), ``psi_j`` is the pseudo-derivative of the spike nonlinearity evaluated on the
    PRE-RESET membrane (the surrogate gradient, a triangular/Gaussian bump at threshold),
    and ``LP_kappa`` is a slow output filter (the eligibility time constant, played here
    by ``tau_leak`` so the three methods are compared at a matched retention). The signed
    LEARNING SIGNAL (R - b) supplies credit/blame at reward; e-prop does NOT use a signed
    coincidence (that is the R-STDP/device mechanism), so its drive is the unsigned
    product psi*zbar. This is the shallow, single-decision-layer reduction of reward-based
    e-prop: faithful to the per-synapse eligibility, but it does not exercise e-prop's
    recurrent-credit machinery (a property of the feedforward task, stated plainly).
    """

    def __init__(self, B, S, A, tau_leak=10.0, dt=5e-3, tau_m=TAU_M, v_th=V_TH,
                 psi_width=0.3, dampening=0.3, **_ignored):
        self.B, self.S, self.A, self.dt = B, S, A, dt
        self.tau_m, self.v_th = tau_m, v_th
        self.psi_width = psi_width          # surrogate-gradient width
        self.gamma = dampening              # surrogate-gradient peak (Bellec's gamma)
        self.tau = tau_leak                 # output (eligibility) filter time constant
        self.Vnmax = 1.0                    # interface parity with the device gate
        self.zbar = np.zeros((B, S))        # filtered presynaptic trace per state-line
        self.elig = np.zeros((B, S, A))     # filtered eligibility per synapse
        self.vn = np.zeros((B, S, A, 1))    # mirror of device .vn[...,-1] for readout

    def reset(self):
        self.zbar[:] = 0.0
        self.elig[:] = 0.0
        self.vn[:] = 0.0

    def step_eprop(self, pre, v_pre, active):
        """Advance one dt-tick. ``pre`` (B,S) presynaptic spikes this tick, ``v_pre``
        (B,A) pre-reset membrane, ``active`` (B,) episode mask. Returns nothing; the
        eligibility is read from ``self.vn[...,-1]`` for parity with the device gate."""
        dt = self.dt
        # presynaptic membrane-filtered trace zbar_i
        self.zbar += (-self.zbar / self.tau_m + pre) * dt
        # pseudo-derivative psi_j on the pre-reset membrane (triangular surrogate):
        # gamma * max(0, 1 - |v - v_th| / width)
        psi = self.gamma * np.maximum(0.0, 1.0 - np.abs(v_pre - self.v_th) / self.psi_width)
        # per-synapse instantaneous eligibility = psi_j * zbar_i  (outer over S x A)
        inst = self.zbar[:, :, None] * psi[:, None, :]          # (B,S,A)
        # slow output filter -> the eligibility trace (time constant tau)
        self.elig += (-self.elig / self.tau + inst) * dt
        self.elig *= active[:, None, None]
        self.vn[..., 0] = self.elig


def _relax(bank, zero_drive, n, stride, abstract):
    """Advance an undriven gate for ``n`` blocks of ``stride`` dt-ticks each, using a
    temporarily coarsened timestep for speed. With zero drive the cascade/leaky
    dynamics are smooth, so a coarse forward-Euler step is accurate provided
    stride*dt/tau << 1 (checked: stride=10, dt=5e-3, tau>=2 s -> step ratio <=0.025).
    Restores the gate's original dt afterwards."""
    if n <= 0:
        return
    dt0 = bank.dt
    bank.dt = dt0 * stride
    try:
        for _ in range(n):
            bank.step(zero_drive)
    finally:
        bank.dt = dt0


def _relax_eprop(bank, B, S, A, n, stride):
    """Coarse-stepped zero-input relaxation of the e-prop eligibility filters (mirror
    of :func:`_relax` for the EpropTrace, whose decay is equally smooth with no input)."""
    if n <= 0:
        return
    zp, zv, act = np.zeros((B, S)), np.zeros((B, A)), np.ones(B)
    dt0 = bank.dt
    bank.dt = dt0 * stride
    try:
        for _ in range(n):
            bank.step_eprop(zp, zv, act)
    finally:
        bank.dt = dt0

def train_sequential(env, *, B=20, tau_leak=10.0, D=2.0, episodes=1500,
                     max_steps=None, dt=5e-3, step_dur=0.4, eta=0.2, V=1.5,
                     in_rate=200.0, ltd=LTD_BIAS, tau_m=TAU_M, v_th=V_TH,
                     sigma0=0.15, sigma1=None, abstract=False, no_trace=False,
                     eprop=False, beta_pol=1.0, beta_leak=1.0, weight_fault=None, seed0=0,
                     return_weights=False):
    """Train the sequential policy on ``env`` for ``B`` parallel seeds.

    The network is a state x action grid of device synapses ``w[state, action]``.
    Each EPISODE runs a trajectory: at every visited state the active state-line
    drives the ``A`` action neurons (Kirchhoff bitline sum), the spiking winner is the
    action, and the signed pre--post coincidence on the chosen (state, action) synapse
    drives the eligibility gate. The gate integrates CONTINUOUSLY across the whole
    trajectory (it is not reset between steps), so eligibility from early steps decays
    while later steps are taken. A single goal reward, delivered ``D`` seconds after
    the goal-entry action, then gates the surviving trace into every synapse:
    ``dw = eta (R - b) e``. Returns rewards ``(B, episodes)``.

    ``abstract`` swaps the device gate for the exponential R-STDP kernel;
    ``no_trace`` zeroes eligibility (necessity control). ``weight_fault`` applies the
    device-fault prior at READ time (once per episode, since weights change only per
    episode): the array reads a faulted conductance while learning still targets the
    clean weight -- a fixed physical realisation, not a learnable parameter, exactly as
    in :func:`mrl_trace.deep.train_deep`. The maze weights are NON-NEGATIVE
    (``[0, W_MAX]``) single conductances, so the signed-pair SiO_x stack cannot be used
    verbatim; pass a fault built for the non-negative weight space (see
    ``device_faults.maze_fault_stack``). ``step_dur`` is the wall-clock time the agent
    dwells in each state (so a trajectory of ``L`` steps spans
    ``L*step_dur`` seconds, against which ``tau_leak`` must be long enough).
    """
    rng = np.random.default_rng(seed0)
    S, A = env.n_states, env.n_actions
    if max_steps is None:
        max_steps = 3 * env.L + 2
    if eprop:
        bank = EpropTrace(B, S, A, tau_leak=tau_leak, dt=dt, tau_m=tau_m, v_th=v_th)
    elif abstract:
        bank = AbstractTrace(B, S, A, tau_leak=tau_leak, dt=dt)
    else:
        bank = GateBankBatched(B, S, A, tau_leak=tau_leak, dt=dt, V=V, beta_leak=beta_leak)
    w = np.full((B, S, A), W_INIT)
    baseline = np.full(B, 1.0 / A)
    bidx = np.arange(B)
    steps_per_state = max(1, int(round(step_dur / dt)))
    reward_lag = int(round(D / dt))
    rewards = np.zeros((B, episodes))

    for ep in range(episodes):
        sigma = sigma0 if sigma1 is None else sigma0 + (sigma1 - sigma0) * ep / episodes
        # Read-time device faults applied ONCE per episode: the array reads the faulted
        # conductance (wr) while learning updates the clean weight (w). Computing this per
        # episode (not per dt-tick) is what keeps the run tractable.
        wr = weight_fault(w) if weight_fault is not None else w
        bank.reset()
        if isinstance(env, TMaze):
            pos = env.start(B, rng)
        else:
            pos = env.start(B)
        v = np.zeros((B, A))
        done = np.zeros(B, dtype=bool)
        got_reward = np.zeros(B, dtype=bool)
        # accumulated eligibility snapshot, captured at each seed's reward instant
        e_rew = np.zeros((B, S, A))
        reward_due = np.full(B, -1)         # step index at which reward gates in
        chosen_final = np.zeros(B, dtype=int)   # decision-point action per seed (e-prop)
        pol_final = np.full((B, A), 1.0 / A)     # decision-point policy pi (e-prop)
        # --- run the trajectory ---
        for st in range(max_steps):
            state = env.encode(pos) if isinstance(env, TMaze) else pos
            chosen = np.zeros(B, dtype=int)
            # dwell in this state for steps_per_state dt-ticks: integrate spikes + trace
            spk = np.zeros((B, A))
            for _ in range(steps_per_state):
                pre = np.zeros((B, S))
                active = ~done
                pre[bidx[active], state[active]] = (
                    rng.random(active.sum()) < in_rate * dt).astype(float)
                charge = np.einsum('bsa,bs->ba', wr, pre)
                v, sp, v_pre = lif_step_batched(v, charge, dt, rng, tau_m=tau_m,
                                                v_th=v_th, noise=sigma, return_pre=True)
                spk += sp * active[:, None]
                if eprop:
                    # e-prop computes eligibility from network dynamics: the pre-reset
                    # membrane (pseudo-derivative) and the presynaptic trace. No signed
                    # coincidence -- credit/blame comes from (R - b) at reward.
                    bank.step_eprop(pre, v_pre, active)
                else:
                    # signed leak-dominant coincidence on the CURRENT state's lines
                    # (device gate and R-STDP share this drive)
                    drive = np.zeros((B, S, A))
                    drv = pre[bidx, state][:, None] * np.where(sp, 1.0, -ltd)
                    drive[bidx, state, :] = drv * active[:, None]
                    if no_trace:
                        drive[:] = 0.0
                    bank.step(drive)
            if eprop:
                # Reward-based e-prop is a policy-gradient method: the action is SAMPLED
                # from a softmax policy over the action-neuron spike counts (this sampling
                # is the exploration), and the policy pi is stored so the per-action
                # learning signal (1[a=chosen] - pi_a) can be formed at reward.
                logits = beta_pol * (spk - spk.mean(1, keepdims=True))
                pol_p = np.exp(logits - logits.max(1, keepdims=True))
                pol_p /= pol_p.sum(1, keepdims=True)
                cdf = np.cumsum(pol_p, axis=1)
                u = rng.random((B, 1))
                chosen = (u > cdf).sum(1)
                chosen = np.clip(chosen, 0, A - 1)
            else:
                # action = spiking winner over the dwell (ties -> random)
                tie = spk.max(1) == spk.min(1)
                chosen = np.argmax(spk, 1)
                chosen[tie] = rng.integers(A, size=int(tie.sum()))
            chosen[done] = 0
            nxt, reached, ep_done = env.step(pos, chosen)
            pos = np.where(done, pos, nxt)
            newly = (~done) & ep_done
            # record the decision-point action and policy (the step that ended the
            # episode) for the e-prop per-action learning signal (1[a=chosen] - pi_a)
            chosen_final = np.where(newly, chosen, chosen_final)
            if eprop:
                pol_final = np.where(newly[:, None], pol_p, pol_final)
            got_reward |= (newly & reached)
            # schedule reward gating reward_lag ticks after goal entry (approx: capture
            # the trace reward_lag dt-ticks later by letting the gate relax that long)
            reward_due = np.where(newly & (reward_due < 0), st, reward_due)
            done |= ep_done
            if done.all():
                break
        # Let the trace relax for the action->reward delay, then snapshot eligibility.
        # During relaxation there is NO drive and no spikes, so the dynamics are smooth;
        # integrate with a coarser effective step (stride x dt) to keep long delays
        # tractable. Equivalent to bank.step(0) repeated reward_lag times, to first
        # order, but ~stride x cheaper -- the relaxation is numerically smooth so this
        # introduces negligible error at stride=10 (verified against the tick loop).
        zero = np.zeros((B, S, A))
        stride = 10 if reward_lag > 200 else 1
        n_relax = reward_lag // stride
        rem = reward_lag - n_relax * stride
        if eprop:
            # e-prop's slow output filter keeps decaying with no input; relax it the
            # same way (coarse-stepped), reading eligibility from vn[...,-1].
            _relax_eprop(bank, B, S, A, n_relax, stride)
            for _ in range(rem):
                bank.step_eprop(np.zeros((B, S)), np.zeros((B, A)), np.ones(B))
            e_rew = bank.vn[..., -1]
        elif no_trace:
            e_rew = np.zeros((B, S, A))
        elif abstract:
            for _ in range(reward_lag):
                bank.step(zero)
            e_rew = bank.e
        else:
            _relax(bank, zero, n_relax, stride, abstract)
            for _ in range(rem):
                bank.step(zero)
            e_rew = bank.vn[..., -1] / bank.Vnmax
        R = got_reward.astype(float)
        adv = eta * (R - baseline)
        if eprop:
            # Reward-based e-prop = policy gradient with an e-prop eligibility. The
            # per-action learning signal is the REINFORCE log-policy gradient
            # L_j = (R - b) * (1[j = chosen] - pi_j), which supplies the arm selectivity
            # that e-prop's unsigned eligibility lacks: the chosen action is pushed up by
            # (1 - pi) and the others down by (-pi), scaled by the advantage. Delta w_ij
            # = eta * (R-b) * (1[j=chosen]-pi_j) * e_ij. eta absorbs the small eligibility
            # magnitude (e-prop eligibility ~1e-4 here; see eta scaling).
            onehot = np.zeros((B, A))
            onehot[bidx, chosen_final] = 1.0
            Lsig = (R - baseline)[:, None] * (onehot - pol_final)     # (B, A)
            w = np.clip(w + eta * Lsig[:, None, :] * e_rew, 0.0, W_MAX)
        else:
            w = np.clip(w + adv[:, None, None] * e_rew, 0.0, W_MAX)
        baseline += 0.02 * (R - baseline)
        rewards[:, ep] = R
    if return_weights:
        return rewards, w          # w: (B, S, A) final learned weights
    return rewards


# ----------------------------------------------------------------------------
# Reductions (reuse the bandit's conventions)
# ----------------------------------------------------------------------------
def reward_rate(rewards, window=100):
    rewards = np.asarray(rewards)
    if rewards.shape[-1] < window:
        return rewards.mean(axis=-1)
    return rewards[..., -window:].mean(axis=-1)


def trials_to_criterion(rew_1d, crit, window=100):
    rew_1d = np.asarray(rew_1d)
    if len(rew_1d) < window:
        return None
    csum = np.cumsum(rew_1d)
    rr = (csum[window:] - csum[:-window]) / window
    if not np.any(rr >= crit):
        return None
    return int(np.argmax(rr >= crit) + window)


def policy_correct(env, w):
    """Fraction of states whose argmax-weight action equals the correct action.

    A policy-quality probe independent of stochastic spiking: for each state, the
    action with the largest learned weight should be the progress/arm-correct action.
    ``w`` is a single seed's ``(S, A)`` weight grid.
    """
    S = env.n_states
    correct = 0
    total = 0
    for s in range(S):
        # correct action per state, derived directly from the state index (no episode
        # state needed): stem states want FORWARD (action 0); a junction-context state
        # s = L + arm wants that arm.
        if isinstance(env, TMaze):
            ca = 0 if s < env.L else (s - env.L)
        else:
            ca = 0                          # LinearTrack: always forward
        if w[s].max() == w[s].min():
            continue                       # untrained / tied row -> skip
        if int(np.argmax(w[s])) == int(ca):
            correct += 1
        total += 1
    return correct / total if total else 0.0


# =============================================================================
# Experiment cores (the sequential-MDP credit-assignment studies that compose
# ``TMaze`` + ``train_sequential``)
#
# Each ``run_*`` returns the result grid as a plain dict -- no file I/O, no
# plotting, no stdout.  Notebooks call these at a small (quick) seed/episode
# count and render the figures inline; ``main()`` (below) calls them at the
# published scale, parallelising the coarse axis over a process pool, and writes
# each grid under ``data/results/`` via ``paths.save_result`` for the notebooks
# to replay.  The original driver scripts each computed one grid:
#   * exp6_sequential.npy            -- sequential T-maze + R-STDP/e-prop/no-trace
#                                       benchmark and the retention x delay grid;
#   * exp8_dmax_law.npy              -- densely-sampled D_max(tau_leak) scaling law;
#   * exp15_long_horizon.npy         -- longer-horizon (trajectory length L) credit
#                                       feasibility vs the established traces;
#   * exp15_long_horizon_faults.npy  -- the same under the SiO_x device-fault prior.
# =============================================================================


# -----------------------------------------------------------------------------
# Experiment 6 -- sequential distal-reward task + R-STDP / e-prop benchmark.
#
# Extends the one-step contextual bandit to a multi-step MDP (T-maze), where reward
# is delivered only at the goal after a trajectory and credit must propagate back
# across intermediate steps -- the regime an eligibility trace exists for. Four
# conditions share ONE harness for a fair comparison:
#   device   : the SiOx trap-cascade eligibility gate (this work);
#   rstdp    : an abstract exponential eligibility trace at matched timescale
#              (the hand-set kernel of Florian 2007 / Izhikevich 2007 / Legenstein 2008);
#   eprop    : network-computed eligibility (reward-modulated policy gradient, Bellec 2020);
#   no_trace : eligibility zeroed (device-necessity control).
# Operating point V=1.5 (rise time tau_r=1.9 s) so the eligibility snapshot reflects
# RETENTION (tau_leak), not the sigmoidal rise.  Pre-registered criteria C1-C4 / kill K1
# are evaluated in the summary.
# -----------------------------------------------------------------------------
EPROP_ETA = 2000.0        # e-prop eligibility magnitude ~1e-4 << device/R-STDP; eta scaled up


def running_rate(r1d, window=50):
    """Running ``window``-episode reward rate of a single mean-over-seeds curve."""
    c = np.cumsum(np.insert(r1d, 0, 0.0))
    return (c[window:] - c[:-window]) / window


def run_sequential(*, seeds=20, episodes=800, V=1.5, D0=4.0, TAU0=10.0,
                   taus=(2.0, 5.0, 20.0), delays=(2, 4, 8, 16, 32, 64), seed0=0):
    """Experiment 6: sequential T-maze learning curves + retention x delay grid.

    Runs the four conditions (device / rstdp / eprop / no_trace) as learning curves
    at the reference ``(D0, TAU0)`` operating point, a device policy-correctness probe
    from the ACTUAL learned weights, and a ``taus x delays`` device reward-rate grid
    from which ``D_max(tau_leak)`` (the prediction-arm datapoints) is read off.

    Every cell is an independent, seed-deterministic ``train_sequential`` call
    (``B=seeds``, ``seed0=0``), so the result is identical whether the cells are run
    serially (here) or dispatched across a process pool (``main``).  Returns the result
    dict (per-condition per-seed finals + running curves, policy correctness, the
    retention x delay grid with bootstrap CIs, ``D_max`` per tau, and the pre-registered
    criteria); no file I/O.
    """
    from .stats import bootstrap_ci
    taus = list(taus); delays = list(delays)
    env = TMaze(L=3, A_goal=2)
    A = env.A_goal
    chance = 1.0 / A
    crit = 0.5 * (1 + chance)                 # manuscript convention -> 0.75 for A=2
    conds = {"device": dict(), "rstdp": dict(abstract=True),
             "eprop": dict(eprop=True, eta=EPROP_ETA),
             "no_trace": dict(no_trace=True)}

    # ---------- (1) learning curves ----------
    curves, finals, finals_ci = {}, {}, {}
    for name, kw in conds.items():
        r = train_sequential(env, B=seeds, tau_leak=TAU0, D=D0, episodes=episodes,
                             V=V, seed0=seed0, **kw)
        finals[name] = reward_rate(r, window=100)
        curves[name] = running_rate(r.mean(0), window=50)
        finals_ci[name] = bootstrap_ci(finals[name])

    # ---------- (2) policy correctness from ACTUAL learned weights ----------
    _, w_dev = train_sequential(env, B=seeds, tau_leak=TAU0, D=D0, episodes=episodes,
                                V=V, seed0=seed0, return_weights=True)
    pol = np.array([policy_correct(env, w_dev[b]) for b in range(seeds)])

    # ---------- (3) retention x delay grid (device) ----------
    grid = np.zeros((len(taus), len(delays)))
    grid_ci = np.zeros((len(taus), len(delays), 2))
    for i, tau in enumerate(taus):
        for j, D in enumerate(delays):
            f = reward_rate(train_sequential(env, B=seeds, tau_leak=tau, D=D,
                                             episodes=episodes, V=V, seed0=seed0), 100)
            grid[i, j] = f.mean()
            grid_ci[i, j] = bootstrap_ci(f)

    # D_max(tau) = largest delay reaching criterion (the prediction-arm datapoints)
    dmax = {}
    for i, tau in enumerate(taus):
        ok = [D for j, D in enumerate(delays) if grid[i, j] >= crit]
        dmax[tau] = max(ok) if ok else 0

    # ---------- pre-registered criteria ----------
    c1 = bool(finals["no_trace"].mean() <= chance + 0.10)
    c2 = bool(finals["device"].mean() >= crit)
    c3 = bool(pol.mean() >= crit)
    c4 = bool((dmax[20.0] >= dmax[5.0] >= dmax[2.0]) and (dmax[20.0] > dmax[2.0])) \
        if set((2.0, 5.0, 20.0)).issubset(dmax) else None
    dev_beats_nt = bool(any(grid[i, j] > finals["no_trace"].mean() + 0.1
                            for i in range(len(taus)) for j in range(len(delays))))
    criteria = {"C1_notrace_at_chance": c1, "C2_device_criterion": c2,
                "C3_device_policy": c3, "C4_dmax_grows_with_tau": c4,
                "K1_device_exceeds_notrace": dev_beats_nt}

    return {
        "finals": {k: v for k, v in finals.items()},
        "finals_ci": finals_ci,
        "curves": curves, "policy": pol,
        "policy_ci": bootstrap_ci(pol),
        "taus": taus, "delays": delays, "grid": grid, "grid_ci": grid_ci,
        "dmax": dmax, "chance": chance, "crit": crit,
        "D0": D0, "TAU0": TAU0, "V": V, "seeds": seeds, "episodes": episodes,
        "criteria": criteria,
    }


# -----------------------------------------------------------------------------
# Experiment 8 -- densely-sampled retention->delay scaling, D_max(tau_leak).
#
# Sharpens the D_max ~= k*tau relation from the 3-point version (exp6) to ~13 retention
# values on a fine delay grid, so the FUNCTIONAL FORM can be read off rather than assumed:
# is it linear in the trace-dominated regime, and where does it saturate (the inverted-U
# the conditioning literature predicts)? D_max is interpolated as the largest delay whose
# reward rate is >= criterion (threshold crossing), not just the coarsest passing grid
# point. Device condition only (the law is a device-physics claim).
#
# HONEST SCOPE: this improves the resolution of the FORM. It does NOT remove the
# in-simulation near-tautology (tau_leak sets the trace decay, which by construction
# bounds the bridgeable delay). The non-tautological tests are the hardware tau_leak sweep
# and the biological arm (future work). This is the "proposed scaling relation, sampled"
# -- not a claim of a tested cross-domain law.
# -----------------------------------------------------------------------------
# Retention values spanning the measured band (~1-40 s). The trace-dominated regime
# (tau<=8 s) is densified with interpolating points so the linear fit rests on 9 points
# rather than 5; the longer tau (>=12 s) are kept for context but remain right-censored at
# the delay-grid ceiling.
DMAX_TAUS = [1.0, 1.5, 2.0, 2.5, 3.0, 4.0, 5.0, 6.5, 8.0, 12.0, 20.0, 30.0, 40.0]
# Delay grid extended (adding 420, 560, 700) so the tau=30 and 40 s points can resolve
# their D_max crossings; with k~14 the tau=40 point is expected near 560 s.
DMAX_DELAYS = [1, 2, 3, 5, 8, 12, 20, 32, 48, 72, 100, 140, 180, 240, 300, 420, 560, 700]
DMAX_PLATEAU_TAU = 20.0


def interp_dmax(delays, rates, crit):
    """Largest delay at which the (monotone-from-peak) reward rate is >= crit, linearly
    interpolated in log-delay between the last passing and first failing grid points.
    Returns 0 if never reached, or the max delay if always above."""
    d = np.asarray(delays, float); r = np.asarray(rates, float)
    above = r >= crit
    if not above.any():
        return 0.0
    # take the falling edge AFTER the peak (D_max is where it drops below crit going up
    # in delay), so scan from the peak rightwards for the last contiguous >=crit delay
    pk = int(np.argmax(r))
    last = pk
    for j in range(pk, len(d)):
        if r[j] >= crit:
            last = j
        else:
            break
    if last == len(d) - 1:
        return float(d[last])              # still above at the largest delay tested
    # interpolate in log-delay between d[last] (>=crit) and d[last+1] (<crit)
    x0, x1 = np.log(d[last]), np.log(d[last + 1])
    y0, y1 = r[last], r[last + 1]
    frac = (y0 - crit) / (y0 - y1) if y0 != y1 else 0.0
    return float(np.exp(x0 + frac * (x1 - x0)))


def _dmax_origin_fit(taus, dvals, mask):
    """Origin-forced linear fit ``D_max = k*tau`` over the masked points -> ``(k, R^2)``."""
    if mask.sum() < 2:
        return float("nan"), float("nan")
    kk = float(np.sum(taus[mask] * dvals[mask]) / np.sum(taus[mask] ** 2))
    res = dvals[mask] - kk * taus[mask]
    sst = float(np.sum((dvals[mask] - dvals[mask].mean()) ** 2))
    return kk, (1 - float(np.sum(res ** 2)) / sst if sst > 0 else float("nan"))


def _dmax_cell(spec, episodes, V, seeds):
    """One (tau, D) device cell -> (tau, D, mean reward rate, lo, hi). Top-level for
    multiprocessing; deterministic by seed0=0."""
    from .stats import bootstrap_ci
    tau, D = spec
    env = TMaze(L=3, A_goal=2)
    r = train_sequential(env, B=seeds, tau_leak=tau, D=D, episodes=episodes, V=V, seed0=0)
    f = reward_rate(r, 100)
    lo, hi = bootstrap_ci(f)
    return (tau, D, float(f.mean()), lo, hi)


def _dmax_assemble(results, taus, delays, crit, seeds, episodes, V):
    """Package raw ``(tau, D, mean, lo, hi)`` cells into the exp8 grid dict (per-tau
    reward-rate rows, interpolated ``D_max``, the origin-forced headline fit + the
    asymptotic plateau slope, and the right-censoring flags)."""
    grid = {tau: {} for tau in taus}
    for tau, D, m, lo, hi in results:
        grid[tau][D] = (m, lo, hi)
    dmax = {tau: interp_dmax(delays, [grid[tau][D][0] for D in delays], crit)
            for tau in taus}
    taus_a = np.array(taus); dvals = np.array([dmax[t] for t in taus])
    censored = dvals >= max(delays) - 1e-6
    fitmask = (~censored) & (dvals > 0)
    k, r2 = _dmax_origin_fit(taus_a, dvals, fitmask)                     # all resolved points
    platmask = fitmask & (taus_a >= DMAX_PLATEAU_TAU)
    k_plateau, r2_plateau = _dmax_origin_fit(taus_a, dvals, platmask)    # asymptotic slope
    return {"taus": list(taus), "delays": list(delays), "grid": grid, "dmax": dmax,
            "k": k, "r2": r2, "k_plateau": k_plateau, "r2_plateau": r2_plateau,
            "plateau_tau": DMAX_PLATEAU_TAU, "crit": crit,
            "censored": censored.tolist(), "seeds": seeds, "episodes": episodes,
            "V_op": V}


def run_dmax_law(*, seeds=20, episodes=800, V=1.5, taus=None, delays=None, crit=0.75):
    """Experiment 8: dense D_max(tau_leak) scaling law (device condition only).

    Runs every ``(tau, D)`` device cell serially and reads off the interpolated
    ``D_max`` per retention (threshold crossing of reward rate >= ``crit``), then
    fits ``D_max = k*tau`` through the origin over the resolved (non-censored) points
    and separately over the trace-dominated plateau (``tau >= DMAX_PLATEAU_TAU``).
    Each cell is deterministic (``seed0=0``); serial here, pooled in ``main``. Returns
    the result dict; no file I/O.
    """
    taus = DMAX_TAUS if taus is None else list(taus)
    delays = DMAX_DELAYS if delays is None else list(delays)
    specs = [(tau, D) for tau in taus for D in delays]
    results = [_dmax_cell(s, episodes, V, seeds) for s in specs]
    return _dmax_assemble(results, taus, delays, crit, seeds, episodes, V)


# -----------------------------------------------------------------------------
# Experiment 15 -- longer-horizon temporal-credit feasibility (Tier 2).
#
# Answers the *task-depth* objection ("your task is too easy") on the axis the device
# primitive is actually for: temporal credit assignment across a multi-step trajectory
# with a delayed goal reward. NOT a perception/benchmark study -- the input stays
# low-dimensional; what grows is the temporal DEPTH (trajectory length L) of the credit
# problem. Shallow state x action policy (no homeostasis). Per cell the device trace is
# compared against the established baselines on the SAME task/network/seeds/budget:
#   device / eprop / abstract / no_trace.
# Scaling axes: trajectory length L (T-maze stem) and arm count A_goal (chance = 1/A).
# tau_leak is set generously so retention can span the longest trajectory tested.
# -----------------------------------------------------------------------------
LONG_HORIZON_CONDS = {
    "device":   dict(),
    "eprop":    dict(eprop=True),
    "abstract": dict(abstract=True),
    "no_trace": dict(no_trace=True),
}


def _lh_final_rate(r):
    """Final-window (100-episode) reward rate averaged over the seed batch."""
    return float(np.mean(reward_rate(r, window=100)))


def _lh_run(job):
    """One (L, A, cond, seed) long-horizon cell -> (L, A, cond, seed, final rate).
    Top-level for multiprocessing; ``B=1`` and ``seed0=seed`` so seeds are independent."""
    L, A, cond, seed, episodes, D, tau_leak = job
    env = TMaze(L=L, A_goal=A)
    r = train_sequential(env, B=1, tau_leak=tau_leak, D=D, episodes=episodes,
                         seed0=seed, **LONG_HORIZON_CONDS[cond])
    return (L, A, cond, seed, _lh_final_rate(r))


def _lh_boot_ci(vals, seed=0, n_boot=10000):
    """Percentile bootstrap CI over the per-seed final rates (matches the driver's own
    resampler; ``bootstrap_ci`` in ``stats`` uses the same convention)."""
    rng = np.random.default_rng(seed)
    b = [vals[rng.integers(0, len(vals), len(vals))].mean() for _ in range(n_boot)]
    return float(np.percentile(b, 2.5)), float(np.percentile(b, 97.5))


def _lh_summarize(results, L_grid, A_grid, episodes, D, tau_leak, seeds):
    """Package raw long-horizon cells into the exp15 grid dict (per-(L, A, cond) mean +
    bootstrap CI)."""
    acc = {}
    for L, A, cond, s, val in results:
        acc.setdefault((L, A, cond), []).append(val)
    summary = {}
    for L in L_grid:
        for A in A_grid:
            for cond in LONG_HORIZON_CONDS:
                v = np.array(acc[(L, A, cond)])
                lo, hi = _lh_boot_ci(v, seed=L * 100 + A * 10)
                summary[(L, A, cond)] = (float(v.mean()), lo, hi)
    return {"L": L_grid, "A": A_grid, "conds": list(LONG_HORIZON_CONDS),
            "summary": summary, "episodes": episodes, "D": D,
            "tau_leak": tau_leak, "seeds": seeds}


def run_long_horizon(*, L_grid=(3, 5, 8, 12), A_grid=(2, 3), seeds=12,
                     episodes=2500, D=2.0, tau_leak=10.0):
    """Experiment 15: longer-horizon temporal-credit feasibility, serial.

    Sweeps trajectory length ``L`` x arm count ``A_goal`` x condition x seed and reports
    the final reward rate (mean + bootstrap CI over seeds) for the device trace against
    the eprop / abstract / no_trace baselines. Each seed is an independent ``B=1``,
    ``seed0=seed`` run, so the serial sweep here matches the pooled sweep in ``main``.
    Returns the result dict; no file I/O.
    """
    L_grid = list(L_grid); A_grid = list(A_grid)
    jobs = [(L, A, cond, s, episodes, D, tau_leak)
            for L in L_grid for A in A_grid
            for cond in LONG_HORIZON_CONDS for s in range(seeds)]
    results = [_lh_run(j) for j in jobs]
    return _lh_summarize(results, L_grid, A_grid, episodes, D, tau_leak, seeds)


# -----------------------------------------------------------------------------
# Experiment 15-faults -- longer-horizon temporal credit UNDER the SiO_x device-fault
# prior. The fault-robustness companion to ``run_long_horizon``.
#
# The maze weights are a SINGLE NON-NEGATIVE conductance per synapse in ``[0, W_MAX]``
# (not the signed differential pair of the deep net), so ``device_faults.maze_fault_stack``
# applies the SAME fault classes (stuck-off, lognormal D2D; Poole-Frenkel I-V excluded as a
# read-nonlinearity out of scope for a programmed-weight learning claim) through the correct
# ``[0, W_MAX] <-> [G_off, G_on]`` mapping. Per (L, p) cell we report the device trace's
# final reward rate (mean + bootstrap CI) under faults beside the clean no_trace control;
# the headline question is graceful degradation as the stuck fraction p rises.
# -----------------------------------------------------------------------------
LONG_HORIZON_SIGMA_G = 0.5                    # lognormal D2D width (matches exp14)


def _lhf_run(job):
    """One (L, p, seed, trace) fault cell -> (L, p, seed, trace, final rate). Top-level
    for multiprocessing; ``B=1`` and ``seed0=seed`` so seeds are independent. ``trace``
    True = device trace under the SiO_x programmed-conductance fault prior (stuck-off +
    D2D, read-time; PF off); False = no-trace control (faults moot with no trace)."""
    from .device_faults import maze_fault_stack
    L, p, seed, trace, A, episodes, D, tau_leak = job
    env = TMaze(L=L, A_goal=A)
    if trace:
        fault = maze_fault_stack(p_stuck=p, sigma_g=LONG_HORIZON_SIGMA_G, pf_on=False,
                                 stuck_kind="off", w_max=W_MAX,
                                 seed=1000 * seed + int(100 * p))
        r = train_sequential(env, B=1, tau_leak=tau_leak, D=D, episodes=episodes,
                             seed0=seed, weight_fault=fault)
    else:
        r = train_sequential(env, B=1, tau_leak=tau_leak, D=D, episodes=episodes,
                             seed0=seed, no_trace=True)
    return (L, p, seed, trace, _lh_final_rate(r))


def _lhf_summarize(results, L_grid, p_grid, A, episodes, D, tau_leak, seeds):
    """Package raw fault cells into the exp15-faults grid dict (device mean + bootstrap
    CI grid, clean no-trace control grid)."""
    dev = {(L, p): [] for L in L_grid for p in p_grid}
    ctl = {(L, p): [] for L in L_grid for p in p_grid}
    for L, p, s, trace, val in results:
        (dev if trace else ctl)[(L, p)].append(val)
    chance = 1.0 / A
    grid = np.zeros((len(L_grid), len(p_grid), 3))
    cgrid = np.zeros((len(L_grid), len(p_grid)))
    for i, L in enumerate(L_grid):
        for j, p in enumerate(p_grid):
            d = np.array(dev[(L, p)]); c = np.array(ctl[(L, p)])
            lo, hi = _lh_boot_ci(d, seed=i * 10 + j)
            grid[i, j] = (d.mean(), lo, hi); cgrid[i, j] = c.mean()
    return {"L": L_grid, "p": p_grid, "A": A, "grid": grid, "ctrl": cgrid,
            "chance": chance,
            "faults": "stuck-off + D2D(0.5) [PF excluded] (maze mapping)",
            "episodes": episodes, "D": D, "tau_leak": tau_leak, "seeds": seeds}


def run_long_horizon_faults(*, L_grid=(3, 5, 8, 12), p_grid=(0, 0.1, 0.2, 0.5),
                            A=2, seeds=12, episodes=2500, D=2.0, tau_leak=12.0):
    """Experiment 15-faults: longer-horizon temporal credit under the SiO_x fault prior.

    Sweeps trajectory length ``L`` x stuck-device fraction ``p``, comparing the device
    eligibility trace under the SiO_x programmed-conductance fault prior (stuck-off +
    lognormal D2D at read time; Poole-Frenkel excluded) against the clean no-trace
    control. Each seed is an independent ``B=1``, ``seed0=seed`` run; serial here, pooled
    in ``main``. Returns the result dict; no file I/O.
    """
    L_grid = list(L_grid); p_grid = list(p_grid)
    jobs = [(L, p, s, trace, A, episodes, D, tau_leak)
            for L in L_grid for p in p_grid
            for s in range(seeds) for trace in (True, False)]
    results = [_lhf_run(j) for j in jobs]
    return _lhf_summarize(results, L_grid, p_grid, A, episodes, D, tau_leak, seeds)


def main(argv=None):
    """Full-scale reproduction CLI for the sequential-maze grids (writes ``data/results``).

    ``python -m mrl_trace.maze [--exp6] [--exp8] [--exp15] [--exp15-faults] [--full|--quick]``

    With no experiment flag, runs all four. ``--full`` = the published scale (exp6/exp8:
    20 seeds; exp15/exp15-faults: 12 seeds); ``--quick`` = a fast few-seed smoke run.
    ``main`` parallelises the coarse axis over a process pool (it runs as ``python -m``,
    a real ``__main__``); the ``run_*`` cores themselves stay serial for in-notebook use.
    """
    import argparse
    import os
    from functools import partial
    from multiprocessing import Pool
    from . import paths

    ap = argparse.ArgumentParser(description="Sequential-maze credit-assignment reproductions")
    ap.add_argument("--exp6", action="store_true",
                    help="sequential T-maze + benchmark grid -> exp6_sequential.npy")
    ap.add_argument("--exp8", action="store_true",
                    help="dense D_max(tau_leak) law -> exp8_dmax_law.npy")
    ap.add_argument("--exp15", action="store_true",
                    help="longer-horizon feasibility -> exp15_long_horizon.npy")
    ap.add_argument("--exp15-faults", dest="exp15_faults", action="store_true",
                    help="longer-horizon under the SiOx fault prior -> exp15_long_horizon_faults.npy")
    ap.add_argument("--quick", action="store_true", help="fast few-seed smoke run")
    ap.add_argument("--full", action="store_true", help="published-scale run (default)")
    a = ap.parse_args(argv)
    run_all = not (a.exp6 or a.exp8 or a.exp15 or a.exp15_faults)

    n_proc = max(1, (os.cpu_count() or 4) - 2)

    # ---- Experiment 6: sequential T-maze + benchmark grid ----
    if a.exp6 or run_all:
        seeds = 6 if a.quick else 20
        episodes = 400 if a.quick else 800
        V, D0, TAU0 = 1.5, 4.0, 10.0
        taus = [2.0, 5.0, 20.0]
        delays = [2, 4, 8, 16, 32, 64]
        print(f"=== exp6 sequential T-maze (N={seeds}, {episodes} ep, V={V}) ===", flush=True)
        conds = {"device": dict(), "rstdp": dict(abstract=True),
                 "eprop": dict(eprop=True, eta=EPROP_ETA), "no_trace": dict(no_trace=True)}
        specs = []
        for name, kw in conds.items():                       # learning-curve runs
            specs.append((("curve", name), TAU0, D0, kw, False))
        specs.append((("weights", "device"), TAU0, D0, dict(), True))   # policy run
        for tau in taus:                                     # grid cells
            for D in delays:
                specs.append((("grid", tau, D), tau, D, dict(), False))
        with Pool(min(len(specs), n_proc)) as pool:
            res = dict(pool.map(partial(_exp6_worker, seeds=seeds, episodes=episodes, V=V),
                                specs))
        grid = _exp6_assemble(res, conds, taus, delays, seeds, episodes, V, D0, TAU0)
        paths.save_result("exp6_sequential.npy", grid)
        print(f"  wrote exp6_sequential.npy  criteria={grid['criteria']}", flush=True)

    # ---- Experiment 8: dense D_max(tau_leak) law ----
    if a.exp8 or run_all:
        seeds = 6 if a.quick else 20
        episodes = 400 if a.quick else 800
        V, crit = 1.5, 0.75
        taus, delays = DMAX_TAUS, DMAX_DELAYS
        print(f"=== exp8 dense D_max(tau_leak): {len(taus)} tau x {len(delays)} delays, "
              f"{seeds} seeds, {episodes} ep, V={V}, crit={crit} ===", flush=True)
        specs = [(tau, D) for tau in taus for D in delays]
        with Pool(min(len(specs), n_proc)) as pool:
            results = pool.map(partial(_dmax_cell, episodes=episodes, V=V, seeds=seeds),
                               specs)
        grid = _dmax_assemble(results, taus, delays, crit, seeds, episodes, V)
        paths.save_result("exp8_dmax_law.npy", grid)
        print(f"  wrote exp8_dmax_law.npy  D_max = {grid['k']:.2f} tau (R^2={grid['r2']:.3f}), "
              f"plateau slope ~{grid['k_plateau']:.2f}", flush=True)

    # ---- Experiment 15: longer-horizon feasibility ----
    if a.exp15 or run_all:
        seeds = 4 if a.quick else 12
        episodes = 400 if a.quick else 2500
        L_grid = [3, 5] if a.quick else [3, 5, 8, 12]
        A_grid = [2] if a.quick else [2, 3]
        D, tau_leak = 2.0, 10.0
        jobs = [(L, A, cond, s, episodes, D, tau_leak)
                for L in L_grid for A in A_grid
                for cond in LONG_HORIZON_CONDS for s in range(seeds)]
        print(f"=== exp15 long-horizon: {len(jobs)} runs; L={L_grid} A={A_grid} "
              f"conds={list(LONG_HORIZON_CONDS)} seeds={seeds} episodes={episodes} "
              f"D={D}s tau_leak={tau_leak}s ===", flush=True)
        with Pool(min(len(jobs), n_proc)) as pool:
            results = pool.map(_lh_run, jobs, chunksize=1)
        grid = _lh_summarize(results, L_grid, A_grid, episodes, D, tau_leak, seeds)
        paths.save_result("exp15_long_horizon.npy", grid)
        print("  wrote exp15_long_horizon.npy", flush=True)

    # ---- Experiment 15-faults: longer-horizon under the SiOx fault prior ----
    if a.exp15_faults or run_all:
        seeds = 4 if a.quick else 12
        episodes = 400 if a.quick else 2500
        L_grid = [3, 5] if a.quick else [3, 5, 8, 12]
        p_grid = [0, 0.2] if a.quick else [0, 0.1, 0.2, 0.5]
        A, D, tau_leak = 2, 2.0, 12.0
        jobs = [(L, p, s, trace, A, episodes, D, tau_leak)
                for L in L_grid for p in p_grid
                for s in range(seeds) for trace in (True, False)]
        print(f"=== exp15-faults long-horizon under SiOx fault prior: {len(jobs)} runs; "
              f"L={L_grid} p={p_grid} A={A} seeds={seeds} episodes={episodes} D={D}s "
              f"tau_leak={tau_leak}s (stuck-off + D2D, non-negative maze mapping) ===",
              flush=True)
        with Pool(min(len(jobs), n_proc)) as pool:
            results = pool.map(_lhf_run, jobs, chunksize=1)
        grid = _lhf_summarize(results, L_grid, p_grid, A, episodes, D, tau_leak, seeds)
        paths.save_result("exp15_long_horizon_faults.npy", grid)
        print("  wrote exp15_long_horizon_faults.npy", flush=True)


# --- exp6 pooled workers (module-level so they pickle for the process pool) ---
def _exp6_worker(spec, seeds, episodes, V):
    """Run one independent exp6 cell for the pool. ``spec`` = (key, tau_leak, D, kwargs,
    want_weights). Returns ``(key, rewards[, weights])``. Deterministic given seed0=0."""
    key, tau_leak, D, kw, want_w = spec
    env = TMaze(L=3, A_goal=2)
    out = train_sequential(env, B=seeds, tau_leak=tau_leak, D=D, episodes=episodes,
                           V=V, seed0=0, return_weights=want_w, **kw)
    return (key, out)


def _exp6_assemble(res, conds, taus, delays, seeds, episodes, V, D0, TAU0):
    """Package pooled exp6 results (keyed dict) into the exp6 grid dict -- identical
    structure to ``run_sequential`` output, built from the parallel cells."""
    from .stats import bootstrap_ci
    A = 2
    chance = 1.0 / A
    crit = 0.5 * (1 + chance)
    curves, finals, finals_ci = {}, {}, {}
    for name in conds:
        r = res[("curve", name)]
        finals[name] = reward_rate(r, window=100)
        curves[name] = running_rate(r.mean(0), window=50)
        finals_ci[name] = bootstrap_ci(finals[name])
    _, w_dev = res[("weights", "device")]
    env = TMaze(L=3, A_goal=2)
    pol = np.array([policy_correct(env, w_dev[b]) for b in range(seeds)])
    grid = np.zeros((len(taus), len(delays)))
    grid_ci = np.zeros((len(taus), len(delays), 2))
    for i, tau in enumerate(taus):
        for j, D in enumerate(delays):
            f = reward_rate(res[("grid", tau, D)], 100)
            grid[i, j] = f.mean()
            grid_ci[i, j] = bootstrap_ci(f)
    dmax = {}
    for i, tau in enumerate(taus):
        ok = [D for j, D in enumerate(delays) if grid[i, j] >= crit]
        dmax[tau] = max(ok) if ok else 0
    c1 = bool(finals["no_trace"].mean() <= chance + 0.10)
    c2 = bool(finals["device"].mean() >= crit)
    c3 = bool(pol.mean() >= crit)
    c4 = bool((dmax[20.0] >= dmax[5.0] >= dmax[2.0]) and (dmax[20.0] > dmax[2.0])) \
        if set((2.0, 5.0, 20.0)).issubset(dmax) else None
    dev_beats_nt = bool(any(grid[i, j] > finals["no_trace"].mean() + 0.1
                            for i in range(len(taus)) for j in range(len(delays))))
    criteria = {"C1_notrace_at_chance": c1, "C2_device_criterion": c2,
                "C3_device_policy": c3, "C4_dmax_grows_with_tau": c4,
                "K1_device_exceeds_notrace": dev_beats_nt}
    return {"finals": {k: v for k, v in finals.items()}, "finals_ci": finals_ci,
            "curves": curves, "policy": pol, "policy_ci": bootstrap_ci(pol),
            "taus": taus, "delays": delays, "grid": grid, "grid_ci": grid_ci,
            "dmax": dmax, "chance": chance, "crit": crit,
            "D0": D0, "TAU0": TAU0, "V": V, "seeds": seeds, "episodes": episodes,
            "criteria": criteria}


if __name__ == "__main__":
    main()
\end{lstlisting}

\begin{lstlisting}[caption={siox\_eligibility/deep.py},label={lst:siox-src-deep}]
"""Deep (two-layer) all-local spiking RL with a physical eligibility trace -- Arm D.

The bandit (:mod:`mrl_trace.bandit`) and the sequential T-maze
(:mod:`mrl_trace.maze`) both train a SINGLE layer of device synapses with a
single global reward scalar ``(R - b)``. That is enough when the policy is linearly
separable in the state lines, but it has two known limits:

1. A single trained layer cannot represent a non-linearly-separable policy (XOR).
2. A single global scalar gives a *hidden* layer no per-neuron credit -- every hidden
   synapse sees the same third factor, so reward-modulated Hebbian learning at depth
   cannot resolve which hidden unit helped (the structural-credit problem; the depth
   analogue of the Chapter 7 scaling wall).

This module asks whether a device-supplied physical eligibility trace can serve as the
TEMPORAL factor of a DEEP spiking policy that is trained FULLY LOCALLY -- no
backpropagation, no weight transport -- on a task a single trained layer cannot solve
(XOR). The temporal factor (the device trace) is held fixed across conditions; only the
SPATIAL-credit pathway varies, so the comparison isolates structural credit:

- ``shallow`` : one trained layer F->A (proves the task needs depth -- should FAIL XOR).
- ``elm``     : deep, hidden layer FIXED RANDOM, only the output trained (the
                random-features shortcut; bounds how much depth-without-learning buys).
- ``global``  : deep, BOTH layers trained by the pure global scalar ``(R - b)``
                (reward-modulated Hebb at the hidden layer; the structural-credit
                failure mode).
- ``dfa``     : deep, BOTH layers trained by DIRECT FEEDBACK ALIGNMENT -- a per-neuron
                learning signal broadcast to the hidden layer through a FIXED RANDOM
                feedback matrix (no W2 transpose, no gradient transport). All-local.
- ``no_trace``: deep DFA with the eligibility zeroed (device-necessity control).

Every condition uses the same :class:`GateBankBatched` device physics for eligibility,
the same signed leak-dominant coincidence drive, and ``dw = eta * L * e`` updates,
where ``L`` is the (global or per-layer) learning signal. See
``experiments/PREREGISTRATION_deep_local.md`` for the pre-registered criteria.
"""
from __future__ import annotations

import numpy as np

from .bandit import GateBankBatched, W_INIT, W_MAX
from .neurons import lif_step_batched, TAU_M, V_TH
from .learning import LTD_BIAS

__all__ = ["xor_inputs", "train_deep", "reward_rate",
           "run_dms", "final_rate", "run_dms_all",
           "run_deep_local", "run_deep_dms", "run_array_scale", "main"]


def _relax_gate(bank, n_relax, stride, rem):
    """Advance an undriven :class:`GateBankBatched` for ``n_relax*stride + rem`` ticks,
    using a coarsened step (``stride*dt``) for the bulk and single ``dt`` steps for the
    remainder. The drive is zero, so the dynamics are smooth and the coarse step is
    accurate (see the inline note at the call site). Restores the gate's ``dt``."""
    B, S, A = bank.B, bank.S, bank.A
    z = np.zeros((B, S, A))
    if n_relax > 0:
        dt0 = bank.dt
        bank.dt = dt0 * stride
        try:
            for _ in range(n_relax):
                bank.step(z)
        finally:
            bank.dt = dt0
    for _ in range(rem):
        bank.step(z)


# ----------------------------------------------------------------------------
# Task: XOR contextual bandit (depth genuinely required)
# ----------------------------------------------------------------------------
def xor_inputs(rng, B):
    """Draw ``B`` XOR trials.

    Two binary features ``b0, b1`` per seed; the correct action is ``XOR(b0, b1)``.
    The features are one-hot-per-feature encoded onto ``F = 4`` input lines
    ``[b0=0, b0=1, b1=0, b1=1]`` so that EXACTLY TWO lines are active every trial --
    a constant total input drive across all four states, which removes any input-rate
    confound between states. XOR is not linearly separable in these 4 lines, so a single
    trained linear layer F->A cannot solve it (this is what forces depth).

    Returns ``(feat_lines (B,4) in {0,1}, correct_action (B,) in {0,1})``.
    """
    b0 = rng.integers(2, size=B)
    b1 = rng.integers(2, size=B)
    lines = np.zeros((B, 4))
    bidx = np.arange(B)
    lines[bidx, 0 + b0] = 1.0          # b0=0 -> line 0, b0=1 -> line 1
    lines[bidx, 2 + b1] = 1.0          # b1=0 -> line 2, b1=1 -> line 3
    correct = (b0 ^ b1).astype(int)
    return lines, correct


# ----------------------------------------------------------------------------
# Training harness
# ----------------------------------------------------------------------------
def train_deep(*, mode="dfa", B=20, H=8, tau_leak=10.0, D=5.0, trials=2000,
               dt=5e-3, cue_dur=1.0, eta=0.2, eta_hidden=None, in_rate=200.0,
               ltd=LTD_BIAS, tau_m=TAU_M, v_th=V_TH, V=1.5, sigma0=0.15, sigma1=0.05,
               fb_scale=1.0, w_scale1=0.6, w_scale2=0.35, w_max=W_MAX,
               bias_o=0.35, homeo=0.0, homeo_target=0.35, homeo_tau=200.0,
               t_distract=None, distract_dur=0.3, weight_fault=None, early_stop=None,
               reward_pools=None, state_sampler=None, n_features=None, n_actions=None,
               seed0=0, return_weights=False, log_align=False):
    """Train the XOR contextual bandit with a two-layer spiking device-synapse network.

    Architecture: ``F=4`` input lines -> ``H`` hidden LIF neurons -> ``A=2`` action LIF
    neurons, with device-synapse weight matrices ``W1 (F,H)`` and ``W2 (H,A)``. Each
    matrix has its own :class:`GateBankBatched` eligibility (same trap-cascade physics,
    retention ``tau_leak``, signed leak-dominant drive). The action is the
    higher-spiking output neuron over the cue epoch (membrane noise = exploration); the
    reward ``R in {0,1}`` is delivered after delay ``D``, contingent on the chosen action
    matching ``XOR``.

    ``mode`` selects the SPATIAL-credit pathway (the eligibility/temporal factor is the
    device trace in all modes); see module docstring and the pre-registration.

    ``state_sampler`` swaps the input STATE without touching the learning rule. By default
    (``None``) the state is the built-in XOR cue (``xor_inputs``, ``F=4``, ``A=2``). Pass a
    callable ``state_sampler(rng, B) -> (lines (B,F) >=0, correct (B,) in [0,A))`` to drive
    the SAME deep all-local stack from an arbitrary state source -- e.g. the low-dimensional
    readout of a frozen perception front-end (the exp16 hybrid). ``lines`` are interpreted as
    per-line input rates exactly as the one-hot cue is, so real-valued (noisy, overlapping)
    readout vectors slot in unchanged; only the policy the network must learn changes, never
    the rule. ``n_features`` and ``n_actions`` are then required (the ``F``/``A`` the sampler
    emits). The eligibility, DFA, homeostasis, fault and reward machinery are all identical,
    so a feasibility conclusion drawn here transfers to the built-in cue and vice versa.

    Returns rewards ``(B, trials)`` (and final weights if ``return_weights``).
    """
    if mode not in {"shallow", "elm", "global", "dfa", "no_trace"}:
        raise ValueError(f"unknown mode {mode!r}")
    rng = np.random.default_rng(seed0)
    if state_sampler is None:
        F, A = 4, 2
        sampler = xor_inputs
    else:
        if n_features is None or n_actions is None:
            raise ValueError("state_sampler requires n_features and n_actions")
        F, A = int(n_features), int(n_actions)
        sampler = state_sampler
    eta_h = eta if eta_hidden is None else eta_hidden
    deep = mode != "shallow"
    no_trace = mode == "no_trace"
    train_w1 = deep and mode != "elm"     # ELM keeps the hidden layer fixed random

    # --- weights ---
    # SIGNED weights in [-w_max, w_max]: each synapse is a differential conductance pair
    # w = G+ - G- (the manuscript's crossbar realisation), so a synapse can be excitatory
    # or inhibitory. This is required, not cosmetic: with the one-hot-per-feature XOR
    # encoding the four patterns carry equal total positive charge, so a positive-only
    # readout cannot separate them -- a signed readout (subtracting the AND-like feature)
    # can. Weights are RANDOM, not uniform: an identical init makes every hidden neuron
    # receive the same drive, collapsing the hidden layer to one effective unit (zero
    # capacity); independent random init breaks that symmetry. Per-layer scale (w_scale1,
    # w_scale2) sets the summed drive near the LIF threshold rather than saturating it
    # (output fan-in ~H/2 active hidden neurons -> smaller w_scale2).
    def _init(shape, scale):
        return np.clip(scale * rng.standard_normal(shape), -w_max, w_max)
    if deep:
        if mode == "elm":
            # fixed random hidden projection (ELM / random features); never updated.
            W1 = _init((B, F, H), w_scale1)
        else:
            W1 = _init((B, F, H), w_scale1)
        W2 = _init((B, H, A), w_scale2)
    else:
        # shallow: a single trained layer F->A (the hidden layer is absent)
        W2 = _init((B, F, A), w_scale1)
        W1 = None

    # --- device eligibility gates (one bank per trained matrix) ---
    g1 = GateBankBatched(B, F, H, tau_leak=tau_leak, dt=dt, V=V) if deep else None
    g_out = GateBankBatched(B, (H if deep else F), A, tau_leak=tau_leak, dt=dt, V=V)

    # --- DFA feedback matrix: fixed random, drawn ONCE, independent of W2 (no transport) ---
    B_fix = fb_scale * rng.standard_normal((A, H)) if deep else None

    baseline = np.full(B, 1.0 / A)
    bidx = np.arange(B)
    cue = (0.3, 0.3 + cue_dur)
    reward_lag = int(round(D / dt))
    nsteps = int(round(cue[1] / dt)) + 2
    rewards = np.zeros((B, trials))

    # Feedback-alignment diagnostic. The credit DFA delivers to the hidden layer is
    # B_fix @ e_out; the credit backprop WOULD deliver is W2 @ e_out. Both are linear in
    # the output error, so per seed the angle between the matrices W2 (H,A) and B_fix.T
    # (H,A) measures alignment: FA theory predicts learning succeeds when W2 evolves to
    # align with B_fix.T (angle -> 0) and stalls when they stay orthogonal (~90 deg).
    # Logged every `align_every` trials when requested; never affects the dynamics.
    align_log = None
    if log_align and deep:
        align_every = max(1, trials // 60)
        align_t, align_a = [], []

    def _align_angle():
        # per-seed angle (deg) between flattened W2 and B_fix.T
        bt = B_fix.T[None]                       # (1,H,A)
        a = W2.reshape(B, -1); b = np.broadcast_to(bt, (B, H, A)).reshape(B, -1)
        num = (a * b).sum(1)
        den = np.linalg.norm(a, axis=1) * np.linalg.norm(b, axis=1) + 1e-12
        return np.degrees(np.arccos(np.clip(num / den, -1.0, 1.0)))

    # Homeostatic activity regulation (optional). Each hidden neuron keeps a slow running
    # estimate of its own firing fraction and is nudged toward a shared target rate. This
    # is LOCAL (a neuron sees only its own activity) and hardware-plausible (an activity
    # integrator per row). It opposes the winner-take-all policy collapse diagnosed as the
    # cause of the DFA bimodality WITHOUT fighting the XOR solution, whose correct policy
    # has balanced marginal activity across actions. Off (homeo=0) reproduces the baseline.
    act_hidden = np.full((B, H), homeo_target) if (deep and homeo > 0) else None

    for tr in range(trials):
        sigma = sigma0 + (sigma1 - sigma0) * tr / trials
        # Read-time device faults applied ONCE per trial (weights change only per trial):
        # the array reads the faulted conductance while learning targets the clean weight,
        # so faults are a fixed physical realisation, not a learnable parameter. Computing
        # this here (not per timestep) is the difference between a tractable run and not.
        W1r = weight_fault(W1) if (weight_fault is not None and deep) else W1
        W2r = weight_fault(W2) if (weight_fault is not None) else W2
        if log_align and deep and (tr % align_every == 0):
            align_t.append(tr); align_a.append(_align_angle())
        if deep:
            g1.reset()
        g_out.reset()
        lines, correct = sampler(rng, B)
        vh = np.zeros((B, H))
        vo = np.zeros((B, A))
        spk_o = np.zeros((B, A))
        spk_h = np.zeros((B, H))

        for n in range(nsteps):
            t = n * dt
            on = cue[0] <= t < cue[1]
            pre_in = ((rng.random((B, F)) < (in_rate * dt) * lines).astype(float)
                      if on else np.zeros((B, F)))

            if deep:
                # layer 1: inputs -> hidden LIF
                ch_h = np.einsum('bfh,bf->bh', W1r, pre_in)
                vh, sp_h = lif_step_batched(vh, ch_h, dt, rng, tau_m=tau_m,
                                            v_th=v_th, noise=sigma)
                spk_h += sp_h
                # layer 2: hidden spikes -> action LIF, with a tonic bias drive during
                # the cue. The bias is the LIF analogue of a bias unit: it keeps both
                # action neurons in a firing regime even on patterns whose correct
                # response relies on hidden INHIBITION (the XOR=1 patterns, which fire no
                # excitatory detector), so the output is never fully silent and the
                # spike-count contrast that selects the action can form. It is a fixed
                # constant, not learned, identical for both actions (no class bias).
                ch_o = np.einsum('bha,bh->ba', W2r, sp_h.astype(float))
                if on:
                    ch_o = ch_o + bias_o
                vo, sp_o = lif_step_batched(vo, ch_o, dt, rng, tau_m=tau_m,
                                            v_th=v_th, noise=sigma)
                spk_o += sp_o
                # signed leak-dominant coincidence drives BOTH eligibility banks
                if no_trace:
                    g1.step(np.zeros((B, F, H)))
                    g_out.step(np.zeros((B, H, A)))
                else:
                    drive1 = pre_in[:, :, None] * np.where(sp_h, 1.0, -ltd)[:, None, :]
                    g1.step(drive1)
                    drive_o = (sp_h.astype(float)[:, :, None]
                               * np.where(sp_o, 1.0, -ltd)[:, None, :])
                    g_out.step(drive_o)
            else:
                # shallow: inputs -> action LIF directly (single trained layer F->A)
                ch_o = np.einsum('bfa,bf->ba', W2r, pre_in)
                vo, sp_o = lif_step_batched(vo, ch_o, dt, rng, tau_m=tau_m,
                                            v_th=v_th, noise=sigma)
                spk_o += sp_o
                drive_o = pre_in[:, :, None] * np.where(sp_o, 1.0, -ltd)[:, None, :]
                g_out.step(drive_o)

        # --- snapshot eligibility after the action->reward delay (undriven relaxation) ---
        # The relaxation has NO drive and no spikes, so the cascade/leaky dynamics are
        # numerically smooth: a coarsened forward-Euler step (stride x dt) is accurate
        # while stride*dt/tau << 1 (stride=10, dt=5e-3, tau>=2 s -> ratio <= 0.025), and
        # is ~stride x cheaper than ticking every dt. This is the same trick as
        # maze._relax, and it dominates runtime (reward_lag ~ D/dt undriven ticks/trial).
        #
        # TEMPORAL DISTRACTOR (when t_distract is set): partway through the delay an
        # UNINFORMATIVE coincidence fires on the eligibility gates -- a signed drive,
        # uncorrelated with the XOR label, on the same lines that carried the cue. It adds
        # fresh eligibility late in the delay, so a recency-weighted trace mis-credits it,
        # while the device's interval-peaked band-pass trace (whose surviving eligibility
        # at reward still favours the EARLIER cue) does not. This makes the deep task hard
        # in TIME (eligibility must bridge the cue->reward gap past interference) AND in
        # DEPTH (DFA + homeostasis must still resolve hidden-layer credit) at once.
        if not no_trace:
            stride = 10 if reward_lag > 200 else 1
            if t_distract is None:
                n_relax, rem = divmod(reward_lag, stride)
                _relax_gate(g_out, n_relax, stride, rem)
                if deep:
                    _relax_gate(g1, n_relax, stride, rem)
            else:
                # relax to the distractor, inject it (single dt ticks), relax to reward
                i_d0 = int(round(t_distract / dt))
                i_d1 = i_d0 + int(round(distract_dur / dt))
                i_d0 = min(i_d0, reward_lag); i_d1 = min(i_d1, reward_lag)
                pre_d, post_d = i_d0, reward_lag - i_d1
                n0, r0 = divmod(pre_d, stride)
                _relax_gate(g_out, n0, stride, r0)
                if deep:
                    _relax_gate(g1, n0, stride, r0)
                for _ in range(i_d1 - i_d0):
                    # distractor: random coincidence on the cue lines, label-independent
                    dfire = (rng.random((B, F)) < in_rate * dt).astype(float)
                    if deep:
                        hfire = (rng.random((B, H)) < in_rate * dt).astype(float)
                        g1.step(dfire[:, :, None] * np.where(hfire, 1.0, -ltd)[:, None, :])
                        ofire = (rng.random((B, A)) < in_rate * dt).astype(float)
                        g_out.step(hfire[:, :, None] * np.where(ofire, 1.0, -ltd)[:, None, :])
                    else:
                        ofire = (rng.random((B, A)) < in_rate * dt).astype(float)
                        g_out.step(dfire[:, :, None] * np.where(ofire, 1.0, -ltd)[:, None, :])
                n1, r1 = divmod(post_d, stride)
                _relax_gate(g_out, n1, stride, r1)
                if deep:
                    _relax_gate(g1, n1, stride, r1)
            eo_rew = g_out.vn[..., -1] / g_out.Vnmax
            e1_rew = (g1.vn[..., -1] / g1.Vnmax) if deep else None
        else:
            eo_rew = np.zeros_like(g_out.vn[..., -1])
            e1_rew = np.zeros((B, F, H)) if deep else None

        # --- action selection (argmax output spikes; ties -> random) ---
        tie = spk_o.max(1) == spk_o.min(1)
        chosen = np.argmax(spk_o, 1)
        chosen[tie] = rng.integers(A, size=int(tie.sum()))
        R_true = (chosen == correct).astype(float)        # task performance (returned)
        if reward_pools is None:
            R = R_true                                     # synthetic clean reward
        else:
            # biosignal reward: per seed draw a decoded reward VALUE from the EEG pool
            # matching the TRUE outcome valence (exactly as bandit.train) -- a real, noisy
            # reward-prediction-error gates the update, while R_true scores performance.
            R = np.array([reward_pools[int(R_true[b])][
                rng.integers(len(reward_pools[int(R_true[b])]))] for b in range(B)])

        # --- learning signals ---
        adv = (R - baseline)
        # OUTPUT layer: the proven three-factor scalar update dw = eta (R-b) e (the
        # contextual-bandit rule). Per-action selectivity comes from the SIGNED
        # eligibility itself -- the action neuron that fired has e>0, the others e<0 --
        # so the scalar advantage suffices and no softmax factor is needed (an explicit
        # (onehot-pol) factor at the output double-counts the action and destabilises the
        # update sign). This is identical across all modes; only the HIDDEN credit differs.
        Lo = adv[:, None, None]                                   # (B,1,1) global scalar

        if deep:
            if mode == "global":
                # pure global scalar to the hidden layer too: every hidden synapse sees
                # the SAME third factor, so the rule cannot tell which hidden unit helped
                # (the structural-credit failure mode this experiment is built to expose).
                Lh = adv[:, None, None]                           # (B,1,1)
            else:
                # DFA: a per-action error e_a = (R-b)(1[a=chosen]-pi_a) projected to the
                # hidden layer through a FIXED RANDOM feedback matrix B_fix (no W2^T, no
                # gradient transport). This gives each hidden unit a DIFFERENT, action-
                # dependent learning signal -- the per-neuron credit the global scalar
                # lacks -- while remaining fully local. pi is a softmax over the output
                # spike counts (used only to form the feedback signal, not the output update).
                logits = spk_o - spk_o.mean(1, keepdims=True)
                pol = np.exp(logits - logits.max(1, keepdims=True))
                pol /= pol.sum(1, keepdims=True)
                onehot = np.zeros((B, A)); onehot[bidx, chosen] = 1.0
                L_a = adv[:, None] * (onehot - pol)               # (B,A)
                Lh = np.einsum('ah,ba->bh', B_fix, L_a)[:, None, :]   # (B,1,H)

        W2 = np.clip(W2 + eta * Lo * eo_rew, -w_max, w_max)
        if train_w1:
            dW1 = eta_h * Lh * e1_rew
            if homeo > 0 and act_hidden is not None:
                # Local homeostatic regulation: update each hidden neuron's slow activity
                # estimate from this trial's firing fraction, then nudge its INCOMING
                # weights to push that activity toward the shared target. A neuron firing
                # above target has its drive scaled down, below target scaled up -- which
                # breaks the winner-take-all lock-in (an over-firing unit is reined in)
                # without touching the per-pattern selectivity the DFA signal learns.
                rate = spk_h / nsteps                              # (B,H) this-trial rate
                act_hidden += (rate - act_hidden) / homeo_tau      # slow running estimate
                # multiplicative scaling of the incoming weights toward balanced activity
                scale = 1.0 + homeo * (homeo_target - act_hidden)  # (B,H)
                dW1 = dW1 + (scale[:, None, :] - 1.0) * W1
            W1 = np.clip(W1 + dW1, -w_max, w_max)

        baseline += 0.02 * (R - baseline)        # baseline tracks the gating reward
        rewards[:, tr] = R_true                   # report TRUE task performance

        # Early stop: once the batch-mean reward over a trailing window has plateaued
        # above the criterion (or the run is clearly stuck near chance late on), the
        # remaining trials add no information -- pad with the converged level and break.
        # Cheap O(1) check; large speed-up on the many runs that converge well before
        # ``trials``. Disabled by ``early_stop=None``.
        if early_stop is not None and tr >= early_stop["min_trials"] and \
                (tr % early_stop["check_every"] == 0):
            w = early_stop["window"]
            recent = rewards[:, tr - w + 1:tr + 1].mean()
            prev = rewards[:, tr - 2 * w + 1:tr - w + 1].mean()
            converged = recent >= early_stop["criterion"] and abs(recent - prev) < early_stop["tol"]
            stuck = tr >= early_stop["stuck_after"] and recent < early_stop["chance"] + early_stop["tol"]
            if converged or stuck:
                rewards[:, tr + 1:] = recent      # pad remaining trials at current level
                break

    if log_align and deep:
        # align_log: dict with trial indices and per-seed angle (deg), shape (n_log, B)
        align_log = {"trials": np.array(align_t), "angle_deg": np.array(align_a)}

    if return_weights and log_align:
        return rewards, (W1, W2), align_log
    if return_weights:
        return rewards, (W1, W2)
    if log_align:
        return rewards, align_log
    return rewards


def reward_rate(rewards, window=100):
    rewards = np.asarray(rewards)
    if rewards.shape[-1] < window:
        return rewards.mean(axis=-1)
    return rewards[..., -window:].mean(axis=-1)


# =============================================================================
# Experiment cores (the deep all-local RL studies that compose ``train_deep``)
#
# Each ``run_*`` is SERIAL and returns the result grid as a plain dict -- no file
# I/O, no plotting, no stdout. Notebooks call these at a small (quick) seed/trial
# count and render the figures inline; ``main()`` (below) calls them at the
# published scale (and MAY use a multiprocessing.Pool over the coarse axis), then
# writes each grid under ``data/results/`` via ``paths.save_result``. The DMS
# temporal-distractor task (Experiment 12) is a shallow ``GateBankBatched`` variant
# whose difficulty is purely temporal, so it lives HERE alongside the deep stack
# rather than in the selectivity module.
# =============================================================================

# Frozen operating points (pilot-tuned; see the pre-registration docs). Kept at
# module scope so both the ``run_*`` cores and ``main()`` share one source of truth.
DEEP_LOCAL_HP = dict(H=32, tau_leak=10.0, D=5.0, eta=0.2, eta_hidden=3.0,
                     fb_scale=2.0, bias_o=0.3, V=1.5, sigma0=0.15, sigma1=0.05)
DEEP_LOCAL_HOMEO = 0.1                 # frozen homeostatic strength (pilot-tuned)
DEEP_LOCAL_MODES = ["shallow", "elm", "global", "dfa", "no_trace", "dfa_homeo"]

# Experiment 13 operating point at the VERIFIED homeostasis regime (H=32 clean:
# homeo=0.1 -> 1.00 [1.00,1.00] vs homeo=0 -> 0.58 [0.47,0.69], gap +0.42, disjoint;
# device faults are exp14's separate axis, not confounded into the homeostasis claim).
DEEP_DMS_HP = dict(H=32, tau_leak=10.0, D=3.0, eta=0.2, eta_hidden=3.0,
                   fb_scale=2.0, bias_o=0.3, V=1.5, sigma0=0.15, sigma1=0.05)
DEEP_DMS_HOMEO = 0.1
DEEP_DMS_T_DISTRACT = 2.0              # distractor 2 s into the 3 s delay (1 s before reward)
DEEP_DMS_DISTRACT_DUR = 0.3
# (label, train_mode, homeo, distractor?)
DEEP_DMS_CONDS = [
    ("dfa_homeo_nodist", "dfa", DEEP_DMS_HOMEO, False),
    ("dfa_homeo_dist",   "dfa", DEEP_DMS_HOMEO, True),
    ("dfa_dist",         "dfa", 0.0,   True),
    ("no_trace_dist",    "no_trace", 0.0, True),
]

# Experiment 14 array-scale early-stop (converged/stuck runs pad the rest at the
# reached level, leaving the final-200 metric unchanged; large speed-up at ceiling).
ARRAY_SCALE_EARLY_STOP = {
    "min_trials": 400, "check_every": 100, "window": 150,
    "criterion": 0.78, "tol": 0.03, "chance": 0.5, "stuck_after": 1200,
}


# ----------------------------------------------------------------------------
# Experiment 12: delayed-match-to-sample with a temporal distractor.
#
# The harder credit-assignment task identified as the single best fit for a SLOW
# (seconds-scale) eligibility trace: a delayed-association task whose difficulty is
# purely TEMPORAL (a seconds-scale gap and an interfering distractor), not a larger
# policy. It stresses the trace itself, sits in the device's 1-20 s regime, stays at
# 2 actions, and is the canonical "fails without an eligibility trace" regime (Bellec
# et al. e-prop, Nat. Commun. 2020, 10.1038/s41467-020-17236-y; DMS origin Miller,
# Erickson & Desimone, J. Neurosci. 1996, 10.1523/JNEUROSCI.16-16-05154.1996).
#
# The device band-pass trace (sigmoidal rise + tau_leak decay) is peaked at a
# non-zero lag, so with tau_leak tuned to the sample->reward interval it credits the
# sample over the later distractor -- the closed-loop, standard-task form of the
# interval-selectivity result (manuscript S=3.96 vs 0.48). The abstract single-
# exponential trace (matched tau) is monotone/recency-weighted (the informative
# comparison); the no-trace control must fail.
# ----------------------------------------------------------------------------
def _relax(bank, n_steps, dt, stride=10):
    """Fast-forward the (drive-free) eligibility gate by ``n_steps`` dt-ticks.

    During the silent delay there is no presynaptic drive and no spikes, so the gate
    just relaxes -- a smooth trajectory faithfully reproduced at a coarser effective
    step ``stride*dt`` (the same long-delay shortcut as ``train_sequential``;
    ``stride*dt`` stays well below ``tau_leak``/``tau_r``, so the cascade is unchanged).
    Returns the readout at the end of the interval.
    """
    if n_steps <= 0:
        return bank.step(np.zeros((bank.e.shape if hasattr(bank, "e") else
                                   bank.vn.shape[:-1])) * 0.0)
    zero = None
    base = bank.dt
    coarse = max(1, n_steps // stride)
    bank.dt = base * (n_steps / coarse)               # exact-coverage coarse step
    out = None
    for _ in range(coarse):
        if zero is None:
            # build a correctly-shaped zero drive once (shape (B,S,A))
            shp = bank.e.shape if hasattr(bank, "e") else bank.vn.shape[:-1]
            zero = np.zeros(shp)
        out = bank.step(zero)
    bank.dt = base
    return out


def run_dms(cond, *, B=20, tau_leak=1.5, G=5.0, t_distract=4.0, trials=2500,
            dt=5e-3, cue_dur=0.3, distract_dur=0.3, eta=0.2, V=1.5, in_rate=200.0,
            ltd=LTD_BIAS, tau_m=TAU_M, v_th=V_TH, sigma=0.15, seed0=0):
    """One DMS-with-distractor condition (Experiment 12). SERIAL, returns rewards
    ``(B, trials)``; no file I/O.

    State lines ``[sample0, sample1, distractor]`` -> 2 action neurons. A trial:
    sample coincidence at ``t in [0, cue_dur)``; distractor coincidence at
    ``t in [t_distract, t_distract+distract_dur)``; decision + reward at ``t = G``.
    The eligibility gate integrates continuously across the whole trial; the surviving
    trace at the reward instant (``t=G``) is what the three-factor rule commits.

    ``cond`` is ``"device"`` (trap-discharge :class:`GateBankBatched`), ``"abstract"``
    (matched-tau single exponential :class:`AbstractTrace`, the recency-trace baseline)
    or ``"no_trace"`` (eligibility zeroed, the necessity control).

    Optimised by stepping the LIF + gate only during the two ACTIVE windows (sample,
    distractor) and fast-forwarding the gate through the two SILENT gaps with a coarse
    stride (:func:`_relax`); ~88% of dt-ticks are silent, so this is ~6x fewer
    iterations with no change to the dynamics that matter.
    """
    from .bandit import GateBankBatched, AbstractTrace, W_INIT, W_MAX
    rng = np.random.default_rng(seed0)
    S, A = 3, 2                       # [sample0, sample1, distractor] x {A, B}
    DIST_LINE = 2
    if cond == "abstract":
        bank = AbstractTrace(B, S, A, tau_elig=tau_leak, dt=dt)
    else:
        bank = GateBankBatched(B, S, A, tau_leak=tau_leak, V=V, dt=dt)
    no_trace = (cond == "no_trace")

    w = np.full((B, S, A), W_INIT)
    baseline = np.full(B, 1.0 / A)
    bidx = np.arange(B)
    n_cue = int(round(cue_dur / dt))
    n_gap1 = int(round((t_distract - cue_dur) / dt))
    n_dist = int(round(distract_dur / dt))
    n_gap2 = int(round((G - t_distract - distract_dur) / dt))
    rewards = np.zeros((B, trials))

    def active_window(v, spk, sample_line, count_spikes, distractor):
        """Step LIF + gate for one active window; return (v, e_last)."""
        e = None
        nsteps = n_cue if not distractor else n_dist
        for _ in range(nsteps):
            pre = np.zeros((B, S))
            if not distractor:
                pre[bidx, sample_line] = (rng.random(B) < in_rate * dt).astype(float)
            else:
                pre[:, DIST_LINE] = (rng.random(B) < in_rate * dt).astype(float)
            charge = np.einsum('bsa,bs->ba', w, pre)
            v, sp = lif_step_batched(v, charge, dt, rng, tau_m=tau_m, v_th=v_th,
                                     noise=sigma)
            if count_spikes:
                spk += sp
            if no_trace:
                e = bank.step(np.zeros((B, S, A)))
            else:
                drive = pre[:, :, None] * np.where(sp, 1.0, -ltd)[:, None, :]
                e = bank.step(drive)
        return v, e

    for tr in range(trials):
        bank.reset()
        cls = rng.integers(2, size=B)                  # sample class this trial
        sample_line = np.where(cls == 0, 0, 1)
        v = np.zeros((B, A))
        spk = np.zeros((B, A))                          # decision-window spike count
        # SAMPLE window (informative; reads the decision)
        v, _ = active_window(v, spk, sample_line, count_spikes=True, distractor=False)
        # GAP 1 -- silent relaxation
        _relax(bank, n_gap1, dt)
        # DISTRACTOR window (uninformative; spikes not counted)
        v, _ = active_window(v, spk, sample_line, count_spikes=False, distractor=True)
        # GAP 2 -- silent relaxation to the reward instant; snapshot surviving trace
        e_rew = _relax(bank, n_gap2, dt)
        # action = spiking winner over the sample window (ties -> random); membrane
        # noise provides exploration
        tie = spk.max(1) == spk.min(1)
        chosen = np.argmax(spk, 1)
        chosen[tie] = rng.integers(A, size=int(tie.sum()))
        r = (chosen == cls).astype(float)              # match-to-sample reward
        w = np.clip(w + (eta * (r - baseline))[:, None, None] * e_rew, 0.0, W_MAX)
        baseline += 0.02 * (r - baseline)
        rewards[:, tr] = r
    return rewards


def final_rate(rw, window=300):
    """Per-seed reward rate over the last ``window`` trials (DMS final metric)."""
    return rw[:, -window:].mean(1)


def _summarize_dms(raw, conds, *, trials, chance, crit, tau_leak=1.5, G=5.0,
                   t_distract=4.0):
    """Package raw per-condition :func:`run_dms` output into the exp12 grid dict
    (seed-mean curves, per-seed finals + bootstrap CIs, pre-registered H1-H3/K1).

    Pre-registered (fixed before running):
      H1 device learns DMS-with-distractor:   device final >= crit
      H2 trace is necessary:                   no-trace <= chance + 0.10
      H3 band-pass beats recency at crediting the sample over the distractor:
         device CI lower bound > abstract CI upper bound (genuine separation)
      K1 (kill) device <= chance+0.10 -> the device-trace claim fails on this task
    """
    from .stats import bootstrap_ci
    curves, finals, ci = {}, {}, {}
    for cond in conds:
        rw = raw[cond]
        f = final_rate(rw)
        curves[cond] = rw.mean(0)
        finals[cond] = f
        ci[cond] = bootstrap_ci(f)
    dev = finals["device"].mean()
    nt = finals["no_trace"].mean()
    h1 = bool(dev >= crit)
    h2 = bool(nt <= chance + 0.10)
    h3 = bool(ci["device"][0] > ci["abstract"][1])       # genuine separation (CIs disjoint)
    k1 = bool(dev <= chance + 0.10)
    B = raw[conds[0]].shape[0]
    return {"curves": curves, "raw": {k: raw[k] for k in conds},
            "finals": {k: v for k, v in finals.items()},
            "ci": ci, "seeds": B, "trials": trials, "chance": chance, "crit": crit,
            "tau_leak": tau_leak, "G": G, "t_distract": t_distract,
            "criteria": {"H1": h1, "H2": h2, "H3": h3, "K1": k1}}


def run_dms_all(*, seeds=20, trials=2500, conds=("device", "abstract", "no_trace")):
    """Experiment 12 core: DMS-with-distractor over all conditions (SERIAL, returns
    the packaged grid dict; no file I/O). Notebooks call this in quick mode."""
    raw = {c: run_dms(c, B=seeds, trials=trials) for c in conds}
    return _summarize_dms(raw, conds, trials=trials, chance=0.5, crit=0.75)


# ----------------------------------------------------------------------------
# Experiment 7: deep all-local RL with a physical eligibility trace (XOR).
#
# Five conditions isolate the SPATIAL-credit pathway (the device trace is the
# temporal factor in all of them); "dfa_homeo" adds the local homeostatic activity
# regulator (the diagnosed fix for the DFA policy-collapse bimodality). See
# experiments/PREREGISTRATION_deep_local.md for the frozen operating point + criteria.
# ----------------------------------------------------------------------------
def _deep_local_one(mode, *, seeds, trials, hp, homeo):
    """Run one deep-local condition over all seeds. Returns (mode, finals[seeds],
    curve[trials-100]). "dfa_homeo" is the dfa rule plus the homeostatic regulator."""
    kw = dict(hp)
    train_mode = mode
    if mode == "dfa_homeo":                           # dfa rule + homeostatic regulator
        train_mode = "dfa"
        kw["homeo"] = homeo
    rew = train_deep(mode=train_mode, B=seeds, trials=trials, seed0=0, **kw)
    finals = reward_rate(rew, window=200)             # per-seed final reward rate
    # running-mean learning curve (window 100), averaged over seeds
    win = 100
    csum = np.cumsum(rew, axis=1)
    run = (csum[:, win:] - csum[:, :-win]) / win      # (B, trials-win)
    curve = run.mean(0)
    return mode, finals, curve


def _summarize_deep_local(res, *, seeds, trials, hp, homeo, chance, crit,
                          modes=None):
    """Package raw ``(mode, finals, curve)`` triples into the exp7 grid dict with
    bootstrap CIs and the pre-registered criteria C1-C4/C6/K4.

    C1 shallow fails (<= chance+0.10) -> depth is genuinely required
    C2 no-trace fails                 -> eligibility necessity
    C3 DFA learns (>= crit)
    C4 DFA > global (CIs disjoint)    -> per-neuron credit beats the global scalar
    C6 homeostasis fix (>= crit and DFA+homeo CI-above DFA)
    K4 (kill/bound) ELM shortcut solves XOR? (bounds the depth-without-learning claim)
    """
    from .stats import bootstrap_ci
    modes = list(modes) if modes is not None else DEEP_LOCAL_MODES
    finals = {m: f for m, f, _ in res}
    curves = {m: c for m, _, c in res}
    ci = {m: bootstrap_ci(finals[m]) for m in modes}
    dfa, homeo_f = finals["dfa"], finals["dfa_homeo"]
    c1 = bool(finals["shallow"].mean() <= chance + 0.10)
    c2 = bool(finals["no_trace"].mean() <= chance + 0.10)
    c3 = bool(dfa.mean() >= crit)
    c4 = bool(ci["dfa"][0] > ci["global"][1])            # CI gap excludes 0
    k4 = bool(finals["elm"].mean() >= crit)              # ELM shortcut solves it?
    c6 = bool(homeo_f.mean() >= crit and ci["dfa_homeo"][0] > ci["dfa"][1])
    return {"finals": finals, "curves": curves, "ci": ci, "HP": hp, "homeo": homeo,
            "seeds": seeds, "trials": trials, "chance": chance, "crit": crit,
            "criteria": {"C1": c1, "C2": c2, "C3": c3, "C4": c4, "C6": c6, "K4": k4}}


def run_deep_local(*, seeds=20, trials=3000, hp=None, homeo=DEEP_LOCAL_HOMEO,
                   modes=None, chance=0.5, crit=0.75):
    """Experiment 7 core: deep all-local XOR over all conditions (SERIAL, returns the
    packaged grid dict; no file I/O). Notebooks call this in quick mode."""
    hp = DEEP_LOCAL_HP if hp is None else hp
    modes = list(modes) if modes is not None else DEEP_LOCAL_MODES
    res = [_deep_local_one(m, seeds=seeds, trials=trials, hp=hp, homeo=homeo)
           for m in modes]
    return _summarize_deep_local(res, seeds=seeds, trials=trials, hp=hp, homeo=homeo,
                                 chance=chance, crit=crit, modes=modes)


# ----------------------------------------------------------------------------
# Experiment 13: the convergence test -- deep XOR with a temporal distractor.
#
# One task where the device eligibility trace must do TEMPORAL credit assignment
# (bridge a seconds-scale cue->reward delay past an interfering distractor) AND DFA +
# local homeostasis must do SPATIAL credit assignment through a hidden layer (XOR).
# ----------------------------------------------------------------------------
def _deep_dms_one(spec, *, seeds, trials, hp, t_distract, distract_dur):
    """Run one exp13 condition over all seeds. ``spec`` is (label, mode, homeo, dist).
    Returns (label, rewards (B, trials))."""
    label, mode, homeo, dist = spec
    kw = dict(hp)
    kw["homeo"] = homeo
    if dist:
        kw["t_distract"] = t_distract
        kw["distract_dur"] = distract_dur
    rew = train_deep(mode=mode, B=seeds, trials=trials, seed0=0, **kw)
    return label, rew                       # full (B, trials) for curve + bar


def _summarize_deep_dms(raw, conds, *, seeds, trials, hp, homeo, t_distract,
                        chance, crit):
    """Package raw exp13 output into the grid dict with per-seed finals, bootstrap
    CIs, seeds-solved robustness, and the pre-registered H1-H3/K1.

    H1 full stack survives the distractor:   dfa_homeo+dist >= crit
    H2 eligibility necessary:                no_trace+dist <= chance + 0.10
    H3 homeostasis still helps under interference: dfa_homeo+dist > dfa+dist (CIs)
    K1 (kill) dfa_homeo+dist collapses to chance -> the convergence claim fails
    """
    from .stats import bootstrap_ci
    finals, ci, curves, seeds_solved = {}, {}, {}, {}
    for label, *_ in conds:
        rw = raw[label]
        f = reward_rate(rw, window=200)
        finals[label] = f
        ci[label] = bootstrap_ci(f)
        curves[label] = rw.mean(0)                      # seed-mean per trial
        seeds_solved[label] = int((f >= crit).sum())    # per-seed robustness
    fh = finals["dfa_homeo_dist"].mean()
    nt = finals["no_trace_dist"].mean()
    h1 = bool(fh >= crit)
    h2 = bool(nt <= chance + 0.10)
    h3 = bool(ci["dfa_homeo_dist"][0] > ci["dfa_dist"][1])   # homeostasis helps (CIs disjoint)
    k1 = bool(fh <= chance + 0.10)
    return {"finals": {k: v for k, v in finals.items()}, "ci": ci,
            "curves": curves, "seeds_solved": seeds_solved,
            "HP": hp, "homeo": homeo, "seeds": seeds, "trials": trials,
            "t_distract": t_distract, "chance": chance, "crit": crit,
            "criteria": {"H1": h1, "H2": h2, "H3": h3, "K1": k1}}


def run_deep_dms(*, seeds=20, trials=3000, hp=None, homeo=DEEP_DMS_HOMEO,
                 conds=None, t_distract=DEEP_DMS_T_DISTRACT,
                 distract_dur=DEEP_DMS_DISTRACT_DUR, chance=0.5, crit=0.75):
    """Experiment 13 core: deep XOR + temporal distractor over all conditions (SERIAL,
    returns the packaged grid dict; no file I/O). Notebooks call this in quick mode."""
    hp = DEEP_DMS_HP if hp is None else hp
    conds = DEEP_DMS_CONDS if conds is None else conds
    raw = dict(_deep_dms_one(spec, seeds=seeds, trials=trials, hp=hp,
                             t_distract=t_distract, distract_dur=distract_dur)
               for spec in conds)
    return _summarize_deep_dms(raw, conds, seeds=seeds, trials=trials, hp=hp,
                               homeo=homeo, t_distract=t_distract,
                               chance=chance, crit=crit)


# ----------------------------------------------------------------------------
# Experiment 14: array-scale feasibility of the all-local rule under measured faults.
#
# Does the device-eligibility three-factor rule (DFA + homeostasis, the deep all-local
# stack) keep learning as the network grows (hidden width H) AND as the synaptic array
# is corrupted by the SAME measured SiO_x non-idealities (stuck-off + lognormal D2D +
# optionally the Poole-Frenkel I-V nonlinearity)? PASS = beats its own no-trace control
# with disjoint CIs AND > 0.75. KILL: K1 large-H fails at p=0 (no width composition);
# K2 high-p collapses everywhere; K3 control learns (artefact).
#
# The two provenance grids differ only in the fault prior:
#   exp14_array_scale       : pf_on=True,  sigma_g=0.5 symmetric  (the "full PF" grid);
#   exp14_array_scale_sweep : pf_on=False, sigma_g=0.5/sigma_g_on=0.05 asymmetric
#     (PF is a READ-path nonideality, excluded when crediting the PROGRAMMED weight).
# ``run_array_scale`` is the single core; the two are selected by its fault kwargs.
# ----------------------------------------------------------------------------
def _array_scale_run(job):
    """One independent training run -> (H, p, seed, trace, final reward rate).
    Picklable module-level worker for the ``main()`` Pool. ``job`` carries the fault-
    stack parameters so children need no shared config.
    """
    from .device_faults import siox_fault_stack
    H, p, seed, trace, fkw = job
    fault = siox_fault_stack(p_stuck=p, sigma_g=fkw["sigma_g"],
                             sigma_g_on=fkw["sigma_g_on"], pf_on=fkw["pf_on"],
                             stuck_kind="off", seed=1000 * seed + int(100 * p))
    rewards = train_deep(
        mode="dfa" if trace else "no_trace",
        H=H, tau_leak=10.0, D=5.0, t_distract=fkw["t_distract"], distract_dur=0.3,
        homeo=1.0 if trace else 0.0, trials=fkw["trials"], seed0=seed,
        weight_fault=fault, early_stop=ARRAY_SCALE_EARLY_STOP,
    )
    return (H, p, seed, trace, float(rewards[:, -200:].mean()))


def _array_scale_reduce(results, H_grid, p_grid, *, seeds, sigma_g, sigma_g_on,
                        pf_on, trials):
    """Reduce a list of ``_array_scale_run`` outputs into the exp14 grid dict
    (device mean/lo/hi per (H,p), control mean, fault label)."""
    from .stats import bootstrap_ci
    dev = {(H, p): [] for H in H_grid for p in p_grid}
    ctl = {(H, p): [] for H in H_grid for p in p_grid}
    for H, p, s, trace, val in results:
        (dev if trace else ctl)[(H, p)].append(val)
    grid = np.zeros((len(H_grid), len(p_grid), 3))       # mean, lo, hi (device)
    cgrid = np.zeros((len(H_grid), len(p_grid)))         # control mean
    passes = {}
    for i, H in enumerate(H_grid):
        for j, p in enumerate(p_grid):
            d = np.array(dev[(H, p)]); c = np.array(ctl[(H, p)])
            lo, hi = bootstrap_ci(d, seed=i * 10 + j)
            grid[i, j] = (d.mean(), lo, hi); cgrid[i, j] = c.mean()
            passes[(H, p)] = bool(lo > c.mean() and d.mean() > 0.75)
    if pf_on:
        faults = "stuck-off + D2D(0.5) + PF-nonlinearity"
    else:
        faults = ("stuck-off + D2D({:g},{:g}) [PF excluded: read-path nonideality]"
                  .format(sigma_g_on if sigma_g_on is not None else sigma_g, sigma_g))
    return {"H": list(H_grid), "p": list(p_grid), "grid": grid, "ctrl": cgrid,
            "faults": faults, "seeds": seeds, "trials": trials, "passes": passes}


def run_array_scale(*, H_grid=(8, 32, 128, 512), p_grid=(0.0, 0.05, 0.20, 0.50),
                    seeds=12, trials=2000, sigma_g=0.5, sigma_g_on=None,
                    pf_on=True, pool=None):
    """Experiment 14 core: array-scale fault-tolerance feasibility (SERIAL by default,
    returns the packaged grid dict; no file I/O). Unifies ``array_scale_faults`` (the
    ``pf_on=True`` full-PF grid, ``exp14_array_scale.npy``) and ``entry_array_scale``
    (the ``pf_on=False`` asymmetric-D2D sweep, ``exp14_array_scale_sweep.npy``) -- the
    fault prior is selected by the ``sigma_g``/``sigma_g_on``/``pf_on`` kwargs.

    Every (H, p, seed) x {device, control} cell is an independent :func:`train_deep`
    call. In-notebook this runs serially; ``main()`` passes a ``multiprocessing.Pool``
    via ``pool`` to fan the jobs out across the coarse (cell x seed) axis.
    """
    H_grid = list(H_grid); p_grid = list(p_grid)
    fkw = dict(sigma_g=sigma_g, sigma_g_on=sigma_g_on, pf_on=pf_on,
               t_distract=3.0, trials=trials)
    jobs = [(H, p, s, trace, fkw)
            for H in H_grid for p in p_grid for s in range(seeds)
            for trace in (True, False)]
    if pool is None:
        results = [_array_scale_run(j) for j in jobs]
    else:
        results = pool.map(_array_scale_run, jobs, chunksize=1)
    return _array_scale_reduce(results, H_grid, p_grid, seeds=seeds, sigma_g=sigma_g,
                               sigma_g_on=sigma_g_on, pf_on=pf_on, trials=trials)


def main(argv=None):
    """Full-scale reproduction CLI for the deep all-local + DMS grids (writes ``data/results``).

    ``python -m mrl_trace.deep [--exp7] [--exp12] [--exp13] [--exp14]
    [--exp14-sweep] [--full|--quick]``

    With no experiment flag, runs all of exp7/exp12/exp13/exp14 (the sweep variant is
    opt-in via ``--exp14-sweep`` since it shares exp14's compute). ``--full`` = the
    published seed/trial scale; ``--quick`` = a fast few-seed smoke run. The exp14 grids
    fan their independent cell x seed runs across a ``multiprocessing.Pool``.
    """
    import argparse
    import os
    import multiprocessing as mp
    from . import paths
    ap = argparse.ArgumentParser(description="Deep all-local + DMS RL reproductions")
    ap.add_argument("--exp7", action="store_true",
                    help="deep all-local XOR -> exp7_deep_local.npy")
    ap.add_argument("--exp12", action="store_true",
                    help="DMS + temporal distractor -> exp12_dms.npy")
    ap.add_argument("--exp13", action="store_true",
                    help="deep XOR + temporal distractor (convergence) -> exp13_deep_dms.npy")
    ap.add_argument("--exp14", action="store_true",
                    help="array-scale faults (full PF) -> exp14_array_scale.npy")
    ap.add_argument("--exp14-sweep", action="store_true",
                    help="array-scale faults (asymmetric D2D, no PF) -> exp14_array_scale_sweep.npy")
    ap.add_argument("--quick", action="store_true", help="fast few-seed smoke run")
    ap.add_argument("--full", action="store_true", help="published-scale run (default)")
    ap.add_argument("--workers", type=int, default=(os.cpu_count() or 4),
                    help="Pool workers for the exp14 grid(s)")
    a = ap.parse_args(argv)
    any_flag = a.exp7 or a.exp12 or a.exp13 or a.exp14 or a.exp14_sweep
    run_default = not any_flag           # no flag -> exp7/12/13/14 (sweep opt-in)

    if a.exp7 or run_default:
        seeds = 6 if a.quick else 20
        trials = 600 if a.quick else 3000
        print(f"=== exp7 deep all-local XOR (N={seeds}, {trials} trials, H={DEEP_LOCAL_HP['H']}) ===")
        grid = run_deep_local(seeds=seeds, trials=trials)
        paths.save_result("exp7_deep_local.npy", grid)
        print(f"  wrote exp7_deep_local.npy  criteria={grid['criteria']}")

    if a.exp12 or run_default:
        seeds = 6 if a.quick else 20
        trials = 800 if a.quick else 2500
        print(f"=== exp12 DMS + distractor (N={seeds}, {trials} trials) ===")
        grid = run_dms_all(seeds=seeds, trials=trials)
        paths.save_result("exp12_dms.npy", grid)
        print(f"  wrote exp12_dms.npy  criteria={grid['criteria']}")

    if a.exp13 or run_default:
        seeds = 6 if a.quick else 20
        trials = 600 if a.quick else 3000
        print(f"=== exp13 deep XOR + distractor (N={seeds}, {trials} trials, "
              f"distractor@{DEEP_DMS_T_DISTRACT}s) ===")
        grid = run_deep_dms(seeds=seeds, trials=trials)
        paths.save_result("exp13_deep_dms.npy", grid)
        print(f"  wrote exp13_deep_dms.npy  criteria={grid['criteria']}")

    if a.exp14 or a.exp14_sweep or run_default:
        try:
            ctx = mp.get_context("fork")     # pure-NumPy single-threaded runs; fork is safe/fast
        except ValueError:                   # pragma: no cover (non-fork OS)
            ctx = mp.get_context()
        if a.quick:
            H_grid, p_grid, seeds, trials = (32, 512), (0.0, 0.5), 4, 400
        else:
            H_grid, p_grid, seeds, trials = (8, 32, 128, 512), (0.0, 0.05, 0.20, 0.50), 12, 2000
        n_cells = len(H_grid) * len(p_grid) * seeds * 2
        with ctx.Pool(min(a.workers, max(1, n_cells))) as pool:
            if a.exp14 or run_default:
                print(f"=== exp14 array-scale faults [full PF] (N={seeds}, H={list(H_grid)}, "
                      f"p={list(p_grid)}) ===")
                grid = run_array_scale(H_grid=H_grid, p_grid=p_grid, seeds=seeds,
                                       trials=trials, sigma_g=0.5, pf_on=True, pool=pool)
                paths.save_result("exp14_array_scale.npy", grid)
                print(f"  wrote exp14_array_scale.npy  faults={grid['faults']!r}")
            if a.exp14_sweep:
                print(f"=== exp14 array-scale faults [asymmetric D2D, no PF] (N={seeds}) ===")
                grid = run_array_scale(H_grid=H_grid, p_grid=p_grid, seeds=seeds,
                                       trials=trials, sigma_g=0.5, sigma_g_on=0.05,
                                       pf_on=False, pool=pool)
                paths.save_result("exp14_array_scale_sweep.npy", grid)
                print(f"  wrote exp14_array_scale_sweep.npy  faults={grid['faults']!r}")


if __name__ == "__main__":
    main()
\end{lstlisting}

\begin{lstlisting}[caption={siox\_eligibility/deep\_sequential.py},label={lst:siox-src-deep_sequential}]
"""Deep sequential all-local spiking RL with a physical eligibility trace -- Arm E.

This module is the composition of the two working reductions in the package:

- :func:`mrl_trace.deep.train_deep` -- a DEEP (hidden-layer) all-local policy
  trained by Direct Feedback Alignment (DFA) with optional local homeostasis, on a
  non-linearly-separable XOR contextual bandit a single trained layer cannot solve.
  It supplies the depth + DFA + homeostasis machinery and the correct update signs.
- :func:`mrl_trace.maze.train_sequential` (the T-maze) -- a SHALLOW policy that
  bridges a cue->reward gap across a multi-step trajectory, with the device eligibility
  integrating CONTINUOUSLY across the steps and a single delayed goal reward gating the
  surviving trace into every used synapse.

The question this module asks is whether a device-supplied physical eligibility trace
lets a FULLY-LOCAL (no backprop, no weight transport) DEEP policy bridge the cue->reward
gap across a genuine multi-step trajectory whose final decision is non-linearly
separable (XOR of two cue bits). It therefore demands all three at once:

  (a) a longer-horizon trajectory: ``L >= 4`` stem steps actually traversed, with the
      goal reward DELAYED by ``D >= 2`` s after the junction decision;
  (b) a HIDDEN layer trained by DFA, required because the cued decision is XOR (a single
      trained layer cannot solve it);
  (c) LOCAL HOMEOSTASIS stabilising hidden-unit firing.

Task ("deep T-maze"). At episode start a two-bit cue ``(b0, b1)`` is drawn; the correct
arm at the junction is ``XOR(b0, b1)``. The agent then AUTO-ADVANCES through ``L`` stem
states (a one-way corridor, as in :class:`mrl_trace.maze.TMaze`, so the only
learned decision is the arm choice -- this keeps the reward distal and removes the
stem-advance reinforcement artefact). Input lines carry the cue on dedicated cue lines
AND the current stem position on dedicated one-hot position lines, so each step presents
a genuinely different state (a real trajectory, not a longer cue). The eligibility gates
of BOTH layers integrate continuously across all ``L`` steps -- never reset between
steps -- so the trace from early steps has decayed more than from late steps when the
single goal reward lands ``D`` seconds after the junction action. The reward gates the
surviving trace into both layers via the same three-factor / DFA update as
:func:`train_deep`.

All conventions follow the package: ``B`` seeds as a vectorised batch, device gate via
:class:`GateBankBatched`, signed leak-dominant coincidence, ``dw = eta * L * e`` with
``L`` the (output scalar / DFA hidden) learning signal, ``dt = 5e-3`` s, coarsened
undriven relaxation for the delay.
"""
from __future__ import annotations

import numpy as np

from .bandit import GateBankBatched, W_MAX
from .neurons import lif_step_batched, TAU_M, V_TH
from .learning import LTD_BIAS
from .maze import reward_rate

__all__ = ["train_deep_sequential", "reward_rate"]


def _relax_gate(bank, n_relax, stride, rem):
    """Advance an undriven :class:`GateBankBatched` for ``n_relax*stride + rem`` ticks,
    using a coarsened step (``stride*dt``) for the bulk and single ``dt`` steps for the
    remainder. Identical to :func:`mrl_trace.deep._relax_gate`: with zero drive
    the dynamics are smooth, so the coarse step is accurate while ``stride*dt/tau << 1``.
    Restores the gate's ``dt``."""
    B, S, A = bank.B, bank.S, bank.A
    z = np.zeros((B, S, A))
    if n_relax > 0:
        dt0 = bank.dt
        bank.dt = dt0 * stride
        try:
            for _ in range(n_relax):
                bank.step(z)
        finally:
            bank.dt = dt0
    for _ in range(rem):
        bank.step(z)


def _cue(rng, B):
    """Draw ``B`` XOR cues. Two binary features ``(b0, b1)``; correct arm = ``b0 ^ b1``.
    The features are one-hot-per-feature encoded onto ``Fc = 4`` cue lines
    ``[b0=0, b0=1, b1=0, b1=1]`` so exactly two cue lines are active every episode (a
    constant cue drive across all four states; no input-rate confound). XOR is not
    linearly separable in these lines, so a single trained layer cannot solve it -- this
    is what forces the hidden layer. Returns ``(cue_lines (B,4), correct (B,))``."""
    b0 = rng.integers(2, size=B)
    b1 = rng.integers(2, size=B)
    lines = np.zeros((B, 4))
    bidx = np.arange(B)
    lines[bidx, 0 + b0] = 1.0
    lines[bidx, 2 + b1] = 1.0
    correct = (b0 ^ b1).astype(int)
    return lines, correct


def train_deep_sequential(*, mode="dfa", B=20, H=16, L=4, A=2, tau_leak=20.0, D=3.0,
                          trials=2000, dt=5e-3, step_dur=0.4, dec_dur=1.0, eta=0.5,
                          eta_hidden=0.4,
                          in_rate=200.0, ltd=LTD_BIAS, tau_m=TAU_M, v_th=V_TH, V=1.5,
                          sigma0=0.15, sigma1=0.05, fb_scale=1.0, w_scale1=0.6,
                          w_scale2=0.3, w_max=W_MAX, bias_o=0.35, homeo=1.0,
                          homeo_target=0.35, homeo_tau=200.0, weight_fault=0.0,
                          early_stop=None, seed0=0, return_weights=False):
    """Deep sequential XOR T-maze, ``B`` parallel seeds; returns rewards ``(B, trials)``.

    Architecture (extends :func:`train_deep` with a stem-position input block):
      ``F = 4 + L`` input lines -> ``H`` hidden LIF neurons -> ``A`` action LIF neurons,
      with device-synapse matrices ``W1 (F,H)`` and ``W2 (H,A)``. The first 4 lines carry
      the XOR cue (persistent context, active on every step so the junction can decide on
      it); the next ``L`` lines are a one-hot of the current stem position (so each step
      is a genuinely different state). Each weight matrix has its own
      :class:`GateBankBatched` eligibility (retention ``tau_leak``, signed leak-dominant
      drive) that integrates CONTINUOUSLY across all ``L`` stem steps -- never reset
      between steps.

    Trajectory. The agent auto-advances through ``L`` stem states (a one-way corridor:
    the action does not move it, exactly as :class:`mrl_trace.maze.TMaze`), so the
    only learned decision is the arm chosen at the junction (the last state). Reward
    ``R in {0,1}``, contingent on the junction action matching ``XOR(b0, b1)``, is
    delivered after the action->reward delay ``D``. The surviving eligibility then gates
    into both layers.

    ``mode`` selects the spatial-credit pathway (eligibility is the device trace in all):
      ``shallow``  one trained layer F->A (should FAIL XOR -- depth-necessity control);
      ``elm``      hidden layer fixed random, only the output trained;
      ``global``   both layers trained by the pure global scalar (R-b) (structural-credit
                   failure mode);
      ``dfa``      both layers trained by Direct Feedback Alignment (all-local; the
                   device arm of the demonstration);
      ``no_trace`` deep DFA with the eligibility zeroed (device-necessity control).

    ``homeo > 0`` enables local homeostatic regulation of hidden-unit firing.
    ``weight_fault`` randomly stuck-at-zeroes that fraction of W1/W2 synapses (a device
    yield stressor; 0 disables). ``early_stop`` (a reward-rate threshold) stops a seed's
    updates once its trailing reward rate exceeds it (the policy is then run frozen);
    None disables. The signature is batch-sweep compatible.
    """
    if mode not in {"shallow", "elm", "global", "dfa", "no_trace"}:
        raise ValueError(f"unknown mode {mode!r}")
    rng = np.random.default_rng(seed0)
    Fc = 4
    F = Fc + L
    eta_h = eta if eta_hidden is None else eta_hidden
    deep = mode != "shallow"
    no_trace = mode == "no_trace"
    train_w1 = deep and mode != "elm"

    def _init(shape, scale):
        return np.clip(scale * rng.standard_normal(shape), -w_max, w_max)

    if deep:
        W1 = _init((B, F, H), w_scale1)
        W2 = _init((B, H, A), w_scale2)
    else:
        W2 = _init((B, F, A), w_scale1)
        W1 = None

    # --- weight-fault mask: stuck-at-zero synapses (device yield stressor) ---
    m1 = m2 = None
    if weight_fault > 0:
        if deep:
            m1 = (rng.random((B, F, H)) >= weight_fault).astype(float)
            W1 *= m1
        m2 = (rng.random(W2.shape) >= weight_fault).astype(float)
        W2 *= m2

    # --- device eligibility gates (one bank per trained matrix) ---
    g1 = GateBankBatched(B, F, H, tau_leak=tau_leak, dt=dt, V=V) if deep else None
    g_out = GateBankBatched(B, (H if deep else F), A, tau_leak=tau_leak, dt=dt, V=V)

    # --- DFA feedback matrix: fixed random, drawn ONCE, independent of W2 (no transport) ---
    B_fix = fb_scale * rng.standard_normal((A, H)) if deep else None

    baseline = np.full(B, 1.0 / A)
    bidx = np.arange(B)
    steps_per_state = max(1, int(round(step_dur / dt)))
    dec_dwell = max(1, int(round(dec_dur / dt)))     # longer dwell at the junction
    reward_lag = int(round(D / dt))
    rewards = np.zeros((B, trials))

    # Local homeostatic activity estimate per hidden neuron (see train_deep).
    act_hidden = np.full((B, H), homeo_target) if (deep and homeo > 0) else None
    frozen = np.zeros(B, dtype=bool)        # early-stop: seeds whose updates are frozen

    for tr in range(trials):
        sigma = sigma0 + (sigma1 - sigma0) * tr / trials
        if deep:
            g1.reset()
        g_out.reset()
        cue_lines, correct = _cue(rng, B)
        vh = np.zeros((B, H))
        vo = np.zeros((B, A))
        # junction-decision accumulators (the LAST stem step is the decision point)
        spk_o_dec = np.zeros((B, A))
        spk_h_dec = np.zeros((B, H))

        # --- traverse the L-step stem trajectory ---
        # The cue lines are presented on every step (persistent context); the position
        # line for the current state is active too, so each step is a distinct state. The
        # eligibility gates integrate across ALL steps (never reset). Only the LAST step
        # is the junction where the arm choice is read out (auto-advance corridor).
        for st in range(L):
            decision_step = (st == L - 1)
            lines = np.zeros((B, F))
            lines[:, :Fc] = cue_lines
            lines[:, Fc + st] = 1.0                      # one-hot stem position
            # The agent dwells LONGER at the junction (the decision point) than at the
            # auto-advanced stem states: more spike accumulation -> a cleaner argmax arm
            # choice and a stronger (less noisy) output-layer coincidence to credit. The
            # extra dwell also lengthens the cue->reward gap, which the device must bridge.
            dwell = dec_dwell if decision_step else steps_per_state
            for _ in range(dwell):
                pre_in = (rng.random((B, F)) < (in_rate * dt) * lines).astype(float)
                if deep:
                    ch_h = np.einsum('bfh,bf->bh', W1, pre_in)
                    vh, sp_h = lif_step_batched(vh, ch_h, dt, rng, tau_m=tau_m,
                                                v_th=v_th, noise=sigma)
                    ch_o = np.einsum('bha,bh->ba', W2, sp_h.astype(float)) + bias_o
                    vo, sp_o = lif_step_batched(vo, ch_o, dt, rng, tau_m=tau_m,
                                                v_th=v_th, noise=sigma)
                    if decision_step:
                        spk_h_dec += sp_h
                        spk_o_dec += sp_o
                    # Signed leak-dominant coincidence. The INPUT->HIDDEN bank g1 is
                    # driven on EVERY step: its cue+position rows differ per step, so it
                    # accumulates the trajectory eligibility (early steps decay more than
                    # late ones by reward time -- the bridging the device must support).
                    # The HIDDEN->OUTPUT bank g_out shares the same (H,A) synapses across
                    # steps, so a stem-step output coincidence (made by an as-yet random
                    # arm decision) would POLLUTE the junction credit; the arm decision --
                    # and the only credit-relevant output coincidence -- happens at the
                    # junction, so g_out is driven only on the decision step (the maze
                    # analogue: each step's coincidence lands on its own state row, the
                    # junction decision on its own row).
                    if no_trace:
                        g1.step(np.zeros((B, F, H)))
                        g_out.step(np.zeros((B, H, A)))
                    else:
                        drive1 = pre_in[:, :, None] * np.where(sp_h, 1.0, -ltd)[:, None, :]
                        g1.step(drive1)
                        if decision_step:
                            drive_o = (sp_h.astype(float)[:, :, None]
                                       * np.where(sp_o, 1.0, -ltd)[:, None, :])
                            g_out.step(drive_o)
                        else:
                            g_out.step(np.zeros((B, H, A)))
                else:
                    ch_o = np.einsum('bfa,bf->ba', W2, pre_in)
                    vo, sp_o = lif_step_batched(vo, ch_o, dt, rng, tau_m=tau_m,
                                                v_th=v_th, noise=sigma)
                    if decision_step:
                        spk_o_dec += sp_o
                    drive_o = pre_in[:, :, None] * np.where(sp_o, 1.0, -ltd)[:, None, :]
                    g_out.step(drive_o)

        # --- relax for the action->reward delay, then snapshot eligibility ---
        if not no_trace:
            stride = 10 if reward_lag > 200 else 1
            n_relax, rem = divmod(reward_lag, stride)
            _relax_gate(g_out, n_relax, stride, rem)
            if deep:
                _relax_gate(g1, n_relax, stride, rem)
            eo_rew = g_out.vn[..., -1] / g_out.Vnmax
            e1_rew = (g1.vn[..., -1] / g1.Vnmax) if deep else None
        else:
            eo_rew = np.zeros_like(g_out.vn[..., -1])
            e1_rew = np.zeros((B, F, H)) if deep else None

        # --- arm selection at the junction (argmax decision-step spikes; ties -> rand) ---
        tie = spk_o_dec.max(1) == spk_o_dec.min(1)
        chosen = np.argmax(spk_o_dec, 1)
        chosen[tie] = rng.integers(A, size=int(tie.sum()))
        R = (chosen == correct).astype(float)

        # --- learning signals (identical math/signs to train_deep) ---
        adv = (R - baseline)
        Lo = adv[:, None, None]                                   # output: global scalar
        if deep:
            if mode == "global":
                Lh = adv[:, None, None]
            else:
                logits = spk_o_dec - spk_o_dec.mean(1, keepdims=True)
                pol = np.exp(logits - logits.max(1, keepdims=True))
                pol /= pol.sum(1, keepdims=True)
                onehot = np.zeros((B, A)); onehot[bidx, chosen] = 1.0
                L_a = adv[:, None] * (onehot - pol)               # (B,A)
                Lh = np.einsum('ah,ba->bh', B_fix, L_a)[:, None, :]   # (B,1,H)

        upd = (~frozen).astype(float)[:, None, None]              # freeze early-stopped seeds
        W2 = np.clip(W2 + upd * eta * Lo * eo_rew, -w_max, w_max)
        if train_w1:
            dW1 = eta_h * Lh * e1_rew
            if homeo > 0 and act_hidden is not None:
                rate = spk_h_dec / max(1, dec_dwell)              # decision-step firing rate
                act_hidden += (rate - act_hidden) / homeo_tau
                scale = 1.0 + homeo * (homeo_target - act_hidden)  # (B,H)
                dW1 = dW1 + (scale[:, None, :] - 1.0) * W1
            W1 = np.clip(W1 + upd * dW1, -w_max, w_max)
        if m1 is not None:
            W1 *= m1
        if m2 is not None:
            W2 *= m2

        baseline += 0.02 * (R - baseline)
        rewards[:, tr] = R

        if early_stop is not None and tr >= 150:
            rr = rewards[:, max(0, tr - 149):tr + 1].mean(axis=1)
            frozen |= (rr >= early_stop)

    if return_weights:
        return rewards, (W1, W2)
    return rewards
\end{lstlisting}

\begin{lstlisting}[caption={siox\_eligibility/distal\_cue.py},label={lst:siox-src-distal_cue}]
"""Distal-cue single-deposit task -- retention as the credit-assignment asset.

This module replaces the two temporally-weak experiments on the package's temporal
axis (the shallow T-maze of :func:`mrl_trace.maze.train_sequential` and the deep
T-maze of :func:`mrl_trace.deep_sequential.train_deep_sequential`) with a task
that is *designed to stress the eligibility trace* so that long device retention is a
measurable asset rather than a decoration.

WHY THE OLD TASKS WERE WEAK. In both T-mazes the informative cue is re-presented on
*every* stem step (it is a persistent context line), so the cue->action eligibility is
effectively *refreshed* the whole way down the corridor. A short-retention device then
bridges any gap (verified: TMaze device tau=1.5 s scores ~0.99 even at a 12 s delay),
because the trace it must carry is topped up continuously and never has to *survive* on
its own. The flat device-vs-length curve therefore reads as "the task never stressed the
trace", not "long retention bridges a gap a short trace cannot".

THE FIX -- A SINGLE DEPOSIT. Here the informative cue drives the credited synapses'
eligibility coincidence in ONE early window only. After that window the cue is never
re-presented and the credited eligibility is never driven again: it can only DECAY
(relax) while the agent traverses the rest of the trajectory / waits out the delay. A
single delayed reward then gates whatever eligibility has SURVIVED on those synapses.
As the trajectory-time / delay gap ``T_gap`` grows past the device retention
``tau_leak``, the surviving deposit on a SHORT-tau device collapses toward zero (no
credit, chance behaviour) while a LONG-tau device still carries it (credit assigned,
above-chance behaviour). That crossover -- device-LONG beating device-SHORT only once
the gap exceeds the short retention -- is the sequential form of the ``D_max ~ k*tau``
law: the learnable credit horizon scales with the device retention, a fabrication /
defect-engineering knob.

THE CONTROL THAT MUST FAIL IS A SHORT-RETENTION DEVICE, NOT AN ABSTRACT TRACE. A linear
exponential eligibility decaying over the gap rescales *every* credited synapse by the
SAME factor, so it preserves the relative credit pattern and still learns the policy
(abstract-short stays high even at a large gap). Only the device-long vs device-SHORT
contrast isolates retention as the asset, because the trap-cascade gate's *nonlinear*
rise+leak means a short ``tau_leak`` deposit decays below the noise floor and loses the
*pattern*, not merely its scale. ``abstract`` is reported here only as a baseline; the
headline contrast is device-long vs device-short (the intra-device STC fabrication-knob
story).

Two variants share the design:

- :func:`train_distal_shallow` -- the SHALLOW replacement for the shallow T-maze: a single
  trained ``context x action`` device-synapse layer; cue deposited once, trajectory/delay,
  then a single gated reward.
- :func:`train_distal_deep` -- the DEEP strengthening of
  :func:`mrl_trace.deep_sequential.train_deep_sequential`: it reuses that module's
  hidden-layer + DFA + homeostasis machinery but makes the XOR cue a SINGLE early deposit
  (the cue lines are silent after the deposit window) so the gap genuinely stresses the
  trace.

Conventions follow the package: pure NumPy, ``B`` seeds as a vectorised batch (leading
axis ``B``), device gate via :class:`mrl_trace.bandit.GateBankBatched` (``tau_leak``
the retention), signed leak-dominant coincidence, three-factor update
``dw = eta (R - b) e``, ``dt = 5e-3`` s, coarsened undriven relaxation for the gap.
"""
from __future__ import annotations

import numpy as np

from .bandit import GateBankBatched, AbstractTrace, W_INIT, W_MAX
from .neurons import lif_step_batched, TAU_M, V_TH
from .learning import LTD_BIAS
from .maze import reward_rate, trials_to_criterion
from . import deep_sequential

__all__ = ["train_distal_shallow", "train_distal_deep", "reward_rate",
           "trials_to_criterion", "run_distal_cue", "main"]


# ----------------------------------------------------------------------------
# Shared: coarsened undriven relaxation of a GateBankBatched / AbstractTrace
# ----------------------------------------------------------------------------
def _relax(bank, reward_lag, S, A):
    """Advance an undriven eligibility bank for ``reward_lag`` dt-ticks.

    With zero drive the cascade / leaky dynamics are smooth, so the bulk is integrated
    with a coarsened step (``stride*dt``) for tractability on long gaps and the
    remainder with single ``dt`` steps (the same scheme as
    :func:`mrl_trace.maze._relax`). ``bank`` may be a :class:`GateBankBatched`
    or an :class:`AbstractTrace`; both expose ``step`` and ``dt``. Restores ``dt``.
    """
    if reward_lag <= 0:
        return
    z = np.zeros((bank.B if hasattr(bank, "B") else bank.e.shape[0], S, A))
    stride = 10 if reward_lag > 200 else 1
    n_relax, rem = divmod(reward_lag, stride)
    if n_relax > 0:
        dt0 = bank.dt
        bank.dt = dt0 * stride
        try:
            for _ in range(n_relax):
                bank.step(z)
        finally:
            bank.dt = dt0
    for _ in range(rem):
        bank.step(z)


def _read_elig(bank, abstract):
    """Read the per-synapse eligibility snapshot (device cascade tip, or abstract e)."""
    if abstract:
        return np.clip(bank.e, -1.0, 1.0)
    return bank.vn[..., -1] / bank.Vnmax


# ----------------------------------------------------------------------------
# SHALLOW variant -- replacement for the shallow T-maze
# ----------------------------------------------------------------------------
def train_distal_shallow(*, B=20, tau_leak=20.0, T_gap=4.0, C=4, A=2, trials=1500,
                         dt=5e-3, cue_dur=1.0, dec_dur=1.0, eta=0.2, V=0.9,
                         in_rate=200.0, ltd=LTD_BIAS, tau_m=TAU_M, v_th=V_TH,
                         sigma0=0.15, sigma1=None, abstract=False, no_trace=False,
                         weight_fault=None, seed0=0, return_weights=False):
    """Shallow distal-cue single-deposit task; ``B`` seeds; returns rewards ``(B, trials)``.

    Architecture: ``C`` context lines -> ``A`` action LIF neurons through a single trained
    device-synapse grid ``w (C, A)``. Rewarded mapping: context ``c`` -> action ``c % A``.

    Each trial is a trajectory with three phases:

      1. DEPOSIT (``cue_dur`` s). The context line ``c`` for this trial fires; the action
         neurons spike; the signed leak-dominant pre x post coincidence drives the
         eligibility gate ONLY on the active context row. This is the single, literal
         deposit -- the only informative drive onto the credited synapses.
      2. GAP (``T_gap`` s). The cue is REMOVED and the credited eligibility is never driven
         again -- it only relaxes. (Implemented as a coarsened undriven relaxation, the
         trajectory-time / delay the trace must survive.)
      3. DECISION (``dec_dur`` s). The context line is re-presented PURELY to read out an
         action (Kirchhoff bitline sum -> spiking argmax); the eligibility is NOT driven
         here, so the credit gated by the reward comes only from the SURVIVING early
         deposit, not from any decision-time refresh.

    A single reward ``R in {0,1}`` (action == ``c % A``) then gates the surviving deposit:
    ``dw = eta (R - b) e``. As ``T_gap`` grows past ``tau_leak`` a short-tau device's
    deposit decays away (chance) while a long-tau device's survives (above chance).

    ``abstract`` swaps the device gate for the exponential trace (baseline); ``no_trace``
    zeroes the eligibility (device-necessity control). Signature is batch-sweep
    compatible (B, tau_leak, T_gap, trials, seed0, plus the sweep params).
    """
    rng = np.random.default_rng(seed0)
    bank = (AbstractTrace(B, C, A, tau_elig=tau_leak, dt=dt) if abstract
            else GateBankBatched(B, C, A, tau_leak=tau_leak, dt=dt, V=V))
    w = np.full((B, C, A), W_INIT)
    correct = np.array([c % A for c in range(C)])
    baseline = np.full(B, 1.0 / A)
    bidx = np.arange(B)
    dep_ticks = max(1, int(round(cue_dur / dt)))
    dec_ticks = max(1, int(round(dec_dur / dt)))
    gap_lag = int(round(T_gap / dt))
    rewards = np.zeros((B, trials))

    for tr in range(trials):
        sigma = sigma0 if sigma1 is None else sigma0 + (sigma1 - sigma0) * tr / trials
        bank.reset()
        context = rng.integers(C, size=B)
        v = np.zeros((B, A))
        # Read-time device faults applied ONCE per trial (weights change only per trial):
        # the array reads the faulted conductance wr while learning targets the clean w --
        # a fixed physical realisation, not a learnable parameter (as in deep.py/maze.py).
        # The maze weights are non-negative, so weight_fault should be a maze_fault_stack.
        wr = weight_fault(w) if weight_fault is not None else w

        # --- 1. DEPOSIT: cue fires, single signed coincidence onto context row ---
        for _ in range(dep_ticks):
            pre = np.zeros((B, C))
            pre[bidx, context] = (rng.random(B) < in_rate * dt).astype(float)
            charge = np.einsum('bca,bc->ba', wr, pre)
            v, sp = lif_step_batched(v, charge, dt, rng, tau_m=tau_m, v_th=v_th,
                                     noise=sigma)
            drive = np.zeros((B, C, A))
            drive[bidx, context, :] = pre[bidx, context][:, None] * np.where(sp, 1.0, -ltd)
            if no_trace:
                drive[:] = 0.0
            bank.step(drive)

        # --- 2. GAP: cue removed; credited eligibility only relaxes (single deposit) ---
        if not no_trace:
            _relax(bank, gap_lag, C, A)

        # snapshot the SURVIVING deposit -- this is what the reward will gate
        e_rew = np.zeros((B, C, A)) if no_trace else _read_elig(bank, abstract)

        # --- 3. DECISION: re-present cue to READ OUT an action (no eligibility drive) ---
        v[:] = 0.0
        spk = np.zeros((B, A))
        for _ in range(dec_ticks):
            pre = np.zeros((B, C))
            pre[bidx, context] = (rng.random(B) < in_rate * dt).astype(float)
            charge = np.einsum('bca,bc->ba', wr, pre)
            v, sp = lif_step_batched(v, charge, dt, rng, tau_m=tau_m, v_th=v_th,
                                     noise=sigma)
            spk += sp
            # NOTE: bank is NOT stepped here -- the decision must not refresh the deposit.

        tie = spk.max(1) == spk.min(1)
        chosen = np.argmax(spk, 1)
        chosen[tie] = rng.integers(A, size=int(tie.sum()))
        R = (chosen == correct[context]).astype(float)

        adv = eta * (R - baseline)
        w = np.clip(w + adv[:, None, None] * e_rew, 0.0, W_MAX)
        baseline += 0.02 * (R - baseline)
        rewards[:, tr] = R

    if return_weights:
        return rewards, w
    return rewards


# ----------------------------------------------------------------------------
# DEEP variant -- strengthening of deep_sequential.train_deep_sequential
# ----------------------------------------------------------------------------
def train_distal_deep(*, mode="dfa", B=20, H=16, L=4, A=2, tau_leak=20.0, T_gap=4.0,
                      trials=2000, dt=5e-3, step_dur=0.4, dec_dur=1.0, eta=0.5,
                      eta_hidden=0.4, in_rate=200.0, ltd=LTD_BIAS, tau_m=TAU_M,
                      v_th=V_TH, V=1.5, sigma0=0.15, sigma1=0.05, fb_scale=1.0,
                      w_scale1=0.6, w_scale2=0.3, w_max=W_MAX, bias_o=0.35, homeo=1.0,
                      homeo_target=0.35, homeo_tau=200.0, weight_fault=None, seed0=0,
                      return_weights=False):
    """Deep distal-cue single-deposit XOR task; ``B`` seeds; returns rewards ``(B, trials)``.

    This reuses the hidden-layer + DFA + homeostasis machinery of
    :func:`mrl_trace.deep_sequential.train_deep_sequential` (``F = 4 + L`` input
    lines -> ``H`` hidden LIF (DFA + homeostasis) -> ``A`` action LIF; device gates
    ``g1, g_out``; junction-only output coincidence) but fixes the temporal weakness:
    the XOR CUE IS A SINGLE EARLY DEPOSIT.

    Difference from ``train_deep_sequential``: the 4 cue lines are active ONLY during the
    first stem step (the deposit). For the remaining ``L-1`` stem steps and the junction
    the cue lines are SILENT -- only the one-hot stem-position line for the current step
    fires. The input->hidden eligibility ``g1`` therefore receives its informative
    cue->hidden coincidence only at the deposit; thereafter the cue rows of ``g1`` only
    decay across the trajectory and the ``T_gap`` action->reward delay. The surviving cue
    deposit -- not a re-presentation at the junction -- must carry the XOR association into
    the credited synapses. The position lines still make each step a distinct state (a real
    trajectory), and the junction-step output coincidence still credits the arm choice, so
    the deep all-local credit path is unchanged; only the cue's *refresh* is removed.

    As ``T_gap`` grows past ``tau_leak`` the surviving cue deposit on a short-tau device
    collapses (the XOR association is lost; chance) while a long-tau device retains it
    (above chance). ``mode`` (``shallow``/``elm``/``global``/``dfa``/``no_trace``),
    ``homeo``, and all other kwargs mirror ``train_deep_sequential``; the signature is
    batch-sweep compatible (B, H, tau_leak, T_gap, trials, seed0, mode, homeo, ...).
    """
    if mode not in {"shallow", "elm", "global", "dfa", "no_trace"}:
        raise ValueError(f"unknown mode {mode!r}")
    rng = np.random.default_rng(seed0)
    _cue = deep_sequential._cue
    _relax_gate = deep_sequential._relax_gate
    Fc = 4
    F = Fc + L
    eta_h = eta if eta_hidden is None else eta_hidden
    deep = mode != "shallow"
    no_trace = mode == "no_trace"
    train_w1 = deep and mode != "elm"

    def _init(shape, scale):
        return np.clip(scale * rng.standard_normal(shape), -w_max, w_max)

    if deep:
        W1 = _init((B, F, H), w_scale1)
        W2 = _init((B, H, A), w_scale2)
    else:
        W2 = _init((B, F, A), w_scale1)
        W1 = None

    g1 = GateBankBatched(B, F, H, tau_leak=tau_leak, dt=dt, V=V) if deep else None
    g_out = GateBankBatched(B, (H if deep else F), A, tau_leak=tau_leak, dt=dt, V=V)
    B_fix = fb_scale * rng.standard_normal((A, H)) if deep else None

    baseline = np.full(B, 1.0 / A)
    bidx = np.arange(B)
    steps_per_state = max(1, int(round(step_dur / dt)))
    dec_dwell = max(1, int(round(dec_dur / dt)))
    gap_lag = int(round(T_gap / dt))
    rewards = np.zeros((B, trials))

    act_hidden = np.full((B, H), homeo_target) if (deep and homeo > 0) else None

    for tr in range(trials):
        sigma = sigma0 + (sigma1 - sigma0) * tr / trials
        if deep:
            g1.reset()
        g_out.reset()
        cue_lines, correct = _cue(rng, B)
        # Read-time device faults applied ONCE per trial: the array reads faulted weights
        # (W1r/W2r) while learning targets the clean W1/W2 (as in deep.py). Signed weights,
        # so weight_fault should be a siox_fault_stack.
        W1r = weight_fault(W1) if (weight_fault is not None and deep) else W1
        W2r = weight_fault(W2) if weight_fault is not None else W2
        vh = np.zeros((B, H))
        vo = np.zeros((B, A))
        spk_o_dec = np.zeros((B, A))
        spk_h_dec = np.zeros((B, H))

        # --- traverse the L-step stem; cue ELIGIBILITY is deposited ONLY on step 0 ---
        # Cue->hidden eligibility row mask: 1 during the deposit step, 0 afterwards. The
        # cue lines stay present in the FORWARD pass on every step (so the junction can
        # READ OUT the XOR -- the decision must be able to depend on the cue), but the
        # eligibility coincidence onto the cue rows of g1 is driven ONLY at the deposit.
        # The credit the reward later gates onto the cue->hidden synapses therefore comes
        # solely from the SURVIVING early deposit, never a junction-time refresh. The
        # position rows drive g1 every step (a genuine per-step trajectory state).
        for st in range(L):
            decision_step = (st == L - 1)
            deposit_step = (st == 0)
            lines = np.zeros((B, F))
            lines[:, :Fc] = cue_lines            # cue present every step (forward readout)
            lines[:, Fc + st] = 1.0              # one-hot stem position (every step)
            elig_row = np.ones(F)                # gate cue ROWS of the g1 drive
            if not deposit_step:
                elig_row[:Fc] = 0.0              # single deposit: cue eligibility early only
            dwell = dec_dwell if decision_step else steps_per_state
            for _ in range(dwell):
                pre_in = (rng.random((B, F)) < (in_rate * dt) * lines).astype(float)
                if deep:
                    ch_h = np.einsum('bfh,bf->bh', W1r, pre_in)
                    vh, sp_h = lif_step_batched(vh, ch_h, dt, rng, tau_m=tau_m,
                                                v_th=v_th, noise=sigma)
                    ch_o = np.einsum('bha,bh->ba', W2r, sp_h.astype(float)) + bias_o
                    vo, sp_o = lif_step_batched(vo, ch_o, dt, rng, tau_m=tau_m,
                                                v_th=v_th, noise=sigma)
                    if decision_step:
                        spk_h_dec += sp_h
                        spk_o_dec += sp_o
                    # g1: position rows driven every step; cue rows driven only at deposit
                    # (elig_row masks them off afterwards) -> a single cue->hidden deposit
                    # that then only decays. g_out is driven only at the arm decision.
                    if no_trace:
                        g1.step(np.zeros((B, F, H)))
                        g_out.step(np.zeros((B, H, A)))
                    else:
                        drive1 = ((pre_in * elig_row)[:, :, None]
                                  * np.where(sp_h, 1.0, -ltd)[:, None, :])
                        g1.step(drive1)
                        if decision_step:
                            drive_o = (sp_h.astype(float)[:, :, None]
                                       * np.where(sp_o, 1.0, -ltd)[:, None, :])
                            g_out.step(drive_o)
                        else:
                            g_out.step(np.zeros((B, H, A)))
                else:
                    ch_o = np.einsum('bfa,bf->ba', W2r, pre_in)
                    vo, sp_o = lif_step_batched(vo, ch_o, dt, rng, tau_m=tau_m,
                                                v_th=v_th, noise=sigma)
                    if decision_step:
                        spk_o_dec += sp_o
                    drive_o = pre_in[:, :, None] * np.where(sp_o, 1.0, -ltd)[:, None, :]
                    g_out.step(drive_o)

        # --- relax for the action->reward delay (the gap), then snapshot eligibility ---
        if not no_trace:
            stride = 10 if gap_lag > 200 else 1
            n_relax, rem = divmod(gap_lag, stride)
            _relax_gate(g_out, n_relax, stride, rem)
            if deep:
                _relax_gate(g1, n_relax, stride, rem)
            eo_rew = g_out.vn[..., -1] / g_out.Vnmax
            e1_rew = (g1.vn[..., -1] / g1.Vnmax) if deep else None
        else:
            eo_rew = np.zeros_like(g_out.vn[..., -1])
            e1_rew = np.zeros((B, F, H)) if deep else None

        tie = spk_o_dec.max(1) == spk_o_dec.min(1)
        chosen = np.argmax(spk_o_dec, 1)
        chosen[tie] = rng.integers(A, size=int(tie.sum()))
        R = (chosen == correct).astype(float)

        adv = (R - baseline)
        Lo = adv[:, None, None]
        if deep:
            if mode == "global":
                Lh = adv[:, None, None]
            else:
                logits = spk_o_dec - spk_o_dec.mean(1, keepdims=True)
                pol = np.exp(logits - logits.max(1, keepdims=True))
                pol /= pol.sum(1, keepdims=True)
                onehot = np.zeros((B, A)); onehot[bidx, chosen] = 1.0
                L_a = adv[:, None] * (onehot - pol)
                Lh = np.einsum('ah,ba->bh', B_fix, L_a)[:, None, :]

        W2 = np.clip(W2 + eta * Lo * eo_rew, -w_max, w_max)
        if train_w1:
            dW1 = eta_h * Lh * e1_rew
            if homeo > 0 and act_hidden is not None:
                rate = spk_h_dec / max(1, dec_dwell)
                act_hidden += (rate - act_hidden) / homeo_tau
                scale = 1.0 + homeo * (homeo_target - act_hidden)
                dW1 = dW1 + (scale[:, None, :] - 1.0) * W1
            W1 = np.clip(W1 + dW1, -w_max, w_max)

        baseline += 0.02 * (R - baseline)
        rewards[:, tr] = R

    if return_weights:
        return rewards, (W1, W2)
    return rewards


# =============================================================================
# Experiment core + full-scale driver -- the headline TEMPORAL result.
#
# Absorbed from experiments/sweeps/entry_distal_cue.py (the strengthened exp15
# shallow + exp17 deep temporal experiments).  ``run_distal_cue`` is the SERIAL,
# I/O-free core a notebook calls in quick mode; ``main()`` is the full-scale
# python -m driver that fans the (gap x arm x seed x p) grid over a Pool, bootstraps
# CIs, and writes the grid via ``paths.save_result`` (preserving the original
# ``distal_cue_<variant>.npy`` filename).
#
# THE SCIENCE (preserved verbatim from the driver's provenance). The distal-cue
# task deposits the informative cue into the eligibility trace ONCE, early, then
# never refreshes it: the agent traverses a trajectory / waits out a gap ``T_gap``
# and a single delayed reward gates whatever eligibility SURVIVED. As ``T_gap`` grows
# past the device retention ``tau_leak``, a SHORT-retention device's deposit decays
# away (chance) while a LONG-retention device still carries it (above chance). That
# crossover -- device-LONG beating device-SHORT only once the gap exceeds the short
# retention, while TIED at small gaps -- is the sequential form of the ``D_max ~
# k*tau`` law: the learnable credit horizon scales with retention, a fabrication /
# defect-engineering knob.
#
# CONTROL: the arm that must fail is a SHORT-tau DEVICE (same GateBankBatched
# physics, small tau_leak), NOT the abstract exponential trace -- a linear trace
# rescales every synapse equally and still learns. device-long vs device-short
# isolates retention as the asset.
#
# FAULTS (optional ``p`` > 0): the SiO_x programmed-conductance fault prior (stuck-at
# + realistic asymmetric D2D) is applied at read time -- maze_fault_stack for the
# non-negative shallow weights, siox_fault_stack for the signed deep weights. The
# Poole-Frenkel I-V term is excluded (a read/inference nonideality, out of scope for
# a programmed-weight rule). p=0 is the clean crossover; p>0 asks whether the
# retention asset survives device faults.
# =============================================================================

# Bootstrap resamples and the device-to-device (D2D) spreads for the SiO_x fault
# prior, plus the two retentions the crossover contrasts -- all preserved exactly
# from the original driver.
_DISTAL_N_BOOT = 10000
_DISTAL_SIGMA_G = 0.5        # D2D R_off-end log-std
_DISTAL_SIGMA_G_ON = 0.05    # D2D R_on-end log-std (realistic asymmetric spread)
_DISTAL_TAU_LONG = 20.0
_DISTAL_TAU_SHORT = 1.5

# Full-scale driver config, mutated inside ``main()`` before the fork Pool spawns
# (fork inherits it, matching the original entry point's module-global ``_CFG``).
# ``run_distal_cue`` does NOT read this -- it takes every parameter explicitly so a
# notebook can call it standalone; only the Pool worker consults it.
_DISTAL_CFG = {"trials": 2000, "variant": "deep", "window": 150}


def _distal_final(r, window):
    """Mean-over-seeds final reward rate of a ``(B, trials)`` reward array."""
    return float(np.mean(reward_rate(r, window=window)))


def _distal_make_fault(p, variant, seed):
    """SiO_x read-time programmed-conductance fault stack for a run, or ``None`` at p=0.

    ``maze_fault_stack`` for the non-negative shallow single-conductance weights,
    ``siox_fault_stack`` for the signed deep weights; the Poole-Frenkel I-V term is
    excluded (a read/inference nonideality, out of scope for a programmed-weight rule).
    """
    from .device_faults import siox_fault_stack, maze_fault_stack
    if p <= 0:
        return None
    if variant == "shallow":   # non-negative single-conductance weights
        return maze_fault_stack(p_stuck=p, sigma_g=_DISTAL_SIGMA_G, pf_on=False,
                                stuck_kind="off", w_max=W_MAX, seed=seed)
    return siox_fault_stack(p_stuck=p, sigma_g=_DISTAL_SIGMA_G,
                            sigma_g_on=_DISTAL_SIGMA_G_ON, pf_on=False,
                            stuck_kind="off", seed=seed)


def _distal_one_run(job):
    """Single (gap, p, arm, seed) run; returns ``(gap, p, arm, seed, final_reward)``.

    A module-level function (not a closure) so it is picklable for the fork Pool in
    ``main()``. Reads ``_DISTAL_CFG`` for the variant/trials/window the driver set.
    """
    gap, p, arm, seed = job
    variant = _DISTAL_CFG["variant"]
    trials = _DISTAL_CFG["trials"]
    window = _DISTAL_CFG["window"]
    tau = _DISTAL_TAU_LONG if arm in ("long", "abstract") else _DISTAL_TAU_SHORT
    abstract = (arm == "abstract")
    no_trace = (arm == "no_trace")
    fault = _distal_make_fault(p, variant,
                               seed=1000 * seed + int(100 * p) + int(10 * gap))
    if variant == "shallow":
        r = train_distal_shallow(B=1, tau_leak=tau, T_gap=gap, trials=trials,
                                 seed0=seed, abstract=abstract, no_trace=no_trace,
                                 weight_fault=fault)
    else:
        mode = "no_trace" if no_trace else "dfa"
        # the deep variant has no abstract gate; skip abstract arm for deep (handled in main)
        r = train_distal_deep(mode=mode, B=1, H=16, tau_leak=tau, T_gap=gap,
                              trials=trials, seed0=seed, homeo=1.0,
                              weight_fault=fault)
    return (gap, p, arm, seed, _distal_final(r, window))


def _distal_boot_ci(vals, seed=0):
    """Percentile bootstrap 95% CI for the mean of ``vals`` (matches the driver's own
    resample count so reproduced numbers are bit-identical)."""
    vals = np.asarray(vals, float)
    rng = np.random.default_rng(seed)
    b = [vals[rng.integers(0, len(vals), len(vals))].mean() for _ in range(_DISTAL_N_BOOT)]
    return float(np.percentile(b, 2.5)), float(np.percentile(b, 97.5))


def run_distal_cue(*, variant="deep", gaps=(1.0, 4.0, 10.0), p_grid=(0.0,), seeds=12,
                   trials=2000, window=150, workers=1):
    """Distal-cue retention-as-asset grid; returns the result dict (no file I/O, no plots).

    SERIAL by default (``workers=1``) so a notebook can call it in-kernel at a small
    ``seeds``/``trials``; ``main()`` sets ``workers=os.cpu_count()`` for the published run.

    Sweeps ``gaps x arms x p_grid x seeds`` where the arms are device-LONG
    (``tau=20 s``), device-SHORT (``tau=1.5 s``) and ``no_trace`` always, plus an
    ``abstract`` (matched-tau linear-trace) baseline for the SHALLOW variant only (the
    deep variant has no abstract eligibility gate). Reports each arm's final reward rate
    (mean + bootstrap 95% CI) per (gap, p), and the device-long minus device-short gap.

    Returns the grid dict with the SAME keys the original driver saved:
    ``{"variant", "gaps", "p", "arms", "tau_long", "tau_short", "summary", "faults",
    "seeds", "trials"}`` where ``summary[(gap, p, arm)] = (mean, lo, hi)``.
    """
    gaps = [float(x) for x in gaps]
    p_grid = [float(x) for x in p_grid]
    _DISTAL_CFG.update(trials=trials, variant=variant, window=window)

    # arms: device-long, device-short, no_trace always; abstract-long only for the
    # shallow variant (the deep variant has no abstract eligibility gate).
    arms = ["long", "short", "no_trace"] + (["abstract"] if variant == "shallow" else [])
    jobs = [(g, p, arm, s) for g in gaps for p in p_grid for arm in arms
            for s in range(seeds)]

    if workers and workers > 1:
        import multiprocessing as mp
        try:
            ctx = mp.get_context("fork")
        except ValueError:                                   # pragma: no cover
            ctx = mp.get_context()
        with ctx.Pool(workers) as pool:
            results = pool.map(_distal_one_run, jobs, chunksize=1)
    else:
        results = [_distal_one_run(job) for job in jobs]

    acc = {}
    for g, p, arm, s, val in results:
        acc.setdefault((g, p, arm), []).append(val)

    summary = {}
    for g in gaps:
        for p in p_grid:
            for arm in arms:
                v = np.array(acc[(g, p, arm)])
                lo, hi = _distal_boot_ci(v, seed=int(10 * g) + int(100 * p))
                summary[(g, p, arm)] = (float(v.mean()), lo, hi)

    return {"variant": variant, "gaps": gaps, "p": p_grid, "arms": arms,
            "tau_long": _DISTAL_TAU_LONG, "tau_short": _DISTAL_TAU_SHORT,
            "summary": summary,
            "faults": "none" if p_grid == [0.0] else "stuck-off + D2D(0.05,0.5) [no PF]",
            "seeds": seeds, "trials": trials}


def _distal_figure(grid):
    """Optional convenience figure: crossover (device-long vs device-short vs T_gap) at
    the clean ``p``. Written under ``paths.results_dir()``; skipped silently if
    matplotlib is unavailable. Returns the saved path or ``None``."""
    from . import paths
    try:
        import matplotlib
        matplotlib.use("Agg")
        import matplotlib.pyplot as plt
    except Exception as e:                                   # plotting is optional
        print(f"[distal/{grid['variant']}] figure skipped "
              f"({type(e).__name__}: {e})", flush=True)
        return None
    gaps, arms, summary = grid["gaps"], grid["arms"], grid["summary"]
    tau_long, tau_short = grid["tau_long"], grid["tau_short"]
    p0 = grid["p"][0]
    fig, ax = plt.subplots(figsize=(5.4, 4.0))
    col = {"long": "#3aa07a", "short": "#c0392b", "no_trace": "#9aa6b2",
           "abstract": "#2f4b8f"}
    for arm in arms:
        m = [summary[(g, p0, arm)][0] for g in gaps]
        lo = [summary[(g, p0, arm)][1] for g in gaps]
        hi = [summary[(g, p0, arm)][2] for g in gaps]
        lbl = {"long": f"device tau={tau_long:g}s", "short": f"device tau={tau_short:g}s",
               "no_trace": "no trace", "abstract": f"abstract tau={tau_long:g}s"}[arm]
        ax.plot(gaps, m, "o-", color=col[arm], label=lbl, lw=1.9, ms=4)
        ax.fill_between(gaps, lo, hi, color=col[arm], alpha=0.15)
    ax.axhline(0.5, color="0.6", ls=":", lw=1.0, label="chance")
    ax.set_xlabel("cue->reward gap $T_{gap}$ (s)")
    ax.set_ylabel("final reward rate")
    ax.set_title(f"Retention is the asset ({grid['variant']}, p={p0:g})", fontsize=10)
    ax.set_ylim(0.4, 1.02)
    ax.legend(fontsize=8, frameon=False, loc="lower left")
    for sp in ("top", "right"):
        ax.spines[sp].set_visible(False)
    fig.tight_layout()
    out = paths.results_dir() / f"fig_distal_cue_{grid['variant']}.png"
    fig.savefig(out, dpi=200, facecolor="white")
    plt.close(fig)
    print(f"[distal/{grid['variant']}] wrote {out.name}", flush=True)
    return out


def main(argv=None):
    """Full-scale reproduction CLI for the distal-cue temporal grid.

    ``python -m mrl_trace.distal_cue [--variant deep|shallow] [--gaps 1,4,10]``
    ``[-p 0] [--seeds N] [--trials N] [--workers N] [--quick|--full]``

    Fans the (gap x arm x seed x p) grid over a fork Pool (``__main__``-only), bootstraps
    95% CIs per (gap, p, arm), prints the per-row table + device-long-minus-short gap, and
    writes the grid via ``paths.save_result("distal_cue_<variant>.npy", obj)`` -- the SAME
    filename the original ``entry_distal_cue.py`` wrote. ``--full`` = the published 12-seed
    / 2000-trial run; ``--quick`` = a fast few-seed / few-trial smoke run.
    """
    import argparse
    import os
    from . import paths
    ap = argparse.ArgumentParser(description="Distal-cue retention-as-asset temporal grid")
    ap.add_argument("--variant", choices=["shallow", "deep"], default="deep")
    ap.add_argument("--gaps", default="1,4,10", help="T_gap values (s) to sweep")
    ap.add_argument("-p", "--p", default="0",
                    help="stuck-fault fractions (0 = clean crossover)")
    ap.add_argument("--seeds", type=int, default=None)
    ap.add_argument("--trials", type=int, default=None)
    ap.add_argument("--workers", type=int, default=int(os.cpu_count() or 4))
    ap.add_argument("--quick", action="store_true", help="fast few-seed smoke run")
    ap.add_argument("--full", action="store_true",
                    help="published 12-seed / 2000-trial run (default)")
    a = ap.parse_args(argv)

    seeds = a.seeds if a.seeds is not None else (3 if a.quick else 12)
    trials = a.trials if a.trials is not None else (200 if a.quick else 2000)
    gaps = [float(x) for x in a.gaps.split(",")]
    p_grid = [float(x) for x in a.p.split(",")]
    window = _DISTAL_CFG["window"]

    arms = ["long", "short", "no_trace"] + (["abstract"] if a.variant == "shallow" else [])
    n_jobs = len(gaps) * len(p_grid) * len(arms) * seeds
    print(f"[distal/{a.variant}] {n_jobs} runs over {a.workers} workers; "
          f"gaps={gaps} p={p_grid} arms={arms} seeds={seeds} trials={trials}; "
          f"tau_long={_DISTAL_TAU_LONG} tau_short={_DISTAL_TAU_SHORT}; "
          f"{'clean' if p_grid == [0.0] else 'stuck+realistic-D2D (no PF)'}", flush=True)

    grid = run_distal_cue(variant=a.variant, gaps=gaps, p_grid=p_grid, seeds=seeds,
                          trials=trials, window=window, workers=a.workers)

    summary = grid["summary"]
    print(f"  {'gap':>5} {'p':>5} " + " ".join(f"{arm:>10}" for arm in arms) +
          "   long-short", flush=True)
    for g in gaps:
        for p in p_grid:
            line = f"  {g:5.1f} {p:5.2f} "
            means = {}
            for arm in arms:
                m, lo, hi = summary[(g, p, arm)]
                means[arm] = m
                line += f" {m:.2f}[{lo:.2f},{hi:.2f}]"
            line += f"   {means['long'] - means['short']:+.2f}"
            print(line, flush=True)

    out = paths.save_result(f"distal_cue_{a.variant}.npy", grid)
    print(f"[distal/{a.variant}] wrote {out.name} to {out.parent}", flush=True)

    _distal_figure(grid)


if __name__ == "__main__":
    main()
\end{lstlisting}

\begin{lstlisting}[caption={siox\_eligibility/probselect.py},label={lst:siox-src-probselect}]
"""Frank probabilistic-selection task (PST) -- a canonical RL paradigm for the device trace.

The XOR bandit of :mod:`mrl_trace.deep` is the minimal test of depth, but a toy.
This module replaces it with the probabilistic-selection task of Frank, Seeberger &
O'Reilly (2004) -- the SAME task the ds003474 EEG subjects performed -- so the device-trace
agent learns a real, canonical reinforcement-learning paradigm whose behaviour can be
compared to human data from the same dataset.

Task. Six stimuli A--F in three TRAIN pairs with fixed reinforcement probabilities
  AB: A 80% / B 20%;  CD: C 70% / D 30%;  EF: E 60% / F 40%.
Each trial shows one pair (the two stimuli in randomised left/right slots); the agent picks
a slot; the chosen slot's stimulus is rewarded stochastically at its probability. Learning
the better stimulus of each pair requires integrating reward OVER TRIALS -- the regime an
eligibility trace exists for. After training, a no-feedback TEST phase shows novel
recombinations; the two canonical Frank measures are
  choose-A : tendency to pick A (the most-positive stimulus) against C,D,E,F -- learning
             from POSITIVE feedback;
  avoid-B  : tendency to avoid B (the most-negative stimulus) against C,D,E,F -- learning
             from NEGATIVE feedback.

Each stimulus is a fixed random sparse binary code; the cue presented to the network is the
two codes placed in two position-tagged input blocks. The mapping from cue to the
better-slot action is a comparison of learned stimulus values, made non-linear by using
low-dimensional codes, so a single trained layer cannot in general solve every pair (a
``shallow`` control checks this empirically). The network, device eligibility, signed
coincidence, DFA feedback, and homeostatic stabiliser are exactly those of
:mod:`mrl_trace.deep`; only the task wrapper differs.
"""
from __future__ import annotations

import numpy as np

from .bandit import GateBankBatched, W_MAX
from .neurons import lif_step_batched, TAU_M, V_TH
from .learning import LTD_BIAS
from .deep import _relax_gate

__all__ = ["PST_PROBS", "make_codes", "train_pst", "reward_rate"]

#: Reinforcement probability of each stimulus (Frank PST).
PST_PROBS = {"A": 0.8, "B": 0.2, "C": 0.7, "D": 0.3, "E": 0.6, "F": 0.4}
#: The three training pairs.
TRAIN_PAIRS = [("A", "B"), ("C", "D"), ("E", "F")]
STIM = ["A", "B", "C", "D", "E", "F"]


def make_codes(rng, n_bits=4, n_active=2):
    """Six fixed random sparse binary codes (one per stimulus) over ``n_bits`` lines.

    Low ``n_bits`` (default 4 for six stimuli) makes the six learned values not in general
    a linear function of the code, so a hidden layer is needed; the ``shallow`` control
    verifies this empirically. Returns ``codes`` of shape ``(6, n_bits)``."""
    codes = np.zeros((6, n_bits))
    for s in range(6):
        on = rng.choice(n_bits, size=n_active, replace=False)
        codes[s, on] = 1.0
    return codes


def _present(rng, B, codes, n_bits):
    """Sample one TRAIN trial per seed. Returns the input field (B, 2*n_bits) with the two
    stimuli in position-tagged slot blocks, the per-seed (slot0_stim, slot1_stim) ids, and
    the better slot (the higher-probability stimulus' slot) for accuracy scoring."""
    pair_idx = rng.integers(len(TRAIN_PAIRS), size=B)
    swap = rng.integers(2, size=B).astype(bool)              # randomise left/right
    s0 = np.empty(B, int); s1 = np.empty(B, int)
    for b in range(B):
        a, c = TRAIN_PAIRS[pair_idx[b]]
        ia, ic = STIM.index(a), STIM.index(c)
        s0[b], s1[b] = (ic, ia) if swap[b] else (ia, ic)     # slot0, slot1 stimulus ids
    field = np.zeros((B, 2 * n_bits))
    field[:, :n_bits] = codes[s0]
    field[:, n_bits:] = codes[s1]
    probs = np.array([PST_PROBS[STIM[i]] for i in range(6)])
    better = (probs[s1] > probs[s0]).astype(int)             # better slot (0 or 1)
    return field, s0, s1, better, probs


def train_pst(*, mode="dfa_homeo", B=20, H=16, n_bits=4, tau_leak=10.0, D=5.0,
              trials=4000, dt=5e-3, cue_dur=1.0, eta=0.2, eta_hidden=3.0, in_rate=200.0,
              ltd=LTD_BIAS, tau_m=TAU_M, v_th=V_TH, V=1.5, sigma0=0.15, sigma1=0.05,
              fb_scale=2.0, w_scale1=0.6, w_scale2=0.35, w_max=W_MAX, bias_o=0.3,
              homeo=0.1, homeo_target=0.35, homeo_tau=200.0, reward_pools=None,
              seed0=0, return_test=False):
    """Train the PST policy. ``mode`` in {shallow, dfa, dfa_homeo, no_trace}; semantics as
    in :func:`mrl_trace.deep.train_deep`. Device eligibility on every plastic layer.

    Returns per-trial TRAIN correctness ``(B, trials)`` (chose the higher-probability
    stimulus). With ``return_test`` also returns ``(choose_A, avoid_B)`` arrays of shape
    ``(B,)`` from a no-feedback test phase over the novel recombinations.
    """
    if mode == "dfa_homeo":
        mode_deep, hm = "dfa", homeo
    elif mode == "no_trace":
        mode_deep, hm = "no_trace", homeo
    else:
        mode_deep, hm = mode, 0.0
    rng = np.random.default_rng(seed0)
    codes = make_codes(rng, n_bits=n_bits)
    F, A = 2 * n_bits, 2
    deep = mode_deep != "shallow"
    no_trace = mode_deep == "no_trace"
    train_w1 = deep

    def _init(shape, scale):
        return np.clip(scale * rng.standard_normal(shape), -w_max, w_max)
    if deep:
        W1 = _init((B, F, H), w_scale1)
        W2 = _init((B, H, A), w_scale2)
    else:
        W2 = _init((B, F, A), w_scale1); W1 = None

    g1 = GateBankBatched(B, F, H, tau_leak=tau_leak, dt=dt, V=V) if deep else None
    g_out = GateBankBatched(B, (H if deep else F), A, tau_leak=tau_leak, dt=dt, V=V)
    B_fix = fb_scale * rng.standard_normal((A, H)) if deep else None
    baseline = np.full(B, 0.5)
    bidx = np.arange(B)
    cue = (0.3, 0.3 + cue_dur)
    reward_lag = int(round(D / dt))
    nsteps = int(round(cue[1] / dt)) + 2
    rewards = np.zeros((B, trials))
    act_hidden = np.full((B, H), homeo_target) if (deep and hm > 0) else None

    def _run_cue(field, learn=True, rng_local=None):
        """Run one cue presentation; return chosen slot (B,) and (for learning) cached
        spike/eligibility state. If learn, integrates eligibility and is used in the loop."""
        r = rng if rng_local is None else rng_local
        vh = np.zeros((B, H)); vo = np.zeros((B, A))
        spk_o = np.zeros((B, A)); spk_h = np.zeros((B, H))
        if deep and learn:
            g1.reset()
        if learn:
            g_out.reset()
        for n in range(nsteps):
            on = cue[0] <= n * dt < cue[1]
            pre_in = ((r.random((B, F)) < (in_rate * dt) * field).astype(float)
                      if on else np.zeros((B, F)))
            if deep:
                ch_h = np.einsum('bfh,bf->bh', W1, pre_in)
                vh, sp_h = lif_step_batched(vh, ch_h, dt, r, tau_m=tau_m, v_th=v_th, noise=sigma)
                spk_h += sp_h
                ch_o = np.einsum('bha,bh->ba', W2, sp_h.astype(float))
                if on:
                    ch_o = ch_o + bias_o
                vo, sp_o = lif_step_batched(vo, ch_o, dt, r, tau_m=tau_m, v_th=v_th, noise=sigma)
                spk_o += sp_o
                if learn:
                    if no_trace:
                        g1.step(np.zeros((B, F, H))); g_out.step(np.zeros((B, H, A)))
                    else:
                        g1.step(pre_in[:, :, None] * np.where(sp_h, 1.0, -ltd)[:, None, :])
                        g_out.step(sp_h.astype(float)[:, :, None]
                                   * np.where(sp_o, 1.0, -ltd)[:, None, :])
            else:
                ch_o = np.einsum('bfa,bf->ba', W2, pre_in)
                if on:
                    ch_o = ch_o + bias_o
                vo, sp_o = lif_step_batched(vo, ch_o, dt, r, tau_m=tau_m, v_th=v_th, noise=sigma)
                spk_o += sp_o
                if learn:
                    g_out.step(pre_in[:, :, None] * np.where(sp_o, 1.0, -ltd)[:, None, :])
        tie = spk_o.max(1) == spk_o.min(1)
        chosen = np.argmax(spk_o, 1)
        chosen[tie] = r.integers(A, size=int(tie.sum()))
        return chosen, spk_h, spk_o

    for tr in range(trials):
        sigma = sigma0 + (sigma1 - sigma0) * tr / trials
        field, s0, s1, better, probs = _present(rng, B, codes, n_bits)
        chosen, spk_h, spk_o = _run_cue(field, learn=True)
        # stochastic reward at the chosen stimulus' reinforcement probability
        chosen_stim = np.where(chosen == 0, s0, s1)
        p_chosen = probs[chosen_stim]
        R_true = (chosen == better).astype(float)                # chose higher-prob stim
        outcome = (rng.random(B) < p_chosen).astype(int)         # actual stochastic reward
        if reward_pools is None:
            R = outcome.astype(float)
        else:
            R = np.array([reward_pools[int(outcome[b])][
                rng.integers(len(reward_pools[int(outcome[b])]))] for b in range(B)])

        if not no_trace:
            stride = 10 if reward_lag > 200 else 1
            n_relax, rem = divmod(reward_lag, stride)
            _relax_gate(g_out, n_relax, stride, rem)
            if deep:
                _relax_gate(g1, n_relax, stride, rem)
            eo_rew = g_out.vn[..., -1] / g_out.Vnmax
            e1_rew = (g1.vn[..., -1] / g1.Vnmax) if deep else None
        else:
            eo_rew = np.zeros_like(g_out.vn[..., -1])
            e1_rew = np.zeros((B, F, H)) if deep else None

        adv = (R - baseline)
        Lo = adv[:, None, None]
        if deep:
            logits = spk_o - spk_o.mean(1, keepdims=True)
            pol = np.exp(logits - logits.max(1, keepdims=True)); pol /= pol.sum(1, keepdims=True)
            onehot = np.zeros((B, A)); onehot[bidx, chosen] = 1.0
            L_a = adv[:, None] * (onehot - pol)
            Lh = np.einsum('ah,ba->bh', B_fix, L_a)[:, None, :]
        W2 = np.clip(W2 + eta * Lo * eo_rew, -w_max, w_max)
        if train_w1:
            dW1 = eta_hidden * Lh * e1_rew
            if hm > 0 and act_hidden is not None:
                rate = spk_h / nsteps
                act_hidden += (rate - act_hidden) / homeo_tau
                scale = 1.0 + hm * (homeo_target - act_hidden)
                dW1 = dW1 + (scale[:, None, :] - 1.0) * W1
            W1 = np.clip(W1 + dW1, -w_max, w_max)
        baseline += 0.02 * (R - baseline)
        rewards[:, tr] = R_true

    if not return_test:
        return rewards

    # --- TEST phase: novel recombinations, no learning, greedy (low-noise) choice ---
    rng_t = np.random.default_rng(seed0 + 99991)
    saved = sigma
    def _choice_prob(slot0_stim, slot1_stim):
        """For each seed, fraction of repeats it picks slot 1 (the second listed)."""
        field = np.zeros((B, F))
        field[:, :n_bits] = codes[slot0_stim]; field[:, n_bits:] = codes[slot1_stim]
        picks = np.zeros(B)
        reps = 21
        nonlocal sigma
        sigma = 0.05                                          # near-greedy at test
        for _ in range(reps):
            chosen, _, _ = _run_cue(field, learn=False, rng_local=rng_t)
            picks += chosen
        sigma = saved
        return picks / reps                                   # P(choose slot1) per seed
    iA, iB = STIM.index("A"), STIM.index("B")
    others = [STIM.index(s) for s in ("C", "D", "E", "F")]
    # choose-A: A vs each other; pick rate of A (A placed in slot0 -> P(choose A)=1-P(slot1))
    cA = np.zeros(B)
    for o in others:
        cA += 1.0 - _choice_prob(np.full(B, iA), np.full(B, o))
    cA /= len(others)
    # avoid-B: B vs each other; avoid rate of B (B in slot0 -> avoid = P(choose slot1))
    aB = np.zeros(B)
    for o in others:
        aB += _choice_prob(np.full(B, iB), np.full(B, o))
    aB /= len(others)
    return rewards, cA, aB


def reward_rate(rewards, window=200):
    rewards = np.asarray(rewards)
    if rewards.shape[-1] < window:
        return rewards.mean(axis=-1)
    return rewards[..., -window:].mean(axis=-1)


# =============================================================================
# Experiment core (Experiment 10, Arm F -- the probabilistic-selection study)
#
# ``run_probselect`` returns one condition's result as a plain dict -- no file I/O,
# no plotting, no stdout -- so a notebook can call it in-kernel at a small (quick)
# trial count.  ``main()`` (below) runs the published 20-seed grid across a process
# Pool over conditions, computes bootstrap CIs + the pre-registered F1--F5 criteria,
# and writes ``exp10_probselect.npy`` under ``data/results/`` for the notebook to
# replay.  Provenance: experiments/05_biological_grounding/probabilistic_selection.py
# (Arm F), pre-registered criteria F1--F5 in PREREGISTRATION_probselect.md.
# =============================================================================

#: Fixed hyper-parameters of the published Arm-F grid (H, code width, retention,
#: reward lag, learning rates, DFA feedback scale, output bias) -- the shape shared
#: by every condition so only mode/reward-source varies.
PST_HP = dict(H=16, n_bits=12, tau_leak=10.0, D=5.0, eta=0.2, eta_hidden=3.0,
              fb_scale=2.0, bias_o=0.3)
#: Homeostatic-stabiliser strength used by the ``dfa_homeo`` family.
PST_HOMEO = 0.1
#: Published scale + task references.
PST_SEEDS = 20
PST_TRIALS = 5000
PST_CHANCE = 0.5
PST_CRIT = 0.75


def run_probselect(cond, *, trials=PST_TRIALS, seeds=PST_SEEDS, pools=None, shuf=None,
                   hp=None, homeo=PST_HOMEO, seed0=0):
    """One independent Arm-F cell: train the PST for a single condition ``cond`` on a
    vectorised ``seeds``-seed batch and return its behavioural summary as a plain dict.

    ``cond`` selects mode + reward source (semantics as in the Arm-F driver):
      ``shallow``          one trained device-synapse layer (the linearly-separable
                           control; solves the task -- headline Frank signature);
      ``dfa``              deep, DFA feedback, no homeostasis;
      ``dfa_homeo``        deep, DFA feedback + homeostatic stabiliser (the device agent);
      ``no_trace``         eligibility zeroed (device-necessity control);
      ``dfa_homeo_eeg``    ``dfa_homeo`` gated by the subjects' OWN decoded real-EEG
                           reward (requires ``pools``; the human non-invasive RPE);
      ``dfa_homeo_shuf``   ``dfa_homeo`` gated by the shuffled-reward control
                           (requires ``shuf``).

    The test phase (choose-A / avoid-B) is run only for the conditions the Frank
    figure reports (``shallow``, ``dfa``, ``dfa_homeo``).  Returns
    ``dict(cond, finals (seeds,), curve (trials-200,), chooseA, avoidB)`` where
    ``finals`` is the last-300-trial train accuracy per seed, ``curve`` is the
    seed-mean 200-trial running P(better stimulus), and ``chooseA``/``avoidB`` are the
    test-phase Frank measures per seed (``None`` for the non-test conditions).
    """
    hp = PST_HP if hp is None else hp
    rp = None
    mode = cond
    if cond == "dfa_homeo_eeg":
        mode, rp = "dfa_homeo", {0: pools[0], 1: pools[1]}
    elif cond == "dfa_homeo_shuf":
        mode, rp = "dfa_homeo", shuf
    want_test = cond in ("dfa_homeo", "shallow", "dfa")
    out = train_pst(mode=mode, B=seeds, trials=trials, seed0=seed0, homeo=homeo,
                    reward_pools=rp, return_test=want_test, **hp)
    if want_test:
        rew, cA, aB = out
    else:
        rew, cA, aB = out, None, None
    finals = reward_rate(rew, window=300)
    win = 200
    csum = np.cumsum(rew, axis=1)
    curve = ((csum[:, win:] - csum[:, :-win]) / win).mean(0)
    return {"cond": cond, "finals": finals, "curve": curve, "chooseA": cA, "avoidB": aB}


def _load_eeg_pools(cache="/tmp/eeg_pools.npy", seed=0):
    """Load the cached EEG out-of-fold reward pools (built by the EEG-capstone driver
    from ds003474) and derive the shuffled-reward control, degrading gracefully.

    Returns ``(pools, shuf, meta)``; ``(None, None, {})`` when the cache is absent (the
    subjects' raw EEG is not bundled), so the EEG conditions are skipped exactly as the
    original Arm-F driver does when ``pools`` is ``None``."""
    import os
    if not os.path.exists(cache):
        return None, None, {}
    pools = np.load(cache, allow_pickle=True).item()
    meta = pools.get("_meta", {})
    rng = np.random.default_rng(seed)
    allv = np.concatenate([pools[1], pools[0]])
    shuf = {1: rng.permutation(allv), 0: rng.permutation(allv)}
    return pools, shuf, meta


def _summarize_probselect(res_by_cond, conds, *, seeds, trials, meta, hp,
                          chance=PST_CHANCE, crit=PST_CRIT):
    """Package the per-condition :func:`run_probselect` output into the saved grid dict
    (the exact schema the Arm-F driver wrote to ``exp10_probselect.npy``) plus the
    pre-registered criteria F1--F5.

    Criteria (no goalpost moving; from PREREGISTRATION_probselect.md):
      F1  device learns the train phase             (dfa_homeo mean >= ``crit``);
      F2  depth is needed                            (homeo CI lower > shallow CI upper);
      F3  eligibility is necessary                   (no-trace mean <= 0.60);
      F4  Frank asymmetry present                    (choose-A and avoid-B CI lower > 0.5);
      F5  the real-EEG loop beats its shuffled control (EEG CI lower > shuf CI upper) --
          only evaluated when the EEG pools are available.
    """
    from .stats import bootstrap_ci
    finals = {c: res_by_cond[c]["finals"] for c in conds}
    curves = {c: res_by_cond[c]["curve"] for c in conds}
    tests = {c: (res_by_cond[c]["chooseA"], res_by_cond[c]["avoidB"])
             for c in conds if res_by_cond[c]["chooseA"] is not None}
    ci = {c: bootstrap_ci(finals[c]) for c in finals}

    dh = finals["dfa_homeo"]
    cA, aB = tests["dfa_homeo"]
    criteria = {
        "F1": bool(dh.mean() >= crit),
        "F2": bool(ci["dfa_homeo"][0] > ci["shallow"][1]),
        "F3": bool(finals["no_trace"].mean() <= 0.60),
        "F4": bool(bootstrap_ci(cA)[0] > 0.5 and bootstrap_ci(aB)[0] > 0.5),
    }
    if "dfa_homeo_eeg" in finals and "dfa_homeo_shuf" in finals:
        criteria["F5"] = bool(ci["dfa_homeo_eeg"][0] > ci["dfa_homeo_shuf"][1])

    return {
        "finals": finals, "curves": curves,
        "tests": {k: {"chooseA": v[0], "avoidB": v[1]} for k, v in tests.items()},
        "ci": {k: ci[k] for k in ci}, "meta": meta, "HP": hp,
        "seeds": seeds, "trials": trials, "chance": chance, "crit": crit,
        "criteria": criteria,
    }


def main(argv=None):
    """Full-scale reproduction CLI for the probabilistic-selection grid (Experiment 10,
    Arm F): the Frank PST -- the SAME task the ds003474 EEG subjects performed --
    learned by the deep all-local device-trace agent, lifting the RL result off the XOR toy.

    ``python -m mrl_trace.probselect [--probselect] [--full|--quick]``
    ``--full`` = 20 seeds x 5000 trials (published); ``--quick`` = 6 seeds x 1000 trials.
    Each condition (mode/reward source) is an independent cell run across a process Pool.
    The synthetic conditions always run; the two real-EEG conditions run only when the
    cached EEG reward pools (``/tmp/eeg_pools.npy``, built by the EEG-capstone driver from
    ds003474) are present -- absent, they are skipped with a printed note, exactly as the
    original Arm-F driver degrades.  Writes ``exp10_probselect.npy`` under ``data/results``.
    """
    import argparse
    import os
    from functools import partial
    from multiprocessing import Pool
    from . import paths
    from .stats import bootstrap_ci

    ap = argparse.ArgumentParser(description="Frank probabilistic-selection RL reproduction")
    ap.add_argument("--probselect", action="store_true",
                    help="run the Arm-F grid -> exp10_probselect.npy (default)")
    ap.add_argument("--quick", action="store_true", help="fast few-seed smoke run")
    ap.add_argument("--full", action="store_true", help="published 20-seed run (default)")
    a = ap.parse_args(argv)

    seeds = 6 if a.quick else PST_SEEDS
    trials = 1000 if a.quick else PST_TRIALS

    pools, shuf, meta = _load_eeg_pools()
    conds = ["shallow", "dfa", "dfa_homeo", "no_trace"]
    if pools is not None:
        conds += ["dfa_homeo_eeg", "dfa_homeo_shuf"]
    else:
        print("  (no EEG reward pools at /tmp/eeg_pools.npy -- skipping the real-EEG "
              "conditions; run the EEG-capstone driver first to include them)")

    print(f"Probabilistic-selection task | {seeds} seeds, {trials} trials, "
          f"H={PST_HP['H']}, n_bits={PST_HP['n_bits']}")

    worker = partial(run_probselect, trials=trials, seeds=seeds, pools=pools, shuf=shuf)
    with Pool(processes=min(len(conds), (os.cpu_count() or 4) - 1)) as pool:
        res = pool.map(worker, conds)
    res_by_cond = {r["cond"]: r for r in res}

    grid = _summarize_probselect(res_by_cond, conds, seeds=seeds, trials=trials,
                                 meta=meta, hp=PST_HP)

    print("\n=== final train accuracy (chose higher-prob stimulus, last 300 trials) ===")
    for c in conds:
        lo, hi = grid["ci"][c]
        print(f"  {c:22s} {grid['finals'][c].mean():.3f} [{lo:.3f}, {hi:.3f}]")

    print("\n=== test-phase Frank signature (choose-A / avoid-B) ===")
    for c, t in grid["tests"].items():
        cA, aB = t["chooseA"], t["avoidB"]
        clo, chi = bootstrap_ci(cA); alo, ahi = bootstrap_ci(aB)
        print(f"  {c:22s} choose-A {cA.mean():.3f} [{clo:.2f},{chi:.2f}]  "
              f"avoid-B {aB.mean():.3f} [{alo:.2f},{ahi:.2f}]")

    print(f"\n=== pre-registered criteria ===  {grid['criteria']}")
    paths.save_result("exp10_probselect.npy", grid)
    print("  wrote exp10_probselect.npy")


if __name__ == "__main__":
    main()
\end{lstlisting}

\begin{lstlisting}[caption={siox\_eligibility/selectivity.py},label={lst:siox-src-selectivity}]
r"""Timing/selectivity constructions on the device eligibility cascade.

Two experiments that treat the SiO\ :sub:`x` trap-cascade (see
:mod:`mrl_trace.device`) as a *temporal* primitive rather than as a scalar
credit-assignment window.  Both compose the existing gates (the trace kernel of
:mod:`mrl_trace.distal_reward` and the vectorised
:class:`mrl_trace.bandit.GateBankBatched`) with the LIF decision layer, but
neither wraps a single existing training primitive -- they build new timing tasks on
top of the shared physics, so they live in their own module.

- **Interval selectivity (Experiment 10).**  The device trace is *peaked* at a
  non-zero lag ``t*`` (Eqs. 7.1--7.2), so it is a band-pass in time: it credits a
  coincidence at a *preferred interval* and suppresses a coincidence at lag~0
  (reward time), whereas a matched-decay exponential is monotone (recency-biased)
  and credits the lag-0 confound.  The preferred interval is fabrication-tunable via
  ``tau_leak``.  :func:`run_interval_selectivity` measures the *learned* selectivity
  ratio ``S = w_pref/w_late`` from a competitive three-factor loop and returns the
  design-point comparison plus the ``S(D)`` crossover sweep.  Pre-registered criteria
  P1--P4 and kills K1--K2 (PREREGISTRATION_interval_selectivity.md).

- **Vector timer (Experiment 20).**  The cascade occupancy *vector* ``v^1..v^k`` is a
  distributed clockless code for elapsed time: the stage occupancies peak at
  progressively later lags, so the full vector separates elapsed times that the
  non-monotone last-stage *scalar* aliases.  :func:`run_vector_timer` runs a
  closed-loop interval-discrimination bandit whose policy is driven by either the full
  cascade vector (this work), the last stage only (the CET/current-paper baseline), or
  a zeroed trace (necessity control), at a scalar-aliasing interval pair computed from
  the device model.  Pre-registered criteria H1--H4 and kill K1.  (Physical-
  observability caveat: the individual stage occupancies are not exposed by a
  two-terminal cell, so the vector result *motivates* a staged-readout device rather
  than being a capability of the present one -- thesis Ch7 exploratory.)

Each ``run_*`` is serial and import-light (a notebook calls it in-kernel at a small
seed/trial count); ``main()`` runs the published-scale grids, parallelising the
coarse axis with a process pool, and writes each grid via
``paths.save_result`` under ``data/results/`` (``exp10_interval.npy``,
``exp20_vector_timer.npy``).
"""
from __future__ import annotations

import numpy as np

from .device import TransientGate, tau_r
from .distal_reward import trace_kernel, abstract_kernel, W_INIT
from .bandit import GateBankBatched
from .neurons import lif_step_batched, TAU_M, V_TH
from .learning import LTD_BIAS
from .stats import bootstrap_ci

__all__ = [
    "run_interval_selectivity",
    "peak_lag",
    "selectivity_ratio",
    "run_vector_timer",
    "aliasing_pair",
    "run_interval",
]

# =============================================================================
# Experiment 10 -- the dual-phase trace as a tunable interval-selective eligibility.
#
# Pre-registered in PREREGISTRATION_interval_selectivity.md (criteria P1-P4, kills
# K1-K2). Claim: the device trace (Eqs. 7.1-7.2) is PEAKED at a non-zero lag t*, so it
# implements a band-pass in time -- it credits a coincidence at a *preferred interval*
# and SUPPRESSES a coincidence at lag~0 (reward time). A matched-decay exponential is
# monotone (recency-biased) and credits the lag-0 confound instead. The preferred
# interval is fabrication-tunable via tau_leak.
#
# Task (trace level, competition for credit):
#   - syn_pref : reward-correlated coincidence at lag D_rew (the preferred interval).
#   - syn_late : reward-correlated coincidence at lag ~0.3 s (the reward-time confound).
#     BOTH are perfectly reward-correlated, so the test is purely about TIMING.
#   - N-2 reward-UNcorrelated distractors at random lags (the LTD wing handles them).
# Metric: learned selectivity ratio S = w[syn_pref] / w[syn_late]. S>1 => interval-
# selective; S<1 => recency-biased. S is a LEARNED outcome of a competitive 1500-trial
# three-factor loop with a shared adapting baseline, not a direct read of the kernel.
# =============================================================================

LATE_LAG = 0.3       # s, the reward-time confound lag
LTD = 0.6
JITTER = 0.15        # per-trial fractional lag jitter (biological timing is noisy)
DECAY = 0.08         # weight-decay alpha: the leaky three-factor rule Dw=eta((R-b)e-alpha w)
                     # (the form of Ch2 Eq. 2.17 / Bianchi 2020); makes the equilibrium
                     # weight eligibility-proportional rather than clip-saturated.


def peak_lag(tau_leak, V=0.9, T=80.0, dt=0.05):
    """Preferred interval t* = argmax of the (unnormalised) device trace."""
    g = TransientGate(V=V, tau_leak=tau_leak, dt=dt)
    t = np.arange(0, T, dt)
    raw = g.trace(t, coincidence_at=0.0, coincidence_dur=0.3, normalise=False)
    return float(t[np.argmax(raw)])


def _kernel(kind, t, tau_leak, V):
    if kind == "device":
        return trace_kernel(t, tau_leak, V=V)
    if kind == "exp":
        return abstract_kernel(t, tau_leak)
    raise ValueError(kind)


def train_selectivity(kind, tau_leak, D_rew, *, V=0.9, N=8, trials=1500,
                      T=160.0, dt=0.1, eta=0.05, seed=0):
    """One seed. Returns final weights w (w[0]=syn_pref @ lag D_rew,
    w[1]=syn_late @ lag LATE_LAG, w[2:]=reward-uncorrelated distractors).

    Both syn_pref and syn_late are PERFECTLY reward-correlated: when a trial is
    rewarded, both emit a coincidence (at their respective lags before reward).
    The only thing distinguishing them is WHEN they fire relative to reward, so the
    learned ratio w_pref/w_late isolates the trace's timing selectivity.
    """
    rng = np.random.default_rng(seed)
    t = np.arange(0, T, dt)
    nt = len(t)
    kern = _kernel(kind, t, tau_leak, V)
    # eligibility a synapse carries at reward time, given its coincidence LAG before
    # reward: e = kern[lag_index]. (kern is e(t) after a coincidence at t=0.)
    def elig_at_lag(lag):
        li = int(round(lag / dt))
        return kern[li] if 0 <= li < nt else 0.0

    w = np.full(N, W_INIT)
    baseline = 0.5
    for _ in range(trials):
        rewarded = rng.random() < 0.5
        e = np.zeros(N)
        if rewarded:
            # per-trial multiplicative lag jitter: biological action--reward timing is
            # not exact, so neither synapse sits at a fixed point lag. This injects
            # genuine seed/trial variance (the degenerate-CI fix) and makes the test
            # robust to a synapse exploiting one exact grid point.
            j_pref = D_rew * (1.0 + JITTER * rng.standard_normal())
            j_late = max(0.0, LATE_LAG + 0.1 * abs(rng.standard_normal()))
            e[0] = elig_at_lag(j_pref)       # preferred-interval cue
            e[1] = elig_at_lag(j_late)       # reward-time confound
        # reward-uncorrelated distractors at random lags (LTD wing nets them down).
        for i in range(2, N):
            lag = rng.uniform(LATE_LAG, 0.9 * T)
            sign = 1.0 if rng.random() < 0.5 else -LTD
            e[i] = sign * elig_at_lag(lag)
        R = 1.0 if rewarded else 0.0
        baseline += 0.02 * (R - baseline)
        # leaky three-factor update: the -alpha*w term bounds weights by DECAY so the
        # equilibrium w* is proportional to the mean eligibility a synapse accrues,
        # NOT pinned to a hard clip. This is what makes w_pref/w_late report the
        # trace's timing selectivity rather than a saturation artefact.
        w = np.clip(w + eta * ((R - baseline) * e - DECAY * (w - W_INIT)), 0.0, 2.0)
    return w


def selectivity_ratio(w):
    """S = w_pref / w_late, guarding the denominator away from 0."""
    return w[0] / max(w[1], 1e-6)


def _sweep_selectivity(kind, tau_leak, D_rew, *, V=0.9, seeds=20, trials=1500):
    """Per-seed selectivity ratios over ``seeds`` seeds (serial)."""
    return np.array([
        selectivity_ratio(train_selectivity(kind, tau_leak, D_rew, V=V,
                                             trials=trials, seed=s))
        for s in range(seeds)
    ])


def run_interval_selectivity(*, seeds=20, trials=1500, V=0.9):
    """Experiment 10 core: device vs exponential interval selectivity + tunability.

    Serial; returns the result grid as a plain dict, no file I/O / no plotting /
    no stdout (the notebook renders the figure inline).  Preserves the science of
    ``experiments/04_temporal_selectivity/interval_selectivity.py`` exactly.

    Computes, at the design point ``tau_leak=10 s, D_rew=t*`` (the last-stage peak),
    the per-seed learned selectivity ``S = w_pref/w_late`` for the device trace and
    the matched-tau exponential control, with bootstrap CIs; and the ``S(D)`` curve
    over a log-spaced design-interval grid for three retentions ``tau_leak in
    {5,10,20} s`` (the pre-registered CROSSOVER: a short design interval is best
    served by short retention, a long one by long retention -- the peak sliding right
    with ``tau_leak``).  Pre-registered criteria (fixed before running):
      P1  device interval-selective     : device S CI lower bound > 1.
      P2  exponential recency-biased     : exp S CI upper bound < 1.
      P3  preferred interval shifts w/ tau_leak : S(short D) decreases with tau_leak
          AND S(long D) increases with it (crossover).
      P4  device/exp CIs disjoint at design point.
      K1  (kill) device CI includes 1.
      K2  (kill) device/exp CIs overlap (= not P4).
    """
    taus = [5.0, 10.0, 20.0]
    tstar = {tl: peak_lag(tl, V=V) for tl in taus}

    tl0 = 10.0
    D0 = round(tstar[tl0], 1)
    dev_S = _sweep_selectivity("device", tl0, D0, V=V, seeds=seeds, trials=trials)
    exp_S = _sweep_selectivity("exp", tl0, D0, V=V, seeds=seeds, trials=trials)
    dlo, dhi = bootstrap_ci(dev_S)
    elo, ehi = bootstrap_ci(exp_S)

    # P3: preferred interval SHIFTS RIGHT with tau_leak (crossover). Selectivity S(D)
    # plateaus once D clears the rise, so an argmax is ill-defined on the plateau. The
    # robust, non-tautological signature of a *moving* preferred interval is a
    # CROSSOVER: a short design interval is best served by short retention, a long one
    # by long retention. We test that S(short D) decreases with tau_leak while S(long D)
    # increases with it -- the peak sliding rightward.
    D_grid = [1, 2, 4, 8, 16, 32, 48]
    Scurve = {}
    for tl in taus:
        Scurve[tl] = np.array([
            _sweep_selectivity("device", tl, D, V=V, seeds=seeds, trials=trials).mean()
            for D in D_grid
        ])
    D_short, D_long = 2, 32
    si = D_grid.index(D_short); li = D_grid.index(D_long)
    short_decreases = Scurve[5.0][si] > Scurve[20.0][si]
    long_increases = Scurve[20.0][li] > Scurve[5.0][li]

    P1 = bool(dlo > 1.0)
    P2 = bool(ehi < 1.0)
    P3 = bool(short_decreases and long_increases)
    P4 = bool((dlo > ehi) or (dhi < elo))

    return {
        "V": V, "taus": taus, "tstar": tstar, "D0": D0, "late_lag": LATE_LAG,
        "dev_S": dev_S, "exp_S": exp_S, "dev_ci": (dlo, dhi), "exp_ci": (elo, ehi),
        "D_grid": D_grid, "Scurve": Scurve,
        "P1": P1, "P2": P2, "P3": P3, "P4": P4, "seeds": seeds, "trials": trials,
    }


# =============================================================================
# Experiment 20 -- the cascade occupancy VECTOR as a clockless elapsed-time code.
#
# THE LEARNING CONTRIBUTION (the one slice the novelty search left open). Every
# eligibility trace in this work, and in the nearest prior art (Ralambomihanta et al.,
# Cascading Eligibility Traces, arXiv:2506.14598), reads a SINGLE SCALAR -- the last
# cascade stage -- as the trace. But the device's internal trap-cascade is a
# MULTI-STATE object: the occupancies of its k stages (v^1..v^k) peak at progressively
# later lags (verified: at tau_leak=10s the per-stage peaks sit at ~0, 9, 19 s), so the
# FULL VECTOR is a distributed code for TIME-SINCE-EVENT that a single scalar cannot be.
# Because the scalar readout is non-monotonic (it rises then decays), two different
# elapsed times on opposite sides of its peak ALIAS to the same scalar value -- a rule
# reading the scalar literally cannot tell them apart. The full vector breaks aliasing.
#
# This experiment tests, in a CLOSED LOOP, whether a three-factor rule that reads the
# FULL stage vector solves an interval-discrimination task that the scalar readout
# cannot -- i.e. whether the device's intrinsic cascade state is a usable temporal
# basis for credit assignment, with no clock and no recurrence.
#
# TASK (interval-discrimination bandit, 2 actions). Each trial: a cue fires a pre--post
# coincidence at t=0, driving the eligibility gate. After an elapsed interval drawn as
# one of two values (tA or tB), a decision is taken and reward is contingent on the
# action matching WHICH interval elapsed (act 0 if tA, act 1 if tB). The only signal
# distinguishing the two is the SHAPE of the surviving cascade state at decision time.
# tA and tB STRADDLE the scalar readout's peak at EQUAL last-stage value, so the scalar
# genuinely aliases them while the vector separates them. The pair is COMPUTED from the
# device model (aliasing_pair()): locate the last-stage peak t*, take a rising-flank tA
# and the falling-flank tB of matching last-stage value -- an earlier hard-coded
# tA=3,tB=15 sat on the rising flank, so the scalar was never aliased and the
# aliasing-rescue test was vacuous.
#
# READOUT CONDITIONS (the comparison that isolates the contribution):
#   vector   : the rule reads all k cascade occupancies through a learned per-stage
#              combination c_m (k extra local parameters); credit = (R-b)*sum_m c_m
#              v^(m).  [THIS WORK]
#   scalar   : the rule reads only the last stage v^(k) -- the current-paper/CET
#              readout.  [BASELINE]
#   no_trace : eligibility zeroed (necessity control).
#
# CHARACTERISATION SWEEP (turns "bounded resolution" into a result): vector-vs-scalar
# accuracy across (interval separation) x (cascade depth k). Establishes (i) the
# vector's advantage band, (ii) that resolution is limited by readout NOISE not by k
# (high-k plateaus/regresses), so the measured k~3 already sits near the useful optimum
# -- not a limitation appealing to a deeper device that does not exist.
#
# PRE-REGISTRATION (fixed before running; reported as-is):
#   H1  vector readout learns the interval-discrimination task (straddling pair):
#       vector >= CRIT.
#   H2  scalar readout FAILS the same task (aliasing): scalar <= chance + 0.10, AND
#       vector CI lower bound disjoint above scalar CI upper bound.
#   H3  trace necessary: no_trace <= chance + 0.10.
#   H4  (DESCRIPTIVE) resolution characterisation: vector-minus-scalar advantage as a
#       function of interval separation, and the k-dependence at fixed separation
#       (expect a sweet spot ~k=3-5 then plateau/regress -> noise-limited, not
#       depth-limited).
#   K1  (kill) vector ~ scalar on the straddling pair (vector adds nothing) -> dead.
# =============================================================================

CHANCE = 0.5
CRIT = 0.75


def run_interval(readout, *, B=20, k=3, tau_leak=10.0, tA=3.0, tB=15.0, trials=4000,
                 dt=5e-3, cue_dur=0.5, eta=0.3, in_rate=200.0, ltd=LTD_BIAS,
                 tau_m=TAU_M, v_th=V_TH, sigma=0.4, read_noise=0.10, seed0=0):
    """One condition of the interval-discrimination bandit, STATE-DRIVEN policy.

    A single cue fires a coincidence at t=0, driving the device gate. After an elapsed
    interval (tA or tB, chosen per trial) a decision is taken and reward is contingent
    on the action matching WHICH interval elapsed. The decision must therefore depend
    on the elapsed interval, and the ONLY trace of the interval at decision time is the
    SHAPE of the surviving cascade state -- so the cascade readout DRIVES the action
    neurons (a state-driven policy), and the readout mode is exactly what determines
    whether the two intervals are separable:
      readout='vector'  : action input = M . [v^1..v^k]  (full occupancy vector, k inputs/action)
      readout='scalar'  : action input = m . v^(k)        (last stage only -- CET/current readout)
      readout='no_trace': action input = 0                (necessity control)
    The readout->action weights (M or m) are learned online by the same reward-modulated
    rule (REINFORCE-style: nudged toward the chosen action's input when reward beats
    baseline), local and without weight transport. Because the scalar last stage is
    NON-MONOTONE, tA and tB on opposite sides of its peak ALIAS to one value, so no
    linear scalar->action map can separate them; the vector breaks the aliasing. A small
    multiplicative read noise on the cascade state is what makes the aliasing bite and
    the resolution finite.
    Returns rewards (B, trials)."""
    rng = np.random.default_rng(seed0)
    A = 2
    no_trace = (readout == "no_trace")
    nin = 1 if readout == "scalar" else k                  # readout dimensionality feeding action
    bank = GateBankBatched(B, 1, 1, tau_leak=tau_leak, k=k, dt=dt)
    # readout->action weights, (B, nin, A); start small/random so the policy must LEARN the map.
    M = 0.01 * rng.standard_normal((B, nin, A))
    baseline = np.full(B, 1.0 / A)
    cue = (0.3, 0.3 + cue_dur)
    riA = int((cue[1] + tA) / dt); riB = int((cue[1] + tB) / dt)
    nsteps = max(riA, riB) + 2
    rewards = np.zeros((B, trials))
    bidx = np.arange(B)

    for tr in range(trials):
        bank.reset()
        which = rng.integers(2, size=B)                    # 0 -> tA, 1 -> tB
        ri = np.where(which == 0, riA, riB)
        captured = np.zeros(B, bool)
        e_state = np.zeros((B, k))                          # cascade vector snapshot at decision
        for n in range(nsteps):
            t = n * dt
            d = (rng.random(B) < in_rate * dt).astype(float) if cue[0] <= t < cue[1] \
                else np.zeros(B)
            drive = np.zeros((B, 1, 1)); drive[:, 0, 0] = d * 1.0   # causal coincidence on the cue
            if no_trace:
                drive[:] = 0.0
            bank.step(drive)
            due = (n == ri) & (~captured)
            if due.any():
                e_state[due] = bank.vn[due, 0, 0, :]; captured[due] = True
        # read cascade state (device read noise) and form the action input per readout mode
        noisy = e_state / bank.Vnmax * (1.0 + read_noise * rng.standard_normal((B, k)))
        if no_trace:
            x = np.zeros((B, nin))
        elif readout == "scalar":
            x = noisy[:, -1:]                               # last stage only (aliased)
        else:
            x = noisy                                       # full vector
        # state-driven action: logits = x . M (+ exploration noise), softmax choice
        logits = np.einsum('bi,bia->ba', x, M) + sigma * rng.standard_normal((B, A))
        chosen = np.argmax(logits, 1)
        r_true = (chosen == which).astype(float)
        adv = (r_true - baseline)
        # REINFORCE-style local update of the readout->action map: push the chosen action's
        # weights toward the input when reward beats baseline (no gradient transport, local to x).
        oh = np.zeros((B, A)); oh[bidx, chosen] = 1.0
        M = M + eta * (adv[:, None, None] * x[:, :, None] * oh[:, None, :])
        M = np.clip(M, -5.0, 5.0)
        baseline += 0.02 * (r_true - baseline)
        rewards[:, tr] = r_true
    return rewards


def _final(rw, window=400):
    return rw[:, -window:].mean(1)


def aliasing_pair(*, k=3, tau_leak=10.0, dt=5e-3, cue_dur=0.5, in_rate=200.0,
                  reps=4000, tmax=40.0, seed=0):
    """Locate the last-stage occupancy peak t* and return a pair (tA, tB) that ALIASES
    the scalar readout: tA on the rising flank and tB on the falling flank, chosen to
    have an equal last-stage value. This is the condition the aliasing-rescue test
    REQUIRES -- the scalar reads the same value at both intervals, so only the full
    vector can separate them. Computed from the device model (not hand-set), so the
    pair tracks tau_leak and k."""
    rng = np.random.default_rng(seed)
    cue = (0.3, 0.3 + cue_dur)
    ts = np.arange(1.0, tmax, 0.5)
    last = np.empty(len(ts))
    for i, te in enumerate(ts):
        bank = GateBankBatched(reps, 1, 1, tau_leak=tau_leak, k=k, dt=dt)
        ri = int((cue[1] + te) / dt)
        bank.reset()
        for n in range(ri + 1):
            t = n * dt
            d = (rng.random(reps) < in_rate * dt).astype(float) if cue[0] <= t < cue[1] \
                else np.zeros(reps)
            drive = np.zeros((reps, 1, 1)); drive[:, 0, 0] = d
            bank.step(drive)
        last[i] = bank.vn[:, 0, 0, -1].mean() / bank.Vnmax
    pk = int(np.argmax(last)); tstar = float(ts[pk])
    tA = max(ts[0], tstar - 8.0)                       # a point on the rising flank
    vA = float(np.interp(tA, ts, last))
    fall = ts[ts > tstar]; lastfall = last[ts > tstar]
    tB = float(fall[np.argmin(np.abs(lastfall - vA))]) # falling-flank point of equal value
    return tA, tB, tstar


def _worker_main(readout, B, trials, tA, tB):
    return readout, run_interval(readout, B=B, trials=trials, tA=tA, tB=tB)


def _worker_sweep(job, B, trials):
    sep, k = job
    centre = 9.0
    tA, tB = centre - sep / 2, centre + sep / 2
    accs = {}
    for readout in ("scalar", "vector"):
        rw = run_interval(readout, B=B, k=k, trials=trials, tA=tA, tB=tB)
        accs[readout] = float(_final(rw).mean())
    return (sep, k), accs


def run_vector_timer(*, seeds=20, trials=3000, quick=False, pool=None):
    """Experiment 20 core: the cascade occupancy VECTOR as a clockless elapsed-time code.

    Serial by default (a notebook calls it in-kernel); returns the result grid as a
    plain dict, no file I/O / no plotting / no stdout.  Preserves the science of
    ``experiments/04_temporal_selectivity/vector_timer.py`` exactly.

    Runs the vector/scalar/no_trace comparison at the MEASURED scalar-aliasing pair
    (:func:`aliasing_pair`, so the pair tracks ``tau_leak``/``k``), then the
    (interval separation) x (cascade depth ``k``) characterisation sweep.  Evaluates
    the pre-registered H1--H4/K1.  ``quick`` shrinks the sweep grid exactly as the
    original ``--quick`` did.  An optional ``pool`` (a ``multiprocessing.Pool``) is
    used to parallelise the coarse axes when called from :func:`main`; when ``None``
    (the notebook path) everything runs serially in-kernel.
    """
    from functools import partial

    # --- main comparison at the MEASURED scalar-aliasing pair, vector vs scalar vs no_trace ---
    # The aliasing test requires tA, tB to give the SAME last-stage value on opposite sides of
    # the last-stage peak (else the scalar is not aliased and the test is vacuous). Earlier runs
    # hard-set tA=3,tB=15, but both sit on the last-stage rising flank (peak ~18.5s), so the
    # scalar was never aliased and H2 could not pass. The pair is now computed from the device
    # model so it tracks tau_leak/k.
    tA, tB, tstar = aliasing_pair(k=3, tau_leak=10.0)
    conds = ("vector", "scalar", "no_trace")
    if pool is not None:
        res = dict(pool.map(partial(_worker_main, B=seeds, trials=trials, tA=tA, tB=tB),
                            conds))
    else:
        res = {c: run_interval(c, B=seeds, trials=trials, tA=tA, tB=tB) for c in conds}
    finals = {c: _final(res[c]) for c in conds}
    ci = {c: bootstrap_ci(finals[c]) for c in conds}

    vec, sca, nt = finals["vector"].mean(), finals["scalar"].mean(), finals["no_trace"].mean()
    h1 = vec >= CRIT
    h2 = (sca <= CHANCE + 0.10) and (ci["vector"][0] > ci["scalar"][1])
    h2_weak = ci["vector"][0] > ci["scalar"][1]            # vector disjoint-above scalar (the real claim)
    h3 = nt <= CHANCE + 0.10
    k1 = ci["vector"][0] <= ci["scalar"][1]                # vector not disjoint above scalar

    # --- characterisation sweep: (interval separation) x (k); resolution noise- not k-limited ---
    seps = [12.0, 6.0, 3.0, 1.5] if not quick else [12.0, 3.0]
    ks = [3, 5, 8] if not quick else [3, 8]
    jobs = [(sep, k) for sep in seps for k in ks]
    if pool is not None:
        sweep_res = dict(pool.map(partial(_worker_sweep, B=seeds, trials=trials), jobs))
    else:
        sweep_res = dict(_worker_sweep(job, seeds, trials) for job in jobs)
    sweep = {}
    for sep in seps:
        for k in ks:
            sweep[(sep, k)] = sweep_res[(sep, k)]

    return {
        "finals": {c: finals[c] for c in conds}, "ci": ci,
        "sweep": {f"{s}_{k}": v for (s, k), v in sweep.items()},
        "seps": seps, "ks": ks, "tau_leak": 10.0, "tA": tA, "tB": tB, "tstar": tstar,
        "seeds": seeds, "trials": trials, "chance": CHANCE, "crit": CRIT,
        "criteria": {"H1": bool(h1), "H2": bool(h2), "H3": bool(h3), "K1": bool(k1)},
        # extra diagnostics preserved from the original stdout report (H2' weak claim)
        "h2_weak": bool(h2_weak),
    }


def main(argv=None):
    """Full-scale reproduction CLI for the timing/selectivity grids (writes ``data/results``).

    ``python -m mrl_trace.selectivity [--exp10] [--exp20] [--full|--quick]``
    With no experiment flag, runs both.  ``--full`` = 20 seeds (published); ``--quick``
    = a fast few-seed smoke run.  The vector-timer sweep parallelises its coarse axes
    with a process pool (only here, in the real ``__main__``); the ``run_*`` cores stay
    serial so a notebook can call them in-kernel.
    """
    import argparse
    import os
    from multiprocessing import Pool
    from . import paths

    ap = argparse.ArgumentParser(description="Timing/selectivity reproductions")
    ap.add_argument("--exp10", action="store_true",
                    help="interval selectivity -> exp10_interval.npy")
    ap.add_argument("--exp20", action="store_true",
                    help="cascade vector timer -> exp20_vector_timer.npy")
    ap.add_argument("--quick", action="store_true", help="fast few-seed smoke run")
    ap.add_argument("--full", action="store_true", help="published 20-seed run (default)")
    a = ap.parse_args(argv)
    run_all = not (a.exp10 or a.exp20)
    nproc = max(1, min(6, (os.cpu_count() or 4) - 2))

    if a.exp10 or run_all:
        seeds = 6 if a.quick else 20
        trials = 600 if a.quick else 1500
        print("=== Experiment 10: interval-selectivity ===")
        print(f"  V=0.9, tau_r={tau_r(0.9):.1f}s | {seeds} seeds, {trials} trials, "
              f"late-lag={LATE_LAG}s")
        grid = run_interval_selectivity(seeds=seeds, trials=trials, V=0.9)
        print(f"  preferred interval t* (= argmax device trace): "
              + ", ".join(f"tau={tl:g}->t*={grid['tstar'][tl]:.2f}s" for tl in grid["taus"]))
        (dlo, dhi), (elo, ehi) = grid["dev_ci"], grid["exp_ci"]
        print(f"  design point tau_leak=10s, D_rew=t*={grid['D0']}s: "
              f"device S={grid['dev_S'].mean():.3f} CI[{dlo:.3f},{dhi:.3f}]  "
              f"exp S={grid['exp_S'].mean():.3f} CI[{elo:.3f},{ehi:.3f}]")
        print(f"  criteria P1={grid['P1']} P2={grid['P2']} P3={grid['P3']} P4={grid['P4']}")
        paths.save_result("exp10_interval.npy", grid)
        print("  wrote exp10_interval.npy")

    if a.exp20 or run_all:
        seeds = 6 if a.quick else 20
        trials = 1500 if a.quick else 3000
        print("=== Experiment 20: cascade occupancy VECTOR as a clockless timer ===")
        print(f"  interval-discrimination bandit, tau_leak=10s, k=3; {seeds} seeds, "
              f"{trials} trials; chance={CHANCE}, criterion={CRIT}")
        with Pool(nproc) as pool:
            grid = run_vector_timer(seeds=seeds, trials=trials, quick=a.quick, pool=pool)
        print(f"  last-stage peak t*={grid['tstar']:.1f}s -> aliasing pair "
              f"tA={grid['tA']:.1f}s, tB={grid['tB']:.1f}s")
        for c in ("vector", "scalar", "no_trace"):
            lo, hi = grid["ci"][c]
            print(f"    {c:9s}: final {grid['finals'][c].mean():.3f}  CI[{lo:.3f}, {hi:.3f}]")
        print(f"  criteria={grid['criteria']}")
        paths.save_result("exp20_vector_timer.npy", grid)
        print("  wrote exp20_vector_timer.npy")


if __name__ == "__main__":
    main()
\end{lstlisting}

\begin{lstlisting}[caption={siox\_eligibility/extensions.py},label={lst:siox-src-extensions}]
r"""Chapter-8 cross-cutting studies: how the two slow device primitives (the
eligibility trace and eligibility-magnitude homeostasis) compose, and how far the
device-native reinforcement-learning story stretches.

Four experiments live here, all built on the validated ``GateBankBatched`` cascade
gate (Section 5.3.2 physics) and the sequential ``LinearTrack``/``TMaze`` MDPs -- none
re-derives a primitive, each COMPOSES the ones already in :mod:`mrl_trace.bandit`,
:mod:`mrl_trace.maze`, :mod:`mrl_trace.neurons` and
:mod:`mrl_trace.learning`:

- **Experiment 19 -- multi-timescale credit from the MEASURED retention spread**
  (:func:`run_multitimescale`).  A real ITO array does not have one ``tau_leak``, it has
  a DISTRIBUTION of them (the n=53 measured trap-discharge fits, median ~1.3 s, range
  0.56--7.7 s).  A mixed-delay task needs more than one credit horizon, so the measured
  retention POPULATION is the resource, not the defect.  NAIVELY the spread fails (the
  cascade readout magnitude scales ~linearly with ``tau_leak``, so a shared learning rate
  starves the short-tau synapses); a LOCAL, activity-driven, tau-BLIND eligibility-
  magnitude homeostasis -- the same CLASS of set-point negative feedback used for
  firing-rate stability -- equalises the gain and recovers full multi-timescale credit,
  matching a tau-aware oracle.  One material, two slow primitives interacting.

- **Experiment 21 -- beta_leak sensitivity of the load-bearing results**
  (:func:`run_beta_sensitivity`).  Every learning result integrates a SINGLE-RATE
  (mono-exponential, ``beta_leak=1``) leak; the device measures a DISPERSIVE
  (stretched-exponential) discharge, ``beta_leak ~ 0.85`` field-free (~0.54 held-bias).
  This re-runs the two load-bearing results -- the ``D_max ~= k*tau_leak`` law and the
  exp19 multi-timescale coupling -- at the measured ``beta_leak`` and checks they survive,
  closing the "measured 0.85 but simulated 1.0" gap.

- **Experiment 22 -- WM and STC on ONE device: a NULL on the two-site hypothesis**
  (:func:`run_wm_stc`, :func:`wm_isolated`).  Working memory (inference-time cue hold) and
  STC/eligibility (learning-time credit) occupy different brain sites.  Does a device
  circuit therefore NEED two co-fabricated traces, or can one serve both?  A DMS task with
  a stimulus-free WM delay and a distractor-bearing reward gap shows a SINGLE shared trace
  does not collapse (97.6% at 6 s, 100% at 10 s -- the retention the rule already runs at),
  so two devices are NOT needed; the credit role, not the hold, sets the required retention.

- **Experiment 23 -- a device-native temporal-difference / actor-critic rule**
  (:func:`run_device_td`).  REINFORCE-with-baseline (``dw = eta (R-b) e``) is extended to a
  value-based learner that BOOTSTRAPS.  Two device-native identifications make it native:
  the discount ``gamma = exp(-step_dur/tau_leak)`` is the trap-discharge decay over one
  step (one material constant sets both trace lifetime and RL horizon), and the critic
  ``V(s)`` is a second on-device weight bank read by the same crossbar product.  The TD
  error then gates the same signed three-factor write, ``dw = eta delta e``.

Every ``run_*`` core is SERIAL, import-light and returns plain arrays/dicts with no file
I/O -- notebooks call them in-kernel at a small seed/trial count.  The module-level
:func:`main` is the full-scale driver: it may parallelise the coarse axis with a
``multiprocessing.Pool`` (it runs as ``python -m``), computes the bootstrap CIs and the
pre-registered criteria, and writes each grid via :func:`mrl_trace.paths.save_result`.
NOTE: reversal learning (Experiment 18) is a ``train`` variant and lives in
:mod:`mrl_trace.bandit`, not here.
"""
from __future__ import annotations

import math
from functools import partial

import numpy as np

from .bandit import GateBankBatched, W_INIT, W_MAX
from .maze import LinearTrack, TMaze, train_sequential, reward_rate
from .neurons import lif_step_batched, TAU_M, V_TH
from .learning import LTD_BIAS
from .stats import bootstrap_ci

__all__ = [
    # exp19 -- multi-timescale credit from the measured retention spread
    "load_measured_tau",
    "run_multitimescale",
    # exp22 -- WM + STC on one device (null on two-site hypothesis)
    "run_wm_stc",
    "wm_isolated",
    # exp23 -- device-native TD / actor-critic
    "run_device_td",
    # exp21 -- beta_leak sensitivity of the load-bearing results
    "dmax_law",
    "run_beta_sensitivity",
    # shared constants
    "CHANCE",
    "CRIT",
    "TAU_BAND",
    "BETAS",
    "V_INIT",
    "V_MAX",
    "ETA_V",
]

# =============================================================================
# Shared constants
# =============================================================================
CHANCE = 0.5
CRIT = 0.75

#: exp22 measured ITO retention band (s), swept in the WM/STC study.
TAU_BAND = [0.8, 1.3, 2.0, 3.6, 6.0, 10.0]

#: exp21 dispersion exponents: single-rate / field-free operating / held-bias stress.
BETAS = [1.0, 0.85, 0.54]

#: exp23 on-device critic value store: one scalar V per state (a second weight bank).
V_INIT = 0.0
V_MAX = 1.0
ETA_V = 0.1                    # critic learning rate (value TD update)


# =============================================================================
# Experiment 19 -- multi-timescale credit from the MEASURED trap-discharge spread
#
# The scaling/fault studies (exp14) treat device-to-device spread in retention as a
# NON-IDEALITY to tolerate. This inverts that: the measured retention DISTRIBUTION is
# the resource for a mixed-delay task, not the defect. No competitor can claim this --
# a single-time-constant cell (Ming 2026) has one horizon; an engineered cascade of
# traces (Ralambomihanta 2025) DESIGNS the spread; here the spread is MEASURED, drawn
# from the device population itself.
# =============================================================================


def load_measured_tau():
    """The n=53 measured ITO trap-discharge time constants (s), from the device population.

    Reads ``data/device_model/ito_decay_data.npz`` (the measured trap-discharge fits;
    resolved through :func:`mrl_trace.paths.device_model_dir` so it works from any
    working directory). Falls back to the documented summary (median 1.34 s, range
    0.56--7.72 s) -- logged, never silent -- only if the npz is absent.
    """
    from . import paths
    try:
        d = np.load(paths.device_model_dir() / "ito_decay_data.npz")
        tau = np.asarray(d["tau"], float)
        tau = tau[np.isfinite(tau) & (tau > 0)]
        if tau.size >= 10:
            return tau, "measured (ito_decay_data.npz, n=%d)" % tau.size
    except Exception as exc:                                  # noqa: BLE001
        print(f"  [tau] measured npz unavailable ({type(exc).__name__}); using summary")
    # documented summary fallback (logged, never silent)
    rng = np.random.default_rng(0)
    tau = np.clip(np.exp(rng.normal(np.log(1.34), 0.6, 53)), 0.56, 7.72)
    return tau, "summary-proxy (median 1.34, range 0.56-7.72)"


def _peak_gain_per_tau(tau_row, dt, cue=0.5, Tmax=20.0):
    """Analytic peak eligibility for each retention in ``tau_row`` (1-D, length S+1): drive a
    single cascade gate over a ``cue``-long causal coincidence, then relax, and take the max
    readout. This is the tau-dependent magnitude (the nuisance gain) the oracle reference
    divides out. The homeostatic condition does NOT use this -- it estimates the gain from the
    synapse's own running activity instead."""
    gains = np.empty(len(tau_row))
    for i, tau in enumerate(tau_row):
        g = GateBankBatched(1, 1, 1, tau_leak=float(tau), dt=dt)
        for _ in range(int(cue / dt)):
            g.step(np.ones((1, 1, 1)))
        m = g.vn[0, 0, 0, -1]
        for _ in range(int(Tmax / dt)):
            g.step(np.zeros((1, 1, 1))); m = max(m, g.vn[0, 0, 0, -1])
        gains[i] = m
    return gains


def _assign_tau(kind, S, A, tau_pool, short_states, long_states, rng):
    """Per-synapse (S, A) retention array for each condition.

    hetero_measured : draw from the measured pool, then route the LONGER draws to the
                      long-delay states and the SHORTER draws to the short-delay states
                      (a programmed array can place a device of the right retention at the
                      right crosspoint; the claim under test is that the measured spread
                      SPANS both horizons).
    best_single     : a single homogeneous tau (value passed in tau_pool as a scalar).
    mean_tau        : homogeneous mean of the pool.
    """
    if kind == "best_single" or kind == "mean_tau":
        return np.full((S, A), float(tau_pool))
    # hetero_measured: sort the pool, give long states high-tau draws, short states low-tau
    draws = rng.choice(tau_pool, size=S, replace=True)
    order = np.argsort(draws)                                # ascending tau
    tau_state = np.empty(S)
    # short-delay states get the smallest draws, long-delay states the largest
    ns = len(short_states)
    tau_state[np.asarray(short_states)] = np.sort(draws)[:ns]
    tau_state[np.asarray(long_states)] = np.sort(draws)[ns:]
    return np.repeat(tau_state[:, None], A, axis=1)          # (S, A), same tau across actions


def run_multitimescale(kind, *, S=4, A=2, B=20, tau_arg=None, D_short=2.0, D_long=8.0,
                       t_distract_short=1.5, trials=3000, dt=5e-3, cue_dur=0.5, eta=0.2,
                       in_rate=200.0, ltd=LTD_BIAS, tau_m=TAU_M, v_th=V_TH, sigma=0.15,
                       seed0=0, tau_pool=None, elig_norm="none", tau_homeo=300.0,
                       beta_leak=1.0):
    """One condition on the TWO-HORIZON contextual bandit (Experiment 19). SERIAL.

    The diagnostic that motivates the design: on a pure survival task a long tau bridges every
    delay (highest surviving eligibility at every D), so heterogeneity cannot help. A genuine
    two-horizon task must make NO single tau optimal. We achieve that by giving each delay
    group a failure mode for the WRONG tau:

      * SHORT group (states 0..S/2-1): cue at t=0, reward at D_short, and an UNINFORMATIVE
        distractor coincidence near reward (t_distract_short). A too-LONG tau is low-pass /
        recency-weighted, so at reward the recent distractor dominates its trace and it
        MIS-CREDITS the distractor (the exp12 result). Only a tau whose band-pass peak sits
        near the cue lag credits the cue over the near-reward distractor. -> long tau FAILS.
      * LONG group (states S/2..S-1): cue at t=0, reward at D_long, NO distractor. The cue
        sits a long lag before reward, beyond a short trace's reach (survival diagnostic), so
        a too-SHORT tau has no surviving eligibility at reward. -> short tau FAILS.

    No single homogeneous tau serves both: short tau loses the long group (reach), long tau
    loses the short group (distractor). A heterogeneous population -- short-tau synapses on the
    short-group states, long-tau synapses on the long-group states -- serves both, which is
    exactly what the measured ITO retention SPREAD supplies. The reward-prediction-error gate
    and signed rule are unchanged from the base bandit.

    ``kind`` selects the tau assignment ("hetero_measured", "best_single", "mean_tau",
    "no_trace"); ``elig_norm`` selects the eligibility-magnitude normalisation
    ("none"/"homeo"/"oracle"). ``beta_leak`` sets the discharge dispersion (1 = single-rate,
    the exp21 sensitivity knob). Returns
    ``(rewards (B,trials), grp_rw dict of (B,trials) reward-or-nan-by-group,
    short_states, long_states)``."""
    rng = np.random.default_rng(seed0)
    DIST = S                                   # one extra distractor input line (not a state)
    short_states = list(range(0, S // 2))
    long_states = list(range(S // 2, S))
    is_long = np.zeros(S, bool); is_long[long_states] = True
    correct = np.array([s % A for s in range(S)])

    no_trace = (kind == "no_trace")
    if no_trace:
        tau_grid = np.full((S + 1, A), 10.0)
    elif kind == "hetero_measured":
        tau_grid = _assign_tau("hetero_measured", S, A, tau_pool, short_states, long_states, rng)
        tau_grid = np.vstack([tau_grid, tau_grid.mean(0, keepdims=True)])     # distractor row
    else:
        tau_grid = _assign_tau(kind, S, A, tau_arg, short_states, long_states, rng)
        tau_grid = np.vstack([tau_grid, tau_grid.mean(0, keepdims=True)])
    bank = GateBankBatched(B, S + 1, A, tau_leak=tau_grid, dt=dt, beta_leak=beta_leak)

    # Eligibility-magnitude normalisation (the fix for the tau-dependent gain).
    #   "none"   : raw eligibility (the naive heterogeneous-spread result).
    #   "homeo"  : divide each synapse's snapshotted eligibility by a SLOW running estimate of
    #              its OWN eligibility magnitude (gbar), updated once per trial. This is
    #              synaptic-scaling / divisive homeostasis -- a local, set-point negative
    #              feedback of the SAME CLASS as the firing-rate homeostasis used elsewhere,
    #              but on eligibility amplitude. It is BLIND to tau and to the task/reward: it
    #              only sees the synapse's own activity, so equalising the per-synapse gain is
    #              an emergent activity-driven property, NOT a tau lookup.
    #   "oracle" : divide by the precomputed analytic peak gain(tau) -- a tau-AWARE calibration,
    #              included only as the upper-bound reference the homeostasis should approach
    #              WITHOUT knowing tau.
    gbar = np.full((B, S + 1), 1e-3)             # per-row running |e| estimate (homeo state)
    oracle_gain = None
    if elig_norm == "oracle":
        oracle_gain = np.maximum(_peak_gain_per_tau(tau_grid[:, 0], dt), 1e-6)  # (S+1,)

    w = np.full((B, S + 1, A), W_INIT)
    w[:, DIST, :] = W_INIT                       # distractor synapses plastic too (can be mis-credited)
    baseline = np.full(B, 1.0 / A)
    cue = (0.3, 0.3 + cue_dur)
    ri_short = int((cue[1] + D_short) / dt)
    ri_long = int((cue[1] + D_long) / dt)
    nsteps = ri_long + 2
    ds0 = int((cue[1] + t_distract_short) / dt)         # distractor window (short trials only)
    ds1 = ds0 + int(0.3 / dt)
    rewards = np.zeros((B, trials))
    grp_rw = {"short": np.full((B, trials), np.nan), "long": np.full((B, trials), np.nan)}
    bidx = np.arange(B)

    for tr in range(trials):
        bank.reset()
        state = rng.integers(S, size=B)
        long_trial = is_long[state]
        ri = np.where(long_trial, ri_long, ri_short)
        v = np.zeros((B, A)); spk = np.zeros((B, A))
        e_rew = np.zeros((B, S + 1, A))
        captured = np.zeros(B, bool)
        for n in range(nsteps):
            t = n * dt
            pre = np.zeros((B, S + 1))
            if cue[0] <= t < cue[1]:
                pre[bidx, state] = (rng.random(B) < in_rate * dt).astype(float)
            # near-reward distractor on SHORT-group trials only (uninformative)
            if ds0 <= n < ds1:
                fire = (~long_trial) & (rng.random(B) < in_rate * dt)
                pre[fire, DIST] = 1.0
            charge = np.einsum('bsa,bs->ba', w, pre)
            v, sp = lif_step_batched(v, charge, dt, rng, tau_m=tau_m, v_th=v_th, noise=sigma)
            spk += sp
            drive = pre[:, :, None] * np.where(sp, 1.0, -ltd)[:, None, :]
            if no_trace:
                drive[:] = 0.0
            e = bank.step(drive)
            due = (n == ri) & (~captured)
            if due.any():
                e_rew[due] = e[due]; captured[due] = True
        # eligibility-magnitude normalisation (see header): equalise the tau-dependent gain
        if elig_norm == "homeo":
            row_mag = np.abs(e_rew).max(axis=2)                  # (B, S+1) this trial's row mag
            e_used = e_rew / (gbar[:, :, None] + 1e-6)
            active = row_mag > 1e-9                              # only adapt rows that fired
            gbar = np.where(active, gbar + (row_mag - gbar) / tau_homeo, gbar)
        elif elig_norm == "oracle":
            e_used = e_rew / oracle_gain[None, :, None]
        else:
            e_used = e_rew
        tie = spk.max(1) == spk.min(1)
        chosen = np.argmax(spk, 1); chosen[tie] = rng.integers(A, size=int(tie.sum()))
        r_true = (chosen == correct[state]).astype(float)
        adv = eta * (r_true - baseline)
        w = np.clip(w + adv[:, None, None] * e_used, 0.0, W_MAX)
        baseline += 0.02 * (r_true - baseline)
        rewards[:, tr] = r_true
        # record reward into the correct group's array (nan elsewhere -> clean group stats)
        grp_rw["long"][long_trial, tr] = r_true[long_trial]
        grp_rw["short"][~long_trial, tr] = r_true[~long_trial]
    return rewards, grp_rw, short_states, long_states


def _mt_final(rw, window=400):
    return rw[:, -window:].mean(1)


def _mt_group_final(grp_rw_g, window=800):
    """Per-seed mean reward over this group's trials in the last ``window`` trials.

    ``grp_rw_g`` is (B, trials) with the reward on this group's trials and nan elsewhere, so a
    nanmean over the tail window selects exactly the group's recent trials."""
    B, T = grp_rw_g.shape
    tail = grp_rw_g[:, T - window:]
    with np.errstate(invalid="ignore"):
        return np.array([np.nanmean(tail[b]) if np.any(~np.isnan(tail[b])) else np.nan
                         for b in range(B)])


def _mt_worker(cond, B, trials, tau_pool, best_single, mean_tau):
    """cond = (label, kind, elig_norm). ``kind`` picks the tau assignment, ``elig_norm`` the
    eligibility-magnitude normalisation (none/homeo/oracle). Top-level so it pickles for a Pool."""
    label, kind, elig_norm = cond
    arg = {"best_single": best_single, "mean_tau": mean_tau}.get(kind)
    rw, grp_rw, ss, ls = run_multitimescale(
        kind, B=B, trials=trials, tau_arg=arg, tau_pool=tau_pool, elig_norm=elig_norm)
    return label, rw, grp_rw


def _mt_grid_one(tau, trials):
    """Score one homogeneous tau (small fixed B -- only needs to RANK, not give CIs)."""
    rw, *_ = run_multitimescale("best_single", B=6, trials=trials, tau_arg=float(tau), seed0=0)
    return float(tau), _mt_final(rw).mean()


def _mt_grid_best_single(tau_pool, trials, pool=None):
    """Find the single homogeneous tau that maximises mean reward on the two-horizon task --
    the strongest one-horizon competitor. Ranking only, so a small B and short run; run the
    grid points in PARALLEL when a Pool is provided (the serial grid was the main bottleneck)."""
    grid = np.round(np.geomspace(max(0.5, tau_pool.min()), max(tau_pool.max(), 14.0), 6), 2)
    gt = min(trials, 1200)
    if pool is not None:
        scored = pool.map(partial(_mt_grid_one, trials=gt), [float(t) for t in grid])
    else:
        scored = [_mt_grid_one(float(t), gt) for t in grid]
    best_tau = max(scored, key=lambda x: x[1])[0]
    return best_tau, grid


# =============================================================================
# Experiment 22 -- WM and STC on ONE device: a NULL on the two-site hypothesis
# =============================================================================


def _wm_relax(bank, n_steps, dt, stride=10):
    """Drive-free coarse relaxation of a gate over a silent delay (exp12 helper, verbatim)."""
    if n_steps <= 0:
        shp = bank.e.shape if hasattr(bank, "e") else bank.vn.shape[:-1]
        return bank.step(np.zeros(shp))
    base = bank.dt; coarse = max(1, n_steps // stride)
    bank.dt = base * (n_steps / coarse)
    shp = bank.e.shape if hasattr(bank, "e") else bank.vn.shape[:-1]
    zero = np.zeros(shp); out = None
    for _ in range(coarse):
        out = bank.step(zero)
    bank.dt = base
    return out


def run_wm_stc(tau_wm, tau_stc, *, B=20, trials=2500, dt=5e-3, D_wm=3.0, D_stc=3.0,
               cue_dur=0.3, distract_dur=0.3, in_rate=200.0, eta=0.2, V=1.5,
               sigma=0.15, ltd=LTD_BIAS, seed0=0, shared_device=False):
    """DMS with a WM delay before the decision (Experiment 22). SERIAL. Returns per-seed
    reward ``(B, trials)``.

    Two-device (default): a dedicated cue-hold device (tau_wm, positive drive) supplies the
    WM readout, and a separate eligibility device (tau_stc, exp12 loop) supplies STC credit.

    shared_device=True -- the GENUINE one-device collapse test: a SINGLE trace bank per
    (sample,action) synapse must serve BOTH roles. It receives ONE drive (the STC pre*post
    coincidence, LTD-biased, as physics dictates -- a synapse's trace is set by its own
    plasticity), and the WM decision must read cue identity off that SAME trace. There is
    no separate positive cue-hold trace, because one device is one physical state. tau_wm is
    ignored (only tau_stc, the single device's retention, applies)."""
    rng = np.random.default_rng(seed0)
    S, A = 3, 2; DIST = 2
    stc = GateBankBatched(B, S, A, tau_leak=tau_stc, V=V, dt=dt)   # eligibility device
    wm = None if shared_device else GateBankBatched(B, 2, 1, tau_leak=tau_wm, V=V, dt=dt)
    w = np.full((B, S, A), W_INIT)
    baseline = np.full(B, 1.0 / A)
    bidx = np.arange(B)
    n_cue = int(round(cue_dur / dt)); n_wm = int(round(D_wm / dt))
    n_gap1 = int(round(0.5 * D_stc / dt)); n_dist = int(round(distract_dur / dt))
    n_gap2 = int(round((D_stc - 0.5 * D_stc - distract_dur) / dt))
    rewards = np.zeros((B, trials))

    for tr in range(trials):
        stc.reset()
        if wm is not None:
            wm.reset()
        cls = rng.integers(2, size=B)                 # sample class
        sline = np.where(cls == 0, 0, 1)
        v = np.zeros((B, A)); spk = np.zeros((B, A))
        e_cue = None                                   # shared-device cue snapshot (from stc)
        # --- SAMPLE window: drive the eligibility trace (exp12-style) ---
        for _ in range(n_cue):
            pre = np.zeros((B, S)); firing = (rng.random(B) < in_rate * dt).astype(float)
            pre[bidx, sline] = firing
            charge = np.einsum('bsa,bs->ba', w, pre)
            v, sp = lif_step_batched(v, charge, dt, rng, tau_m=TAU_M, v_th=V_TH, noise=sigma)
            # STC eligibility: pre*post coincidence with LTD bias (exp12 verbatim)
            e_stc = stc.step(pre[:, :, None] * np.where(sp, 1.0, -ltd)[:, None, :])
            if wm is not None:
                # SEPARATE WM device: positive cue coincidence only (clean hold, no LTD)
                dwm = np.zeros((B, 2, 1)); dwm[bidx, sline, 0] = firing
                wm.step(dwm)
        # --- WM DELAY: trace(s) relax silently; the cue must SURVIVE to the decision ---
        if wm is not None:
            wm_held = _wm_relax(wm, n_wm, dt)          # dedicated cue trace
            cuevec = wm_held[:, :, 0]                  # (B,2)
            _wm_relax(stc, n_wm, dt)
        else:
            # ONE device: the cue can only be read off the SAME eligibility trace. Its value
            # at each (sample,action) synapse is the surviving eligibility, whose sign/scale
            # were set by the plasticity rule during the sample window -- NOT a clean cue.
            e_cue = _wm_relax(stc, n_wm, dt)           # (B,S,A) surviving eligibility
            cuevec = e_cue[:, :2, :].max(axis=2)       # best available cue readout off the trace
        # --- DECISION: read from the held cue (dedicated device or shared trace) ---
        logit = np.einsum('bs,bsa->ba', cuevec, w[:, :2, :]) + sigma * rng.standard_normal((B, A))
        chosen = np.argmax(logit, axis=1)
        # tag the chosen action's eligibility at the decision
        dstc = np.zeros((B, S, A)); dstc[bidx, sline, chosen] = 1.0
        stc.step(dstc)
        # --- ACTION->REWARD gap with a DISTRACTOR competing for credit ---
        _wm_relax(stc, n_gap1, dt)
        for _ in range(n_dist):
            pd = (rng.random(B) < in_rate * dt).astype(float)
            dd = np.zeros((B, S, A)); dd[:, DIST, :] = pd[:, None]
            stc.step(dd)
        e_rew = _wm_relax(stc, n_gap2, dt)              # surviving eligibility at reward
        # --- three-factor update ---
        r = (chosen == cls).astype(float)
        w = np.clip(w + (eta * (r - baseline))[:, None, None] * e_rew, 0.0, W_MAX)
        baseline += 0.02 * (r - baseline)
        rewards[:, tr] = r
    return rewards


def wm_isolated(tau_wm, *, B=20, trials=600, dt=5e-3, D_wm=3.0, cue_dur=0.3,
                in_rate=200.0, sigma=0.05, seed=1):
    """WM alone: weights FROZEN to the correct bijection; only WM retention varies. SERIAL.
    Returns per-seed accuracy (length B)."""
    rng = np.random.default_rng(seed); S = 2; A = 2
    wm = GateBankBatched(B, S, 1, tau_leak=tau_wm, V=1.5, dt=dt)
    cmap = np.stack([rng.permutation(A) for _ in range(B)])
    w = np.zeros((B, S, A))
    for b in range(B):
        for s in range(S): w[b, s, cmap[b, s]] = 1.0
    n_cue = int(round(cue_dur / dt)); n_wm = int(round(D_wm / dt)); bidx = np.arange(B)
    hit = np.zeros((B, trials))
    for tr in range(trials):
        wm.reset(); sample = rng.integers(S, size=B)
        for _ in range(n_cue):
            dwm = np.zeros((B, S, 1)); dwm[bidx, sample, 0] = (rng.random(B) < in_rate * dt).astype(float)
            wm.step(dwm)
        held = _wm_relax(wm, n_wm, dt)
        logit = np.einsum('bs,bsa->ba', held[:, :, 0], w) + sigma * rng.standard_normal((B, A))
        hit[:, tr] = (np.argmax(logit, 1) == cmap[bidx, sample]).astype(float)
    return hit.mean(axis=1)                             # per-seed accuracy (length B)


def _wm_final(rw, window=500):
    """Per-seed final performance (mean over the last ``window`` trials) -> length-B."""
    return rw[:, -window:].mean(axis=1)


def _wm_fmt(vec):
    """Mean [lo, hi] percentage string from a per-seed vector, bootstrap 95% CI."""
    lo, hi = bootstrap_ci(vec)
    return f"{100*vec.mean():5.1f}% [{100*lo:4.1f}, {100*hi:4.1f}]"


# =============================================================================
# Experiment 23 -- a device-native temporal-difference / actor-critic rule
# =============================================================================


def _td_relax(bank, n_steps, dt, stride=10):
    """Drive-free coarse relaxation over a silent delay (exp12/21 helper, verbatim)."""
    if n_steps <= 0:
        shp = bank.vn.shape[:-1]
        return bank.step(np.zeros(shp))
    base = bank.dt
    coarse = max(1, n_steps // stride)
    bank.dt = base * (n_steps / coarse)
    shp = bank.vn.shape[:-1]
    zero = np.zeros(shp)
    out = None
    for _ in range(coarse):
        out = bank.step(zero)
    bank.dt = base
    return out


def run_device_td(scheme, *, L=5, B=20, episodes=2000, dt=5e-3, D=2.0, step_dur=0.4,
                  tau_leak=10.0, eta=0.2, V=1.5, in_rate=200.0, sigma=0.15, ltd=LTD_BIAS,
                  tau_homeo=200.0, max_steps=None, seed0=0):
    """Train the LinearTrack corridor under one of three reward-modulated schemes
    (Experiment 23). SERIAL.

    ``scheme`` in {"reinforce", "td_actor_critic", "td_no_homeo", "no_trace"}. Returns
    ``(rewards (B, episodes), V(s) (B, S))``. At each of ``L`` states the agent chooses
    forward/back; the eligibility gate is snapshotted per decision so each step carries
    its own trace. REINFORCE gates every step's write by the whole-episode return minus a
    scalar baseline; TD gates step ``t`` by the bootstrapped error
    ``delta_t = r_t + gamma V(s_{t+1}) - V(s_t)`` with an on-device critic ``V(s)``. The
    device-native discount ``gamma = exp(-step_dur / tau_leak)`` is the trap-discharge decay
    over one decision step -- the SAME tau_leak that sets the eligibility trace, so one
    material constant sets both the trace lifetime and the RL horizon. ``td_no_homeo`` is the
    same TD rule with the eligibility-magnitude homeostasis of exp19 removed (the deadly-triad
    check); ``no_trace`` zeroes eligibility (device-necessity control, anchors chance).
    """
    rng = np.random.default_rng(seed0)
    env = LinearTrack(L=L)
    S, A = env.n_states, env.n_actions
    bidx = np.arange(B)
    bank = GateBankBatched(B, S, A, tau_leak=tau_leak, V=V, dt=dt)
    w = np.full((B, S, A), W_INIT)                 # policy weights (actor)
    Vval = np.full((B, S), V_INIT)                 # on-device critic V(s), per state
    baseline = np.full(B, 1.0 / A)                 # scalar baseline (reinforce only)
    ghom = np.ones((B, S, A))                      # eligibility-magnitude homeostasis EMA
    # device-native discount: the trap-discharge decay over one decision step
    gamma = math.exp(-step_dur / tau_leak)
    n_cue = int(round(0.3 / dt))
    reward_lag = int(round(D / dt))
    if max_steps is None:
        max_steps = 3 * L                          # allow back-steps, bounded
    rewards = np.zeros((B, episodes))

    n_step = int(round(step_dur / dt))              # inter-decision interval in ticks
    for ep in range(episodes):
        pos = env.start(B)
        done = np.zeros(B, bool)
        got_reward = np.zeros(B, bool)
        bank.reset()                                 # ONE gate carries all decisions' traces
        traj_state, traj_active = [], []
        for st in range(max_steps):
            state = pos.copy()
            active = (~done).astype(float)
            v = np.zeros((B, A)); spk = np.zeros((B, A))
            for n in range(n_cue):
                pre = np.zeros((B, S))
                pre[bidx, state] = (rng.random(B) < in_rate * dt).astype(float)
                charge = np.einsum('bsa,bs->ba', w, pre)
                v, sp = lif_step_batched(v, charge, dt, rng, tau_m=TAU_M, v_th=V_TH, noise=sigma)
                spk += sp
                drive = np.zeros((B, S, A))
                if scheme != "no_trace":
                    drv = pre[bidx, state][:, None] * np.where(sp, 1.0, -ltd)
                    drive[bidx, state, :] = drv * active[:, None]
                bank.step(drive)
            tie = spk.max(1) == spk.min(1)
            chosen = np.argmax(spk, 1)
            chosen[tie] = rng.integers(A, size=int(tie.sum()))
            chosen[done] = 0
            traj_state.append(state); traj_active.append(active)
            nxt, reached, ep_done = env.step(pos, chosen)
            got_reward |= ((~done) & reached)
            pos = np.where(done, pos, nxt)
            done |= ep_done
            if done.all():
                break
            _td_relax(bank, n_step, dt)              # decay one inter-decision interval only
        # Goal reward gates the surviving trace, read ONCE after the action->reward delay.
        # Earlier decisions have decayed longer (they are further from the goal), so the
        # per-(state,action) eligibility at reward carries the correct distance-weighting;
        # this replaces the per-step full-window relax at a fraction of the cost.
        _td_relax(bank, reward_lag, dt)
        e_rew = (bank.vn[..., -1] / bank.Vnmax)      # surviving trace on ALL used synapses
        R = got_reward.astype(float)
        T = len(traj_state)

        # Each decision t owns the (state, action) rows of the shared gate; its surviving
        # eligibility is e_rew masked to that decision's state. One goal reward gates them.
        if scheme == "no_trace":
            pass                                                # trace zeroed -> no learning (chance)
        elif scheme == "reinforce":
            # whole-episode return minus a scalar baseline, applied to every used synapse
            adv = eta * (R - baseline)
            w = np.clip(w + adv[:, None, None] * e_rew, 0.0, W_MAX)
            baseline += 0.02 * (R - baseline)
        else:
            # device-native TD(lambda) / actor-critic, swept along the trajectory. r_t = R
            # only on the terminal (goal) step, 0 otherwise; V(s_{t+1}) bootstraps (0 past
            # the terminal step). The SAME tau_leak sets both the trace and the discount
            # gamma. delta_t gates the per-decision write, masked to that decision's state.
            if scheme == "td_actor_critic":
                ghom += (np.abs(e_rew) - ghom) / tau_homeo
                e_use = e_rew / np.maximum(ghom, 1e-3)
            else:                                               # td_no_homeo
                e_use = e_rew
            for t in range(T):
                s_t = traj_state[t]
                v_st = Vval[bidx, s_t]
                v_next = Vval[bidx, traj_state[t + 1]] if t + 1 < T else np.zeros(B)
                r_t = R if t == T - 1 else np.zeros(B)          # reward only at the goal step
                delta = r_t + gamma * v_next - v_st             # bootstrapped TD error
                mask = np.zeros((B, S))                         # this decision's state rows
                mask[bidx, s_t] = traj_active[t]
                gate = (eta * delta)[:, None, None] * mask[:, :, None]
                w = np.clip(w + gate * e_use, 0.0, W_MAX)
                Vval[bidx, s_t] = np.clip(v_st + ETA_V * delta * traj_active[t], 0.0, V_MAX)
        rewards[:, ep] = R
    return rewards, Vval


def _td_final(rw, window=200):
    return rw[:, -window:].mean(axis=1)


def _td_fmt(v):
    lo, hi = bootstrap_ci(v)
    return f"{v.mean():.3f} [{lo:.3f}, {hi:.3f}]"


def _td_worker(spec):
    """Run one independent (scheme, L) cell. Top-level so it pickles for the Pool; each
    cell is seed-deterministic (seed0=0), so the parallel result is identical to serial."""
    sc, L, B, episodes = spec
    rw, Vval = run_device_td(sc, L=L, B=B, episodes=episodes)
    return (sc, L), _td_final(rw), rw.mean(0), Vval


# =============================================================================
# Experiment 21 -- beta_leak sensitivity of the two load-bearing learning results
#
# All learning results integrate a SINGLE-RATE (mono-exponential, beta_leak=1) leak; the
# device measures a DISPERSIVE (stretched-exponential) discharge, beta_leak ~ 0.85 field-free
# (~0.54 held-bias). Re-run the D_max ~= k*tau law (A) and the exp19 coupling (B) at the
# measured beta_leak and check they survive -- closing the "measured 0.85 but simulated 1.0" gap.
# =============================================================================


def _dmax_one(job):
    """One (tau, beta) cell: D_max = largest delay whose reward rate >= criterion.
    Top-level so it pickles for the Pool."""
    tau, beta, seeds, episodes, delays = job
    crit = 0.75
    dmax = 0.0
    for D in delays:
        env = TMaze(L=3, A_goal=2)
        r = train_sequential(env, B=seeds, tau_leak=tau, D=D, episodes=episodes,
                             beta_leak=beta, seed0=0)
        if reward_rate(r, window=100).mean() >= crit:
            dmax = D
        else:
            break                                   # first failure = the crossing (monotone)
    return (tau, beta), dmax


def dmax_law(betas, seeds, episodes, pool=None):
    """(A) The D_max ~= k*tau_leak scaling law re-fit at each beta_leak (Experiment 21).

    For each beta, sweep the retention taus x delays, take D_max = largest delay whose reward
    rate reaches criterion, and origin-fit D_max = k*tau (report k, R^2, monotonicity). A
    ``multiprocessing.Pool`` may be passed to parallelise the (tau, beta) cells; otherwise the
    cells run serially. Returns a dict keyed by beta."""
    taus = [1.0, 2.0, 5.0, 10.0, 20.0]
    delays = [2, 5, 10, 20, 40, 80, 160]
    jobs = [(tau, b, seeds, episodes, delays) for b in betas for tau in taus]
    if pool is not None:
        res = dict(pool.map(_dmax_one, jobs))
    else:
        res = dict(_dmax_one(j) for j in jobs)
    out = {}
    for b in betas:
        d = np.array([res[(t, b)] for t in taus]); t = np.array(taus)
        # origin fit D_max = k*tau
        k = float(np.sum(t * d) / np.sum(t * t)) if np.sum(t * t) > 0 else float("nan")
        pred = k * t
        ss = 1 - np.sum((d - pred) ** 2) / max(np.sum((d - d.mean()) ** 2), 1e-9)
        mono = bool(np.all(np.diff(d) >= -1e-9))
        out[b] = {"taus": taus, "dmax": d.tolist(), "k": k, "r2": float(ss), "monotone": mono}
    return out


def _exp19_at_beta(job):
    """Run exp19's hetero_raw / hetero_homeo / best_single at a given beta_leak by passing
    ``beta_leak`` straight through :func:`run_multitimescale` (both live in this module now,
    so the old importlib/monkeypatch construction is unnecessary). Top-level so it pickles
    for the Pool. Returns ``(beta, {label: (mean, (lo, hi))})``."""
    beta, seeds, trials, tau_pool, best_tau = job
    res = {}
    for label, kind, norm in [("hetero_raw", "hetero_measured", "none"),
                              ("hetero_homeo", "hetero_measured", "homeo"),
                              ("best_single", "best_single", "none")]:
        arg = best_tau if kind == "best_single" else None
        rw, *_ = run_multitimescale(kind, B=seeds, trials=trials, tau_arg=arg,
                                    tau_pool=tau_pool, elig_norm=norm, beta_leak=beta)
        f = rw[:, -400:].mean(1)
        res[label] = (float(f.mean()), bootstrap_ci(f))
    return beta, res


def run_beta_sensitivity(*, betas=BETAS, seeds=12, episodes=2000, trials=3000, pool=None):
    """The full beta_leak sensitivity study (Experiment 21): (A) the D_max law and (B) the
    exp19 multi-timescale coupling, each re-run at every ``beta`` in ``betas``, plus the
    pre-registered S1/S2 criteria evaluated at the field-free operating beta=0.85. SERIAL by
    default; pass a ``multiprocessing.Pool`` to parallelise the coarse cells (main() does).

    PRE-REGISTRATION (fixed before running; reported as-is):
      S1  D_max law SURVIVES at beta=0.85: still monotone D_max(tau), origin-fit R^2 >= 0.90,
          and the slope k within ~30% of the beta=1 value (the law is not a single-rate artefact).
      S2  exp19 coupling SURVIVES at beta=0.85: hetero+homeo still >= 0.75 AND still
          disjoint-above naive hetero (the multi-timescale homeostasis is not a single-rate
          artefact).
      S3  (DESCRIPTIVE) report the same at the strongly-dispersive held-bias beta=0.54 -- the
          stress point, expected to degrade (that regime is NOT the operating point).

    Returns the grid dict (dmax law per beta, exp19 per beta, betas, seeds, tau_source, S1/S2).
    """
    tau_pool, src = load_measured_tau()

    # (A) D_max ~= k*tau at each beta_leak
    dmax = dmax_law(betas, seeds, episodes, pool=pool)

    # (B) exp19 multi-timescale + homeostasis at each beta_leak. Rank the best single tau
    # first (ranking only, small B / short run -- run serially, it is cheap).
    grid = np.round(np.geomspace(max(0.5, tau_pool.min()), max(tau_pool.max(), 14.0), 6), 2)
    scored = [(float(t), run_multitimescale("best_single", B=6, trials=min(trials, 1200),
                                            tau_arg=float(t), tau_pool=tau_pool)[0]
               [:, -400:].mean()) for t in grid]
    best_tau = max(scored, key=lambda x: x[1])[0]
    jobs = [(b, seeds, trials, tau_pool, best_tau) for b in betas]
    if pool is not None:
        e19res = dict(pool.map(_exp19_at_beta, jobs))
    else:
        e19res = dict(_exp19_at_beta(j) for j in jobs)

    # criteria at the OPERATING beta=0.85 (present only if 0.85 is in the sweep)
    s1 = s2 = None
    if 1.0 in dmax and 0.85 in dmax:
        k1, k085 = dmax[1.0]["k"], dmax[0.85]["k"]
        s1 = bool(dmax[0.85]["r2"] >= 0.90 and dmax[0.85]["monotone"]
                  and abs(k085 - k1) <= 0.30 * k1)
    if 0.85 in e19res:
        h = e19res[0.85]
        s2 = bool(h["hetero_homeo"][0] >= 0.75 and h["hetero_homeo"][1][0] > h["hetero_raw"][1][1])

    return {"dmax": dmax, "exp19": {str(b): e19res[b] for b in betas},
            "betas": list(betas), "seeds": seeds, "tau_source": src,
            "best_tau": best_tau,
            "criteria": {"S1": s1, "S2": s2}}


# =============================================================================
# Full-scale reproduction CLI (writes data/results). Pool ONLY here.
# =============================================================================


def main(argv=None):
    """Full-scale reproduction CLI for the Chapter-8 extension grids (writes ``data/results``).

    ``python -m mrl_trace.extensions [--exp19] [--exp21] [--exp22] [--exp23] [--quick|--full]``
    With no experiment flag, runs all four. ``--full`` = published seed count; ``--quick`` =
    a fast few-seed smoke run. Each writes its published grid filename via
    :func:`mrl_trace.paths.save_result`:
      exp19 -> exp19_multitimescale.npy   exp21 -> exp21_beta_sensitivity.npy
      exp22 -> exp22_wm_stc.npy           exp23 -> exp23_device_td.npy
    """
    import argparse
    import os
    import time
    from multiprocessing import Pool
    from . import paths

    ap = argparse.ArgumentParser(description="Chapter-8 cross-cutting RL extension studies")
    ap.add_argument("--exp19", action="store_true",
                    help="multi-timescale credit from measured spread -> exp19_multitimescale.npy")
    ap.add_argument("--exp21", action="store_true",
                    help="beta_leak sensitivity of load-bearing results -> exp21_beta_sensitivity.npy")
    ap.add_argument("--exp22", action="store_true",
                    help="WM+STC one-vs-two-device null test -> exp22_wm_stc.npy")
    ap.add_argument("--exp23", action="store_true",
                    help="device-native TD/actor-critic -> exp23_device_td.npy")
    ap.add_argument("--quick", action="store_true", help="fast few-seed smoke run")
    ap.add_argument("--full", action="store_true", help="published seed count (default)")
    a = ap.parse_args(argv)
    run_all = not (a.exp19 or a.exp21 or a.exp22 or a.exp23)
    quick = a.quick
    nproc = max(1, min(6, (os.cpu_count() or 4) - 2))

    # ---------------- Experiment 19 ----------------
    if a.exp19 or run_all:
        B = 6 if quick else 20
        trials = 1500 if quick else 3000
        S, A = 4, 2
        # Delays chosen to STRADDLE the measured tau band-pass peaks (tau~1s peaks at ~1.75s,
        # tau~4-8s peaks at ~7-15s), so the short delay is best served by a short measured tau
        # and the long delay by a long one, and no single tau matches both.
        D_short, D_long = 2.0, 8.0
        tau_pool, src = load_measured_tau()
        mean_tau = float(np.mean(tau_pool))
        print("Experiment 19: multi-timescale credit from the measured trap-discharge spread")
        print(f"  tau source = {src}; median={np.median(tau_pool):.2f}s mean={mean_tau:.2f}s "
              f"range=[{tau_pool.min():.2f},{tau_pool.max():.2f}]s")
        print(f"  mixed-delay bandit: S={S} (half D_short={D_short}s, half D_long={D_long}s), A={A}")
        print(f"  {B} seeds, {trials} trials; chance={CHANCE} crit={CRIT}; pre-registered H1-H5/K1\n")
        # (label, kind, elig_norm). Headline: hetero_raw (naive spread, fails) vs hetero_homeo
        # (spread + eligibility-magnitude homeostasis, recovers) vs hetero_oracle (tau-aware
        # upper bound). best_single is the strongest single-horizon device; no_trace the control.
        conds = [
            ("hetero_raw",    "hetero_measured", "none"),
            ("hetero_homeo",  "hetero_measured", "homeo"),
            ("hetero_oracle", "hetero_measured", "oracle"),
            ("best_single",   "best_single",     "none"),
            ("no_trace",      "no_trace",        "none"),
        ]
        labels = [c[0] for c in conds]
        with Pool(nproc) as pool:
            print("  grid-searching the best single homogeneous tau (parallel) ...")
            best_tau, grid = _mt_grid_best_single(tau_pool, trials, pool=pool)
            print(f"  best_single tau = {best_tau}s  (grid {list(grid)})\n")
            res = pool.map(partial(_mt_worker, B=B, trials=trials, tau_pool=tau_pool,
                                   best_single=best_tau, mean_tau=mean_tau), conds)
        raw = {label: (rw, grp_rw) for label, rw, grp_rw in res}

        finals, ci, grp = {}, {}, {}
        for label in labels:
            rw, grp_rw = raw[label]
            f = _mt_final(rw); finals[label] = f; ci[label] = bootstrap_ci(f)
            gs = _mt_group_final(grp_rw["short"]); gl = _mt_group_final(grp_rw["long"])
            grp[label] = (float(np.nanmean(gs)), float(np.nanmean(gl)))
            print(f"  {label:14s}: overall {f.mean():.3f}[{ci[label][0]:.3f},{ci[label][1]:.3f}]  "
                  f"short-grp {grp[label][0]:.3f}  long-grp {grp[label][1]:.3f}")

        raw_m = finals["hetero_raw"].mean()
        hom_m = finals["hetero_homeo"].mean()
        ora_m = finals["hetero_oracle"].mean()
        bst_m = finals["best_single"].mean()
        nt_m = finals["no_trace"].mean()
        # H1: naive measured spread does NOT solve the two-horizon task (the problem to fix).
        h1 = raw_m < CRIT
        # H2: eligibility-magnitude HOMEOSTASIS recovers it -- reaches criterion AND disjoint
        #     above the naive spread (the mechanism genuinely helps, blind to tau).
        h2 = (hom_m >= CRIT) and (ci["hetero_homeo"][0] > ci["hetero_raw"][1])
        # H3: tau-blind homeostasis matches OR exceeds the tau-aware oracle (within a small margin).
        h3 = ci["hetero_homeo"][0] >= ci["hetero_oracle"][0] - 0.05
        # H4: the homeostasis-equipped spread beats the best single tau (the multi-timescale payoff).
        h4 = ci["hetero_homeo"][0] > ci["best_single"][1]
        # H5: trace necessary.
        h5 = nt_m <= CHANCE + 0.10
        k1c = hom_m <= raw_m                                 # kill: homeostasis fails to help
        print("\n  === pre-registered criteria ===")
        print(f"  H1 naive measured spread fails the two-horizon task (< {CRIT}): "
              f"{'PASS' if h1 else 'FAIL'} ({raw_m:.3f})")
        print(f"  H2 eligibility HOMEOSTASIS recovers it (>= {CRIT}, disjoint > naive): "
              f"{'PASS' if h2 else 'FAIL'} ({hom_m:.3f} vs naive {raw_m:.3f})")
        print(f"  H3 tau-blind homeostasis matches tau-aware oracle (CIs overlap): "
              f"{'PASS' if h3 else 'FAIL'} ({hom_m:.3f} vs oracle {ora_m:.3f})")
        print(f"  H4 homeostasis-spread beats best single tau (disjoint): "
              f"{'PASS' if h4 else 'FAIL'} ({hom_m:.3f} vs best {bst_m:.3f})")
        print(f"  H5 no-trace fails (<= {CHANCE+0.10:.2f}): {'PASS' if h5 else 'FAIL'} ({nt_m:.3f})")
        print(f"  K1 homeostasis does NOT help (homeo <= naive): {'KILL' if k1c else 'ok'} "
              f"({hom_m:.3f} vs {raw_m:.3f})")
        print("\n  HEADLINE: the measured tau spread is unusable raw (the tau-dependent eligibility")
        print("  gain starves the short-tau synapses), but a LOCAL, ACTIVITY-DRIVEN eligibility-")
        print("  magnitude homeostasis -- the same CLASS of set-point negative feedback used for")
        print("  firing-rate stability, here on eligibility amplitude and BLIND to tau -- equalises")
        print("  the gain and recovers full multi-timescale credit, matching the tau-aware oracle.")

        paths.save_result("exp19_multitimescale.npy",
                          {"finals": {k: finals[k] for k in labels}, "ci": ci, "grp": grp,
                           "tau_source": src, "tau_pool": tau_pool, "best_tau": best_tau,
                           "mean_tau": mean_tau, "S": S, "A": A, "D_short": D_short,
                           "D_long": D_long, "seeds": B, "trials": trials, "chance": CHANCE,
                           "crit": CRIT,
                           "criteria": {"H1": bool(h1), "H2": bool(h2), "H3": bool(h3),
                                        "H4": bool(h4), "H5": bool(h5), "K1": bool(k1c)}})
        print("\n  wrote exp19_multitimescale.npy")

    # ---------------- Experiment 21 ----------------
    if a.exp21 or run_all:
        seeds = 6 if quick else 12
        episodes = 1200 if quick else 2000
        trials = 1500 if quick else 3000
        betas = [1.0, 0.85] if quick else BETAS
        print("\nExperiment 21: beta_leak sensitivity of the load-bearing learning results")
        print(f"  betas={betas}; D_max law + exp19 coupling; {seeds} seeds\n")
        with Pool(nproc) as pool:
            print("  (A) D_max ~= k*tau at each beta_leak ...")
            grid = run_beta_sensitivity(betas=betas, seeds=seeds, episodes=episodes,
                                        trials=trials, pool=pool)
        dmax = grid["dmax"]
        for b in betas:
            r = dmax[b]
            print(f"    beta={b:.2f}: D_max={r['dmax']}  k={r['k']:.1f}  R2={r['r2']:.3f}  "
                  f"monotone={r['monotone']}")
        print(f"\n  best_single tau = {grid['best_tau']}s")
        print("  (B) exp19 multi-timescale + homeostasis at each beta_leak ...")
        for b in betas:
            r = grid["exp19"][str(b)]
            print(f"    beta={b:.2f}: hetero_raw={r['hetero_raw'][0]:.3f}  "
                  f"hetero_homeo={r['hetero_homeo'][0]:.3f} {r['hetero_homeo'][1]}  "
                  f"best_single={r['best_single'][0]:.3f}")
        s1, s2 = grid["criteria"]["S1"], grid["criteria"]["S2"]
        print("\n  === pre-registered criteria (at the field-free operating beta=0.85) ===")
        if s1 is not None:
            k1v, k085 = dmax[1.0]["k"], dmax[0.85]["k"]
            print(f"  S1 D_max law survives (monotone, R2>=0.90, k within 30% of beta=1): "
                  f"{'PASS' if s1 else 'FAIL'} (k {k1v:.1f}->{k085:.1f}, R2={dmax[0.85]['r2']:.3f})")
        if s2 is not None:
            h = grid["exp19"]["0.85"]
            print(f"  S2 exp19 coupling survives (homeo>=0.75 & disjoint above naive): "
                  f"{'PASS' if s2 else 'FAIL'} (homeo {h['hetero_homeo'][0]:.3f} vs naive "
                  f"{h['hetero_raw'][0]:.3f})")
        print(f"  S3 (descriptive) held-bias beta=0.54 reported above as the stress point.")
        # store beta keys as strings for the dmax dict too (npy pickle round-trips fine either way)
        paths.save_result("exp21_beta_sensitivity.npy",
                          {"dmax": grid["dmax"], "exp19": grid["exp19"],
                           "betas": grid["betas"], "seeds": grid["seeds"],
                           "tau_source": grid["tau_source"],
                           "criteria": {"S1": bool(s1) if s1 is not None else None,
                                        "S2": bool(s2) if s2 is not None else None}})
        print("\n  wrote exp21_beta_sensitivity.npy")
        print("  HEADLINE: if S1+S2 PASS, the single-rate (beta=1) idealisation is SAFE at the")
        print("  measured field-free beta~0.85 -- the learning results are not artefacts of assuming")
        print("  a mono-exponential discharge, closing the 'measured 0.85 but simulated 1.0' gap.")

    # ---------------- Experiment 22 ----------------
    if a.exp22 or run_all:
        B = 8 if quick else 20
        trials = 800 if quick else 2500
        t0 = time.time()
        print(f"\n=== exp22: WM + STC on the device trace -- one-vs-two-device NULL test, DMS "
              f"chance 50% | B={B}, trials={trials}, bootstrap 95% CI ===")
        res = {"B": B, "trials": trials, "tau_band": TAU_BAND, "chance": 0.5,
               "wm": {}, "stc": {}, "two_cotuned": {}, "one_shared": {}, "two_device": {}}

        print("\n[diag 1] WM isolated (weights frozen correct) -- perf vs tau_wm:")
        for tw in TAU_BAND:
            v = wm_isolated(tw, B=B, trials=max(400, trials // 4))
            res["wm"][tw] = v
            print(f"    tau_wm={tw:4.1f}: {_wm_fmt(v)}   ({time.time()-t0:.0f}s)")

        print("\n[diag 2] STC via exp12 loop, WM long (tau_wm=10) -- learning perf vs tau_stc:")
        for ts in TAU_BAND:
            v = _wm_final(run_wm_stc(10.0, ts, B=B, trials=trials))
            res["stc"][ts] = v
            print(f"    tau_stc={ts:4.1f}: {_wm_fmt(v)}   ({time.time()-t0:.0f}s)")

        print("\n[two co-tuned devices] tau_wm=tau_stc, but TWO separate traces (control):")
        for t in TAU_BAND:
            v = _wm_final(run_wm_stc(t, t, B=B, trials=trials, shared_device=False))
            res["two_cotuned"][t] = v
            print(f"    tau={t:4.1f}: {_wm_fmt(v)}   ({time.time()-t0:.0f}s)")

        print("\n[COLLAPSE TEST] ONE shared trace serving BOTH roles (shared_device=True):")
        for t in TAU_BAND:
            v = _wm_final(run_wm_stc(t, t, B=B, trials=trials, shared_device=True))
            res["one_shared"][t] = v
            print(f"    tau={t:4.1f}: {_wm_fmt(v)}   ({time.time()-t0:.0f}s)")

        print("\n[head-to-head at matched tau] one-shared vs two-cotuned vs two-split:")
        for tw, ts in [(6.0, 6.0), (10.0, 10.0)]:
            one = _wm_final(run_wm_stc(tw, ts, B=B, trials=trials, shared_device=True))
            two = _wm_final(run_wm_stc(tw, ts, B=B, trials=trials, shared_device=False))
            res["two_device"][(tw, ts)] = two
            print(f"    tau={tw:4.1f}: one-shared {_wm_fmt(one)}  |  two-device {_wm_fmt(two)}"
                  f"   ({time.time()-t0:.0f}s)")

        paths.save_result("exp22_wm_stc.npy", res)
        print(f"\n  wrote exp22_wm_stc.npy  (total {time.time()-t0:.0f}s)")

    # ---------------- Experiment 23 ----------------
    if a.exp23 or run_all:
        B = 6 if quick else 20
        episodes = 600 if quick else 2500
        t0 = time.time()
        # LinearTrack corridor: reward = goal reached. The credit-assignment difficulty grows
        # with corridor length L, since the goal reward must propagate back across L forward/back
        # decisions. Pre-registered question: does the bootstrapped device-native TD rule
        # overtake the whole-episode-return REINFORCE baseline once L is long enough that the
        # scalar baseline fails, and is that crossover monotone in L? A no-trace control anchors
        # the chance (random-walk goal-reach) level at each L.
        L_GRID = [5, 10, 15, 20]
        crit = 0.75
        schemes = ["no_trace", "reinforce", "td_actor_critic", "td_no_homeo"]
        print(f"\n=== exp23 device-native TD/actor-critic vs REINFORCE | LinearTrack L-sweep "
              f"{L_GRID} | B={B}, {episodes} ep ===")
        print(f"    device-native discount gamma = exp(-step_dur/tau_leak); critic V(s) "
              f"on-device; delta gates the three-factor write. Reward = goal reached.")
        # every (scheme, L) cell is independent and seed-deterministic -> dispatch across a Pool
        specs = [(sc, L, B, episodes) for L in L_GRID for sc in schemes]
        n_proc = min(len(specs), max(1, (os.cpu_count() or 4) - 2))
        print(f"    dispatching {len(specs)} cells across {n_proc} processes...\n")
        with Pool(n_proc) as pool:
            out = pool.map(_td_worker, specs)
        cells = {key: (fin, curve, V) for key, fin, curve, V in out}

        res = {"B": B, "episodes": episodes, "L_grid": L_GRID, "crit": crit,
               "tau_leak": 10.0, "final": {}, "curve": {}}
        for (sc, L), (fin, curve, V) in cells.items():
            res["final"][(sc, L)] = fin
            res["curve"][(sc, L)] = curve
            res[f"V_{sc}_{L}"] = V

        print(f"  {'L':>3} | {'no_trace':>20} | {'reinforce':>20} | {'td_actor_critic':>20} | "
              f"{'td_no_homeo':>20} | TD-RE")
        for L in L_GRID:
            r = {sc: res["final"][(sc, L)] for sc in schemes}
            d = r["td_actor_critic"].mean() - r["reinforce"].mean()
            disj = bootstrap_ci(r["td_actor_critic"])[0] > bootstrap_ci(r["reinforce"])[1]
            print(f"  {L:>3} | {_td_fmt(r['no_trace']):>20} | {_td_fmt(r['reinforce']):>20} | "
                  f"{_td_fmt(r['td_actor_critic']):>20} | {_td_fmt(r['td_no_homeo']):>20} | "
                  f"{d:+.3f}{' *' if disj else ''}")

        # crossover: smallest L at which TD is disjoint-above REINFORCE
        cross = None
        for L in L_GRID:
            rt, rr = res["final"][("td_actor_critic", L)], res["final"][("reinforce", L)]
            if bootstrap_ci(rt)[0] > bootstrap_ci(rr)[1]:
                cross = L; break
        print("\n  === verdict ===")
        if cross is not None:
            print(f"  device-native TD overtakes REINFORCE (disjoint CIs) at L>={cross}: "
                  f"bootstrapping helps once credit must bridge that many decisions.")
        else:
            print(f"  no crossover in L{L_GRID}: device-native TD machinery runs and is stable "
                  f"but does not beat the return-baseline REINFORCE on this task/range (honest negative).")

        paths.save_result("exp23_device_td.npy", res)
        print(f"\n  wrote exp23_device_td.npy  (total {time.time()-t0:.0f}s)")


if __name__ == "__main__":
    main()
\end{lstlisting}

\begin{lstlisting}[caption={siox\_eligibility/dopamine.py},label={lst:siox-src-dopamine}]
r"""Real measured-dopamine reward signal for the biosignal reinforcement-learning
experiment -- the direct-neuromodulator analogue of :mod:`biosignal` (EEG).

Where :mod:`biosignal` decodes a reward-prediction error from the human EEG reward
positivity (a *cortical correlate* of dopaminergic RPE), this module decodes it from a
**directly recorded dopamine transient**: NAcc dLight1.3b fiber photometry during
Pavlovian conditioning (Jeong et al. 2022, Science; DANDI 000351). Dopamine is the literal
third factor the three-factor theory names, so this is a strictly stronger biological
grounding of the reward term.

Pipeline, mirroring :mod:`biosignal` exactly so the learning code is unchanged:
  1. ``load_session`` reads a cached session (dF/F trace @100 Hz + cue/reward event times).
  2. ``epoch_outcome`` epochs the dF/F around the expected-outcome time of each cue and
     labels each cue reward (CS+, a solenoid followed within the window) or omission (CS-).
  3. ``rewp_features`` -> ``decode_reward`` train a cross-validated single-trial decoder of
     reward vs omission; its OUT-OF-FOLD per-trial predictions are the reward gate (R-b).

Honesty: this supplies the THIRD FACTOR (a real, noisy dopaminergic reward-prediction
error). It is offline replay of a recorded animal's dopamine -- the right KIND of signal
(the actual neuromodulator), cross-species, not a live brain in the loop. It is a
biorealism/robustness test, not a new device capability. Whatever the single-trial decode
accuracy is, it is reported as-is.

The cache files under ``DA_CACHE`` are produced by the extractor that streams the dF/F
trace and eventlog out of the (large) NWB files via HTTP range reads; only the small
arrays are kept. See ``experiments/`` for the extractor and provenance.
"""
from __future__ import annotations

import glob
import os

import numpy as np

__all__ = ["load_session", "epoch_outcome", "rewp_features", "decode_reward",
           "build_reward_pools", "DA_CACHE",
           "make_shuffled_pools", "run_dopamine_shallow", "run_dopamine_deep",
           "HP", "HOMEO", "REW"]

#: Default cache directory of extracted per-session dopamine arrays.
DA_CACHE = os.environ.get("DA_CACHE", "/tmp/da_cache")

#: Reward window after a cue: a solenoid in (lo, hi) s marks the cue as rewarded (CS+).
REWARD_WIN = (0.2, 5.0)


def load_session(path):
    """Load one cached session.

    Returns ``dict(dff, dff_t, sound_t, reward_t, fs, sub)`` where ``dff`` is the
    analysis-ready dF/F dopamine trace, ``dff_t`` its timestamps (s), ``sound_t`` the cue
    onset times (s) and ``reward_t`` the solenoid (reward) times (s)."""
    return np.load(path, allow_pickle=True).item()


def _label_cues(sound_t, reward_t, win=REWARD_WIN):
    """Pair each reward to its nearest preceding cue within ``win``; label cues CS+/CS-.

    Returns (rewarded[bool, n_cue], delay) where ``delay`` is the median cue->reward lag
    over CS+ trials (the expected-outcome lag used to centre the omission epochs)."""
    lo, hi = win
    rewarded = np.zeros(len(sound_t), dtype=bool)
    delays = []
    for r in reward_t:
        cand = np.where((sound_t < r - lo) & (sound_t > r - hi))[0]
        if len(cand):
            j = cand[-1]
            if not rewarded[j]:
                rewarded[j] = True
                delays.append(r - sound_t[j])
    delay = float(np.median(delays)) if delays else 0.5 * (lo + hi)
    return rewarded, delay


def epoch_outcome(sess, *, pre=1.0, post=3.0):
    """Outcome-locked, baseline-corrected dF/F epochs for one session.

    Each cue is epoched over ``[-pre, +post]`` s around its expected-outcome time
    (cue + median CS+ delay), baseline-subtracted on ``[-pre, 0]``. Returns
    ``(epochs[n_cue, n_time], t, valence[n_cue])`` with valence 1 = reward (CS+),
    0 = omission (CS-). Cues whose window falls outside the recording are dropped."""
    dff = np.asarray(sess["dff"], float)
    ts = np.asarray(sess["dff_t"], float)
    fs = float(sess["fs"])
    rewarded, delay = _label_cues(sess["sound_t"], sess["reward_t"])
    npre, npost = int(pre * fs), int(post * fs)
    n = npre + npost
    t = np.arange(-npre, npost) / fs
    segs, val = [], []
    for s, rw in zip(sess["sound_t"], rewarded):
        c = s + delay
        i0 = int(np.searchsorted(ts, c)) - npre
        i1 = i0 + n
        if i0 < 0 or i1 > len(dff):
            continue
        seg = dff[i0:i1] - dff[i0:i0 + npre].mean()
        segs.append(seg)
        val.append(int(rw))
    return np.array(segs), t, np.array(val, dtype=int)


def rewp_features(epochs, t):
    """Per-trial features for reward/omission decoding: mean dF/F in successive bins over
    the post-outcome dopamine-response window (0-2 s), parallel to the EEG RewP features.

    Coarse, interpretable bins (avoids overfitting the full time course); the dopamine
    reward response is a transient peak in the first ~1.5 s after a rewarded outcome."""
    bins = [(0.0, 0.5), (0.5, 1.0), (1.0, 1.5), (1.5, 2.0), (0.0, 1.5)]
    feats = []
    for lo, hi in bins:
        w = (t >= lo) & (t < hi)
        feats.append(epochs[:, w].mean(1, keepdims=True))
    return np.concatenate(feats, axis=1)            # (n_trials, len(bins))


def decode_reward(features, valence, *, n_splits=5, seed=0, C=0.1):
    """Cross-validated single-trial reward-vs-omission decoder (logistic regression).

    Returns ``dict(acc, bal_acc, auc, proba, pred, valence)`` with ``proba``/``pred`` the
    OUT-OF-FOLD per-trial predictions, so the reward signal handed to the agent is never
    trained on its own trial. Identical contract to :func:`biosignal.decode_reward`."""
    from sklearn.linear_model import LogisticRegression
    from sklearn.preprocessing import StandardScaler
    from sklearn.model_selection import StratifiedKFold
    from sklearn.metrics import balanced_accuracy_score, roc_auc_score

    X, y = np.asarray(features, float), np.asarray(valence, int)
    proba = np.zeros(len(y)); pred = np.zeros(len(y), dtype=int)
    skf = StratifiedKFold(n_splits=n_splits, shuffle=True, random_state=seed)
    for tr, te in skf.split(X, y):
        sc = StandardScaler().fit(X[tr])
        clf = LogisticRegression(C=C, class_weight="balanced", max_iter=1000)
        clf.fit(sc.transform(X[tr]), y[tr])
        proba[te] = clf.predict_proba(sc.transform(X[te]))[:, 1]
        pred[te] = clf.predict(sc.transform(X[te]))
    return {
        "acc": float((pred == y).mean()),
        "bal_acc": float(balanced_accuracy_score(y, pred)),
        "auc": float(roc_auc_score(y, proba)) if len(np.unique(y)) > 1 else float("nan"),
        "proba": proba, "pred": pred, "valence": y,
    }


def build_reward_pools(cache_dir=DA_CACHE, *, seed=0, min_trials=40):
    """Decode every usable cached session and pool OUT-OF-FOLD predictions by TRUE valence.

    Returns ``(pools, meta)`` where ``pools = {1: array, 0: array}`` are the decoded reward
    VALUES on true-reward and true-omission trials respectively -- the exact
    ``reward_pools`` interface ``bandit.train`` / ``deep.train_deep`` accept. ``meta``
    carries the subject count and decode statistics. Sessions with one valence class or
    too few trials are skipped (not silently mis-shaped)."""
    files = sorted(glob.glob(os.path.join(cache_dir, "*.npy")))
    pool1, pool0, baccs, aucs, subs = [], [], [], [], []
    for fn in files:
        sess = load_session(fn)
        ep, t, val = epoch_outcome(sess)
        if len(np.unique(val)) < 2 or len(val) < min_trials:
            continue
        r = decode_reward(rewp_features(ep, t), val, seed=seed)
        baccs.append(r["bal_acc"]); aucs.append(r["auc"]); subs.append(sess["sub"])
        pool1.extend(r["pred"][r["valence"] == 1].tolist())
        pool0.extend(r["pred"][r["valence"] == 0].tolist())
    pools = {1: np.array(pool1, float), 0: np.array(pool0, float)}
    meta = {"n_subj": len(baccs), "subjects": subs,
            "mean_bal_acc": float(np.mean(baccs)) if baccs else float("nan"),
            "sd_bal_acc": float(np.std(baccs)) if baccs else float("nan"),
            "mean_auc": float(np.mean(aucs)) if aucs else float("nan"),
            "P_R1_given_reward": float(pools[1].mean()) if len(pool1) else float("nan"),
            "P_R1_given_omission": float(pools[0].mean()) if len(pool0) else float("nan")}
    return pools, meta


# =============================================================================
# Experiment 11 / Arm G -- the capstone with a REAL MEASURED DOPAMINE reward
#
# The direct-neuromodulator analogue of the EEG capstone (Arm E / exp9). The third
# factor that gates the device eligibility trace is a per-trial reward-prediction
# error decoded (above) from a *directly recorded dopamine transient* -- NAcc
# dLight1.3b fiber photometry during Pavlovian conditioning (Jeong et al. 2022,
# Science; DANDI 000351) -- rather than the human EEG reward-positivity *correlate*
# of that signal. Dopamine is the literal third factor the three-factor theory names,
# so this is the strongest biological grounding of the reward term in the paper.
#
# Two arms, both at frozen operating points (nothing retuned for the dopamine reward):
#   G-shallow : single-layer contextual bandit (mirrors exp7_biosignal), reward
#               sources {synthetic, dopamine, shuffled}.
#   G-deep    : deep XOR, rules {DFA, DFA+homeostasis} x reward {synthetic, dopamine,
#               shuffled} (mirrors exp9 capstone, frozen Arm-D operating point).
#
# Reported reward rate = TRUE task performance; the decoded dopamine value only gates
# the update. Pre-registered criteria G1-G4 / KILL live in
# PREREGISTRATION_dopamine_capstone.md. Negative results reported as-is; no goalpost
# moving.
#
# The two ``run_*`` cores below are SERIAL and return plain result dicts (no file I/O,
# no plotting, no stdout) so a notebook can call them at a small seed/trial count.
# ``main()`` runs the published 20-seed scale, parallelises the deep grid over specs,
# computes the criteria and writes the grid via ``paths.save_result``.
# =============================================================================

#: Frozen deep operating point -- identical to the exp9 capstone (Arm D/E). Nothing retuned.
HP = dict(H=32, tau_leak=10.0, D=5.0, eta=0.2, eta_hidden=3.0,
          fb_scale=2.0, bias_o=0.3, V=1.5, sigma0=0.15, sigma1=0.05)
#: Homeostasis strength for the ``dfa_homeo`` rule (moderate; the Arm-D value).
HOMEO = 0.1
#: Reward sources, in reporting order.
REW = ["synthetic", "dopamine", "shuffled"]

_CHANCE = 0.5
_CRIT = 0.75


def make_shuffled_pools(pools, *, seed=0):
    """Shuffled-reward control: destroy the valence->reward mapping while preserving the
    marginal distribution of decoded values.

    Concatenates both true-valence pools and returns ``{1: perm, 0: perm}`` -- two
    independent permutations of the SAME pooled values, so a trial's reward no longer
    carries any information about its true outcome. This is the exact ``shuf`` construction
    the original driver used to bound how much of the learning is genuine reward information
    versus the reward's marginal statistics."""
    rng = np.random.default_rng(seed)
    allv = np.concatenate([pools[1], pools[0]])
    return {1: rng.permutation(allv), 0: rng.permutation(allv)}


def _running(r1d, window=50):
    """Running-``window`` mean of a 1-D per-trial reward sequence (the learning curve)."""
    c = np.cumsum(np.insert(np.asarray(r1d, float), 0, 0.0))
    return (c[window:] - c[:-window]) / window


def run_dopamine_shallow(pools, shuf, *, seeds=20, trials=600, S=2, A=2,
                         tau_leak=10.0, D=2.0, seed0=0):
    """G-shallow (Experiment 11): single-layer contextual bandit under the measured
    dopamine reward.

    Trains the 2x2 bandit (:func:`mrl_trace.bandit.train`) with three reward
    sources -- ``synthetic`` (clean {0,1}), ``dopamine`` (the decoded ``pools`` gating the
    three-factor update), and ``shuffled`` (``shuf``, the valence-scrambled control) -- at
    the frozen operating point (nothing retuned for the dopamine reward). The reported
    reward rate is TRUE task performance; the decoded dopamine value only gates the update.

    Serial, returns a plain dict ``{finals, curves, ci}`` (per-seed final reward rate,
    mean running curve, bootstrap CI per source); no file I/O, no plotting, no stdout."""
    from .bandit import train as bandit_train, reward_rate as bandit_rate
    from .stats import bootstrap_ci

    finals, curves = {}, {}
    for name, rp in [("synthetic", None), ("dopamine", pools), ("shuffled", shuf)]:
        r = bandit_train(S, A, B=seeds, tau_leak=tau_leak, D=D, trials=trials,
                         seed0=seed0, reward_pools=rp)
        finals[name] = bandit_rate(r, window=100)
        curves[name] = _running(r.mean(0))
    ci = {n: bootstrap_ci(finals[n]) for n in finals}
    return {"finals": finals, "curves": curves, "ci": ci}


def _deep_worker(spec, trials, nseed, pools, shuf):
    """One deep-XOR condition (rule x reward source) -- module-level so ``main()`` can
    dispatch it over a :class:`multiprocessing.Pool`. Returns
    ``((rule, rname), finals (nseed,), curve)``."""
    from .deep import train_deep, reward_rate as deep_rate
    rule, rname = spec
    homeo = 0.0 if rule == "dfa" else HOMEO
    rp = {"synthetic": None, "dopamine": pools, "shuffled": shuf}[rname]
    rew = train_deep(mode="dfa", B=nseed, trials=trials, seed0=0, homeo=homeo,
                     reward_pools=rp, **HP)
    finals = deep_rate(rew, window=200)
    win = 100
    csum = np.cumsum(rew, axis=1)
    curve = ((csum[:, win:] - csum[:, :-win]) / win).mean(0)
    return (rule, rname), finals, curve


def run_dopamine_deep(pools, shuf, *, seeds=20, trials=3000, pool=None):
    """G-deep (Experiment 11): deep XOR at the frozen Arm-D operating point under the
    measured dopamine reward.

    Runs the full rule x reward grid: rules ``{dfa, dfa_homeo}`` x reward sources
    ``{synthetic, dopamine, shuffled}`` (six conditions). The temporal factor (the device
    trace) and every hyper-parameter (:data:`HP`, :data:`HOMEO`) are frozen at the Arm-D
    point; only the SPATIAL-credit rule (homeostasis on/off) and the reward SOURCE vary.

    Serial by default (a notebook can call it in-kernel); pass an open
    :class:`multiprocessing.Pool` as ``pool`` to dispatch the six conditions in parallel
    (used by ``main()``). Returns a plain dict ``{finals, curves, ci}`` keyed by the
    ``(rule, rname)`` tuple; no file I/O, no plotting, no stdout."""
    from .stats import bootstrap_ci
    specs = [(rule, rname) for rule in ("dfa", "dfa_homeo") for rname in REW]
    args = (specs, [trials] * len(specs), [seeds] * len(specs),
            [pools] * len(specs), [shuf] * len(specs))
    if pool is None:
        res = [_deep_worker(*a) for a in zip(*args)]
    else:
        res = pool.starmap(_deep_worker, list(zip(*args)))
    finals = {k: f for k, f, _ in res}
    curves = {k: c for k, _, c in res}
    ci = {k: bootstrap_ci(finals[k]) for k in finals}
    return {"finals": finals, "curves": curves, "ci": ci}


def _dopamine_criteria(shallow, deep, chance=_CHANCE):
    """Evaluate the pre-registered criteria G1-G3 from the shallow + deep result dicts.

    G1  shallow learns from real dopamine (CI-lo > max(chance, shuffled CI-hi)).
    G2  deep genuine reward info (homeo dopamine CI-lo > its own shuffled CI-hi, disjoint).
    G3  homeostasis robustness (homeo dopamine CI-lo > DFA-alone dopamine CI-hi).
    (G4 is descriptive-only: the dopamine ceiling vs the EEG capstone; reported in
    ``main()``.) Returns ``{"G1": bool, "G2": bool, "G3": bool}``."""
    s_ci = shallow["ci"]
    d_ci = deep["ci"]
    g1 = s_ci["dopamine"][0] > max(chance, s_ci["shuffled"][1])
    hda = d_ci[("dfa_homeo", "dopamine")]
    hsh = d_ci[("dfa_homeo", "shuffled")]
    g2 = hda[0] > hsh[1]
    dda = d_ci[("dfa", "dopamine")]
    g3 = hda[0] > dda[1]
    return {"G1": bool(g1), "G2": bool(g2), "G3": bool(g3)}


def _build_grid(shallow, deep, meta, *, seeds, trials_deep, trials_shallow):
    """Package the shallow + deep result dicts into the exact saved-grid schema the
    original ``exp11_dopamine_capstone.npy`` used (tuple deep-keys stringified for the
    pickle, criteria pre-computed)."""
    finals = shallow["finals"]
    d_finals, d_curves, d_ci = deep["finals"], deep["curves"], deep["ci"]
    return {
        "shallow": {"finals": finals, "curves": shallow["curves"], "ci": shallow["ci"]},
        "deep": {"finals": {str(k): v for k, v in d_finals.items()},
                 "curves": {str(k): v for k, v in d_curves.items()},
                 "ci": {str(k): v for k, v in d_ci.items()}},
        "meta": meta, "HP": HP, "homeo": HOMEO, "seeds": seeds,
        "trials_deep": trials_deep, "trials_shallow": trials_shallow,
        "chance": _CHANCE, "crit": _CRIT,
        "criteria": _dopamine_criteria(shallow, deep),
    }


def _replot_from_cache():
    """Graceful ``--figonly`` path: load the cached grid via ``paths`` and report its
    contents WITHOUT rerunning training (or requiring the DANDI cache to be present).

    The figure itself is rendered by the topic notebook from this same grid; this entry
    point just confirms the cached grid is readable and prints its headline numbers.
    Returns the loaded grid dict, or ``None`` if no cache has been written yet."""
    from . import paths
    fn = "exp11_dopamine_capstone.npy"
    path = paths.results_dir() / fn
    if not path.exists():
        print(f"  [figonly] no cached grid at {path}; run without --figonly first")
        return None
    d = paths.load_result(fn)
    meta = d.get("meta", {})
    print(f"  [figonly] loaded {fn}: {d.get('seeds')} seeds, "
          f"DA bal-acc {meta.get('mean_bal_acc', float('nan')):.3f}, "
          f"criteria={d.get('criteria')}")
    return d


def main(argv=None):
    """Full-scale reproduction CLI for the Arm-G dopamine capstone (writes ``data/results``).

    ``python -m mrl_trace.dopamine [--exp11] [--figonly] [--full|--quick]``

    Decodes the real dopamine reward pools from the DANDI 000351 cache
    (:func:`build_reward_pools`), runs the frozen-operating-point shallow bandit and deep
    XOR grids, computes the pre-registered G1-G4 criteria and writes
    ``exp11_dopamine_capstone.npy`` (the exact filename the original driver produced) via
    ``paths.save_result``.

    Graceful degradation (matching the original): with ``--figonly`` it only re-reports
    the cached grid; and if the DANDI cache under ``DA_CACHE`` is empty/absent (as it is
    wherever the multi-GB raw photometry has not been streamed), it prints a note and
    skips the run without error rather than crashing."""
    import argparse
    import time
    from . import paths

    ap = argparse.ArgumentParser(description="Arm-G dopamine capstone (Experiment 11)")
    ap.add_argument("--exp11", action="store_true",
                    help="dopamine capstone -> exp11_dopamine_capstone.npy (the default)")
    ap.add_argument("--figonly", action="store_true",
                    help="re-report the cached grid without rerunning training")
    ap.add_argument("--quick", action="store_true", help="fast few-seed smoke run")
    ap.add_argument("--full", action="store_true", help="published 20-seed run (default)")
    a = ap.parse_args(argv)

    if a.figonly:
        _replot_from_cache()
        return

    t0 = time.time()
    td = 600 if a.quick else 3000
    ts = 300 if a.quick else 600
    nseed = 6 if a.quick else 20

    print("Dopamine capstone (Arm G): REAL measured-dopamine reward x device trace")
    print(f"  decoding dopamine reward pools from DANDI 000351 (Jeong 2022) [DA_CACHE={DA_CACHE}] ...")
    pools, meta = build_reward_pools()
    if meta["n_subj"] == 0 or pools[1].size == 0 or pools[0].size == 0:
        # Graceful skip: the DANDI cache is empty/absent here (the multi-GB raw photometry
        # has not been streamed via extract_dopamine_cache.py). Report and exit cleanly --
        # do NOT write a degenerate grid over any existing published one.
        print(f"  [skip] no usable dopamine sessions under DA_CACHE={DA_CACHE!r} "
              f"(n_subj={meta['n_subj']}); run extract_dopamine_cache first. "
              f"Nothing written.")
        return
    print(f"  {meta['n_subj']} mice | single-trial decoder bal-acc {meta['mean_bal_acc']:.3f}"
          f" +/- {meta['sd_bal_acc']:.3f}, AUC {meta['mean_auc']:.3f}")
    print(f"  P(R=1 | reward)={meta['P_R1_given_reward']:.3f}  "
          f"P(R=1 | omission)={meta['P_R1_given_omission']:.3f}  (clean would be 1.00/0.00)")

    shuf = make_shuffled_pools(pools, seed=0)

    print(f"\n[G-shallow] single-layer contextual bandit | {nseed} seeds, {ts} trials")
    shallow = run_dopamine_shallow(pools, shuf, seeds=nseed, trials=ts)
    for name in REW:
        lo, hi = shallow["ci"][name]
        print(f"  {name:10s} reward rate = {shallow['finals'][name].mean():.3f} [{lo:.3f}, {hi:.3f}]")

    print(f"\n[G-deep] deep XOR, frozen Arm-D point | {nseed} seeds, {td} trials")
    import multiprocessing as mp
    specs = [(rule, rname) for rule in ("dfa", "dfa_homeo") for rname in REW]
    nproc = min(len(specs), (os.cpu_count() or 4) - 1)
    with mp.Pool(processes=max(1, nproc)) as pool:
        deep = run_dopamine_deep(pools, shuf, seeds=nseed, trials=td, pool=pool)
    d_finals, d_ci = deep["finals"], deep["ci"]
    print(f"  {'rule':18s} {'reward':12s} {'mean':>6} {'95% CI':>16} {'solved':>8}")
    for rule in ("dfa", "dfa_homeo"):
        for rname in REW:
            k = (rule, rname); lo, hi = d_ci[k]
            print(f"  {rule:18s} {rname:12s} {d_finals[k].mean():6.3f} "
                  f"[{lo:5.3f},{hi:5.3f}]  {int((d_finals[k] >= 0.9).sum()):>3}/{nseed}")

    crit = _dopamine_criteria(shallow, deep)
    da_deep = d_finals[("dfa_homeo", "dopamine")].mean()
    da_shal = shallow["finals"]["dopamine"].mean()
    print("\n=== pre-registered criteria (PREREGISTRATION_dopamine_capstone.md) ===")
    print(f"  G1 shallow learns from dopamine (>chance & >shuffled, disjoint): "
          f"{'PASS' if crit['G1'] else 'FAIL'} "
          f"(dopamine {da_shal:.3f}{shallow['ci']['dopamine']} vs "
          f"shuf {shallow['finals']['shuffled'].mean():.3f}{shallow['ci']['shuffled']})")
    print(f"  G2 deep genuine reward info (homeo dopamine > shuffled, disjoint): "
          f"{'PASS' if crit['G2'] else 'FAIL'} "
          f"(homeo {da_deep:.3f}{d_ci[('dfa_homeo', 'dopamine')]} vs "
          f"shuf {d_finals[('dfa_homeo', 'shuffled')].mean():.3f}{d_ci[('dfa_homeo', 'shuffled')]})")
    print(f"  G3 homeostasis robustness (homeo dopamine CI-lo > DFA-alone CI-hi): "
          f"{'PASS' if crit['G3'] else 'FAIL'} "
          f"(homeo {da_deep:.3f}{d_ci[('dfa_homeo', 'dopamine')]} vs "
          f"DFA {d_finals[('dfa', 'dopamine')].mean():.3f}{d_ci[('dfa', 'dopamine')]})")
    print(f"  G4 ceiling vs EEG (descriptive): dopamine deep {da_deep:.3f} / shallow {da_shal:.3f}"
          f" | EEG was 0.62 / 0.66 ; DA decoder {meta['mean_bal_acc']:.2f} vs EEG 0.69")

    grid = _build_grid(shallow, deep, meta, seeds=nseed,
                       trials_deep=td, trials_shallow=ts)
    paths.save_result("exp11_dopamine_capstone.npy", grid)
    print(f"\n  wrote exp11_dopamine_capstone.npy  criteria={grid['criteria']} "
          f"[{time.time() - t0:.0f}s]")


if __name__ == "__main__":
    main()
\end{lstlisting}

\begin{lstlisting}[caption={siox\_eligibility/biosignal.py},label={lst:siox-src-biosignal}]
"""Real-EEG reward signal for the biosignal reinforcement-learning experiment.

Loads the OpenNeuro ds003474 probabilistic-selection EEG (Cavanagh; CC0), epochs the
signal around feedback events, and trains a single-trial decoder of the feedback
valence (Correct vs Incorrect) from the frontocentral reward-positivity / feedback-
related negativity. The decoder's per-trial output is used as the biologically measured
reward gate (R - b) of the three-factor rule, in place of a synthetic reward.

No MNE dependency: EEGLAB ``.set`` is read with scipy.io.loadmat and the ``.fdt`` as a
float32 (pnts x nbchan) array; feedback events come from the BIDS ``events.tsv`` (clean
"Feedback: Correct"/"Feedback: Incorrect" labels with onsets), not the .set event
struct (which carries only numeric trigger codes).

Honesty: this supplies the THIRD FACTOR (a real, noisy reward-prediction-error signal)
to the device-trace bandit; it is a robustness/biorealism test, not a new device
capability, and it is offline replay, not a live brain in the loop.
"""
from __future__ import annotations

import csv
import os

import numpy as np
import scipy.io as sio

__all__ = ["load_subject", "epoch_feedback", "rewp_features", "decode_reward",
           "FRONTOCENTRAL", "build_reward_pools", "run_biosignal_reward",
           "run_eeg_capstone", "EEG_DATA_DEFAULT", "EEG_POOLS_CACHE"]

#: Frontocentral channels carrying the reward-positivity / FRN (uppercased labels).
FRONTOCENTRAL = ("FZ", "FC1", "FCZ", "FC2", "CZ", "C1", "C2")


def load_subject(eeg_set, eeg_fdt, events_tsv):
    """Load one subject: returns (data[nbchan, pnts], srate, labels, fb_onsets,
    fb_labels) where fb_labels is 1 for Correct (reward) and 0 for Incorrect."""
    m = sio.loadmat(eeg_set, squeeze_me=True, struct_as_record=False)["EEG"]
    srate = float(m.srate)
    nbchan, pnts = int(m.nbchan), int(m.pnts)
    raw = np.fromfile(eeg_fdt, dtype=np.float32)
    # The .set header nbchan/pnts can disagree with the .fdt (some ds003474 subjects
    # have interpolated/extra channels; README flags this). Reconcile against the file
    # size: trust pnts, recover the true channel count, and fail loudly if it does not
    # divide evenly so such subjects are skipped rather than silently mis-shaped.
    if raw.size != nbchan * pnts:
        if raw.size % pnts == 0:
            nbchan = raw.size // pnts
        else:
            raise ValueError(f"{os.path.basename(eeg_fdt)}: .fdt size {raw.size} "
                             f"not divisible by pnts {pnts}; channel layout unclear")
    data = raw.reshape((pnts, nbchan)).T
    labels = [str(c.labels).upper() for c in m.chanlocs]
    # if labels list is shorter/longer than recovered channels, pad/trim defensively so
    # frontocentral selection (by name) still works on the channels that ARE labelled
    if len(labels) < nbchan:
        labels = labels + [f"X{i}" for i in range(nbchan - len(labels))]
    elif len(labels) > nbchan:
        labels = labels[:nbchan]
    onsets, valence = [], []
    with open(events_tsv) as f:
        for row in csv.DictReader(f, delimiter="\t"):
            tt = row.get("trial_type", "")
            if tt == "Feedback: Correct":
                onsets.append(float(row["onset"])); valence.append(1)
            elif tt == "Feedback: Incorrect":
                onsets.append(float(row["onset"])); valence.append(0)
    return data, srate, labels, np.array(onsets), np.array(valence)


def epoch_feedback(data, srate, labels, onsets, *, t0=-0.1, t1=0.6,
                   channels=FRONTOCENTRAL, pnts=None):
    """Feedback-locked, baseline-corrected frontocentral epochs.

    Returns ``epochs[n_trials, n_chan, n_time]`` (baseline -100..0 ms removed) and the
    time vector. Trials whose window falls outside the recording are dropped (and the
    caller must drop the matching labels via the returned ``keep`` mask)."""
    if pnts is None:
        pnts = data.shape[1]
    ch = [i for i, l in enumerate(labels) if l in channels]
    a, b = int(t0 * srate), int(t1 * srate)
    base = int(-t0 * srate)
    segs, keep = [], []
    for o in onsets:
        s, e = int(o * srate) + a, int(o * srate) + b
        if 0 <= s and e < pnts:
            seg = data[ch, s:e]
            seg = seg - seg[:, :base].mean(1, keepdims=True)
            segs.append(seg); keep.append(True)
        else:
            keep.append(False)
    t = np.arange(a, b) / srate
    return np.array(segs), t, np.array(keep, dtype=bool)


def rewp_features(epochs, t):
    """Per-trial features for valence decoding: mean amplitude in successive 50 ms bins
    over 150-500 ms post-feedback (the RewP/FRN window), per frontocentral channel,
    flattened. Robust, interpretable, and avoids overfitting the full time course."""
    bins = [(0.15, 0.25), (0.25, 0.35), (0.35, 0.45), (0.20, 0.40)]
    feats = []
    for lo, hi in bins:
        w = (t >= lo) & (t < hi)
        feats.append(epochs[:, :, w].mean(2))      # (n_trials, n_chan)
    return np.concatenate(feats, axis=1)            # (n_trials, n_chan*len(bins))


def decode_reward(features, valence, *, n_splits=5, seed=0, C=0.1):
    """Cross-validated single-trial valence decoder (logistic regression, standardised).

    Returns ``dict(acc, bal_acc, proba, pred, auc)`` where ``proba``/``pred`` are
    OUT-OF-FOLD per-trial predictions (so the reward signal handed to the bandit is never
    trained on its own trial). Honest single-trial ErrP/RewP decoding is ~0.6-0.8
    balanced accuracy; whatever it is, it is reported as-is.
    """
    from sklearn.linear_model import LogisticRegression
    from sklearn.preprocessing import StandardScaler
    from sklearn.model_selection import StratifiedKFold
    from sklearn.metrics import balanced_accuracy_score, roc_auc_score

    X, y = np.asarray(features), np.asarray(valence)
    proba = np.zeros(len(y)); pred = np.zeros(len(y), dtype=int)
    skf = StratifiedKFold(n_splits=n_splits, shuffle=True, random_state=seed)
    for tr, te in skf.split(X, y):
        sc = StandardScaler().fit(X[tr])
        clf = LogisticRegression(C=C, class_weight="balanced", max_iter=1000)
        clf.fit(sc.transform(X[tr]), y[tr])
        proba[te] = clf.predict_proba(sc.transform(X[te]))[:, 1]
        pred[te] = clf.predict(sc.transform(X[te]))
    return {
        "acc": float((pred == y).mean()),
        "bal_acc": float(balanced_accuracy_score(y, pred)),
        "auc": float(roc_auc_score(y, proba)) if len(np.unique(y)) > 1 else float("nan"),
        "proba": proba, "pred": pred, "valence": y,
    }


# =============================================================================
# Experiment cores (the biosignal-reward RL studies that compose ``bandit.train``
# and ``deep.train_deep``)
#
# Each ``run_*`` returns the result grid as a plain dict -- no file I/O, no
# plotting, no stdout.  Notebooks call these at a small (quick) seed/trial count
# and render the figures inline; ``main()`` (below) calls them at the published
# 20-seed scale and writes the grid under ``data/results/`` for the notebooks to
# replay.
#
# The raw OpenNeuro ds003474 EEG is NOT bundled: it is a several-GB external
# download (Cavanagh; CC0), so these cores DEGRADE GRACEFULLY -- when the EEG is
# absent :func:`build_reward_pools` finds no usable subject, returns empty pools,
# and each core returns ``{"skipped": True, ...}`` with a printed note (the
# original drivers behaved the same way).  Provenance for the reported numbers:
#   - exp7 biosignal   : experiments/05_biological_grounding/biosignal_reward.py
#   - exp9 EEG capstone: experiments/05_biological_grounding/eeg_capstone.py
# =============================================================================

#: Default location of the (unbundled) OpenNeuro ds003474 EEG download, overridable
#: via ``SIOX_EEG_DIR`` or the ``--data`` CLI flag.
EEG_DATA_DEFAULT = os.environ.get("SIOX_EEG_DIR", "/tmp/ds003474")

#: Optional cache of the decoded reward pools (the capstone driver's ``POOLS_CACHE``).
EEG_POOLS_CACHE = os.environ.get("SIOX_EEG_POOLS_CACHE", "/tmp/eeg_pools.npy")


def build_reward_pools(data_dir=EEG_DATA_DEFAULT, *, seed=0, min_trials=40,
                       cache=None):
    """Decode per-trial reward from every usable ds003474 subject and pool the
    out-of-fold predictions by TRUE valence.

    Returns ``(pools, meta)`` where ``pools = {1: array, 0: array}`` are the decoded
    reward VALUES (the decoder's 0/1 prediction) on true-correct and true-incorrect
    trials respectively -- the exact ``reward_pools`` interface ``bandit.train`` /
    ``deep.train_deep`` accept -- and ``meta`` records the decode accuracy actually
    achieved (the realistic reward-noise level).  Subjects with a broken channel
    layout, a single valence class, or fewer than ``min_trials`` feedback events are
    skipped (not silently mis-shaped).  Identical to the two drivers' local
    ``build_reward_pools``.

    Graceful degradation: if ``data_dir`` holds no usable subject (e.g. the several-GB
    EEG is not downloaded here), both pools come back EMPTY and ``meta == {}``; callers
    treat that as "skip this experiment" rather than erroring.

    ``cache`` (default ``None``) may be a path to a ``.npy`` reward-pool cache; if it
    exists it is loaded verbatim (the capstone driver's ``POOLS_CACHE`` behaviour),
    otherwise a freshly decoded pool is written to it.
    """
    if cache and os.path.exists(cache):
        d = np.load(cache, allow_pickle=True).item()
        meta = d.get("_meta", {})
        pools = {1: np.asarray(d[1], float), 0: np.asarray(d[0], float)}
        return pools, meta
    import glob
    subs = sorted(glob.glob(os.path.join(data_dir, "sub-*/eeg/"
                  "sub-*_task-ProbabilisticSelection_eeg.fdt")))
    pool1, pool0, baccs = [], [], []
    used = 0
    for fdt in subs:
        pre = fdt[:-len("_eeg.fdt")]
        try:
            data, srate, labels, onsets, val = load_subject(
                pre + "_eeg.set", pre + "_eeg.fdt", pre + "_events.tsv")
        except ValueError:
            continue                                   # skip channel-layout-broken subj
        ep, t, keep = epoch_feedback(data, srate, labels, onsets, pnts=data.shape[1])
        y = val[keep]
        if len(np.unique(y)) < 2 or len(y) < min_trials:
            continue
        r = decode_reward(rewp_features(ep, t), y, seed=seed)
        baccs.append(r["bal_acc"]); used += 1
        # out-of-fold predicted reward, split by the TRUE outcome valence
        pool1.extend(r["pred"][r["valence"] == 1].tolist())   # decoded R | true=correct
        pool0.extend(r["pred"][r["valence"] == 0].tolist())   # decoded R | true=incorrect
    pools = {1: np.array(pool1, dtype=float), 0: np.array(pool0, dtype=float)}
    if used == 0:
        return pools, {}                               # graceful: nothing usable here
    meta = {"n_subj": used, "mean_bal_acc": float(np.mean(baccs)),
            "sd_bal_acc": float(np.std(baccs)),
            "P_R1_given_correct": float(pools[1].mean()),
            "P_R1_given_incorrect": float(pools[0].mean())}
    if cache:
        np.save(cache, {1: pools[1], 0: pools[0], "_meta": meta}, allow_pickle=True)
    return pools, meta


def _running(r1d, window=50):
    """Running-mean reward curve (window trials), matching the exp7 driver."""
    c = np.cumsum(np.insert(np.asarray(r1d, float), 0, 0.0))
    return (c[window:] - c[:-window]) / window


def run_biosignal_reward(pools, meta, *, seeds=20, trials=600, S=2, A=2,
                         tau_leak=10.0, D=2.0, seed0=0):
    """Experiment 7 -- close the RL loop with a REAL biosignal reward (Arm: biorealism).

    The contextual spiking bandit's synthetic reward ``R in {0,1}`` is replaced by a
    per-trial reward-prediction-error DECODED FROM REAL EEG (ds003474 reward positivity /
    FRN).  The device supplies the eligibility trace; biology supplies the third factor.
    Reported reward rate is TASK PERFORMANCE (true outcome); the EEG-decoded value only
    gates the ``dw = eta (R - b) e`` update.

    Conditions (device gate, same task/seeds/budget):
      synthetic : clean ``R in {0,1}`` (reference);
      biosignal : per-trial reward sampled from the EEG decoder's out-of-fold predictions,
                  keyed by the true outcome valence (realistic decode noise);
      shuffled  : the SAME decoded values with valence permuted (a noise-matched control:
                  same reward statistics, no genuine outcome information).

    Pre-registered criteria B1-B4 (BIOSIGNAL_RL_notes.md):
      B1 decoder beats chance (bal-acc > 0.55);
      B2 biosignal reaches criterion (>= 0.5*(1+1/A));
      B3 biosignal beats shuffled by > 0.10;
      B4 biosignal graceful vs synthetic (>= synthetic - 0.25).

    ``pools``/``meta`` come from :func:`build_reward_pools`.  If ``pools`` is empty (EEG
    absent) the experiment is skipped and ``{"skipped": True}`` is returned.  Serial: no
    multiprocessing, no file I/O -- notebook-callable.
    """
    from .bandit import train, reward_rate
    from .stats import bootstrap_ci
    if len(pools.get(1, [])) == 0 or len(pools.get(0, [])) == 0:
        return {"skipped": True, "reason": "no usable ds003474 EEG subjects found"}
    chance = 1.0 / A
    crit = 0.5 * (1 + chance)

    # shuffled control: identical pooled reward values, valence association destroyed
    allvals = np.concatenate([pools[1], pools[0]])
    rng = np.random.default_rng(0)
    shuf = {1: rng.permutation(allvals), 0: rng.permutation(allvals)}

    conds = {
        "synthetic": dict(reward_pools=None),
        "biosignal": dict(reward_pools=pools),
        "shuffled":  dict(reward_pools=shuf),
    }
    finals, curves = {}, {}
    for name, kw in conds.items():
        r = train(S, A, B=seeds, tau_leak=tau_leak, D=D, trials=trials, seed0=seed0, **kw)
        finals[name] = reward_rate(r, window=100)
        curves[name] = _running(r.mean(0))

    b1 = meta.get("mean_bal_acc", float("nan")) > 0.55
    b2 = finals["biosignal"].mean() >= crit
    b3 = finals["biosignal"].mean() > finals["shuffled"].mean() + 0.1
    b4 = finals["biosignal"].mean() >= finals["synthetic"].mean() - 0.25
    return {
        "finals": finals, "curves": curves, "meta": meta,
        "chance": chance, "crit": crit, "seeds": seeds, "trials": trials,
        "S": S, "A": A, "tau_leak": tau_leak, "D": D,
        "criteria": {"B1": bool(b1), "B2": bool(b2), "B3": bool(b3), "B4": bool(b4)},
    }


# --- capstone (Experiment 9, Arm E): EEG reward x deep all-local RL + homeostasis ---
# Frozen Arm-D operating point (PREREGISTRATION_capstone_eeg_deep.md).
_CAP_HP = dict(H=16, tau_leak=10.0, D=5.0, eta=0.2, eta_hidden=3.0,
               fb_scale=2.0, bias_o=0.3, V=1.5, sigma0=0.15, sigma1=0.05)
_CAP_HOMEO = 0.1
_CAP_RULES = [("dfa", 0.0), ("dfa_homeo", _CAP_HOMEO)]
_CAP_REW = ["synthetic", "eeg", "shuffled"]


def _capstone_worker(spec, trials, pools, shuf, seeds):
    """Train one (rule, reward-source) cell of the capstone grid; returns
    ``((rule, rname), finals, curve)``.  Module-level so it is picklable for the
    ``multiprocessing.Pool`` in :func:`main`."""
    from .deep import train_deep, reward_rate
    rule, rname = spec
    mode, homeo = ("dfa", 0.0) if rule == "dfa" else ("dfa", _CAP_HOMEO)
    rp = {"synthetic": None, "eeg": {0: pools[0], 1: pools[1]}, "shuffled": shuf}[rname]
    rew = train_deep(mode=mode, B=seeds, trials=trials, seed0=0, homeo=homeo,
                     reward_pools=rp, **_CAP_HP)
    finals = reward_rate(rew, window=200)
    win = 100
    csum = np.cumsum(rew, axis=1)
    curve = ((csum[:, win:] - csum[:, :-win]) / win).mean(0)
    return (rule, rname), finals, curve


def _summarize_capstone(res, meta, *, seeds, trials, chance, crit, homeo):
    """Package raw per-cell ``_capstone_worker`` output into the saved grid dict
    (str-keyed ``finals``/``curves``/``ci`` as the driver saved them, plus the
    pre-registered E2-E4 criteria)."""
    from .stats import bootstrap_ci
    finals = {k: f for k, f, _ in res}
    curves = {k: c for k, _, c in res}
    ci = {k: bootstrap_ci(finals[k]) for k in finals}

    d_eeg, h_eeg = ci[("dfa", "eeg")], ci[("dfa_homeo", "eeg")]
    h_eeg_m = finals[("dfa_homeo", "eeg")].mean()
    h_shuf = ci[("dfa_homeo", "shuffled")]
    h_syn_m = finals[("dfa_homeo", "synthetic")].mean()
    e2 = bool(h_eeg_m >= crit and h_eeg[0] > d_eeg[1])
    e3 = bool(h_eeg[0] > h_shuf[1])
    e4 = bool(h_eeg_m >= h_syn_m - 0.25)
    return {
        "finals": {str(k): v for k, v in finals.items()},
        "curves": {str(k): v for k, v in curves.items()},
        "ci": {str(k): v for k, v in ci.items()}, "meta": meta,
        "HP": _CAP_HP, "homeo": homeo, "seeds": seeds, "trials": trials,
        "chance": chance, "crit": crit,
        "criteria": {"E2": e2, "E3": e3, "E4": e4},
    }


def run_eeg_capstone(pools, meta, *, seeds=20, trials=3000, chance=0.5, crit=0.75):
    """Experiment 9 (Arm E) -- capstone: EEG reward x deep all-local RL + homeostasis.

    Composes the three biorealistic ingredients of Chapter 7 into ONE experiment:
      - temporal credit : the physical device eligibility trace (on both layers);
      - spatial credit  : a deep two-layer policy trained fully locally by direct
                          feedback alignment (no backprop, no weight transport);
      - the third factor: a REAL human reward-prediction-error decoded from EEG
                          (ds003474), the same biosignal reward as
                          :func:`run_biosignal_reward`.

    Tests the pre-registered hypothesis (PREREGISTRATION_capstone_eeg_deep.md) that the
    local homeostatic stabiliser -- which prevents the policy collapse behind deep-DFA's
    unreliability (Arm D) -- ALSO confers robustness to a noisy biological reward, so
    DFA+homeostasis learns under the EEG reward where DFA-alone does not.

    Design: 2 rules {DFA, DFA+homeostasis} x 3 reward sources {synthetic, EEG, shuffled}
    on deep XOR at the frozen Arm-D operating point, 20 seeds, bootstrap 95% CI.  Reported
    reward rate is TRUE task performance; the EEG-decoded value only gates the update.

    Pre-registered criteria:
      E2 homeostasis robustness (EEG homeo >= crit AND its CI-lo > DFA CI-hi);
      E3 genuine reward info (EEG homeo CI-lo > shuffled CI-hi, disjoint);
      E4 graceful vs clean (EEG homeo >= synthetic homeo - 0.25).

    SERIAL (notebook-callable): iterates the 6 grid cells in-process, no ``Pool``, no file
    I/O.  ``main()`` runs the same cells across a ``multiprocessing.Pool``.  If ``pools`` is
    empty (EEG absent) the experiment is skipped and ``{"skipped": True}`` is returned.
    """
    if len(pools.get(1, [])) == 0 or len(pools.get(0, [])) == 0:
        return {"skipped": True, "reason": "no usable ds003474 EEG subjects found"}
    rng = np.random.default_rng(0)
    allv = np.concatenate([pools[1], pools[0]])
    shuf = {1: rng.permutation(allv), 0: rng.permutation(allv)}
    specs = [(rule, rname) for rule, _ in _CAP_RULES for rname in _CAP_REW]
    res = [_capstone_worker(s, trials, pools, shuf, seeds) for s in specs]
    return _summarize_capstone(res, meta, seeds=seeds, trials=trials,
                               chance=chance, crit=crit, homeo=_CAP_HOMEO)


def main(argv=None):
    """Full-scale reproduction CLI for the biosignal-reward RL grids (writes ``data/results``).

    ``python -m mrl_trace.biosignal [--biosignal] [--capstone] [--full|--quick]
    [--data DIR]``

    With no experiment flag, runs both.  ``--full`` = 20 seeds (published); ``--quick`` =
    a fast few-seed smoke run.  ``--data DIR`` points at the (unbundled) OpenNeuro ds003474
    download (default ``$SIOX_EEG_DIR`` or ``/tmp/ds003474``).  When the EEG is absent the
    reward pools come back empty and each experiment is SKIPPED with a printed note (exit
    0) -- the graceful-degradation path, since the several-GB EEG is not bundled.
    """
    import argparse
    from . import paths
    ap = argparse.ArgumentParser(description="Biosignal-reward RL reproductions")
    ap.add_argument("--biosignal", action="store_true",
                    help="exp7 device bandit under real-EEG reward -> exp7_biosignal.npy")
    ap.add_argument("--capstone", action="store_true",
                    help="exp9 EEG x deep DFA + homeostasis -> exp9_capstone.npy")
    ap.add_argument("--quick", action="store_true", help="fast few-seed smoke run")
    ap.add_argument("--full", action="store_true", help="published 20-seed run (default)")
    ap.add_argument("--data", default=EEG_DATA_DEFAULT,
                    help="OpenNeuro ds003474 EEG directory (default $SIOX_EEG_DIR or /tmp/ds003474)")
    a = ap.parse_args(argv)
    run_all = not (a.biosignal or a.capstone)

    print("=== decoding real-EEG reward pools from ds003474 ===")
    pools, meta = build_reward_pools(a.data)
    if not meta:
        print(f"  no usable ds003474 EEG subjects under {a.data!r} -- "
              f"biosignal experiments SKIPPED (EEG is a several-GB external download; "
              f"set --data / $SIOX_EEG_DIR to reproduce).")
        return
    print(f"  {meta['n_subj']} subjects | decoder bal-acc {meta['mean_bal_acc']:.3f} "
          f"+/- {meta.get('sd_bal_acc', float('nan')):.3f}")
    print(f"  P(R=1 | correct)={pools[1].mean():.3f}  "
          f"P(R=1 | incorrect)={pools[0].mean():.3f}  (clean would be 1.00 / 0.00)")

    if a.biosignal or run_all:
        seeds = 6 if a.quick else 20
        trials = 400 if a.quick else 600
        print(f"\n=== exp7 biosignal-reward bandit ({seeds} seeds, {trials} trials) ===")
        grid = run_biosignal_reward(pools, meta, seeds=seeds, trials=trials)
        if grid.get("skipped"):
            print(f"  skipped: {grid['reason']}")
        else:
            f = grid["finals"]
            from .stats import bootstrap_ci
            for name in ("synthetic", "biosignal", "shuffled"):
                lo, hi = bootstrap_ci(f[name])
                print(f"  {name:10s} reward rate = {f[name].mean():.3f} [{lo:.3f}, {hi:.3f}]")
            paths.save_result("exp7_biosignal.npy", grid)
            print(f"  wrote exp7_biosignal.npy  criteria={grid['criteria']}")

    if a.capstone or run_all:
        import os as _os
        from functools import partial
        from multiprocessing import Pool
        seeds = 6 if a.quick else 20
        trials = 600 if a.quick else 3000
        chance, crit = 0.5, 0.75
        print(f"\n=== exp9 capstone: EEG x deep all-local + homeostasis "
              f"({seeds} seeds, {trials} trials) ===")
        rng = np.random.default_rng(0)
        allv = np.concatenate([pools[1], pools[0]])
        shuf = {1: rng.permutation(allv), 0: rng.permutation(allv)}
        specs = [(rule, rname) for rule, _ in _CAP_RULES for rname in _CAP_REW]
        # Pool over the 6 (rule x reward-source) cells -- this is a real __main__.
        with Pool(processes=min(len(specs), (_os.cpu_count() or 4) - 1)) as pool:
            res = pool.map(partial(_capstone_worker, trials=trials, pools=pools,
                                   shuf=shuf, seeds=seeds), specs)
        grid = _summarize_capstone(res, meta, seeds=seeds, trials=trials,
                                   chance=chance, crit=crit, homeo=_CAP_HOMEO)
        fin = {eval(k): v for k, v in grid["finals"].items()}
        cis = {eval(k): v for k, v in grid["ci"].items()}
        print(f"  {'rule':18s} {'reward':10s} {'mean':>6} {'95% CI':>16}")
        for rule, _ in _CAP_RULES:
            for rname in _CAP_REW:
                k = (rule, rname); lo, hi = cis[k]
                print(f"  {rule:18s} {rname:10s} {fin[k].mean():6.3f} [{lo:5.3f},{hi:5.3f}]")
        paths.save_result("exp9_capstone.npy", grid)
        print(f"  wrote exp9_capstone.npy  criteria={grid['criteria']}")


if __name__ == "__main__":
    main()
\end{lstlisting}

\begin{lstlisting}[caption={siox\_eligibility/capture.py},label={lst:siox-src-capture}]
r"""Two-state ('second-order') device model: can *capture* be memristive?

This module asks whether the reward-gated **capture** step of the three-factor rule
--- committing the surviving eligibility trace into a lasting weight change ---
can be performed by the device itself, rather than by a peripheral active element.
It does so with a deliberately minimal, provenance-labelled model:

  STATE 1 (the TAG, volatile)  -- the eligibility trace ``v(t)``.
      GROUNDED.  Generated by the fitted :class:`mrl_trace.device.TransientGate`
      (trap-cascade rise + ``tau_leak`` leakage; KWW field-acceleration laws from
      Chapter 5 / ``kww_final.json``; ``tau_leak`` from the ITO discharge set, median
      ~1.3 s).  No new parameters.

  STATE 2 (the captured WEIGHT, non-volatile) -- the stored conductance ``G``.
      GROUNDED.  The reachable conductance ladder is the *measured* SiO_x multistate
      set used for inference in vu2024snn (``SiO_x-multistate-data.mat``): 53 states,
      G_off ~ 2.7 uS to G_on ~ 3.5 mS (~1300x range).  A capture event moves ``G`` by
      a fraction of one measured step; weights are bounded by the measured ladder.

  THE COUPLING ``g(v)`` -- how strongly a reward pulse writes, given surviving trace.
      *** ASSUMED (the single free element of this model). ***
      A monotonic, saturating gate: capture is proportional to the trace still present
      when the reward arrives.  Motivated by trap-assisted switching --- the trapped
      charge that *is* the tag lowers the local barrier / seeds the SET --- but the
      coefficient is NOT fitted to a measurement.  This is the one law the
      delta-G-vs-delay experiment (manuscript Limitations) would pin down.

Mechanism class precedent (volatile internal state gating a non-volatile write):
second-order memristors and STP->LTP consolidation, demonstrated across material
systems --- Kim/Lu Nano Lett. 2015 (10.1021/acs.nanolett.5b00697); Chang/Lu ACS
Nano 2011 (10.1021/nn202983n); Ohno et al. Nat. Mater. 2011 (10.1038/nmat3054);
Wang et al. Nat. Mater. 2016 (10.1038/nmat4756).  None is SiO_x: whether *this*
material realises the coupling is the open question the model frames, not a result.

The headline observable is the **capture curve**: the committed weight change as a
function of action->reward delay ``D``.  A correct tag-and-capture device commits
change in proportion to the surviving trace and commits ~nothing once ``D`` exceeds
``tau_leak`` --- the physical signature of a capture window set by the tag lifetime.
"""
from __future__ import annotations

import os
from dataclasses import dataclass

import numpy as np

from .device import TransientGate

__all__ = [
    "load_measured_conductance_ladder",
    "CouplingModel",
    "TwoStateSynapse",
    "capture_curve",
]

# Path to the measured multistate conductance data (shared with mnn-torch / vu2024snn).
_DEFAULT_MAT = os.path.normpath(
    os.path.join(os.path.dirname(__file__), "..", "..", "..",
                 "mnn-torch", "data", "SiO_x-multistate-data.mat")
)


def load_measured_conductance_ladder(mat_path: str | None = None) -> np.ndarray:
    """Return the sorted *measured* SiO_x conductance levels (S).

    GROUNDED.  Reads the same ``SiO_x-multistate-data.mat`` used for inference in
    vu2024snn and reduces each multistate sweep to a single read conductance
    (median of ``I/V`` over the upper-voltage portion of the 0--0.5 V read sweep).
    """
    from scipy.io import loadmat  # local import: keep module import light

    path = mat_path or _DEFAULT_MAT
    raw = loadmat(path)["data"]            # (799, 3, 53) = [points, I/V/_, states]
    # mirror mnn_torch.devices.load_SiOx_multistate exactly (known-good ordering):
    raw = np.flip(raw, axis=2)
    raw = np.transpose(raw, (1, 2, 0))     # -> (3, states, points)
    raw = raw[:2]                          # -> (2, states, points) = [I, V]
    cur, volt = raw[0], raw[1]
    levels = []
    for s in range(cur.shape[0]):
        v, i = volt[s], cur[s]
        m = (v > 0.05) & np.isfinite(v) & np.isfinite(i)
        if m.sum() > 5:
            levels.append(np.median(i[m] / v[m]))
    return np.array(sorted(levels))


@dataclass
class CouplingModel:
    r"""The ASSUMED capture coupling ``g(v)`` --- the one free element of the model.

    Capture rate per reward pulse is ``kappa * (R - b) * g(v)`` where ``g`` is a
    monotonic saturating function of the surviving trace ``v``:

        g(v) = v**p / (v**p + v_half**p)        (Hill form, p >= 1)

    ``p = 1`` is the plain proportional gate; ``p > 1`` gives a soft threshold
    (capture only once the trace clears ``v_half``).  All three parameters are
    ASSUMED, not fitted; they are the targets of the delta-G-vs-delay measurement.
    """

    kappa: float = 0.6      # capture gain (fraction of a measured G-step per unit reward)
    v_half: float = 0.25    # trace level at half-maximal capture (normalised trace units)
    p: float = 1.0          # Hill exponent: 1 = proportional, >1 = soft threshold

    def g(self, v: np.ndarray | float) -> np.ndarray | float:
        v = np.clip(np.abs(v), 0.0, None)
        vp = v ** self.p
        return vp / (vp + self.v_half ** self.p)


class TwoStateSynapse:
    r"""One synapse with a volatile tag and a non-volatile captured weight.

    The tag ``v(t)`` is the fitted device transient (:class:`TransientGate`).  The
    weight ``G`` lives on the measured conductance ladder and only moves when a
    reward pulse arrives, by an amount gated by the surviving tag through the
    assumed :class:`CouplingModel`.  Sign of ``(R - b)`` sets SET vs RESET; the
    weight is clipped to the measured ``[G_off, G_on]`` bounds.
    """

    def __init__(self, V: float = 0.9, tau_leak: float = 2.0, dt: float = 0.05,
                 coupling: CouplingModel | None = None,
                 ladder: np.ndarray | None = None):
        self.gate = TransientGate(V=V, tau_leak=tau_leak, dt=dt)
        self.coupling = coupling or CouplingModel()
        self.ladder = load_measured_conductance_ladder() if ladder is None else ladder
        self.g_off, self.g_on = float(self.ladder[0]), float(self.ladder[-1])
        # one measured step ~ median spacing of the ladder (in S)
        self.g_step = float(np.median(np.diff(self.ladder)))
        self.dt = dt
        self.reset()

    def reset(self) -> None:
        self.gate.reset()
        self.G = 0.5 * (self.g_off + self.g_on)   # start mid-ladder

    def run(self, t_grid: np.ndarray, coincidence_at: float, reward_at: float,
            reward: float = 1.0, coincidence_dur: float = 0.2,
            reward_dur: float = 0.2) -> dict:
        """Drive a coincidence at ``coincidence_at`` and a reward pulse at ``reward_at``.

        Returns time series of the tag ``v`` and weight ``G``, and the committed
        change ``dG`` (final minus initial).  ``reward`` is the signed ``(R - b)``.
        """
        t_grid = np.asarray(t_grid, float)
        v_ts = np.zeros_like(t_grid)
        G_ts = np.zeros_like(t_grid)
        G0 = self.G
        captured = False     # one reward pulse -> one capture event (impulse model)
        for n, tt in enumerate(t_grid):
            drive = 1.0 if coincidence_at <= tt < coincidence_at + coincidence_dur else 0.0
            v = self.gate.step(drive)            # advance the tag (grounded)
            v_ts[n] = v
            # capture: a single write at reward onset, gated by the surviving tag.
            # Applying it once (not per-dt) makes dG independent of the grid spacing,
            # and matches one reward pulse committing one weight change.
            if (not captured) and tt >= reward_at:
                dG = (self.coupling.kappa * reward
                      * self.coupling.g(v) * self.g_step)
                self.G = float(np.clip(self.G + dG, self.g_off, self.g_on))
                captured = True
            G_ts[n] = self.G
        return {"t": t_grid, "v": v_ts, "G": G_ts, "dG": self.G - G0, "G0": G0}


def capture_curve(delays: np.ndarray, V: float = 0.9, tau_leak: float = 2.0,
                  reward: float = 1.0, dt: float = 0.05,
                  coupling: CouplingModel | None = None,
                  ladder: np.ndarray | None = None) -> dict:
    """Committed weight change ``dG`` vs action->reward delay ``D``.

    For each delay the coincidence fires at a fixed time and the reward pulse
    arrives ``D`` seconds later.  The capture window is expected to follow the tag:
    ``dG`` falls off as ``D`` approaches and exceeds ``tau_leak``.
    """
    delays = np.asarray(delays, float)
    dG = np.zeros_like(delays)
    v_at_reward = np.zeros_like(delays)
    coincidence_at = 1.0
    pad = 1.0
    for i, D in enumerate(delays):
        syn = TwoStateSynapse(V=V, tau_leak=tau_leak, dt=dt,
                              coupling=coupling, ladder=ladder)
        reward_at = coincidence_at + D
        t_end = reward_at + max(0.5, pad)
        t_grid = np.arange(0.0, t_end, dt)
        out = syn.run(t_grid, coincidence_at=coincidence_at, reward_at=reward_at,
                      reward=reward)
        dG[i] = out["dG"]
        # tag value at the moment the reward arrives (for the v vs dG check)
        ridx = int(round(reward_at / dt))
        v_at_reward[i] = out["v"][min(ridx, len(out["v"]) - 1)]
    return {"D": delays, "dG": dG, "v_at_reward": v_at_reward,
            "tau_leak": tau_leak, "g_off": syn.g_off, "g_on": syn.g_on}
\end{lstlisting}

\begin{lstlisting}[caption={siox\_eligibility/device\_faults.py},label={lst:siox-src-device_faults}]
r"""Device non-idealities for array-scale feasibility studies.

A faithful NumPy port of the fault catalogue in the nonideality-aware-mnn-training
package (``awarememristor.crossbar.nonidealities``), the engine the SiO_x
inference/homeostasis work trains against. We port the *semantics* rather than import
the package, because that package is TensorFlow and this learning rule is pure NumPy
(running a TF op inside the per-timestep spiking loop would be both a heavy dependency
and prohibitively slow). The fault definitions below mirror those classes exactly, so
the device-fault prior is shared with the inference work while the learning stays here.

Two physical classes (after Joksas et al.):

* Linearity-preserving -- corrupt the conductance, the multiply stays linear:
    - ``StuckAtGOff`` / ``StuckAtGOn`` : a fraction of devices pinned to G_off / G_on.
    - ``StuckDistribution``           : a fraction stuck at values drawn from a KDE of
                                        measured stuck conductances (reflected at 0).
    - ``D2DLognormal``                : lognormal device-to-device programming spread,
                                        width interpolated between R_on and R_off.
* Linearity-NONpreserving -- the per-device I--V law itself is nonlinear:
    - ``IVNonlinearityPF``            : Poole--Frenkel forward map with per-synapse
                                        (c, d_epsilon) sampled from the measured
                                        bivariate regression; an always-on distortion
                                        of every analogue multiply.

Each fault exposes ``disturb(G, rng)`` (linearity-preserving) or
``apply(V_eff, G, rng)`` (nonpreserving), and a ``FaultStack`` composes them into the
single ``weight_fault`` hook ``train_deep`` accepts. Faults are a fixed physical
realisation: drawn once (cached) and reused for the run, as a programmed array behaves.
"""
import numpy as np
import scipy.constants as const

__all__ = ["StuckAtGOff", "StuckAtGOn", "StuckDistribution", "D2DLognormal",
           "IVNonlinearityPF", "FaultStack", "SIOX_PF", "siox_fault_stack",
           "maze_fault_stack"]

# --- Pre-computed SiO_x device constants (fitted ONCE, offline, from the measured
# SiO_x-multistate-data.mat via mnn-torch's PF regression; baked here so no runtime
# fit or TF dependency is needed). These ARE the parameters the inference work uses. ---
SIOX_PF = {
    "G_off": 6.9013e-4,
    "G_on": 3.4506e-3,
    "slopes": [0.236877, 1.091365],
    "intercepts": [-11.535144, -47.585831],
    "res_cov": [[0.0154065, -0.0031134], [-0.0031134, 0.0766984]],
    "k_V": 0.5,
}


class _Fault:
    """Base: draws a fixed realisation on first call, caches it per weight shape."""
    def __init__(self):
        self._cache = {}

    def disturb(self, G, rng):
        raise NotImplementedError


class StuckAtGOff(_Fault):
    """A fraction ``probability`` of devices pinned to G_off (open synapses)."""
    def __init__(self, probability, G_off=0.0):
        super().__init__(); self.p = probability; self.G_off = G_off

    def disturb(self, G, rng):
        key = G.shape
        if key not in self._cache:
            self._cache[key] = rng.random(G.shape) < self.p
        return np.where(self._cache[key], self.G_off, G)


class StuckAtGOn(_Fault):
    """A fraction ``probability`` of devices pinned to G_on (shorted synapses)."""
    def __init__(self, probability, G_on=1.0):
        super().__init__(); self.p = probability; self.G_on = G_on

    def disturb(self, G, rng):
        key = G.shape
        if key not in self._cache:
            self._cache[key] = rng.random(G.shape) < self.p
        return np.where(self._cache[key], self.G_on, G)


class StuckDistribution(_Fault):
    """A fraction stuck at values drawn from a reflected-Gaussian KDE of measured
    stuck conductances (``means``), Scott's-rule bandwidth -- the realistic stuck
    model (cf. ``StuckDistribution`` in awarememristor)."""
    def __init__(self, probability, means, bandwidth=None):
        super().__init__()
        self.p = probability
        self.means = np.asarray(means, float)
        n = len(self.means)
        # Scott's rule (matches KDEpy default used in the reference).
        self.bw = bandwidth if bandwidth is not None else \
            self.means.std(ddof=1) * n ** (-1.0 / 5.0) if n > 1 else 0.1 * abs(self.means.mean() + 1e-9)

    def _sample(self, n, rng):
        # KDE sample: pick a mean, add Gaussian(bw), reflect at 0 (truncated >=0).
        centres = rng.choice(self.means, size=n)
        vals = centres + self.bw * rng.standard_normal(n)
        return np.abs(vals)                       # reflect negatives to keep G >= 0

    def disturb(self, G, rng):
        key = G.shape
        if key not in self._cache:
            mask = rng.random(G.shape) < self.p
            stuck = np.zeros(G.shape)
            stuck[mask] = self._sample(int(mask.sum()), rng)
            self._cache[key] = (mask, stuck)
        mask, stuck = self._cache[key]
        return np.where(mask, stuck, G)


class D2DLognormal(_Fault):
    """Lognormal device-to-device programming spread, the log-std interpolated
    linearly in resistance between R_on and R_off (cf. ``D2DLognormal``)."""
    def __init__(self, R_on_log_std, R_off_log_std, G_on=1.0, G_off=1e-3):
        super().__init__()
        self.s_on, self.s_off = R_on_log_std, R_off_log_std
        self.G_on, self.G_off = G_on, G_off

    def disturb(self, G, rng):
        key = G.shape
        if key not in self._cache:
            Gpos = np.clip(np.abs(G), self.G_off, self.G_on)
            R = 1.0 / Gpos
            R_on, R_off = 1.0 / self.G_on, 1.0 / self.G_off
            # piecewise-linear interp of log-std between R_on and R_off
            frac = np.clip((R - R_on) / (R_off - R_on + 1e-18), 0.0, 1.0)
            log_std = self.s_on + frac * (self.s_off - self.s_on)
            # mean-preserving lognormal factor on G (G = G_target * exp(N(0, sigma)))
            self._cache[key] = np.exp(log_std * rng.standard_normal(G.shape))
        return np.sign(G) * np.abs(G) * self._cache[key]


class IVNonlinearityPF(_Fault):
    """Poole--Frenkel per-device I--V nonlinearity (linearity-NONpreserving).

    Each synapse draws ``(ln c, ln d_epsilon)`` from the measured bivariate regression
    on ``ln R`` (slopes ``m``, intercepts ``b``, residual covariance ``Sigma``), then
    its read current is the PF map I = c V exp[(2e/kT) sqrt(eV / (4 pi d_epsilon))],
    not the linear product G*V. Returned as an *effective conductance* I/V so it drops
    into the existing linear einsum -- i.e. the nonlinearity is folded into a
    voltage-dependent effective weight, evaluated at the read amplitude ``V_read``.
    """
    def __init__(self, slopes, intercepts, res_cov, k_V=0.5, V_read=0.25,
                 G_on=1.0, G_off=1e-3, w_max=1.5):
        super().__init__()
        self.m = np.asarray(slopes, float); self.b = np.asarray(intercepts, float)
        self.L = np.linalg.cholesky(np.asarray(res_cov, float)
                                    + 1e-12 * np.eye(2))
        self.k_V, self.V_read = k_V, V_read
        self.G_on, self.G_off = G_on, G_off
        self.w_max = w_max                       # weight-domain scale for the I -> y map

    def disturb(self, W, rng):
        # Faithful crossbar I-V nonlinearity, following the reference pipeline
        # (awarememristor crossbar.map.w_to_G -> IVNonlinearityPF.compute_I -> map.I_to_y),
        # with the per-device regression RESIDUAL drawn ONCE (fixed device identity) and the
        # weight-dependent regression MEAN recomputed each call so the read current tracks the
        # programmed weight.
        #   1. weight -> conductance: G = k_G*|W| + G_off, k_G = (G_on - G_off)/w_max.
        #   2. PF coefficients from the bivariate regression on ln R (= -ln G); slopes come
        #      from scipy.linregress so the mean is ``+m*lnR + b`` (mnn-torch convention). A
        #      prior port used ``-m*lnR``, driving d_epsilon to ~1e-24 (vs physical ~1e-18),
        #      exploding the exponent and collapsing the output to a saturated sign-only value.
        #   3. PF read current, then map current -> weight: y = I / k_I, k_I = k_V*k_G.
        # NB: at the measured SiO_x constants over a low read window the PF map is close to a
        # mild rescale + per-device spread (a graded, sign-preserving distortion). It is the
        # linearity-NON-preserving term in the prior; the dominant programmed-weight faults are
        # stuck-at and D2D.
        key = W.shape
        if key not in self._cache:
            self._cache[key] = (self.L @ rng.standard_normal((2,) + W.shape).reshape(2, -1)
                                ).reshape((2,) + W.shape)
        eps = self._cache[key]
        Gpos = np.clip(np.abs(W), self.G_off, self.G_on)
        lnR = np.log(1.0 / Gpos)
        c = np.exp(self.m[0] * lnR + self.b[0] + eps[0])
        d_eps = np.exp(self.m[1] * lnR + self.b[1] + eps[1])
        V = self.k_V * self.V_read
        e, kT = const.elementary_charge, const.Boltzmann * (const.zero_Celsius + 20.0)
        expo = np.clip((2.0 * e / kT) * np.sqrt(e * abs(V) / (4 * np.pi * d_eps) + 1e-18),
                       None, 50.0)
        g_eff = c * np.exp(expo)
        # Mean-preserving anchor to the programmed-weight magnitude: keeps the analogue
        # dot-product gain stable (so the LIF stays in regime) while the per-device residual
        # spread + the c(G) warp remain as the PF non-ideality content. Empirically this is the
        # form under which the deep net learns through PF (PF-only ~0.73, full prior graceful);
        # a strict weight->[G_off,G_on] mapping compresses lnR so far that the residual swamps
        # the weight dependence and the net regresses to chance.
        g_eff = g_eff * (np.mean(Gpos) / (np.mean(g_eff) + 1e-30))
        return np.sign(W) * g_eff


class FaultStack:
    """Compose faults into one ``weight_fault`` hook for ``train_deep``.

    Faults apply in order; a single ``rng`` seeded once gives a fixed realisation.
    Linearity-preserving faults are applied first (they set the conductance), then the
    PF nonlinearity maps that conductance through the device I--V law.
    """
    def __init__(self, faults, seed=0):
        self.faults = faults
        self.rng = np.random.default_rng(seed)

    def __call__(self, W):
        G = W
        for f in self.faults:
            G = f.disturb(G, self.rng)
        return G


def siox_fault_stack(p_stuck=0.0, sigma_g=0.5, sigma_g_on=None, pf_on=False,
                     stuck_kind="off", stuck_means=None, seed=0):
    """Assemble the canonical SiO_x fault set from the pre-computed constants.

    Args:
        p_stuck    : fraction of stuck devices (the inference work uses 0.05; the
                     manuscript stress-tests up to 0.5).
        sigma_g    : lognormal D2D log-std at the HIGH-resistance (R_off) end. The
                     awarememristor reference's realistic device is ASYMMETRIC --
                     ``(R_on_log_std, R_off_log_std) = (0.05, 0.5)`` -- and its (0.5, 0.5)
                     symmetric setting is an explicit "high-D2D" STRESS regime. Pass
                     ``sigma_g_on`` to set the R_on end independently; if omitted it defaults
                     to ``sigma_g`` (the legacy symmetric behaviour). Set ``sigma_g=0`` to
                     disable D2D. NB: with the I-V nonlinearity now correctly active (it was
                     previously sign-only), symmetric (0.5, 0.5) drives the deep net to
                     chance, whereas the realistic (0.05, 0.5) leaves it robust -- so the
                     headline grids use the realistic spread, (0.5, 0.5) as a stress point.
        sigma_g_on : R_on-end D2D log-std; defaults to ``sigma_g`` (symmetric) if None.
        pf_on      : include the Poole--Frenkel I--V nonlinearity (always-on distortion).
        stuck_kind : "off" (open synapses, the dominant mode), "on", or "dist" (KDE of
                     measured stuck values, needs ``stuck_means``).
        stuck_means: measured stuck conductances for the KDE model (stuck_kind="dist").
    Returns a ``FaultStack`` callable usable as ``train_deep(weight_fault=...)``.
    """
    G_off, G_on = SIOX_PF["G_off"], SIOX_PF["G_on"]
    s_on = sigma_g if sigma_g_on is None else sigma_g_on
    faults = []
    if p_stuck > 0:
        if stuck_kind == "off":
            faults.append(StuckAtGOff(p_stuck, G_off))
        elif stuck_kind == "on":
            faults.append(StuckAtGOn(p_stuck, G_on))
        elif stuck_kind == "dist":
            faults.append(StuckDistribution(p_stuck,
                          stuck_means if stuck_means is not None else [G_off, G_on]))
    if sigma_g > 0 or s_on > 0:
        faults.append(D2DLognormal(s_on, sigma_g, G_on, G_off))
    if pf_on:
        faults.append(IVNonlinearityPF(SIOX_PF["slopes"], SIOX_PF["intercepts"],
                                       SIOX_PF["res_cov"], k_V=SIOX_PF["k_V"], V_read=0.25,
                                       G_on=G_on, G_off=G_off))
    return FaultStack(faults, seed=seed)


class _MazeFaultStack:
    """Apply the SiO_x fault prior to NON-NEGATIVE single-conductance weights.

    The deep network's weights are signed differential pairs in ``[-w_max, w_max]`` and
    map onto the SiO_x conductance window directly. The sequential-maze (``maze.py``)
    weights are instead a SINGLE non-negative conductance per synapse in ``[0, w_max]``;
    feeding those raw into :func:`siox_fault_stack` is physically wrong, because a weight
    of 0 is NOT ``G_off`` on that stack and ``StuckAtGOn`` (``G_on = 3.45e-3``) would map a
    shorted device to ~0 weight rather than to the maximum. This wrapper fixes the mapping:

      1. affine-map the weight ``[0, w_max] -> [G_off, G_on]`` (0 -> open, w_max -> shorted);
      2. apply the SAME fault classes (stuck-off/on, lognormal D2D, PF I-V nonlinearity)
         in that physical conductance space;
      3. affine-map the faulted conductance back to a non-negative weight, clipped to
         ``[0, w_max]``.

    So stuck-off correctly opens a synapse (weight 0), stuck-on correctly saturates it
    (weight w_max), and D2D/PF distort around the programmed value -- the same measured
    prior as the deep experiments, expressed in the maze's weight space.
    """
    def __init__(self, gc_faults, pf, w_max, G_off, G_on, seed=0):
        # Two fault groups, applied in the right domain for the NON-NEGATIVE maze weight:
        #   gc_faults : linearity-PRESERVING faults (stuck-off/on, D2D). They act in
        #               CONDUCTANCE space, so the wrapper maps weight -> [G_off, G_on],
        #               applies them, maps back to [0, w_max]. Stuck-off -> 0 weight,
        #               stuck-on -> w_max, D2D distorts around the programmed value.
        #   pf        : the PF I-V nonlinearity (or None). PF is SELF-CONTAINED: it maps
        #               weight -> conductance -> current -> weight internally (with its own
        #               w_max), so it is applied DIRECTLY to the (non-negative) maze weight,
        #               NOT pre-mapped (which would double-transform). The earlier min-max
        #               rescale hack is gone -- PF now returns weight-domain values directly.
        self.gc_faults = gc_faults
        self.pf = pf
        self.w_max, self.G_off, self.G_on = w_max, G_off, G_on
        self.rng = np.random.default_rng(seed)

    def __call__(self, W):
        # Map the non-negative maze weight into the SiO_x conductance window, apply ALL faults
        # in that conductance space (stuck/D2D AND the PF I-V map, which expects conductance-
        # scale input -- it clips |.| to [G_off, G_on]), then map back to [0, w_max]. Stuck-off
        # -> 0 weight, stuck-on -> w_max, D2D + PF distort around the programmed value. PF
        # returns a conductance-scale effective value (mean-anchored), so the round-trip stays
        # in-window.
        frac = np.clip(W / self.w_max, 0.0, 1.0)              # [0,w_max] -> [0,1]
        G = self.G_off + frac * (self.G_on - self.G_off)      # -> [G_off, G_on]
        for f in self.gc_faults:
            G = f.disturb(G, self.rng)
        if self.pf is not None:
            G = np.clip(np.abs(self.pf.disturb(G, self.rng)), self.G_off, self.G_on)
        frac_out = (G - self.G_off) / (self.G_on - self.G_off + 1e-18)
        return np.clip(frac_out, 0.0, 1.0) * self.w_max       # back to [0, w_max]


def maze_fault_stack(p_stuck=0.0, sigma_g=0.5, pf_on=False, stuck_kind="off",
                     w_max=1.5, seed=0):
    """SiO_x fault prior for the non-negative single-conductance maze weights.

    Same arguments and physical constants as :func:`siox_fault_stack`, but composed
    through the ``[0, w_max] <-> [G_off, G_on]`` conductance mapping (see
    :class:`_MazeFaultStack`). ``w_max`` must match ``maze.W_MAX`` (default 1.5).
    Returns a callable usable as ``train_sequential(weight_fault=...)``.
    """
    G_off, G_on = SIOX_PF["G_off"], SIOX_PF["G_on"]
    gc_faults = []                                   # linearity-preserving (conductance space)
    if p_stuck > 0:
        if stuck_kind == "off":
            gc_faults.append(StuckAtGOff(p_stuck, G_off))
        elif stuck_kind == "on":
            gc_faults.append(StuckAtGOn(p_stuck, G_on))
    if sigma_g > 0:
        gc_faults.append(D2DLognormal(sigma_g, sigma_g, G_on, G_off))
    pf = None
    if pf_on:
        # PF maps weight<->conductance internally, so it carries the maze w_max.
        pf = IVNonlinearityPF(SIOX_PF["slopes"], SIOX_PF["intercepts"], SIOX_PF["res_cov"],
                              k_V=SIOX_PF["k_V"], V_read=0.25, G_on=G_on, G_off=G_off,
                              w_max=w_max)
    return _MazeFaultStack(gc_faults, pf, w_max, G_off, G_on, seed=seed)
\end{lstlisting}

\begin{lstlisting}[caption={siox\_eligibility/hybrid.py},label={lst:siox-src-hybrid}]
r"""Hybrid architecture -- a BPTT perception front-end + the device-eligibility decision stack.

This module answers the *high-dimensional input* objection to the device eligibility
trace ("real RL has rich perceptual input"), honestly and by DIVISION OF LABOUR: the
slow device trace is NOT a perception engine, so a conventional (gradient-trained,
frozen) perception stage supplies a low-dimensional state, and the all-local device
learning stack learns the policy behind it. The claim is *composition*, never that the
device does perception.

Two complementary studies live here, sharing the same perception front-end:

  * **exp5** (:func:`run_hybrid_decision`, manuscript Fig. 9): a conv-SNN perception
    stage (snntorch ``Leaky`` + surrogate gradient, rate-coded input) classifies noisy
    oriented gratings and is frozen; a contextual spiking bandit then reads either its
    ``C``-dim frozen readout (*hybrid*), the ``P`` raw pixels (*raw*, no front-end), or
    the hybrid wiring with the eligibility zeroed (*no_trace*). Plastic device synapses
    carry the :class:`~mrl_trace.bandit.GateBankBatched` eligibility and are
    updated by the signed three-factor rule at EQUAL RL budget. Pre-registered: C1
    front-end acc >= 0.85; C2 hybrid >= criterion and no-trace at chance; C3 hybrid -
    raw >= 0.15. Reference (MPS): acc 0.87, hybrid 0.62, raw 0.27, no-trace 0.27 -- the
    front-end is the remedy. Writes ``tier6_results.npy``.

  * **exp16** (:func:`run_hybrid_scale`, Tier 3 batch sweep): a WIDTH-SCALING +
    FAULT-TOLERANCE robustness study. STAGE 1 (perception, once, offline, optional GPU)
    trains the SAME conv-SNN and caches its low-dimensional readouts to ``readouts.npz``
    (:func:`prep_readouts`). STAGE 2 (decision, pure NumPy) feeds those cached readouts
    as the STATE of the DEEP all-local agent (:func:`~mrl_trace.deep.train_deep`:
    trace + DFA + homeostasis) via its ``state_sampler`` hook, and sweeps hidden width
    ``H`` x SiO_x stuck-fault fraction ``p`` under the measured device-fault prior. If
    ``readouts.npz`` is ABSENT the sweep falls back to an HONEST proxy front-end (clearly
    logged): class-conditional noisy embeddings whose clusters overlap and whose label is
    a non-linearly-separable XOR of two latent bits, so depth is genuinely required. This
    is a stand-in for the readout *statistics*, not a silent fake of perception -- the
    headline result should use the real cached readouts. Writes ``exp16_hybrid_scale.npy``.

SCOPE (read before citing exp16). exp16 is scoped to width-scaling + fault-tolerance
ONLY: does the all-local device stack, behind a real frozen perception front-end, keep
learning the policy as the hidden layer widens and the stuck-fault fraction rises? It is
NOT a homeostasis experiment and carries NO homeostasis ablation. Homeostasis is
hidden-CAPACITY-dependent (it opposes a winner-take-all hidden-unit collapse), so at this
study's wide hidden layers (H up to 512, overcomplete) it is irrelevant at any strength --
a verified regime search found hybrid-full ~ no-homeo here (gap ~0). The homeostasis
CONTRIBUTION is large and CI-clean only at MODERATE width (H=32) under temporal
interference + device faults, which is exp13's job (gap +0.37 to +0.49, disjoint CIs);
citing exp16's homeostasis-neutrality as evidence about homeostasis would be a category
error, which is why the no-homeo arm is omitted rather than reported as a null.

ACCURACY SCOPE (exp16). On the real A=4 readouts the device policy CONVERGES to ~0.51
reward (flat from 2000 to 4000 trials), well above chance (0.25) and far above the
no-trace control (~0.27), but BELOW the frozen front-end's own classification ceiling
(~0.87). The claim is that the all-local stack learns a SUBSTANTIAL policy behind a real
perception front-end and that this DEGRADES GRACEFULLY with stuck faults across width --
NOT that it matches the front-end's accuracy.

The perception front-end (torch + snntorch) is the OPTIONAL ``hybrid`` extra. It is
imported LAZILY inside the functions that need it, so this module -- and the exp16
Stage-2 sweep on cached/proxy readouts -- imports and runs with NO torch installed. Only
:func:`run_frontend`, :func:`run_hybrid_decision` and :func:`prep_readouts` require torch;
:func:`run_hybrid_scale` never does (it consumes the cached ``.npz`` or the proxy).
Install the extra with ``pip install -e ".[hybrid]"``.
"""
from __future__ import annotations

import numpy as np

from .bandit import GateBankBatched, W_INIT, W_MAX
from .deep import train_deep
from .device_faults import siox_fault_stack
from .stats import bootstrap_ci

__all__ = [
    # exp5 perception front-end + hybrid decision (Fig 9)
    "IMG", "C", "NOISE",
    "make_grating", "make_batch",
    "run_frontend", "run_readout_pool", "run_hybrid_decision",
    # exp16 stage-1 caching + stage-2 sweep
    "prep_readouts", "run_hybrid_scale",
    # analysis helpers
    "final_rate",
]

# =============================================================================
# exp5 -- perception task geometry (noisy oriented gratings) and front-end
# =============================================================================
IMG = 16          # image side (P = 256 pixels)
C = 4             # orientation classes (0, 45, 90, 135 deg)
NOISE = 0.9       # additive noise std (high -> raw-pixel RL is genuinely hard)
_N_SEEDS = 20     # exp5 decision-layer seeds (comp-neuro-grade; batched, so ~free)


def _torch():
    """Lazily import the optional ``hybrid`` extra (torch + snntorch), raising the same
    actionable message as the original driver if it is not installed. Keeping the import
    inside the functions that need it lets this module -- and the pure-NumPy exp16 Stage-2
    sweep -- load and run with no torch present."""
    try:
        import torch
        import torch.nn as nn
        import torch.nn.functional as F
        import snntorch as snn
        from snntorch import surrogate
    except ImportError as exc:  # pragma: no cover
        raise SystemExit(
            "The hybrid perception front-end needs torch + snntorch. Install the "
            'optional extra:\n    pip install -e ".[hybrid]"'
        ) from exc
    dev = torch.device("mps" if torch.backends.mps.is_available()
                       else "cuda" if torch.cuda.is_available() else "cpu")
    return torch, nn, F, snn, surrogate, dev


def make_grating(theta, rng, img=IMG):
    """A single noisy oriented sinusoidal grating at orientation ``theta`` (radians)."""
    x = np.linspace(-1, 1, img)
    xx, yy = np.meshgrid(x, x)
    proj = xx * np.cos(theta) + yy * np.sin(theta)
    g = 0.5 + 0.5 * np.sin(2 * np.pi * 3.0 * proj)
    g = g + NOISE * rng.standard_normal((img, img))
    g = (g - g.min()) / (g.max() - g.min() + 1e-9)
    return g.astype(np.float32)


def make_batch(n, rng):
    """A batch of ``n`` gratings with class labels (one of ``C`` orientations each)."""
    thetas = np.array([0, np.pi / 4, np.pi / 2, 3 * np.pi / 4])
    y = rng.integers(C, size=n)
    X = np.stack([make_grating(thetas[c], rng) for c in y])
    return X[:, None, :, :], y


def _build_convsnn(num_steps=20, beta=0.9):
    """Construct the compact convolutional spiking net (two conv+LIF stages + a readout
    LIF), trained by BPTT with a surrogate gradient on rate-coded input. Defined inside a
    factory so the ``nn.Module`` subclass is created only when torch is present (lazy)."""
    _torch_, nn, F, snn, surrogate, _dev = _torch()

    class ConvSNN(nn.Module):
        def __init__(self, num_steps=num_steps, beta=beta):
            super().__init__()
            sg = surrogate.fast_sigmoid()
            self.num_steps = num_steps
            self.conv1 = nn.Conv2d(1, 8, 3, padding=1)
            self.conv2 = nn.Conv2d(8, 16, 3, padding=1)
            self.lif1 = snn.Leaky(beta=beta, spike_grad=sg)
            self.lif2 = snn.Leaky(beta=beta, spike_grad=sg)
            self.lif3 = snn.Leaky(beta=beta, spike_grad=sg)
            h = IMG // 2 // 2
            self.fc = nn.Linear(16 * h * h, C)

        def forward(self, x):
            m1 = self.lif1.init_leaky()
            m2 = self.lif2.init_leaky()
            m3 = self.lif3.init_leaky()
            memsum = 0.0
            spksum = 0.0
            for _ in range(self.num_steps):
                xt = (_torch_.rand_like(x) < x).float()
                c1 = F.max_pool2d(self.conv1(xt), 2); s1, m1 = self.lif1(c1, m1)
                c2 = F.max_pool2d(self.conv2(s1), 2); s2, m2 = self.lif2(c2, m2)
                c3 = self.fc(s2.flatten(1)); s3, m3 = self.lif3(c3, m3)
                memsum = memsum + m3
                spksum = spksum + s3
            return memsum / self.num_steps, spksum

    return ConvSNN()


def run_frontend(steps=600, bs=128, lr=2e-3, seed=0):
    """Train the perception front-end (conv-SNN, BPTT surrogate gradient) and freeze it.

    Returns ``(net, acc)`` where ``net`` is the frozen ``torch`` module and ``acc`` is its
    held-out classification accuracy on 2000 fresh gratings. Needs the ``hybrid`` extra.
    """
    torch, nn, F, snn, surrogate, dev = _torch()
    torch.manual_seed(seed)
    rng = np.random.default_rng(seed)
    net = _build_convsnn().to(dev)
    opt = torch.optim.Adam(net.parameters(), lr=lr)
    lossf = nn.CrossEntropyLoss()
    net.train()
    for _ in range(steps):
        X, y = make_batch(bs, rng)
        mem, _ = net(torch.tensor(X, device=dev))
        loss = lossf(mem, torch.tensor(y, device=dev))
        opt.zero_grad(); loss.backward(); opt.step()
    net.eval()
    with torch.no_grad():
        Xv, yv = make_batch(2000, rng)
        mv, _ = net(torch.tensor(Xv, device=dev))
        acc = (mv.argmax(1).cpu().numpy() == yv).mean()
    return net, float(acc)


def run_readout_pool(net, per_class=500, seed=1):
    """Extract the frozen front-end's per-class spike-rate readouts + raw-pixel pools.

    Returns ``(feats, raws)`` -- dicts ``class -> (per_class, C)`` normalised readouts and
    ``class -> (per_class, P)`` raw pixels respectively. Needs the ``hybrid`` extra.
    """
    torch, nn, F, snn, surrogate, dev = _torch()
    rng = np.random.default_rng(seed)
    net.eval()
    feats, raws = {}, {}
    thetas = np.array([0, np.pi / 4, np.pi / 2, 3 * np.pi / 4])
    with torch.no_grad():
        for c in range(C):
            X = np.stack([make_grating(thetas[c], rng) for _ in range(per_class)])
            _, spk = net(torch.tensor(X[:, None], device=dev))
            out = np.clip(spk.cpu().numpy(), 0, None)
            out = out / (out.max(1, keepdims=True) + 1e-9)
            feats[c] = out
            raws[c] = X.reshape(per_class, -1)
    return feats, raws


def _decision(state_pool, S, A, B=4, trials=1500, dt=5e-3, cue_dur=1.0, D=2.0,
              eta=0.2, in_rate=200.0, ltd=0.6, tau_m=20e-3, v_th=1.0, noise=0.15,
              tau_leak=10.0, no_trace=False, seed0=0):
    """Contextual bandit over ``C`` classes; ``state_pool`` is a dict ``class -> (n, S)``
    activations (``S=C`` for the hybrid readout, ``S=P`` for raw pixels). Uses the
    package's :class:`~mrl_trace.bandit.GateBankBatched` for the device eligibility
    and returns rewards ``(B, trials)``. Pure NumPy; no torch. This is the exp5 decision
    layer -- its inline current-impulse LIF matches ``neurons.lif_step_batched`` exactly
    (v_reset=0), kept inline here to preserve the RNG call ordering bit-for-bit."""
    rng = np.random.default_rng(seed0)
    bank = GateBankBatched(B, S, A, tau_leak=tau_leak, dt=dt)
    w = np.full((B, S, A), W_INIT)
    correct = np.arange(C) % A
    baseline = np.full(B, 1.0 / A)
    cue = (0.3, 0.3 + cue_dur)
    ri = int((cue[1] + D) / dt)
    nsteps = ri + 2
    rewards = np.zeros((B, trials))
    pools = {c: np.asarray(state_pool[c]) for c in range(C)}
    for tr in range(trials):
        bank.reset()
        cls = rng.integers(C, size=B)
        act = np.zeros((B, S))
        for b in range(B):
            p = pools[cls[b]]
            act[b] = p[rng.integers(len(p))]
        v = np.zeros((B, A)); spk = np.zeros((B, A)); e_rew = np.zeros((B, S, A))
        for n in range(nsteps):
            t = n * dt
            pre = np.zeros((B, S))
            if cue[0] <= t < cue[1]:
                pre = (rng.random((B, S)) < in_rate * dt * act).astype(float)
            charge = np.einsum('bsa,bs->ba', w, pre)
            v = v + charge - dt * v / tau_m + noise * rng.standard_normal((B, A))
            sp = v >= v_th; v = np.where(sp, 0.0, v); spk += sp
            drive = pre[:, :, None] * np.where(sp, 1.0, -ltd)[:, None, :]
            if no_trace:
                drive[:] = 0.0
            e = bank.step(drive)
            if n == ri:
                e_rew = e.copy()
        tie = spk.max(1) == spk.min(1)
        chosen = np.argmax(spk, 1); chosen[tie] = rng.integers(A, size=tie.sum())
        r = (chosen == correct[cls]).astype(float)
        w = np.clip(w + (eta * (r - baseline))[:, None, None] * e_rew, 0.0, W_MAX)
        baseline += 0.02 * (r - baseline)
        rewards[:, tr] = r
    return rewards


def final_rate(rw, window=100):
    """Per-seed reward rate over the last ``window`` trials (axis -1)."""
    rw = np.asarray(rw)
    return rw[:, -window:].mean(1)


def run_hybrid_decision(*, seeds=_N_SEEDS, front_steps=600, front_seed=0,
                        per_class=500, trials=1500, D=2.0):
    """exp5 (Fig 9): train the perception front-end, then run the decision-layer RL at
    ``seeds`` seeds in three conditions -- *hybrid* (reads the C-dim readout), *raw*
    (reads P pixels, no front-end) and *no_trace* (hybrid wiring, eligibility zeroed) --
    at equal RL budget.

    Pre-registered (verbatim): C1 front-end acc >= 0.85; C2 hybrid >= criterion AND
    no-trace at chance (<= chance + 0.10); C3 the front-end is the remedy (hybrid - raw
    >= 0.15). ``criterion = 0.5*(1 + 1/A)``, ``A = C``. Reference (MPS): acc 0.87, hybrid
    0.62, raw 0.27, no-trace 0.27.

    Returns the result grid as a plain dict (no file I/O, no plotting, no stdout); needs
    the ``hybrid`` extra for the front-end. ``main()`` writes it to ``tier6_results.npy``.
    """
    A = C
    chance = 1.0 / A
    crit = 0.5 * (1 + chance)
    net, acc = run_frontend(steps=front_steps, seed=front_seed)
    feats, raws = run_readout_pool(net, per_class=per_class, seed=front_seed + 1)

    per_seed = {
        "hybrid": final_rate(_decision(feats, S=C, A=A, B=seeds, trials=trials, D=D)),
        "raw": final_rate(_decision(raws, S=IMG * IMG, A=A, B=seeds, trials=trials, D=D)),
        "no_trace": final_rate(_decision(feats, S=C, A=A, B=seeds, trials=trials, D=D,
                                         no_trace=True)),
    }
    res = {}
    for k, v in per_seed.items():
        lo, hi = bootstrap_ci(v)
        res[k] = (float(v.mean()), float(v.std()), lo, hi)

    hyb, raw, nt = res["hybrid"][0], res["raw"][0], res["no_trace"][0]
    c2 = (hyb >= crit) and (nt <= chance + 0.10)
    c3 = (hyb >= crit) and (hyb - raw >= 0.15)
    return dict(front_acc=acc, chance=chance, crit=crit, results=res,
                C1=bool(acc >= 0.85), C2=bool(c2), C3=bool(c3),
                C=C, A=A, P=IMG * IMG, noise=NOISE, n_seeds=seeds)


# =============================================================================
# exp16 Stage 1 -- cache the frozen front-end readouts (needs the hybrid extra)
# =============================================================================
def prep_readouts(per_class=800, steps=600, seed=0, out_name="readouts.npz"):
    """exp16 Stage 1: train the SAME conv-SNN front-end ONCE and cache its readouts to
    ``data/results/<out_name>`` in the schema Stage 2 expects::

        readouts.npz : X (N, F) float32 readout vectors, y (N,) int class labels, acc () float

    where ``N = per_class * C``. Rows are shuffled so a sequential sampler sees mixed
    classes. This is the offline perception step that upgrades the exp16 sweep from
    "proxy statistics" to the real frozen-front-end result; it is NOT part of the pure-NumPy
    Stage-2 batch job. Needs the ``hybrid`` extra (torch + snntorch); returns the written
    path. Writes through :mod:`mrl_trace.paths` (never to ``experiments/``).
    """
    from . import paths
    net, acc = run_frontend(steps=steps, seed=seed)
    feats, _raws = run_readout_pool(net, per_class=per_class, seed=seed + 1)
    Xs, ys = [], []
    for c in sorted(feats):
        Xs.append(np.asarray(feats[c], dtype=np.float32))
        ys.append(np.full(len(feats[c]), c, dtype=np.int64))
    X = np.concatenate(Xs, 0)
    y = np.concatenate(ys, 0)
    rng = np.random.default_rng(seed + 2)
    perm = rng.permutation(len(y))
    X, y = X[perm], y[perm]
    out = paths.results_dir() / out_name
    np.savez_compressed(out, X=X, y=y, acc=np.float32(acc))
    return out


# =============================================================================
# exp16 Stage 2 -- pure-NumPy width x fault sweep behind the readouts (or proxy)
# =============================================================================
N_BOOT = 10000
# Realistic ASYMMETRIC lognormal D2D (R_on_log_std=0.05, R_off_log_std=0.5), matching the
# awarememristor reference's realistic SiO_x device; symmetric (0.5,0.5) is its "high-D2D"
# stress regime. With the I-V nonlinearity now correctly active, symmetric 0.5 drives the net
# to chance while the realistic spread stays robust -- so the realistic spread is the headline.
SIGMA_G = 0.5            # R_off-end log-std
SIGMA_G_ON = 0.05        # R_on-end log-std
# Early-stop template; the chance/criterion are set per A in run_hybrid_scale (chance = 1/A,
# converged criterion scaled above chance) so multi-way readout tasks (A>2) are not falsely
# declared "stuck" at a still-improving rate.
EARLY_STOP = {"min_trials": 400, "check_every": 100, "window": 150,
              "criterion": 0.78, "tol": 0.03, "chance": 0.5, "stuck_after": 1200}
# Proxy front-end geometry (used only when no real readouts.npz is supplied).
PROXY_F = 8                    # readout dimensionality (front-end embedding width)
PROXY_NOISE = 0.35             # per-dim Gaussian noise -> overlapping class clusters

# exp16 is a WIDTH + FAULT-TOLERANCE robustness study (the all-local stack behind a real frozen
# front-end), NOT a homeostasis experiment. The no_homeo ablation is DELIBERATELY OMITTED: at
# this study's wide hidden layers (H=128/512, overcomplete) homeostasis is irrelevant at any
# strength (verified regime search: hybrid_full ~ no_homeo, gap ~0), while its contribution is
# large and CI-clean only at MODERATE H (=32) under interference+faults, which is exp13's job.
# Conditions kept: the device stack vs the eligibility-necessity control.
CONDS = {
    "hybrid_full": dict(mode="dfa", homeo=0.1),   # full stack (homeo at its validated strength)
    "no_trace":    dict(mode="no_trace", homeo=0.1),  # eligibility-necessity control
}


def _early_stop_for(A):
    """A-aware early-stop: chance = 1/A, and 'converged' is a high bar above chance so the
    stuck-padding never freezes a run that is still climbing on a >2-way task."""
    chance = 1.0 / A
    es = dict(EARLY_STOP)
    es["chance"] = chance
    es["criterion"] = round(chance + 0.55 * (1.0 - chance), 3)   # A=2 -> 0.78; A=4 -> 0.66
    return es


def _load_readouts(path):
    """Load cached frozen-front-end readouts. Returns ``(X (N,F) float in [0,1], y (N,) int,
    A)``. Readouts are min-max normalised per dimension so they are interpretable as input
    rates (the same [0,1] line-rate convention the built-in one-hot cue uses)."""
    d = np.load(path)
    X = np.asarray(d["X"], dtype=float)
    y = np.asarray(d["y"], dtype=int)
    lo = X.min(0, keepdims=True); hi = X.max(0, keepdims=True)
    X = (X - lo) / (hi - lo + 1e-9)
    return X, y, int(y.max()) + 1


def _make_sampler(ro, F):
    """Build a ``state_sampler(rng, B) -> (lines (B,F) in [0,1], label (B,))`` for
    ``train_deep``.

    With real readouts (``ro = (X, y, A)``): sample rows of the cached ``(X, y)`` -- the
    device stack must learn the policy behind the frozen front-end's actual embedding.
    With the proxy (``ro is None``): class-conditional noisy embeddings (overlapping
    clusters) whose label is XOR of two latent bits, so depth is genuinely required. Either
    way only the STATE source changes; the learning rule, eligibility, DFA, homeostasis and
    fault prior are identical to exp14.
    """
    if ro is not None:
        X, y, _A = ro
        N = len(y)

        def sampler(rng, B):
            idx = rng.integers(0, N, size=B)
            return X[idx], y[idx]
        return sampler

    # fixed class-mean templates (deterministic; the noise gives the spread)
    base = np.array([[1, 0, 1, 0, 0, 1, 0, 1],
                     [0, 1, 0, 1, 1, 0, 1, 0]], dtype=float)
    mu = np.resize(base, (2, F))

    def sampler(rng, B):
        b0 = rng.integers(2, size=B); b1 = rng.integers(2, size=B)
        y = (b0 ^ b1).astype(int)
        x = mu[y] + PROXY_NOISE * rng.standard_normal((B, F))
        return np.clip(x, 0.0, 1.0), y
    return sampler


def _hybrid_scale_run(job):
    """One (H, p, cond, seed) run of the exp16 sweep -> ``(H, p, cond, seed, final_rate)``.

    Module-level (not a closure) so it is picklable by ``multiprocessing.Pool`` in
    ``main()``. Reads the sweep config from the ``job`` tuple's trailing ``cfg`` dict
    ``(F, A, ro, trials)`` so no module-global mutable state has to cross the fork barrier.
    """
    H, p, cond, seed, cfg = job
    F, A, ro, trials = cfg["F"], cfg["A"], cfg["ro"], cfg["trials"]
    # Programmed-conductance fault prior (stuck-at + realistic asymmetric D2D). pf_on=False:
    # the Poole-Frenkel I-V term is a read/inference nonlinearity, out of scope for a rule
    # that credits the programmed weight (see entry_array_scale for the full rationale).
    fault = siox_fault_stack(p_stuck=p, sigma_g=SIGMA_G, sigma_g_on=SIGMA_G_ON, pf_on=False,
                             stuck_kind="off", seed=1000 * seed + int(100 * p))
    sampler = _make_sampler(ro, F)
    kw = CONDS[cond]
    rewards = train_deep(
        H=H, tau_leak=10.0, D=5.0, t_distract=3.0, distract_dur=0.3,
        trials=trials, seed0=seed, weight_fault=fault,
        early_stop=_early_stop_for(A),
        state_sampler=sampler, n_features=F, n_actions=A,
        **kw,
    )
    return (H, p, cond, seed, float(rewards[:, -200:].mean()))


def run_hybrid_scale(*, H_grid=(32, 128, 512), p_grid=(0.0, 0.2, 0.5), seeds=12,
                     trials=2000, readouts=None, pool=None):
    """exp16 Stage 2 (Tier 3): the width-scaling + fault-tolerance sweep behind the frozen
    front-end. Pure NumPy -- NO torch. Serial by default (a notebook can call it in-kernel);
    ``main()`` passes a ``multiprocessing.Pool`` in ``pool`` to parallelise the (H, p, cond,
    seed) grid.

    ``readouts`` is a path to the cached ``readouts.npz`` (Stage 1 output). If it is None or
    the file is absent, the sweep falls back to the HONEST proxy front-end (overlapping class
    clusters, F=8, XOR label, A=2) -- clearly recorded in the returned ``front_end`` field.
    The headline result should pass the real ``readouts.npz``.

    Conditions per (H, p): *hybrid_full* (the device stack: trace + DFA + homeostasis) and
    *no_trace* (eligibility-necessity control). Metric: final reward rate (mean over the last
    200 trials), summarised as mean + bootstrap 95% CI over seeds.

    Returns the result grid dict with the exact keys the driver wrote (H, p, conds, summary,
    readouts, front_end, F, A, seeds, trials, claim); no file I/O. ``main()`` writes it to
    ``exp16_hybrid_scale.npy``.
    """
    import os

    H_grid = [int(x) for x in H_grid]
    p_grid = [float(x) for x in p_grid]
    have_ro = bool(readouts) and os.path.exists(readouts)
    if have_ro:
        X, y, A = _load_readouts(readouts)
        ro, F = (X, y, A), X.shape[1]
    else:
        ro, F, A = None, PROXY_F, 2

    claim = (f"width-scaling + fault-tolerance robustness behind a "
             f"{'real frozen' if ro is not None else 'proxy'} front-end "
             f"(A={A}); homeostasis contribution is exp13's claim, NOT tested here")

    cfg = {"F": F, "A": A, "ro": ro, "trials": trials}
    jobs = [(H, p, cond, s, cfg) for H in H_grid for p in p_grid
            for cond in CONDS for s in range(seeds)]
    if pool is None:
        results = [_hybrid_scale_run(j) for j in jobs]
    else:
        results = pool.map(_hybrid_scale_run, jobs, chunksize=1)

    acc = {}
    for H, p, cond, s, val in results:
        acc.setdefault((H, p, cond), []).append(val)

    summary = {}
    for H in H_grid:
        for p in p_grid:
            for cond in CONDS:
                v = np.array(acc[(H, p, cond)])
                lo, hi = _boot_ci(v, seed=H + int(100 * p))
                summary[(H, p, cond)] = (float(v.mean()), lo, hi)

    return {"H": H_grid, "p": p_grid, "conds": list(CONDS), "summary": summary,
            "readouts": bool(have_ro), "front_end": "real" if have_ro else "proxy",
            "F": F, "A": A, "seeds": seeds, "trials": trials, "claim": claim}


def _boot_ci(vals, seed=0):
    """Percentile bootstrap 95% CI for the mean of ``vals`` (matches the original driver's
    inline bootstrap, N_BOOT resamples, keyed seed for reproducibility)."""
    rng = np.random.default_rng(seed)
    b = [vals[rng.integers(0, len(vals), len(vals))].mean() for _ in range(N_BOOT)]
    return float(np.percentile(b, 2.5)), float(np.percentile(b, 97.5))


# =============================================================================
# Full-scale reproduction CLI
# =============================================================================
def main(argv=None):
    r"""Full-scale reproduction CLI for the hybrid grids (writes ``data/results``).

    ``python -m mrl_trace.hybrid [--exp5] [--exp16] [--prep-readouts]
    [--readouts PATH] [--full|--quick]``

    With no experiment flag, runs exp5 + exp16. ``--full`` = published scale (exp5 20 seeds,
    exp16 H=32,128,512 x p=0,0.2,0.5 x 12 seeds x 2000 trials); ``--quick`` = a fast
    few-seed smoke run. ``--prep-readouts`` (Stage 1) trains the front-end and caches
    ``readouts.npz`` -- needs the ``hybrid`` extra. ``--readouts PATH`` points exp16 Stage 2
    at a cached readouts file; with no path it defaults to ``data/results/readouts.npz`` if
    present, else the honest proxy front-end.
    """
    import argparse
    import multiprocessing as mp
    import os
    from . import paths

    ap = argparse.ArgumentParser(description="Hybrid perception-front-end reproductions")
    ap.add_argument("--exp5", action="store_true",
                    help="hybrid decision (Fig 9) -> tier6_results.npy (needs torch)")
    ap.add_argument("--exp16", action="store_true",
                    help="width x fault sweep behind the front-end -> exp16_hybrid_scale.npy")
    ap.add_argument("--prep-readouts", action="store_true",
                    help="Stage 1: train front-end, cache readouts.npz (needs torch)")
    ap.add_argument("--readouts", default=None,
                    help="path to cached readouts.npz for exp16 Stage 2 "
                         "(default: data/results/readouts.npz if present, else proxy)")
    ap.add_argument("--workers", type=int, default=int(os.cpu_count() or 4),
                    help="worker processes for the exp16 sweep")
    ap.add_argument("--quick", action="store_true", help="fast few-seed smoke run")
    ap.add_argument("--full", action="store_true", help="published-scale run (default)")
    a = ap.parse_args(argv)
    run_all = not (a.exp5 or a.exp16 or a.prep_readouts)

    if a.prep_readouts:
        print("=== exp16 Stage 1: caching frozen front-end readouts ===")
        per_class = 200 if a.quick else 800
        steps = 120 if a.quick else 600
        out = prep_readouts(per_class=per_class, steps=steps)
        print(f"  wrote {out}")

    if a.exp5 or run_all:
        seeds = 6 if a.quick else _N_SEEDS
        trials = 400 if a.quick else 1500
        front_steps = 120 if a.quick else 600
        per_class = 200 if a.quick else 500
        print(f"=== exp5 hybrid decision (Fig 9): {seeds} seeds, {trials} trials ===")
        grid = run_hybrid_decision(seeds=seeds, front_steps=front_steps,
                                   per_class=per_class, trials=trials)
        paths.save_result("tier6_results.npy", grid)
        print(f"  wrote tier6_results.npy  acc={grid['front_acc']:.3f} "
              f"C1={grid['C1']} C2={grid['C2']} C3={grid['C3']}")

    if a.exp16 or run_all:
        # Default the readouts path to the bundled Stage-1 cache if the caller did not
        # supply one; falls back to the honest proxy if that file is absent.
        ro_path = a.readouts
        if ro_path is None:
            cand = paths.results_dir() / "readouts.npz"
            ro_path = str(cand) if cand.exists() else None
        if a.quick:
            H_grid, p_grid, seeds, trials = (32,), (0.0, 0.5), 4, 400
        else:
            H_grid, p_grid, seeds, trials = (32, 128, 512), (0.0, 0.2, 0.5), 12, 2000
        try:
            ctx = mp.get_context("fork")
        except ValueError:                                    # pragma: no cover
            ctx = mp.get_context()
        front = "real" if (ro_path and os.path.exists(ro_path)) else "proxy"
        print(f"=== exp16 hybrid-scale sweep: H={list(H_grid)} p={list(p_grid)} "
              f"seeds={seeds} front_end={front} ===")
        with ctx.Pool(a.workers) as pool:
            grid = run_hybrid_scale(H_grid=H_grid, p_grid=p_grid, seeds=seeds,
                                    trials=trials, readouts=ro_path, pool=pool)
        for H in grid["H"]:
            for p in grid["p"]:
                line = f"  H={H:4d} p={p:.2f}: "
                for cond in grid["conds"]:
                    m, lo, hi = grid["summary"][(H, p, cond)]
                    line += f"{cond}={m:.2f}[{lo:.2f},{hi:.2f}] "
                print(line)
        paths.save_result("exp16_hybrid_scale.npy", grid)
        print(f"  wrote exp16_hybrid_scale.npy  front_end={grid['front_end']}")


if __name__ == "__main__":
    main()
\end{lstlisting}

\section{Reproduction notebooks (\texttt{experiments/})}
\label{app:siox-nb}

\noindent Each notebook below is the code of one topic's cells, in order. Run them against the assembled package; by default they replay the bundled result grids and draw every figure, and \texttt{REPRODUCE\_ALL} runs them all. The markdown narration of each notebook is omitted here (it restates the chapter text); only the executable cells are printed.

\begin{lstlisting}[caption={experiments/01\_trace\_and\_credit\_window.ipynb},label={lst:siox-nb-01-trace-and-credit-window}]
# ---- cell 1 --------------------------------------------------------
import numpy as np
import matplotlib.pyplot as plt
from mrl_trace import paths
from mrl_trace.device import TransientGate, tau_r
from mrl_trace.stats import bootstrap_ci

# QUICK=True re-runs a fast, few-seed version of an experiment in-kernel (serial, no Pool);
# the default (QUICK=False) REPLAYS the committed 20-seed grid so every figure renders
# instantly. Heavy sweeps (exp8's dense 13x18 grid) always document a `--full` shell
# command instead and are never a slow default cell.
QUICK = False

# device retention curves are coloured by a sequential viridis (long->short), matching
# the manuscript trace-window figure; controls/fits use the shared accent palette.
GREEN, INDIGO, RED, GOLD, GREY, INK = "#3aa07a", "#2f4b8f", "#c0392b", "#e0a93b", "#9aa6b2", "#2b2b2b"
VIR = [plt.cm.viridis(x) for x in (0.85, 0.5, 0.15)]   # tau_leak long -> short

def _clean(ax):
    for sp in ("top", "right"):
        ax.spines[sp].set_visible(False)
    ax.set_axisbelow(True); ax.grid(True, color="0.88", lw=0.5)

print("data/results:", paths.results_dir())
print(f"operating point V=1.5 -> fitted rise time tau_r = {tau_r(1.5):.2f} s")

# ---- cell 2 --------------------------------------------------------
# Pure device trace: deterministic, instant, computed inline (no grid replay).
t = np.arange(0.0, 80.0, 0.05)          # s
coincidence_at, coincidence_dur = 1.0, 0.5
taus = [1.0, 5.0, 20.0]                  # the pre-registered retention band (s)
labels = [rf"$\tau_{{\rm leak}}={tl:g}$ s" for tl in taus]

fig, ax = plt.subplots(figsize=(7.0, 3.8))
ax.axvspan(coincidence_at, coincidence_at + coincidence_dur, color=GOLD, alpha=0.30, lw=0)
ax.text(coincidence_at + coincidence_dur / 2, 1.06, "coincidence", color=INK,
        fontsize=8, ha="center")
for tl, col, lab in zip(taus, VIR, labels):
    g = TransientGate(V=1.5, tau_leak=tl, dt=0.05)          # fitted device @ V=1.5
    e = g.trace(t, coincidence_at=coincidence_at, coincidence_dur=coincidence_dur)
    ax.plot(t, e, color=col, lw=2.0, label=lab)
    ax.fill_between(t, 0, e, color=col, alpha=0.10, lw=0)
ax.set_xlabel("time after coincidence (s)"); ax.set_ylabel(r"eligibility trace $e(t)$ (norm.)")
ax.set_xlim(0, 80); ax.set_ylim(0, 1.12)
ax.set_title(r"Device trace: compressed-exponential rise, $\tau_{\rm leak}$-set decay")
ax.legend(frameon=False, fontsize=9, title="deep traps  <->  shallow traps",
          title_fontsize=8); _clean(ax); plt.show()

# half-peak decay delay grows with tau_leak -- the device-physics origin of the window
for tl in taus:
    g = TransientGate(V=1.5, tau_leak=tl, dt=0.05)
    e = g.trace(t, coincidence_at=coincidence_at, coincidence_dur=coincidence_dur)
    pk = t[int(np.argmax(e))]
    post_m = t > pk
    tp, ep = t[post_m], e[post_m]
    d_half = (tp[int(np.argmin(np.abs(ep - 0.5)))] - pk) if (ep < 0.5).any() else float("nan")
    print(f"  tau_leak={tl:5.1f} s: peak @ {pk:4.1f} s, half-peak decay ~ {d_half:5.1f} s")

# ---- cell 3 --------------------------------------------------------
if QUICK:
    from mrl_trace.bandit import run_learning_and_window
    r = run_learning_and_window(seeds=6)
else:
    r = paths.load_result("tier3_results.npy")

delays = np.asarray(r["delays"], float)          # [1, 2, 5, 10, 20]
taus = sorted(r["reward_rate"].keys(), reverse=True)   # [10.0, 2.0, 0.5]
tcol = {10.0: VIR[0], 2.0: VIR[1], 0.5: VIR[2]}
W = 50                                            # running-mean window (matches Fig 6)

def _running(rw_2d, w=W):
    cs = np.cumsum(np.insert(rw_2d, 0, 0.0, axis=1), axis=1)
    rr = (cs[:, w:] - cs[:, :-w]) / w             # (seeds, trials-w)
    return rr.mean(0), np.percentile(rr, 2.5, axis=0), np.percentile(rr, 97.5, axis=0)

fig, (axA, axB) = plt.subplots(1, 2, figsize=(10.0, 3.8))

# (a) delay x retention reward-rate table with 95% CI bands
for tl in taus:
    y = np.asarray(r["reward_rate"][tl], float)
    ci = r["reward_rate_ci"][tl]
    lo = np.array([c[0] for c in ci]); hi = np.array([c[1] for c in ci])
    axA.plot(delays, y, "-o", ms=4, lw=1.7, color=tcol[tl],
             label=rf"$\tau_{{\rm leak}}={tl:g}$ s")
    axA.fill_between(delays, lo, hi, color=tcol[tl], alpha=0.18, lw=0)
axA.axhline(r.get("crit", 0.75), ls="--", color=GREY, lw=1.0)
axA.axhline(0.5, ls=":", color=RED, lw=1.0)
axA.text(delays[-1], 0.505, "chance", fontsize=7.5, ha="right", va="bottom", color=RED)
axA.set_xscale("log"); axA.set_xticks(delays)
axA.set_xticklabels([f"{d:g}" for d in delays])
axA.set_xlabel(r"action$\to$reward delay $D$ (s)"); axA.set_ylabel("final reward rate")
axA.set_ylim(0.45, 1.03); axA.set_title("(a) delay x retention window", fontsize=10)
axA.legend(fontsize=8, frameon=False, loc="upper right"); _clean(axA)

# (b) device vs no-trace learning curve (running reward rate, 95% CI band)
for arr, col, lab in [(r["curve_device"], GREEN, "device trace"),
                      (r["curve_notrace"], GREY, "no-trace")]:
    m, lo, hi = _running(np.asarray(arr, float))
    x = np.arange(W, W + len(m))
    axB.plot(x, m, color=col, lw=1.8, label=lab, zorder=4)
    axB.fill_between(x, lo, hi, color=col, alpha=0.20, lw=0, zorder=2)
axB.axhline(r.get("crit", 0.75), ls="--", color=GREY, lw=1.0)
axB.axhline(0.5, ls=":", color=RED, lw=1.0, label="chance")
axB.set_xlabel("trial"); axB.set_ylabel("reward rate (running)")
axB.set_ylim(0.3, 1.03); axB.set_title(rf"(b) learning curve ($\tau_{{\rm leak}}=10$ s, $D=D_0={r['D0']:g}$ s)",
                                       fontsize=10)
axB.legend(fontsize=8, frameon=False, loc="center right"); _clean(axB)
plt.show()

dev, nt = r["device_final"], r["notrace_final"]
print(f"C2 device learns : final reward rate {dev[0]:.3f}  95% CI [{dev[2]:.3f}, {dev[3]:.3f}]  "
      f"({'PASS' if dev[0] >= r.get('crit', 0.75) else 'fail'} vs crit {r.get('crit', 0.75)})")
print(f"C1 necessity     : no-trace {nt[0]:.3f}  95% CI [{nt[2]:.3f}, {nt[3]:.3f}]  "
      f"({'PASS' if nt[0] <= 0.5 + 0.10 else 'fail'} vs chance+0.10)")
print("max learnable delay index per retention (C4 monotone in tau):",
      {tl: r['max_learn'][tl] for tl in taus})

# ---- cell 4 --------------------------------------------------------
# Replay-only: the dense 13x18 x 20-seed grid is too heavy for an in-kernel default.
r = paths.load_result("exp8_dmax_law.npy")
taus = np.asarray(r["taus"], float)
dvals = np.array([r["dmax"][t] for t in r["taus"]], float)
censored = np.asarray(r["censored"], bool)
k, r2 = r["k"], r["r2"]
k_plateau, plateau_tau, crit = r["k_plateau"], r["plateau_tau"], r["crit"]

# fit / plateau masks (mirror the exp8 driver): resolved = not censored & D_max>0
fitmask = (~censored) & (dvals > 0)
rise = fitmask & (taus < plateau_tau)
plat = fitmask & (taus >= plateau_tau)

fig, ax = plt.subplots(figsize=(6.4, 4.4))
xx = np.linspace(0, taus.max() * 1.05, 100)
# headline origin fit over all resolved points
ax.plot(xx, k * xx, color=INK, lw=1.3, ls="--", zorder=2,
        label=rf"$D_{{\max}} = {k:.0f}\,\tau_{{\rm leak}}$  ($R^2={r2:.2f}$, all points)")
# asymptotic plateau slope (steeper; the trace-dominated regime)
ax.plot(xx, k_plateau * xx, color=RED, lw=1.1, ls=(0, (5, 2)), zorder=2,
        label=rf"plateau slope $\approx{k_plateau:.0f}$ ($\tau_{{\rm leak}}\geq{plateau_tau:g}$ s)")
ax.scatter(taus[rise], dvals[rise], s=58, color=INDIGO, edgecolor="white", zorder=4,
           label="rise-limited / transition")
ax.scatter(taus[plat], dvals[plat], s=58, marker="s", color=GREEN, edgecolor="white",
           zorder=4, label="trace-dominated plateau")
if censored.any():
    ax.scatter(taus[censored], dvals[censored], s=58, marker="^", color=GOLD,
               edgecolor="white", zorder=4, label=r"censored ($D_{\max}\geq$ grid max)")
ax.set_xlabel(r"retention $\tau_{\rm leak}$ (s)"); ax.set_ylabel(r"max learnable delay $D_{\max}$ (s)")
ax.set_title(rf"Densely-sampled retention-delay scaling (crit = {crit:g})", fontsize=10)
ax.legend(frameon=False, fontsize=8, loc="upper left"); _clean(ax); plt.show()

print(f"origin fit over {int(fitmask.sum())} resolved points: D_max = {k:.2f} * tau_leak, R^2 = {r2:.3f}")
print(f"trace-dominated plateau slope (tau_leak >= {plateau_tau:g} s): D_max/tau_leak -> {k_plateau:.2f}")
if censored.any():
    print("right-censored at grid ceiling (D_max >= max delay): tau_leak =",
          [float(t) for t in taus[censored]])
else:
    print("no right-censored points at this grid ceiling")
print("D_max per retention (s):", {float(t): round(float(r['dmax'][t]), 1) for t in r['taus']})

# ---- cell 5 --------------------------------------------------------
# Uncomment to launch the full sweeps as subprocesses (streamed; the exp8 grid is heavy):
# import subprocess, sys
# subprocess.run([sys.executable, "-m", "mrl_trace.bandit", "--tier3", "--full"])
print("see the markdown above for the full-scale commands")
\end{lstlisting}

\begin{lstlisting}[caption={experiments/02\_closed\_loop\_rl.ipynb},label={lst:siox-nb-02-closed-loop-rl}]
# ---- cell 1 --------------------------------------------------------
import numpy as np
import matplotlib.pyplot as plt
from mrl_trace import paths
from mrl_trace.stats import bootstrap_ci

# QUICK=True re-runs a fast, few-seed version of an experiment in-kernel (serial, no Pool);
# the default (QUICK=False) REPLAYS the committed 20-seed grid so every figure renders
# instantly and deterministically. The heavy sweeps document a `--full` shell command instead.
QUICK = False

GREEN, INDIGO, RED, GOLD, GREY = "#3aa07a", "#2f4b8f", "#c0392b", "#e0a93b", "#9aa6b2"
# sequential viridis by retention (long -> short), matching the manuscript window figures
VIR = [plt.cm.viridis(x) for x in (0.85, 0.5, 0.15)]

def _clean(ax):
    for sp in ("top", "right"):
        ax.spines[sp].set_visible(False)
    ax.set_axisbelow(True); ax.grid(True, color="0.88", lw=0.5)

print("data/results:", paths.results_dir())

# ---- cell 2 --------------------------------------------------------
if QUICK:
    from mrl_trace.bandit import run_learning_and_window
    r = run_learning_and_window(seeds=6)
else:
    r = paths.load_result("tier3_results.npy")

delays = np.asarray(r["delays"], float)
D0 = r["D0"]
W = 50                                   # running-mean window (matches gen_rl_curve.py)

def running(rw_2d, w=W):
    cs = np.cumsum(np.insert(rw_2d, 0, 0.0, axis=1), axis=1)
    rr = (cs[:, w:] - cs[:, :-w]) / w    # (seeds, trials-w)
    return rr.mean(0), np.percentile(rr, 2.5, axis=0), np.percentile(rr, 97.5, axis=0)

fig, (axA, axB) = plt.subplots(1, 2, figsize=(9.6, 3.6))

# (a) learning curve: device vs no-trace, running mean + 95% CI band
for arr, c, lab in [(r["curve_device"], GREEN, "device trace"),
                    (r["curve_notrace"], GREY, "no-trace")]:
    m, lo, hi = running(np.asarray(arr, float))
    x = np.arange(W, W + len(m))
    axA.plot(x, m, color=c, lw=1.7, label=lab, zorder=4)
    axA.fill_between(x, lo, hi, color=c, alpha=0.2, lw=0, zorder=2)
axA.axhline(0.5, ls="--", color=RED, lw=1.0, zorder=1, label="chance")
axA.set_xlabel("trial"); axA.set_ylabel("reward rate")
axA.set_ylim(0.3, 1.03); axA.set_title(rf"(a) learning curve ($D_0={D0:g}$ s)", fontsize=10)
axA.legend(fontsize=8, frameon=False, loc="lower right"); _clean(axA)

# (b) delay x retention window
tau_style = [(10.0, VIR[0]), (2.0, VIR[1]), (0.5, VIR[2])]
for tl, c in tau_style:
    y = np.asarray(r["reward_rate"][tl], float)
    ci = r["reward_rate_ci"][tl]
    lo = np.array([b[0] for b in ci]); hi = np.array([b[1] for b in ci])
    axB.plot(delays, y, marker="o", ms=4, lw=1.6, color=c,
             label=rf"$\tau_{{leak}}={tl:g}$ s", zorder=4)
    axB.fill_between(delays, lo, hi, color=c, alpha=0.18, lw=0, zorder=2)
axB.axhline(0.75, ls=":", color=GREY, lw=1.0); axB.axhline(0.5, ls="--", color=RED, lw=0.9)
axB.set_xscale("log"); axB.set_xticks(delays)
axB.set_xticklabels([f"{d:g}" for d in delays], fontsize=8)
axB.set_xlabel(r"action$\to$reward delay $D$ (s)"); axB.set_ylabel("final reward rate")
axB.set_ylim(0.4, 1.03); axB.set_title("(b) delay x retention window", fontsize=10)
axB.legend(fontsize=8, frameon=False, loc="lower left"); _clean(axB)
plt.show()

df, nf = r["device_final"], r["notrace_final"]
print(f"device final reward rate  {df[0]:.3f}  (95% CI {df[2]:.3f}-{df[3]:.3f})")
print(f"no-trace final            {nf[0]:.3f}  (95% CI {nf[2]:.3f}-{nf[3]:.3f})")
print("max learnable delay per tau_leak (s):", {f"{k:g}": v for k, v in r["max_learn"].items()})
print(f"C1 no-trace at chance: {'PASS' if nf[0] <= 0.5 + 0.10 else 'fail'}   "
      f"C2 device >= criterion: {'PASS' if df[0] >= 0.75 else 'fail'}")

# ---- cell 3 --------------------------------------------------------
if QUICK:
    from mrl_trace.bandit import run_scaling
    r = run_scaling(seeds=6)
else:
    r = paths.load_result("tier4_results.npy")

cells_sa = [(2, 2), (4, 2), (4, 4), (8, 4), (8, 8), (12, 8)]
labels = [f"{S}x{A}" for S, A in cells_sa]

def _ci_err(entries):
    m = np.array([e[0] for e in entries])
    lo = np.array([e[2] for e in entries]); hi = np.array([e[3] for e in entries])
    return m, np.vstack([m - lo, hi - m])

dev, dev_err = _ci_err([r[c]["device"] for c in cells_sa])
nt, nt_err = _ci_err([r[c]["no_trace"] for c in cells_sa])
chance = np.array([r[c]["chance"] for c in cells_sa])
crit = np.array([r[c]["crit"] for c in cells_sa])
passed = np.array([r[c]["passed"] for c in cells_sa])
x = np.arange(len(cells_sa))

fig, ax = plt.subplots(figsize=(6.6, 3.8))
ax.plot(x, chance, "--", color=RED, lw=1.1, label=r"chance ($1/A$)", zorder=2)
ax.plot(x, crit, ":", color="0.5", lw=1.2, label="criterion", zorder=2)
ax.errorbar(x, dev, yerr=dev_err, marker="o", ms=5, lw=1.6, color=GREEN,
            capsize=2.5, label="device trace", zorder=5)
ax.errorbar(x, nt, yerr=nt_err, marker="s", ms=4, lw=1.3, color=GREY,
            capsize=2.5, label="no-trace", zorder=4)
for xi, dd, ok in zip(x, dev, passed):
    if ok:
        ax.scatter([xi], [dd], s=110, facecolors="none", edgecolors=GREEN, lw=1.7, zorder=6)
firstfail = int(np.argmax(~passed)) if np.any(~passed) else len(cells_sa)
if firstfail < len(cells_sa):
    ax.axvspan(firstfail - 0.5, len(cells_sa) - 0.5, color=RED, alpha=0.06, zorder=0)
    ax.text((firstfail + len(cells_sa) - 1) / 2, 0.93, "not converged\nin budget",
            fontsize=7.5, ha="center", va="top", color=RED, style="italic")
ax.set_xticks(x); ax.set_xticklabels(labels)
ax.set_xlabel(r"array size $S\times A$"); ax.set_ylabel("final reward rate")
ax.set_ylim(0, 1.02)
ax.legend(fontsize=8, frameon=False, loc="lower left", handlelength=1.6)
ax.set_title("Scaling: device converges up to ~4 actions, then a sample-efficiency wall")
_clean(ax); plt.show()

for c, lab in zip(cells_sa, labels):
    e = r[c]
    ttc = e["trials_to_crit"]
    print(f"  {lab:6s}: device {e['device'][0]:.3f}  crit {e['crit']:.3f}  "
          f"converged {e['n_converged']}/{e['n_seeds']}  "
          f"trials_to_crit {ttc if ttc is not None else 'n/a'}  "
          f"{'PASS' if e['passed'] else 'fail'}")

# ---- cell 4 --------------------------------------------------------
if QUICK:
    from mrl_trace.bandit import run_remedies
    r = run_remedies(seeds=6)
else:
    r = paths.load_result("tier5_results.npy")

keys = list(r.keys())                     # preserves R0..R4 insertion order
short = ["R0\nbaseline", "R1\n4x budget", "R2\ndirected\nexpl.",
         "R3\nabstract\ntrace", "R4\nbudget+\ndirected"]
vals = np.array([r[k]["final"][0] for k in keys])
lo = np.array([r[k]["final"][2] for k in keys]); hi = np.array([r[k]["final"][3] for k in keys])
err = np.vstack([vals - lo, hi - vals])
psd = np.array([r[k]["passed"] for k in keys])
isabs = np.array(["abstract" in k for k in keys])   # abstract-trace row grey, device rows green
bar_c = [GREY if a else GREEN for a in isabs]
xb = np.arange(len(keys))
crit88 = 0.5 * (1 + 1.0 / 8)

fig, ax = plt.subplots(figsize=(6.6, 3.8))
ax.bar(xb, vals, yerr=err, color=bar_c, capsize=3, width=0.66, edgecolor="white", zorder=3)
ax.axhline(crit88, ls=":", color="0.5", lw=1.2, zorder=2, label="criterion")
ax.axhline(1.0 / 8, ls="--", color=RED, lw=1.1, zorder=2, label="chance")
for xi, v, h, ok in zip(xb, vals, hi, psd):
    ax.text(xi, h + 0.03, "PASS" if ok else "fail", fontsize=7.5, ha="center",
            color=(GREEN if ok else RED), fontweight=("bold" if ok else "normal"))
ax.set_xticks(xb); ax.set_xticklabels(short, fontsize=8)
ax.set_ylabel("final reward rate"); ax.set_ylim(0, 1.12)
ax.legend(fontsize=8, frameon=False, loc="upper left", handlelength=1.6)
ax.set_title("Remedies on the failing 8x8 case"); _clean(ax); plt.show()

for k in keys:
    e = r[k]; print(f"  {'PASS' if e['passed'] else 'fail'}  {e['final'][0]:.3f}  {k}")

# ---- cell 5 --------------------------------------------------------
if QUICK:
    from mrl_trace.maze import run_sequential
    r = run_sequential(seeds=6)
else:
    r = paths.load_result("exp6_sequential.npy")

crit, chance = r["crit"], r["chance"]
taus = np.asarray(r["taus"], float)
grid_delays = np.asarray(r["delays"], float)
grid = np.asarray(r["grid"], float)
grid_ci = np.asarray(r["grid_ci"], float)
dmax = r["dmax"]

fig, (axA, axB) = plt.subplots(1, 2, figsize=(10, 3.8))

# (a) learning curves (curves are already the seed-mean running rate)
cond_style = [("device", GREEN, "device trace"), ("rstdp", INDIGO, r"R-STDP (matched $\tau$)"),
              ("eprop", GOLD, "reward-based e-prop"), ("no_trace", GREY, "no-trace")]
for name, c, lab in cond_style:
    if name not in r["curves"]:
        continue
    cur = np.asarray(r["curves"][name], float)
    axA.plot(np.arange(len(cur)), cur, color=c, lw=1.6, label=lab, zorder=4)
axA.axhline(crit, ls=":", color=GREY, lw=1.0); axA.axhline(chance, ls="--", color=RED, lw=0.9)
axA.set_xlabel("episode (running mean)"); axA.set_ylabel("reward rate")
axA.set_ylim(0.3, 1.03)
axA.set_title(rf"(a) learning curves ($D_0={r['D0']:g}$ s, $\tau={r['TAU0']:g}$ s)", fontsize=10)
axA.legend(fontsize=8, frameon=False, loc="lower right"); _clean(axA)

# per-condition finals + CI (committed grid stores per-seed finals only)
for name, _, _ in cond_style:
    if name in r["finals"]:
        f = np.asarray(r["finals"][name], float); lo, hi = bootstrap_ci(f)
        print(f"  {name:9s} final {f.mean():.3f}  (95% CI {lo:.3f}-{hi:.3f})")
pol = np.asarray(r["policy"], float)
print(f"  device policy-correctness {pol.mean():.3f}  (C3 {'PASS' if pol.mean() >= crit else 'fail'})")

# (b) retention x delay grid + D_max crossings
for i, (tau, c) in enumerate(zip(taus, VIR)):
    y = grid[i]; lo = grid_ci[i, :, 0]; hi = grid_ci[i, :, 1]
    axB.plot(grid_delays, y, marker="o", ms=4, lw=1.6, color=c,
             label=rf"$\tau_{{leak}}={tau:g}$ s ($D_{{max}}={dmax[tau]:g}$ s)", zorder=4)
    axB.fill_between(grid_delays, lo, hi, color=c, alpha=0.18, lw=0, zorder=2)
    dm = dmax[tau]
    if dm:
        axB.scatter([dm], [grid[i, list(r["delays"]).index(dm)]], s=90, facecolors="none",
                    edgecolors=c, lw=1.8, zorder=6)
axB.axhline(crit, ls=":", color=GREY, lw=1.0); axB.axhline(chance, ls="--", color=RED, lw=0.9)
axB.set_xscale("log"); axB.set_xticks(grid_delays)
axB.set_xticklabels([f"{d:g}" for d in grid_delays], fontsize=8)
axB.set_xlabel(r"action$\to$reward delay $D$ (s)"); axB.set_ylabel("final reward rate")
axB.set_ylim(0.4, 1.03)
axB.set_title(r"(b) retention law: $D_{max}$ grows with $\tau_{leak}$", fontsize=10)
axB.legend(fontsize=7.5, frameon=False, loc="lower left"); _clean(axB)
plt.show()

nt_mean = np.asarray(r["finals"]["no_trace"], float).mean()
k1 = bool((grid > nt_mean + 0.1).any())
print(f"D_max(tau_leak): {{ {', '.join(f'{k:g}:{v:g}' for k, v in dmax.items())} }} s")
print(f"C1 no-trace at chance: {'PASS' if nt_mean <= chance + 0.10 else 'fail'}   "
      f"C4 D_max grows with tau: {'PASS' if (dmax[20.0] >= dmax[5.0] >= dmax[2.0] and dmax[20.0] > dmax[2.0]) else 'fail'}   "
      f"K1 device exceeds no-trace: {'PASS' if k1 else 'fail'}")

# ---- cell 6 --------------------------------------------------------
# Uncomment to launch the full sweeps as subprocesses (streamed; heavy -- minutes each):
# import subprocess, sys
# subprocess.run([sys.executable, "-m", "mrl_trace.bandit", "--bandit", "--full"])
# subprocess.run([sys.executable, "-m", "mrl_trace.maze", "--exp6", "--full"])
print("see the markdown above for the full-scale commands")
\end{lstlisting}

\begin{lstlisting}[caption={experiments/03\_deep\_local\_learning.ipynb},label={lst:siox-nb-03-deep-local-learning}]
# ---- cell 1 --------------------------------------------------------
import numpy as np
import matplotlib.pyplot as plt
from mrl_trace import paths
from mrl_trace.stats import bootstrap_ci

# QUICK=True re-runs a fast, few-seed version of an experiment in-kernel (serial, no Pool);
# the default (QUICK=False) REPLAYS the committed 20-seed grid so every figure renders
# instantly. Heavy studies (exp14 at H=512; the torch-dependent hybrid front-end) always
# default to replay and document a `--full` shell command instead.
QUICK = False

GREEN, INDIGO, RED, GOLD, GREY = "#3aa07a", "#2f4b8f", "#c0392b", "#e0a93b", "#9aa6b2"
PURPLE, ORANGE, INK = "#b07cc6", "#e0a93b", "#2b2b2b"  # DFA feedback / homeostasis / ink

def _clean(ax):
    for sp in ("top", "right"):
        ax.spines[sp].set_visible(False)
    ax.set_axisbelow(True); ax.grid(True, color="0.88", lw=0.5)

print("data/results:", paths.results_dir())

# ---- cell 2 --------------------------------------------------------
from matplotlib.patches import Circle, FancyArrowPatch, Rectangle

def _compound_node(ax, cx, cy, s=0.085):
    # compound crosspoint: left half = non-volatile weight device, right half = trace device
    ax.add_patch(Rectangle((cx - s, cy - s), s, 2 * s, fc="#cbd6e0", ec=INK, lw=0.9, zorder=4))
    ax.add_patch(Rectangle((cx, cy - s), s, 2 * s, fc="white", ec=INDIGO, lw=1.0, zorder=4))

fig, ax = plt.subplots(figsize=(7.6, 5.0))
ax.set_xlim(0, 10); ax.set_ylim(-0.3, 7.6); ax.axis("off")

# global reward broadcast (top, red dashed)
ax.plot([0.7, 9.4], [7.05, 7.05], color=RED, ls="--", lw=1.6, zorder=2)
ax.text(5.05, 7.32, r"global third factor $(R-b)$ broadcast to every crosspoint",
        color=RED, fontsize=9.5, ha="center")

# layer 1: inputs (rows) x hidden (cols)  ->  W1
in_y = [6.15, 5.5]; h_x = [2.3, 3.1, 3.9]; h_lif_y = 4.05
ax.text(0.15, 5.82, "state\ninputs", color=GREEN, fontsize=9, va="center")
for x in h_x:
    ax.plot([x, x], [6.45, h_lif_y + 0.2], color=INDIGO, lw=1.4, zorder=1)
for y in in_y:
    ax.add_patch(Circle((1.25, y), 0.11, fc="white", ec=GREEN, lw=1.6, zorder=4))
    ax.plot([1.36, 4.25], [y, y], color=GREEN, lw=2.0, zorder=2)
for y in in_y:
    for x in h_x:
        _compound_node(ax, x, y, s=0.085)
ax.text(3.1, 6.62, r"$W_1$ (input$\rightarrow$hidden array)", color=INDIGO, fontsize=8.5, ha="center")

# hidden LIF layer (with homeostasis rings)
for x in h_x:
    ax.add_patch(Circle((x, h_lif_y), 0.32, fc="none", ec=ORANGE, lw=1.3, ls=(0, (2, 1.5)), zorder=4))
    ax.add_patch(Circle((x, h_lif_y), 0.20, fc="#eaf0fb", ec=INDIGO, lw=1.5, zorder=5))
    ax.text(x, h_lif_y, "LIF", color=INDIGO, fontsize=6.5, ha="center", va="center")
ax.text(4.35, h_lif_y, "hidden\nneurons", color=INDIGO, fontsize=8.5, va="center")
ax.annotate("", xy=(1.85, h_lif_y + 0.05), xytext=(1.1, h_lif_y + 0.05),
            arrowprops=dict(arrowstyle="-|>", color=ORANGE, lw=1.5))
ax.text(0.95, h_lif_y + 0.05, "local\nhomeostasis", color=ORANGE, fontsize=8.0, ha="center", va="center")

# layer 2: hidden (rows) x action (cols)  ->  W2
a_x = [6.4, 7.2]; row2_y = [2.45, 2.05, 1.65]; bus_x = 5.55
for x in h_x:
    ax.plot([x, bus_x], [h_lif_y - 0.32, h_lif_y - 0.32], color=INDIGO, lw=1.0, zorder=1)
ax.plot([bus_x, bus_x], [h_lif_y - 0.32, row2_y[-1]], color=INDIGO, lw=1.0, zorder=1)
for y in row2_y:
    ax.plot([bus_x, 7.6], [y, y], color=INDIGO, lw=1.4, zorder=1)
for x in a_x:
    ax.plot([x, x], [2.7, 0.55], color=INDIGO, lw=1.4, zorder=1)
    for y in row2_y:
        _compound_node(ax, x, y, s=0.080)
ax.text(6.8, 2.78, r"$W_2$ (hidden$\rightarrow$action array)", color=INDIGO, fontsize=8.5, ha="center")
ax.add_patch(Rectangle((8.15, 6.05), 0.14, 0.28, fc="#cbd6e0", ec=INK, lw=0.8, zorder=5))
ax.add_patch(Rectangle((8.29, 6.05), 0.14, 0.28, fc="white", ec=INDIGO, lw=0.9, zorder=5))
ax.text(8.55, 6.19, "compound cell:\nweight | trace device", color=INK, fontsize=6.6, va="center")

# action LIF + WTA
for x in a_x:
    ax.add_patch(Circle((x, 0.35), 0.20, fc="#eaf0fb", ec=INDIGO, lw=1.5, zorder=5))
    ax.text(x, 0.35, "LIF", color=INDIGO, fontsize=6.5, ha="center", va="center")
ax.text(7.65, 0.35, "action\nneurons", color=INDIGO, fontsize=8.5, va="center")

# DFA fixed-random feedback (purple) -- no W2^T transport
ax.add_patch(FancyArrowPatch((6.55, 0.15), (3.1, 3.55), connectionstyle="arc3,rad=0.30",
             arrowstyle="-|>", mutation_scale=13, color=PURPLE, lw=1.9, ls=(0, (5, 2)), zorder=2))
ax.text(3.05, 1.15, r"fixed random feedback $B$:  $L_h = B\,L_a$" "\n"
        r"per-neuron credit, no $W_2^{\top}$ transport",
        color=PURPLE, fontsize=8.0, ha="center", va="center")
plt.show()

# ---- cell 3 --------------------------------------------------------
if QUICK:
    from mrl_trace.deep import run_deep_local
    r = run_deep_local(seeds=4, trials=1200)      # fast few-seed recompute; same grid schema
else:
    r = paths.load_result("exp7_deep_local.npy")

ORDER = ["shallow", "elm", "global", "dfa", "no_trace", "dfa_homeo"]
LAB = {"shallow": "shallow\n(1 layer)", "elm": "ELM\n(fixed hidden)",
       "global": "global\nscalar", "dfa": "DFA", "no_trace": "no-trace",
       "dfa_homeo": "DFA +\nhomeostasis"}
COL = {"shallow": GREY, "elm": GOLD, "global": RED, "dfa": INDIGO,
       "no_trace": "#c9ced6", "dfa_homeo": GREEN}
finals, ci, curves = r["finals"], r["ci"], r["curves"]
chance, crit = r["chance"], r["crit"]

fig, (axA, axB) = plt.subplots(1, 2, figsize=(10.2, 3.9),
                               gridspec_kw={"width_ratios": [1.4, 1.0]})
# (a) learning curves (already stored as seed-mean running rate per trial)
for k in ("dfa_homeo", "dfa", "global", "no_trace", "elm", "shallow"):
    c = np.asarray(curves[k], float)
    w = max(1, len(c) // 60)
    rc = np.convolve(c, np.ones(w) / w, mode="valid")
    axA.plot(np.arange(len(rc)), rc, color=COL[k], lw=1.7, label=LAB[k].replace("\n", " "))
axA.axhline(crit, ls=":", color=GREY, lw=1.0); axA.axhline(chance, ls="--", color=RED, lw=1.0)
axA.set_xlabel("trial"); axA.set_ylabel("reward rate (running)"); axA.set_ylim(0.28, 1.04)
axA.set_title("(a) deep-XOR learning curves", fontsize=10, loc="left")
axA.legend(fontsize=6.6, ncol=2, frameon=False, loc="lower right"); _clean(axA)
# (b) final reward rate by condition with bootstrap CI
x = np.arange(len(ORDER))
means = [np.asarray(finals[k]).mean() for k in ORDER]
err = [[m - ci[k][0] for m, k in zip(means, ORDER)], [ci[k][1] - m for m, k in zip(means, ORDER)]]
axB.bar(x, means, 0.66, color=[COL[k] for k in ORDER], yerr=err, capsize=3, error_kw=dict(lw=0.9))
axB.axhline(crit, ls=":", color=GREY, lw=1.0); axB.axhline(chance, ls="--", color=RED, lw=1.0)
axB.set_xticks(x); axB.set_xticklabels([LAB[k] for k in ORDER], fontsize=6.8)
axB.set_ylabel("final reward rate"); axB.set_ylim(0, 1.1)
axB.set_title("(b) final rate", fontsize=10, loc="left"); _clean(axB)
plt.show()

print("finals:", {k: round(float(np.asarray(finals[k]).mean()), 3) for k in ORDER})
print("pre-registered criteria:", {k: ("PASS" if v else "fail") for k, v in r["criteria"].items()})
print("  (C3 fail = plain DFA 0.73 < 0.75, bimodal; C6 PASS = homeostasis fix reaches ceiling)")

# ---- cell 4 --------------------------------------------------------
if QUICK:
    from mrl_trace.deep import run_deep_dms
    r = run_deep_dms(seeds=4, trials=1200)        # fast few-seed recompute; same grid schema
else:
    r = paths.load_result("exp13_deep_dms.npy")

LABELS = {"dfa_homeo_dist": "DFA + homeostasis (distractor)",
          "dfa_homeo_nodist": "DFA + homeostasis (no distractor)",
          "dfa_dist": "DFA, no homeostasis (distractor)",
          "no_trace_dist": "no-trace control (distractor)"}
COL = {"dfa_homeo_dist": GREEN, "dfa_homeo_nodist": INDIGO, "dfa_dist": GOLD, "no_trace_dist": GREY}
ORDER = ["dfa_homeo_dist", "dfa_homeo_nodist", "dfa_dist", "no_trace_dist"]
curves, finals, ci = r["curves"], r["finals"], r["ci"]
solved, seeds = r["seeds_solved"], r.get("seeds", 20)
chance, crit = r.get("chance", 0.5), r.get("crit", 0.75)

fig, (axA, axB) = plt.subplots(1, 2, figsize=(9.6, 3.9),
                               gridspec_kw={"width_ratios": [1.5, 1.0]})
# (a) learning curves (seed-mean per trial -> running mean, window 50)
for k in ORDER:
    c = np.asarray(curves[k], float)
    w = max(1, min(50, len(c) // 20))
    rc = np.convolve(c, np.ones(w) / w, mode="valid")
    axA.plot(np.arange(len(rc)), rc, color=COL[k], lw=1.8, label=LABELS[k])
axA.axhline(chance, ls="--", color=RED, lw=1.1); axA.axhline(crit, ls=":", color=GREY, lw=1.0)
axA.set_xlabel("trial"); axA.set_ylabel("reward rate"); axA.set_ylim(0.3, 1.04)
axA.set_title("(a) deep DMS + distractor", fontsize=10, loc="left")
axA.legend(fontsize=7.0, ncol=1, frameon=False, loc="lower right"); _clean(axA)
# (b) final reward rate bars with CI + seeds-solved
x = np.arange(len(ORDER))
means = [np.asarray(finals[k]).mean() for k in ORDER]
los = [means[i] - ci[k][0] for i, k in enumerate(ORDER)]
his = [ci[k][1] - means[i] for i, k in enumerate(ORDER)]
axB.bar(x, means, 0.62, color=[COL[k] for k in ORDER], yerr=[los, his], capsize=3, error_kw=dict(lw=0.9))
for i, k in enumerate(ORDER):
    axB.text(i, min(means[i] + his[i] + 0.04, 1.05), f"{solved[k]}/{seeds}",
             ha="center", va="bottom", fontsize=8)
axB.axhline(chance, ls="--", color=RED, lw=1.1); axB.axhline(crit, ls=":", color=GREY, lw=1.0)
axB.set_xticks(x)
axB.set_xticklabels(["homeo\n+dist", "homeo\nno-dist", "no-homeo\n+dist", "no-trace"], fontsize=7.5)
axB.set_ylabel("final reward rate"); axB.set_ylim(0, 1.12)
axB.set_title("(b) final rate", fontsize=10, loc="left"); _clean(axB)
plt.show()

print("finals:", {k: round(float(np.asarray(finals[k]).mean()), 3) for k in ORDER})
print("pre-registered criteria:", {k: ("PASS" if v else "fail") for k, v in r["criteria"].items()})

# ---- cell 5 --------------------------------------------------------
r = paths.load_result("exp14_array_scale.npy")            # full PF grid: H={32,512}, p={0,0.5}
rs = paths.load_result("exp14_array_scale_sweep.npy")     # finer sweep: H={8,32,128,512}, p={0,..,0.5}
H, P, grid, ctrl = np.asarray(r["H"]), np.asarray(r["p"]), np.asarray(r["grid"]), np.asarray(r["ctrl"])
Hs, Ps, grids, ctrls = (np.asarray(rs["H"]), np.asarray(rs["p"]),
                        np.asarray(rs["grid"]), np.asarray(rs["ctrl"]))

fig, (axA, axB) = plt.subplots(1, 2, figsize=(10.4, 3.9))
# (a) the coarse full-PF grid: full stack (teal shades by fault level) vs no-trace, per H
pcol = {0.0: GREEN, 0.5: "#7fc4a8"}
xw = np.arange(len(H)); w = 0.8 / (len(P) + 1)
for j, p in enumerate(P):
    m = grid[:, j, 0]; lo = m - grid[:, j, 1]; hi = grid[:, j, 2] - m
    axA.bar(xw + (j - len(P) / 2) * w, m, w, color=pcol.get(float(p), GREEN),
            yerr=[lo, hi], capsize=2.5, error_kw=dict(lw=0.9), label=f"full stack, {p:.0%} stuck")
ntc = ctrl.mean(1)   # no-trace collapses over p (sits at chance regardless)
axA.bar(xw + (len(P) - len(P) / 2) * w, ntc, w, color=GREY, label="no-trace control")
axA.axhline(0.5, ls="--", color=RED, lw=1.1, label="chance $1/A$")
axA.axhline(0.75, ls=":", color=GREY, lw=1.0)
axA.set_xticks(xw); axA.set_xticklabels([f"$H={h}$" for h in H])
axA.set_xlabel("hidden width $H$"); axA.set_ylabel("final reward rate"); axA.set_ylim(0, 1.05)
axA.set_title("(a) full PF fault stack", fontsize=10, loc="left")
axA.legend(fontsize=6.6, ncol=2, frameon=False, loc="lower center", bbox_to_anchor=(0.5, 1.06)); _clean(axA)
# (b) the finer sweep as a heat-map of the full-stack mean over (H, p)
im = axB.imshow(grids[:, :, 0], cmap="viridis", vmin=0.5, vmax=1.0, aspect="auto", origin="lower")
axB.set_xticks(range(len(Ps))); axB.set_xticklabels([f"{p:.0%}" for p in Ps])
axB.set_yticks(range(len(Hs))); axB.set_yticklabels([f"{h}" for h in Hs])
axB.set_xlabel("stuck-off fraction $p$"); axB.set_ylabel("hidden width $H$")
axB.set_title("(b) full-stack mean rate sweep", fontsize=10, loc="left")
for i in range(len(Hs)):
    for j in range(len(Ps)):
        axB.text(j, i, f"{grids[i, j, 0]:.2f}", ha="center", va="center",
                 fontsize=7.5, color="white" if grids[i, j, 0] < 0.8 else "0.1")
cb = fig.colorbar(im, ax=axB, fraction=0.046, pad=0.04); cb.set_label("final reward rate", fontsize=8)
plt.show()

print("full PF grid faults:", r["faults"], f"(seeds={r['seeds']})")
print("sweep faults:", rs["faults"], f"(seeds={rs['seeds']})")
print("no-trace control range (sweep):", round(float(ctrls.min()), 3), "-", round(float(ctrls.max()), 3),
      "-> stays at chance; full stack learns at every H, p with disjoint CIs")

# ---- cell 6 --------------------------------------------------------
r = paths.load_result("tier6_results.npy")
res = r["results"]; chance, crit = r["chance"], r["crit"]
order = ["hybrid", "raw", "no_trace"]
labels = ["hybrid\n(front-end)", "raw\npixels", "no-\ntrace"]
COL = {"hybrid": GREEN, "raw": GREY, "no_trace": "#c9ced6"}
# result tuples are (mean, sd, lo, hi); error bars are the bootstrap 95% CI
vals = np.array([res[k][0] for k in order])
lo = np.array([res[k][2] for k in order]); hi = np.array([res[k][3] for k in order])
err = np.vstack([vals - lo, hi - vals])

fig, ax = plt.subplots(figsize=(5.6, 3.9))
x = np.arange(len(order))
ax.bar(x, vals, 0.6, color=[COL[k] for k in order], yerr=err, capsize=3,
       edgecolor="white", error_kw=dict(lw=0.9), zorder=3)
ax.axhline(crit, ls=":", color=GREY, lw=1.2, label="criterion")
ax.axhline(chance, ls="--", color=RED, lw=1.1, label="chance $1/A$")
ax.text(0, hi[0] + 0.03, "PASS", fontsize=9, ha="center", color=GREEN, fontweight="bold")
ax.set_xticks(x); ax.set_xticklabels(labels, fontsize=9)
ax.set_ylabel("final reward rate"); ax.set_ylim(0, 0.82)
ax.set_title(f"Hybrid stack: front-end acc {r['front_acc']:.2f}  (C={r['C']}, A={r['A']}, P={r['P']})",
             fontsize=9.5)
ax.legend(fontsize=8, frameon=False, loc="upper right"); _clean(ax)
plt.show()

print("finals (mean):", {k: round(float(res[k][0]), 3) for k in order})
print("front-end accuracy:", r["front_acc"], "| seeds:", r["n_seeds"])
print("criteria:", {"C1": r["C1"], "C2": r["C2"], "C3": r["C3"]},
      "-> hybrid reaches criterion; raw-pixel + no-trace stay at chance")

# ---- cell 7 --------------------------------------------------------
# Uncomment to launch a full sweep as a subprocess (streamed; heavy -- minutes each):
# import subprocess, sys
# subprocess.run([sys.executable, "-m", "mrl_trace.deep", "--exp7", "--full"])
print("see the markdown above for the full-scale commands")
\end{lstlisting}

\begin{lstlisting}[caption={experiments/04\_temporal\_selectivity.ipynb},label={lst:siox-nb-04-temporal-selectivity}]
# ---- cell 1 --------------------------------------------------------
import numpy as np
import matplotlib.pyplot as plt
from mrl_trace import paths
from mrl_trace.stats import bootstrap_ci

# QUICK=True re-runs a fast, few-seed version of an experiment in-kernel (serial, no Pool);
# the default (QUICK=False) REPLAYS the committed 20-seed grid so every figure renders
# instantly and deterministically. The full sweeps are documented as `--full` shell
# commands in the last cell.
QUICK = False

GREEN, INDIGO, RED, GOLD, GREY = "#3aa07a", "#2f4b8f", "#c0392b", "#e0a93b", "#9aa6b2"
INK = "#2b2b2b"

def _clean(ax):
    for sp in ("top", "right"):
        ax.spines[sp].set_visible(False)
    ax.set_axisbelow(True); ax.grid(True, color="0.88", lw=0.5)

def _running(x, win=120):
    x = np.asarray(x, float)
    if len(x) < win:
        return x
    c = np.cumsum(np.insert(x, 0, 0.0))
    return (c[win:] - c[:-win]) / win

print("data/results:", paths.results_dir())

# ---- cell 2 --------------------------------------------------------
if QUICK:
    from mrl_trace.deep import run_dms_all
    d = run_dms_all(seeds=4, trials=800)
else:
    d = paths.load_result("exp12_dms.npy")

order = ["device", "abstract", "no_trace"]
disp = {"device": "device (band-pass)", "abstract": "exponential (recency)", "no_trace": "no-trace"}
cols = {"device": GREEN, "abstract": GOLD, "no_trace": GREY}
chance, crit = d["chance"], d["crit"]

fig, (axA, axB) = plt.subplots(1, 2, figsize=(8.4, 3.4),
                               gridspec_kw={"width_ratios": [1.3, 1.0]})

# (a) learning curves with 95% CI band over seeds (from the per-seed raw traces)
raw = d.get("raw")
for k in order:
    if raw is not None:
        seeds = np.apply_along_axis(_running, 1, np.asarray(raw[k]))  # (B, T')
        m = seeds.mean(0); se = seeds.std(0) / np.sqrt(seeds.shape[0])
        t = np.arange(len(m))
        axA.fill_between(t, m - 1.96 * se, m + 1.96 * se, color=cols[k], alpha=0.18, zorder=2)
        axA.plot(t, m, color=cols[k], lw=1.9, label=disp[k], zorder=3)
    else:
        axA.plot(_running(d["curves"][k]), color=cols[k], lw=1.9, label=disp[k])
axA.axhline(chance, color=RED, ls="--", lw=1.1, zorder=1)
axA.axhline(crit, color="0.5", ls=":", lw=1.0, zorder=1)
axA.set_xlabel("trial", fontsize=9.5); axA.set_ylabel("reward rate", fontsize=9.5)
axA.set_ylim(0.35, 1.02); axA.set_title("(a)  Learning curves", fontsize=9.5, loc="left")
axA.legend(fontsize=7.2, loc="lower right", frameon=False); _clean(axA)

# (b) final bars with bootstrap 95% CI
x = np.arange(len(order))
means = [float(np.mean(d["finals"][k])) for k in order]
lo = [means[i] - d["ci"][k][0] for i, k in enumerate(order)]
hi = [d["ci"][k][1] - means[i] for i, k in enumerate(order)]
axB.bar(x, means, color=[cols[k] for k in order], edgecolor=INK, lw=0.8, width=0.62, zorder=3)
axB.errorbar(x, means, yerr=[lo, hi], fmt="none", ecolor=INK, elinewidth=1.2, capsize=3, zorder=4)
axB.axhline(chance, color=RED, ls="--", lw=1.1, zorder=1, label="chance")
axB.axhline(crit, color="0.5", ls=":", lw=1.0, zorder=1, label="criterion")
axB.legend(fontsize=7.0, loc="upper right", frameon=False, handlelength=1.6)
axB.set_xticks(x); axB.set_xticklabels(["device\n(band-pass)", "exp.\n(recency)", "no-trace"], fontsize=7.8)
axB.set_ylim(0, 1.05); axB.set_ylabel("final reward rate", fontsize=9.5)
axB.set_title("(b)  Final rate", fontsize=9.5, loc="left"); _clean(axB)
fig.tight_layout(); plt.show()

print("pre-registered criteria:", {k: ("PASS" if v else "fail") for k, v in d["criteria"].items()})
for k in order:
    print(f"  {k:9s}: final {means[order.index(k)]:.3f}  CI [{d['ci'][k][0]:.3f}, {d['ci'][k][1]:.3f}]")

# ---- cell 3 --------------------------------------------------------
if QUICK:
    from mrl_trace.selectivity import run_interval_selectivity
    d = run_interval_selectivity(seeds=4, trials=400)
else:
    d = paths.load_result("exp10_interval.npy")

dev_S, exp_S = np.asarray(d["dev_S"]), np.asarray(d["exp_S"])
dev_ci, exp_ci = d["dev_ci"], d["exp_ci"]
taus = list(np.asarray(d["taus"], float))
D_grid = np.asarray(d["D_grid"], float)
Scurve = d["Scurve"]

fig, (axA, axB) = plt.subplots(1, 2, figsize=(9.2, 3.5),
                               gridspec_kw={"width_ratios": [0.85, 1.15]})

# (a) design-point selectivity: device band-pass (S>1) vs exponential recency (S<1)
xs = np.arange(2)
mA = [dev_S.mean(), exp_S.mean()]
loA = [mA[0] - dev_ci[0], mA[1] - exp_ci[0]]
hiA = [dev_ci[1] - mA[0], exp_ci[1] - mA[1]]
axA.bar(xs, mA, color=[GREEN, GOLD], edgecolor=INK, lw=0.8, width=0.6, zorder=3)
axA.errorbar(xs, mA, yerr=[loA, hiA], fmt="none", ecolor=INK, elinewidth=1.2, capsize=3, zorder=4)
axA.axhline(1.0, color=RED, ls="--", lw=1.1, zorder=1, label="$S=1$ (no selectivity)")
axA.set_xticks(xs); axA.set_xticklabels(["device\n(band-pass)", "exp.\n(recency)"], fontsize=8)
axA.set_ylabel("selectivity ratio $S = w_{pref}/w_{late}$", fontsize=9)
axA.set_title(f"(a) At design point $D=t^*={d['D0']:.1f}$ s", fontsize=9.5, loc="left")
axA.legend(fontsize=7.2, loc="upper right", frameon=False); _clean(axA)

# (b) S(D) crossover: the preferred interval shifts right with tau_leak
tcol = {5.0: GREEN, 10.0: INDIGO, 20.0: GOLD}
for tl in taus:
    axB.plot(D_grid, np.asarray(Scurve[tl], float), "o-", color=tcol.get(tl, GREY),
             lw=1.8, ms=5, label=rf"$\tau_{{leak}}={tl:g}$ s")
axB.axhline(1.0, color=RED, ls="--", lw=1.0, zorder=1)
axB.axvline(2, color="0.7", ls=":", lw=1.0); axB.axvline(32, color="0.7", ls=":", lw=1.0)
axB.set_xscale("log"); axB.set_xticks([1, 2, 4, 8, 16, 32, 48])
axB.set_xticklabels([1, 2, 4, 8, 16, 32, 48], fontsize=8)
axB.set_xlabel("design interval $D$ (s)"); axB.set_ylabel("selectivity $S$", fontsize=9)
axB.set_title("(b) Crossover: preferred interval shifts with $\\tau_{leak}$", fontsize=9.5, loc="left")
axB.legend(fontsize=7.5, frameon=False, loc="upper right"); _clean(axB)
fig.suptitle("Interval selectivity: a tunable band-pass, not a recency bias", fontsize=10)
fig.tight_layout(); plt.show()

crit = {k: bool(d[k]) for k in ("P1", "P2", "P3", "P4")}
print("pre-registered criteria:", {k: ("PASS" if v else "fail") for k, v in crit.items()})
print(f"  device S = {dev_S.mean():.2f}  CI [{dev_ci[0]:.2f}, {dev_ci[1]:.2f}]  (P1: CI lower > 1)")
print(f"  exp.   S = {exp_S.mean():.2f}  CI [{exp_ci[0]:.2f}, {exp_ci[1]:.2f}]  (P2: CI upper < 1)")
print("  t*(tau_leak):", {k: round(v, 1) for k, v in d["tstar"].items()})

# ---- cell 4 --------------------------------------------------------
if QUICK:
    from mrl_trace.selectivity import run_vector_timer
    d = run_vector_timer(seeds=4, trials=600, quick=True)
else:
    d = paths.load_result("exp20_vector_timer.npy")

fin, ci = d["finals"], d["ci"]
tA, tB = d.get("tA", 10.5), d.get("tB", 30.5)
fig, (axA, axB) = plt.subplots(1, 2, figsize=(9, 3.4))

# (a) aliasing-rescue at the measured scalar-aliasing pair
order = ["vector", "scalar", "no_trace"]
vcols = {"vector": GREEN, "scalar": INDIGO, "no_trace": GREY}
lab = {"vector": "full vector", "scalar": "scalar (last stage)", "no_trace": "no trace"}
m = [float(np.mean(fin[k])) for k in order]
lo = [m[i] - ci[order[i]][0] for i in range(3)]; hi = [ci[order[i]][1] - m[i] for i in range(3)]
axA.bar(np.arange(3), m, color=[vcols[k] for k in order], edgecolor=INK, lw=0.8,
        yerr=np.vstack([lo, hi]), capsize=3, error_kw=dict(lw=1, ecolor="0.3"), zorder=3)
axA.axhline(d.get("chance", 0.5), color=RED, ls="--", lw=1.1, zorder=1, label="chance (0.50)")
axA.axhline(d.get("crit", 0.75), color="0.6", ls=":", lw=1.0, zorder=1, label="criterion (0.75)")
axA.set_xticks(range(3)); axA.set_xticklabels([lab[k] for k in order], fontsize=8)
axA.set_ylabel("reward rate"); axA.set_ylim(0, 1.12)
axA.set_title(f"(a) Aliasing pair $t_A$={tA:.0f}s, $t_B$={tB:.0f}s", fontsize=9.5, loc="left")
axA.legend(fontsize=7, frameon=False, loc="upper right", handlelength=1.6); _clean(axA)

# (b) vector-minus-scalar advantage vs separation x k (noise- not k-bounded)
seps = list(np.asarray(d["seps"], float)); ks = list(np.asarray(d["ks"], int))
kcol = {ks[i]: c for i, c in enumerate([GREEN, INDIGO, GOLD])}
for k in ks:
    adv = [d["sweep"][f"{float(s)}_{k}"]["vector"] - d["sweep"][f"{float(s)}_{k}"]["scalar"]
           for s in seps]
    axB.plot(seps, adv, "o-", color=kcol[k], lw=1.8, ms=5, label=f"$k$={k}")
axB.axhline(0, color="0.5", ls=":", lw=1.0)
axB.set_xlabel("interval separation (s)"); axB.set_ylabel("vector $-$ scalar advantage")
axB.set_title("(b) Advantage is noise-bounded, not $k$-bounded", fontsize=9.5, loc="left")
axB.legend(fontsize=8, frameon=False); _clean(axB)
fig.suptitle("Cascade-vector timer: the vector rescues aliasing the scalar cannot "
             "(simulation; needs stage-level readout)", fontsize=9.5)
fig.tight_layout(); plt.show()

print("pre-registered criteria:", {k: ("PASS" if v else "fail") for k, v in d["criteria"].items()})
for k in order:
    print(f"  {k:9s}: reward rate {m[order.index(k)]:.3f}  CI [{ci[k][0]:.3f}, {ci[k][1]:.3f}]")

# ---- cell 5 --------------------------------------------------------
# Uncomment to launch the full sweeps as subprocesses (streamed; heavy -- minutes each):
# import subprocess, sys
# subprocess.run([sys.executable, "-m", "mrl_trace.selectivity", "--exp10", "--exp20", "--full"])
print("see the markdown above for the full-scale commands")
\end{lstlisting}

\begin{lstlisting}[caption={experiments/05\_biological\_grounding.ipynb},label={lst:siox-nb-05-biological-grounding}]
# ---- cell 1 --------------------------------------------------------
import numpy as np
import matplotlib.pyplot as plt
from mrl_trace import paths
from mrl_trace.stats import bootstrap_ci

# QUICK=True re-runs a fast, few-seed version of an experiment in-kernel (serial, no Pool);
# the default (QUICK=False) REPLAYS the committed 20-seed grid so every figure renders
# instantly. The dopamine/EEG capstones REQUIRE external data absent here, so they are
# always replay-only (a `--full` shell command is documented instead); only the Frank task's
# synthetic-reward device conditions can be recomputed in-kernel.
QUICK = False

GREEN, INDIGO, RED, GOLD, GREY = "#3aa07a", "#2f4b8f", "#c0392b", "#e0a93b", "#9aa6b2"

def _clean(ax):
    for sp in ("top", "right"):
        ax.spines[sp].set_visible(False)
    ax.set_axisbelow(True); ax.grid(True, color="0.88", lw=0.5)

def _running(v, w=50):
    v = np.asarray(v, float)
    c = np.cumsum(np.insert(v, 0, 0.0))
    return (c[w:] - c[:-w]) / w

print("data/results:", paths.results_dir())

# ---- cell 2 --------------------------------------------------------
r = paths.load_result("exp11_dopamine_capstone.npy")   # external DANDI cache absent -> replay only
chance, crit = float(r["chance"]), float(r["crit"])
meta = r["meta"]

fig, (axA, axB) = plt.subplots(1, 2, figsize=(10.4, 3.8),
                               gridspec_kw={"width_ratios": [1.0, 1.4]})

# (a) G-shallow: single-layer final reward rate by reward source, bootstrap 95% CI
sh = r["shallow"]
snames = ["synthetic", "dopamine", "shuffled"]; scols = [GREEN, GOLD, GREY]
sm = [np.asarray(sh["finals"][n]).mean() for n in snames]
sci = [sh["ci"][n] for n in snames]
serr = [[m - lo for m, (lo, hi) in zip(sm, sci)], [hi - m for m, (lo, hi) in zip(sm, sci)]]
axA.bar(range(3), sm, 0.62, color=scols, yerr=serr, capsize=3, edgecolor="white")
axA.axhline(crit, ls=":", color="0.3", lw=1.0); axA.axhline(chance, ls="--", color="0.3", lw=1.0)
axA.set_xticks(range(3)); axA.set_xticklabels(["syn", "DA", "shuf"], fontsize=9)
axA.set_ylabel("final reward rate"); axA.set_ylim(0, 1.05)
axA.set_title("(a) G-shallow | DA bal-acc {:.2f}".format(meta["mean_bal_acc"]),
              loc="left", fontsize=9.5, fontweight="bold"); _clean(axA)

# (b) G-deep: DFA vs DFA+homeostasis x reward source
dp = r["deep"]
order = [("dfa", "synthetic"), ("dfa", "dopamine"), ("dfa", "shuffled"),
         ("dfa_homeo", "synthetic"), ("dfa_homeo", "dopamine"), ("dfa_homeo", "shuffled")]
dcol = {"synthetic": GREEN, "dopamine": GOLD, "shuffled": GREY}
keys = [str(t) for t in order]
dm = [np.asarray(dp["finals"][k]).mean() for k in keys]
dci = [dp["ci"][k] for k in keys]
derr = [[m - lo for m, (lo, hi) in zip(dm, dci)], [hi - m for m, (lo, hi) in zip(dm, dci)]]
x = np.arange(6)
axB.bar(x, dm, 0.7, color=[dcol[rw] for (_, rw) in order], yerr=derr, capsize=3, edgecolor="white")
axB.axhline(crit, ls=":", color="0.3", lw=1.0); axB.axhline(chance, ls="--", color="0.3", lw=1.0)
axB.axvline(2.5, color="0.8", lw=0.8)
axB.set_xticks(x); axB.set_xticklabels(["syn", "DA", "shuf", "syn", "DA", "shuf"], fontsize=8)
axB.set_ylabel("final reward rate"); axB.set_ylim(0, 1.1)
axB.text(1.0, 1.05, "DFA", ha="center", fontsize=8.5)
axB.text(4.0, 1.05, "DFA + homeostasis", ha="center", fontsize=8.5)
axB.set_title("(b) G-deep | homeostasis under measured dopamine",
              loc="left", fontsize=9.5, fontweight="bold"); _clean(axB)
plt.show()

cr = r["criteria"]
print("pre-registered criteria:", {k: ("PASS" if v else "fail") for k, v in cr.items()})
print(f"decoder: n_subj={meta['n_subj']}  bal-acc={meta['mean_bal_acc']:.3f}  AUC={meta['mean_auc']:.3f}"
      f"  P(dec=R|reward)={meta['P_R1_given_reward']:.3f}  P(dec=R|omission)={meta['P_R1_given_omission']:.3f}")
print(f"G4 (descriptive): DA deep DFA+homeo {dp['finals'][str(('dfa_homeo','dopamine'))].mean():.3f}"
      f"  shallow {sh['finals']['dopamine'].mean():.3f}  (cf. EEG deep 0.62 / shallow 0.66)")

# ---- cell 3 --------------------------------------------------------
r = paths.load_result("exp9_capstone.npy")   # external ds003474 EEG absent -> replay only
chance, crit = float(r["chance"]), float(r["crit"])
meta = r["meta"]

rules = ["dfa", "dfa_homeo"]; rewards = ["synthetic", "eeg", "shuffled"]
rcol = {"synthetic": GREEN, "eeg": RED, "shuffled": GREY}
rlab = {"synthetic": "syn", "eeg": "EEG", "shuffled": "shuf"}

fig, (axA, axB) = plt.subplots(1, 2, figsize=(10.6, 3.9),
                               gridspec_kw={"width_ratios": [1.5, 1.2]})

# (a) learning curves, DFA+homeostasis by reward source (the composed deep agent)
for rw in rewards:
    cur = np.asarray(r["curves"][str(("dfa_homeo", rw))], float)
    axA.plot(_running(cur), color=rcol[rw], lw=1.7, label=f"DFA+homeo | {rlab[rw]}")
cur_df = np.asarray(r["curves"][str(("dfa", "eeg"))], float)
axA.plot(_running(cur_df), color=INDIGO, lw=1.4, ls="--", label="DFA (no homeo) | EEG")
axA.axhline(crit, ls=":", color="0.3", lw=1.0); axA.axhline(chance, ls="--", color="0.3", lw=1.0)
axA.set_xlabel("trial"); axA.set_ylabel("reward rate (task)"); axA.set_ylim(0.3, 1.05)
axA.set_title("(a) deep XOR under EEG reward", loc="left", fontsize=9.5, fontweight="bold")
axA.legend(frameon=False, fontsize=7.5, loc="lower right"); _clean(axA)

# (b) final reward rate: 2 rules x 3 reward sources, bootstrap 95% CI
order = [(ru, rw) for ru in rules for rw in rewards]
keys = [str(t) for t in order]
m = [np.asarray(r["finals"][k]).mean() for k in keys]
ci = [r["ci"][k] for k in keys]
err = [[mm - lo for mm, (lo, hi) in zip(m, ci)], [hi - mm for mm, (lo, hi) in zip(m, ci)]]
x = np.arange(6)
axB.bar(x, m, 0.7, color=[rcol[rw] for (_, rw) in order], yerr=err, capsize=3, edgecolor="white")
axB.axhline(crit, ls=":", color="0.3", lw=1.0); axB.axhline(chance, ls="--", color="0.3", lw=1.0)
axB.axvline(2.5, color="0.8", lw=0.8)
axB.set_xticks(x); axB.set_xticklabels([rlab[rw] for (_, rw) in order], fontsize=8)
axB.set_ylabel("final reward rate"); axB.set_ylim(0, 1.1)
axB.text(1.0, 1.05, "DFA", ha="center", fontsize=8.5)
axB.text(4.0, 1.05, "DFA + homeostasis", ha="center", fontsize=8.5)
axB.set_title(f"(b) final rate | EEG bal-acc {meta['mean_bal_acc']:.2f}",
              loc="left", fontsize=9.5, fontweight="bold"); _clean(axB)
plt.show()

cr = r["criteria"]
print("pre-registered criteria:", {k: ("PASS" if v else "fail") for k, v in cr.items()})
print("  E2 = homeostasis-robustness headline; E3 = genuine reward info; E4 = graceful degradation")
print(f"decoder: n_subj={meta['n_subj']}  bal-acc={meta['mean_bal_acc']:.3f}"
      f"  P(R=1|correct)={meta['P_R1_given_correct']:.3f}  P(R=1|incorrect)={meta['P_R1_given_incorrect']:.3f}")
for k in ["('dfa_homeo', 'synthetic')", "('dfa_homeo', 'eeg')", "('dfa_homeo', 'shuffled')",
          "('dfa', 'eeg')"]:
    lo, hi = r["ci"][k]
    print(f"  {k:34s}: {np.asarray(r['finals'][k]).mean():.3f}  CI [{lo:.3f}, {hi:.3f}]")

# ---- cell 4 --------------------------------------------------------
r = paths.load_result("exp7_biosignal.npy")   # external ds003474 EEG absent -> replay only
chance, crit = float(r["chance"]), float(r["crit"])
bacc = r["meta"]["mean_bal_acc"]

fig, (axA, axB) = plt.subplots(1, 2, figsize=(9.6, 3.6),
                               gridspec_kw={"width_ratios": [1.5, 1.0]})

# (a) learning curves by reward source (curves are seed-mean running reward rate)
cur = r["curves"]
tt = np.arange(50, 50 + len(cur["synthetic"]))
axA.plot(tt, cur["synthetic"], color=GREEN, lw=1.8, label="synthetic")
axA.plot(tt, cur["biosignal"], color=RED, lw=1.8, label="real-EEG")
axA.plot(tt, cur["shuffled"], color=GREY, lw=1.6, label="shuffled")
axA.axhline(crit, ls=":", color="0.3", lw=1.0); axA.axhline(chance, ls="--", color="0.3", lw=1.0)
axA.set_xlabel("trial"); axA.set_ylabel("reward rate (task)"); axA.set_ylim(0.3, 1.05)
axA.set_title("(a) single-layer, EEG reward", loc="left", fontsize=9.5, fontweight="bold")
axA.legend(frameon=False, fontsize=8, loc="lower right"); _clean(axA)

# (b) final reward bars with bootstrap CI
names = ["synthetic", "biosignal", "shuffled"]; cols = [GREEN, RED, GREY]
means = [np.asarray(r["finals"][n]).mean() for n in names]
cis = [bootstrap_ci(np.asarray(r["finals"][n])) for n in names]
err = [[m - lo for m, (lo, hi) in zip(means, cis)], [hi - m for m, (lo, hi) in zip(means, cis)]]
axB.bar(range(3), means, 0.62, yerr=err, color=cols, capsize=3, edgecolor="white")
axB.axhline(crit, ls=":", color="0.3", lw=1.0); axB.axhline(chance, ls="--", color="0.3", lw=1.0)
axB.set_xticks(range(3)); axB.set_xticklabels(["syn", "EEG", "shuf"], fontsize=9)
axB.set_ylabel("final reward rate"); axB.set_ylim(0, 1.05)
axB.set_title(f"(b) final rate | bal-acc {bacc:.2f}", loc="left", fontsize=9.5, fontweight="bold")
_clean(axB)
plt.show()

for n in names:
    lo, hi = bootstrap_ci(np.asarray(r["finals"][n]))
    print(f"  {n:10s}: {np.asarray(r['finals'][n]).mean():.3f}  CI [{lo:.3f}, {hi:.3f}]")
print("real-EEG beats shuffled with disjoint CIs -> genuine outcome information at the single layer.")

# ---- cell 5 --------------------------------------------------------
r = paths.load_result("exp10_probselect.npy")   # EEG-reward conditions need absent ds003474 pools
chance, crit = float(r["chance"]), float(r["crit"])

if QUICK:
    from mrl_trace.probselect import run_probselect
    for c in ("shallow", "dfa", "dfa_homeo"):   # synthetic reward -> no external data needed
        out = run_probselect(c, trials=1500, seeds=4)
        r["finals"][c] = out["finals"]
        r["ci"][c] = bootstrap_ci(np.asarray(out["finals"]))
        r["tests"][c] = {"chooseA": out["chooseA"], "avoidB": out["avoidB"]}

fig, (axA, axB, axC) = plt.subplots(1, 3, figsize=(11.4, 3.7),
                                    gridspec_kw={"width_ratios": [1.2, 1.3, 1.1]})

# (a) train-phase learning curves (running P(better stimulus))
ccol = {"shallow": INDIGO, "dfa": GOLD, "dfa_homeo": GREEN, "no_trace": GREY}
clab = {"shallow": "shallow", "dfa": "DFA", "dfa_homeo": "DFA+homeo", "no_trace": "no-trace"}
for c in ["shallow", "dfa", "dfa_homeo", "no_trace"]:
    cur = np.asarray(r["curves"][c], float)
    axA.plot(_running(cur, 100), color=ccol[c], lw=1.7, label=clab[c])
axA.axhline(crit, ls=":", color="0.3", lw=1.0); axA.axhline(chance, ls="--", color="0.3", lw=1.0)
axA.set_xlabel("trial"); axA.set_ylabel("P(better stimulus)"); axA.set_ylim(0.3, 1.02)
axA.set_title("(a) Frank PST -- train", loc="left", fontsize=9.5, fontweight="bold")
axA.legend(frameon=False, fontsize=7.5, loc="lower right"); _clean(axA)

# (b) final train accuracy by condition (bootstrap 95% CI)
order = ["shallow", "dfa", "dfa_homeo", "no_trace", "dfa_homeo_eeg", "dfa_homeo_shuf"]
blab = {"shallow": "shallow", "dfa": "DFA", "dfa_homeo": "DFA\n+homeo", "no_trace": "no-\ntrace",
        "dfa_homeo_eeg": "EEG\nreward", "dfa_homeo_shuf": "EEG\nshuf"}
bcol = {"shallow": INDIGO, "dfa": GOLD, "dfa_homeo": GREEN, "no_trace": GREY,
        "dfa_homeo_eeg": RED, "dfa_homeo_shuf": GREY}
m = [np.asarray(r["finals"][c]).mean() for c in order]
ci = [r["ci"][c] for c in order]
err = [[mm - lo for mm, (lo, hi) in zip(m, ci)], [hi - mm for mm, (lo, hi) in zip(m, ci)]]
x = np.arange(len(order))
axB.bar(x, m, 0.66, color=[bcol[c] for c in order], yerr=err, capsize=3, edgecolor="white")
axB.axhline(crit, ls=":", color="0.3", lw=1.0); axB.axhline(chance, ls="--", color="0.3", lw=1.0)
axB.set_xticks(x); axB.set_xticklabels([blab[c] for c in order], fontsize=7.5)
axB.set_ylabel("final train accuracy"); axB.set_ylim(0, 1.05)
axB.set_title("(b) final accuracy by condition", loc="left", fontsize=9.5, fontweight="bold")
_clean(axB)

# (c) Frank test signature: choose-A vs avoid-B (device conditions only carry the test)
tconds = ["shallow", "dfa", "dfa_homeo"]
tcol = {"shallow": INDIGO, "dfa": GOLD, "dfa_homeo": GREEN}
xw = np.arange(len(tconds)); w = 0.38
cA = [np.asarray(r["tests"][c]["chooseA"]).mean() for c in tconds]
aB = [np.asarray(r["tests"][c]["avoidB"]).mean() for c in tconds]
cA_e = [np.abs(np.subtract(*bootstrap_ci(np.asarray(r["tests"][c]["chooseA"])))) / 2 for c in tconds]
aB_e = [np.abs(np.subtract(*bootstrap_ci(np.asarray(r["tests"][c]["avoidB"])))) / 2 for c in tconds]
axC.bar(xw - w / 2, cA, w, color=GREEN, yerr=cA_e, capsize=3, label="choose-A (positive)")
axC.bar(xw + w / 2, aB, w, color=INDIGO, yerr=aB_e, capsize=3, label="avoid-B (negative)")
axC.axhline(chance, ls="--", color="0.3", lw=1.0)
axC.set_xticks(xw); axC.set_xticklabels(["shallow", "DFA", "DFA\n+homeo"], fontsize=7.5)
axC.set_ylabel("test-phase accuracy"); axC.set_ylim(0, 1.1)
axC.set_title("(c) Frank asymmetry (F4)", loc="left", fontsize=9.5, fontweight="bold")
axC.legend(frameon=False, fontsize=7.5, loc="lower left"); _clean(axC)
plt.show()

for c in order:
    lo, hi = r["ci"][c]
    print(f"  {c:16s}: {np.asarray(r['finals'][c]).mean():.3f}  CI [{lo:.3f}, {hi:.3f}]")
print("Frank test signature (choose-A / avoid-B):",
      {c: (round(float(np.mean(r['tests'][c]['chooseA'])), 2),
           round(float(np.mean(r['tests'][c]['avoidB'])), 2)) for c in tconds})
print("F1/F3/F5 PASS; F2 = pre-registered K2 (12-bit code is linearly separable -> shallow solves it).")

# ---- cell 6 --------------------------------------------------------
# Uncomment to launch the (external-data-free) Frank sweep as a subprocess (heavy -- minutes):
# import subprocess, sys
# subprocess.run([sys.executable, "-m", "mrl_trace.probselect", "--quick"])
print("external DANDI/EEG caches are absent here; see the markdown above for the full-scale commands")
\end{lstlisting}

\begin{lstlisting}[caption={experiments/06\_extensions.ipynb},label={lst:siox-nb-06-extensions}]
# ---- cell 1 --------------------------------------------------------
import numpy as np
import matplotlib.pyplot as plt
from mrl_trace import paths
from mrl_trace.stats import bootstrap_ci

# QUICK=True re-runs a fast, few-seed version of an experiment in-kernel (serial, no Pool);
# the default (QUICK=False) REPLAYS the committed 20-seed grid so every figure renders
# instantly. Heavy sweeps (exp19) always document a `--full` shell command instead.
QUICK = False

GREEN, INDIGO, RED, GOLD, GREY = "#3aa07a", "#2f4b8f", "#c0392b", "#e0a93b", "#9aa6b2"

def _clean(ax):
    for sp in ("top", "right"):
        ax.spines[sp].set_visible(False)
    ax.set_axisbelow(True); ax.grid(True, color="0.88", lw=0.5)

print("data/results:", paths.results_dir())

# ---- cell 2 --------------------------------------------------------
if QUICK:
    from mrl_trace.bandit import run_reversal, _summarize_reversal
    conds = ("device", "abstract", "no_trace")
    raw = {c: run_reversal(c, B=6, trials=4000, n_phases=2, tau_leak=10.0) for c in conds}
    r = _summarize_reversal(raw, conds, 4000, 0.75, 0.5)
else:
    r = paths.load_result("exp18_reversal.npy")

conds = ("device", "abstract", "no_trace")
cmap = {"device": GREEN, "abstract": INDIGO, "no_trace": GREY}
flips, crit, chance = r["flips"], r["crit"], r["chance"]
fig, ax = plt.subplots(figsize=(7.2, 3.6))
for c in conds:
    cur = np.asarray(r["curves"][c], float)
    win = max(1, len(cur) // 120)
    sm = np.convolve(cur, np.ones(win) / win, mode="valid")
    ax.plot(np.arange(len(sm)), sm, color=cmap[c], lw=1.6, label=c)
for f in flips:
    ax.axvline(f, ls="--", color=RED, lw=1.0)
ax.axhline(crit, ls=":", color=GREY, lw=1.0); ax.axhline(chance, ls=":", color=GREY, lw=0.8)
ax.set_xlabel("trial"); ax.set_ylabel("reward rate (running)"); ax.set_ylim(0, 1.03)
ax.set_title("Reversal: device unlearns then re-acquires at the flip"); ax.legend(frameon=False, fontsize=8)
_clean(ax); plt.show()

cr = r["criteria"]
print("pre-registered criteria:", {k: ("PASS" if v else "fail") for k, v in cr.items()})
for c in conds:
    pre, post = np.asarray(r["pre"][c]).mean(), np.asarray(r["post"][c]).mean()
    print(f"  {c:9s}: pre {pre:.3f}  post {post:.3f}")
rc = r["reacq_cost"]
print(f"re-acquisition cost: device {rc['device']} vs abstract {rc['abstract']} trials "
      f"({rc['ratio']:.1f}x -- retention<->flexibility trade-off)"
      if rc.get("ratio") == rc.get("ratio") else "re-acquisition cost: n/a")

# ---- cell 3 --------------------------------------------------------
r = paths.load_result("exp19_multitimescale.npy")
labels = [l for l in ["hetero_raw", "hetero_homeo", "hetero_oracle", "best_single", "no_trace"]
          if l in r["finals"]]
lab_col = {"hetero_raw": RED, "hetero_homeo": GREEN, "hetero_oracle": INDIGO,
           "best_single": GOLD, "no_trace": GREY}
means = [np.asarray(r["finals"][l]).mean() for l in labels]
cis = [r["ci"][l] if "ci" in r else bootstrap_ci(np.asarray(r["finals"][l])) for l in labels]
err = [[m - lo for m, (lo, hi) in zip(means, cis)], [hi - m for m, (lo, hi) in zip(means, cis)]]
fig, ax = plt.subplots(figsize=(6.6, 3.6))
x = np.arange(len(labels))
ax.bar(x, means, color=[lab_col[l] for l in labels], width=0.62)
ax.errorbar(x, means, yerr=err, fmt="none", ecolor="0.3", capsize=3, lw=1.0)
ax.axhline(r.get("crit", 0.75), ls="--", color=GREY, lw=1.0)
ax.axhline(r.get("chance", 0.5), ls=":", color=GREY, lw=0.8)
ax.set_xticks(x); ax.set_xticklabels(labels, rotation=25, ha="right", fontsize=8)
ax.set_ylabel("overall reward rate"); ax.set_ylim(0, 1.03)
ax.set_title("Multi-timescale: homeostasis recovers the naive measured spread"); _clean(ax); plt.show()
if "criteria" in r:
    print("criteria:", {k: ("PASS" if v else "fail") for k, v in r["criteria"].items()})

# ---- cell 4 --------------------------------------------------------
r = paths.load_result("exp22_wm_stc.npy")
taus = np.asarray(r["tau_band"], float)
def _mci(vecs):
    m = np.array([np.asarray(vecs[t]).mean() for t in taus])
    lo = np.array([np.percentile(np.asarray(vecs[t]), 2.5) for t in taus])
    hi = np.array([np.percentile(np.asarray(vecs[t]), 97.5) for t in taus])
    return m, lo, hi
fig, (axA, axB) = plt.subplots(1, 2, figsize=(10, 3.8))
for key, col, lab in [("wm", INDIGO, "cue hold (WM)"), ("stc", GREEN, "credit tag")]:
    m, lo, hi = _mci(r[key]); axA.plot(taus, m, "-o", color=col, lw=1.8, ms=4, label=lab)
    axA.fill_between(taus, lo, hi, color=col, alpha=0.15)
axA.axhline(0.75, ls="--", color=GREY, lw=1.0); axA.set_xscale("log")
axA.set_xticks(taus); axA.set_xticklabels([f"{t:g}" for t in taus])
axA.set_xlabel(r"retention $\tau_{leak}$ (s)"); axA.set_ylabel("task performance")
axA.set_ylim(0.45, 1.02); axA.set_title("(a) per-role requirement", fontsize=10)
axA.legend(fontsize=8, loc="lower right", frameon=False); _clean(axA)
m, lo, hi = _mci(r["one_shared"])
axB.axvspan(1.3, 6.0, color=GREEN, alpha=0.07)
axB.plot(taus, m, "-o", color=GREEN, lw=1.8, ms=4, label="one device, both roles")
axB.fill_between(taus, lo, hi, color=GREEN, alpha=0.15)
axB.axhline(0.75, ls="--", color=GREY, lw=1.0); axB.axhline(0.5, ls=":", color=GREY, lw=1.0)
axB.set_xscale("log"); axB.set_xticks(taus); axB.set_xticklabels([f"{t:g}" for t in taus])
axB.set_xlabel(r"retention $\tau_{leak}$ (s)"); axB.set_ylabel("joint-task performance")
axB.set_ylim(0.45, 1.02); axB.set_title("(b) single substrate, both roles", fontsize=10)
axB.legend(fontsize=8, loc="lower right", frameon=False); _clean(axB)
plt.show()

# ---- cell 5 --------------------------------------------------------
r = paths.load_result("exp23_device_td.npy")
L = np.asarray(r["L_grid"], float); crit = r.get("crit", 0.75)
series = [("reinforce", INDIGO, "policy gradient (REINFORCE)"),
          ("td_actor_critic", GREEN, "device-native TD actor-critic"),
          ("td_no_homeo", GOLD, "TD, no eligibility homeostasis"),
          ("no_trace", GREY, "no-trace (chance)")]
def _mci(final, sc):
    m = np.array([np.asarray(final[(sc, l)]).mean() for l in r["L_grid"]])
    lo = np.array([np.percentile(np.asarray(final[(sc, l)]), 2.5) for l in r["L_grid"]])
    hi = np.array([np.percentile(np.asarray(final[(sc, l)]), 97.5) for l in r["L_grid"]])
    return m, lo, hi
fig, ax = plt.subplots(figsize=(6.6, 4.0))
for sc, col, lab in series:
    if (sc, r["L_grid"][0]) not in r["final"]:
        continue
    m, lo, hi = _mci(r["final"], sc)
    ax.plot(L, m, "-o", color=col, lw=1.7, ms=4, label=lab); ax.fill_between(L, lo, hi, color=col, alpha=0.13)
ax.axhline(crit, ls="--", color=GREY, lw=1.0)
ax.set_xlabel("corridor length $L$"); ax.set_ylabel("final goal-reach rate")
ax.set_xticks(L); ax.set_ylim(0, 1.03)
ax.legend(fontsize=8, loc="upper center", bbox_to_anchor=(0.5, -0.16), ncol=2, frameon=False)
ax.set_title("Device-native TD tracks but does not beat REINFORCE"); _clean(ax); plt.show()

# ---- cell 6 --------------------------------------------------------
r = paths.load_result("exp21_beta_sensitivity.npy")
betas = list(r["betas"])                 # [1.0, 0.85, 0.54]
taus = np.asarray(r["taus_dmax"], float) # [1, 5, 10, 20]
dmax = r["dmax"]                          # {beta: {"dmax":[...per tau...], "k":..., "r2":...}}
bcol = [GREEN, INDIGO, GOLD]
fig, (axA, axB) = plt.subplots(1, 2, figsize=(9.6, 3.6))
# (a) D_max vs tau for each beta: the D_max ~ k*tau retention law, ~unchanged across dispersion
for b, c in zip(betas, bcol):
    dd = np.asarray(dmax[b]["dmax"], float); k = dmax[b]["k"]
    axA.plot(taus, dd, "-o", color=c, lw=1.7, ms=5, label=rf"$\beta={b:g}$  ($k={k:.1f}$)")
axA.set_xlabel(r"retention $\tau_{leak}$ (s)"); axA.set_ylabel(r"$D_{max}$ (s)")
axA.set_title("(a) retention law across dispersion", fontsize=10)
axA.legend(fontsize=8, frameon=False); _clean(axA)
# (b) the exp19 multi-timescale payoff (homeo vs best-single) across beta
exp19 = r["exp19"]; x = np.arange(len(betas)); w = 0.38
homeo = [exp19[b]["homeo"] for b in betas]; best = [exp19[b]["best"] for b in betas]
axB.bar(x - w / 2, homeo, w, color=GREEN, label="hetero + homeostasis")
axB.bar(x + w / 2, best, w, color=GREY, label="best single tau")
axB.set_xticks(x); axB.set_xticklabels([rf"$\beta={b:g}$" for b in betas])
axB.set_ylabel("overall reward rate"); axB.set_ylim(0, 1.03)
axB.set_title("(b) multi-timescale payoff vs dispersion", fontsize=10)
axB.legend(fontsize=8, frameon=False); _clean(axB)
plt.show()
print("D_max slope k across beta:", {b: round(dmax[b]["k"], 2) for b in betas})
print(r.get("verdict", ""))

# ---- cell 7 --------------------------------------------------------
# Uncomment to launch the full sweeps as subprocesses (streamed; heavy -- minutes each):
# import subprocess, sys
# subprocess.run([sys.executable, "-m", "mrl_trace.extensions", "--exp22", "--exp23", "--full"])
print("see the markdown above for the full-scale commands")
\end{lstlisting}

\begin{lstlisting}[caption={experiments/device\_model.ipynb},label={lst:siox-nb-device-model}]
# ---- cell 1 --------------------------------------------------------
import json
import numpy as np
import matplotlib.pyplot as plt
from mrl_trace import device, paths

# QUICK=True re-runs a fast in-kernel version of a study (e.g. device.fit_kww_laws(),
# device.simulate_habituation()); the default (QUICK=False) REPLAYS the committed
# device-model fixtures under data/device_model/ so every figure renders instantly and
# deterministically. Both paths hit the same serial cores in mrl_trace.device.
QUICK = False

GREEN, INDIGO, RED, GOLD, GREY = "#3aa07a", "#2f4b8f", "#c0392b", "#e0a93b", "#9aa6b2"

def _clean(ax):
    for sp in ("top", "right"):
        ax.spines[sp].set_visible(False)
    ax.set_axisbelow(True); ax.grid(True, color="0.88", lw=0.5)

DM = paths.device_model_dir()
print("data/device_model:", DM)

# ---- cell 2 --------------------------------------------------------
VS = list(device.KWW_VOLTAGES)  # [0.8, 0.9, 1.1, 1.2, 1.4, 1.5]

if QUICK:
    res = device.fit_kww_laws()
else:
    res = json.load(open(DM / "kww_final.json"))
LAWS, ROWS = res["laws"], res["rows"]

# the measured gold traces (per-bias averaged) -- the 'replay' data behind the overlay
data = device._kww_load_traces(paths.gold_export_dir())

# free per-trace beta (panel B): the fixed-beta=2 global fit vs each trace's own beta,
# re-fit here from the measured traces (fast, deterministic)
import warnings; warnings.filterwarnings("ignore")
from scipy.optimize import curve_fit
def _free_beta(tg, Ia):
    Im, tmax = Ia.max(), tg[-1]
    def f(t, A, tr, b, td, C):
        return A * (1 - np.exp(-(t / abs(tr)) ** abs(b))) * np.exp(-t / abs(td)) + C
    best = None
    for tr0 in (tmax / 20, tmax / 8, tmax / 4):
        for b0 in (1.0, 2.0):
            try:
                p, _ = curve_fit(f, tg, Ia, p0=[Im, tr0, b0, tmax * 2, Ia[0]], maxfev=40000)
                sse = np.sum((Ia - f(tg, *p)) ** 2)
                if best is None or sse < best[0]:
                    best = (sse, abs(p[2]))
            except Exception:
                pass
    return best[1]
betas = {V: _free_beta(*data[V]) for V in VS}

def kww_global(V, t):
    tr = LAWS["tr0"] * np.exp(-LAWS["cr"] * V)
    td = LAWS["td0"] * np.exp(-LAWS["cd"] * V)
    return (1 - np.exp(-(t / tr) ** LAWS["beta"])) * np.exp(-t / td)

# ============ FIGURE 1: meta-parameter generalisation (tau laws + beta) ============
Vv, Vg = np.array(VS), np.linspace(0.8, 1.5, 100)
fig1, (axA, axB) = plt.subplots(1, 2, figsize=(9.0, 4.0))
tr_pts = np.array([ROWS[f"{V}"]["tau_r"] for V in VS])
td_pts = np.array([ROWS[f"{V}"]["tau_d"] for V in VS])
axA.semilogy(Vv, tr_pts, "o", color=GREEN, ms=6, label=r"$\tau_r$ (rise)")
axA.semilogy(Vg, LAWS["tr0"] * np.exp(-LAWS["cr"] * Vg), "-", color=GREEN, lw=1.8)
axA.semilogy(Vv, td_pts, "s", color=INDIGO, ms=6, label=r"$\tau_d$ (decay)")
axA.semilogy(Vg, LAWS["td0"] * np.exp(-LAWS["cd"] * Vg), "-", color=INDIGO, lw=1.8)
axA.set_xlabel("Voltage (V)"); axA.set_ylabel("Time constant (s)")
axA.set_title("(a) field-acceleration laws", fontsize=10)
axA.legend(fontsize=9, frameon=False); _clean(axA)

axB.plot(Vv, [betas[V] for V in VS], "o", color="#1b1b1b", ms=6, label="per-trace fit")
axB.axhline(2.0, color="#666", ls="--", lw=1.5, label=r"$\beta=2$ (fixed)")
axB.fill_between([0.75, 1.55], 1.86, 2.21, color="0.85", alpha=0.5, zorder=0)
axB.set_xlim(0.75, 1.55); axB.set_ylim(0.8, 3.0)
axB.set_xlabel("Voltage (V)"); axB.set_ylabel(r"Compression exponent $\beta$")
axB.set_title("(b) dispersion is bias-independent", fontsize=10)
axB.text(0.79, 2.04, r"$\approx$ 3 sequential" + "\ntrap stages", fontsize=8.5, color="#444", va="bottom")
axB.legend(fontsize=9, frameon=False, loc="lower right"); _clean(axB)
fig1.tight_layout(); plt.show()

# ============ FIGURE 2: generalised model vs measured data (all biases) ============
fig2, axC = plt.subplots(figsize=(6.6, 4.3))
cmap = plt.cm.viridis(np.linspace(0.1, 0.9, len(VS)))
for c, V in zip(cmap, VS):
    tg, Ia = data[V]
    s = kww_global(V, tg)
    M = np.vstack([s, np.ones_like(s)]).T
    coef, *_ = np.linalg.lstsq(M, Ia, rcond=None)
    f = M @ coef
    mdk = tg <= 200
    axC.plot(tg[mdk], Ia[mdk] * 1e9, ".", color=c, ms=2.5, alpha=0.35)
    axC.plot(tg[mdk], f[mdk] * 1e9, "-", color=c, lw=1.8, label=f"{V:.1f} V")
axC.set_xlim(0, 200); axC.set_xlabel("Time (s)"); axC.set_ylabel("|Current| (nA)")
axC.set_title("Global KWW model (lines) vs measured gold traces (dots)", fontsize=10)
axC.legend(fontsize=8, frameon=False, ncol=2, title="Bias", title_fontsize=8); _clean(axC)
fig2.tight_layout(); plt.show()

print("FINAL KWW global law (beta fixed = {:.0f}):".format(LAWS["beta"]))
print(f"  tau_r(V) = {LAWS['tr0']:.1f} * exp(-{LAWS['cr']:.2f} V)  [s]")
print(f"  tau_d(V) = {LAWS['td0']:.0f} * exp(-{LAWS['cd']:.2f} V)  [s]")
print("  beta<->stages:", res["beta_to_k"], " (beta~2 == ~3 sequential trap stages)")
r2s = [ROWS[f"{V}"]["R2"] for V in VS]
print(f"  per-bias R2: median {np.median(r2s):.3f}  min {min(r2s):.3f}")

# ---- cell 3 --------------------------------------------------------
T, dt = 2.5, 0.05
tg = np.arange(0, T, dt)

def _burst(t, amp, t_on, t_off, period, pw):
    if t < t_on or t >= t_off:
        return 0.0
    return amp if ((t - t_on) % period) < pw else 0.0

def Vfun(t):  # burst 1 = 0.2-0.7 s, burst 2 = 1.2-1.7 s, 60 ms period / 30 ms on
    return _burst(t, 1.0, 0.20, 0.70, 0.06, 0.03) + _burst(t, 1.0, 1.20, 1.70, 0.06, 0.03)

drive = np.array([Vfun(t) for t in tg])

fig = plt.figure(figsize=(8.0, 5.0))
gs = fig.add_gridspec(2, 1, height_ratios=[0.8, 2.6], hspace=0.18)
axV, axI = fig.add_subplot(gs[0]), fig.add_subplot(gs[1])

axV.fill_between(tg, 0, drive, color="0.7", lw=0)
axV.set_ylabel("drive"); axV.set_ylim(0, 1.3); axV.set_yticks([0, 1]); axV.set_xlim(0, T)
axV.text(0.45, 1.08, "burst 1", ha="center", fontsize=9, color="0.3")
axV.text(1.45, 1.08, "burst 2", ha="center", fontsize=9, color="0.3")
for s in ("top", "right"):
    axV.spines[s].set_visible(False)
axV.tick_params(labelbottom=False)

for tl, col, lab in [(0.20, RED, r"$\tau_{leak}=0.2$ s (fast-forgetting)"),
                     (2.00, INDIGO, r"$\tau_{leak}=2$ s (slow)")]:
    # TransientGate: fitted trap-cascade + R_leak relaxation, timescale-rescaled to the
    # RL window via a low-bias-equivalent V that puts tau_r inside [0.1, 2] s. We drive
    # it with the explicit two-burst protocol above rather than the single-coincidence
    # trace() helper.
    g = device.TransientGate(V=0.9, tau_leak=tl, k=device.K_STAGES, dt=dt, vnmax=3.0)
    g.reset()
    e = np.array([g.step(d) for d in drive])
    axI.plot(tg, e / e.max(), color=col, lw=1.8, label=lab)
axI.set_xlabel("Time (s)"); axI.set_ylabel(r"Normalised trace, $e(t)/e_{\max}$")
axI.axhline(0, color="0.85", lw=0.8); axI.set_xlim(0, T); axI.set_ylim(-0.03, 1.12)
axI.legend(fontsize=8.5, frameon=False, loc="upper left", borderaxespad=0.3)
for s in ("top", "right"):
    axI.spines[s].set_visible(False)
axI.annotate("potentiates\nduring burst", xy=(0.62, 0.82), xytext=(0.40, 0.42),
             fontsize=8.5, color="0.3", ha="left", va="center",
             arrowprops=dict(arrowstyle="->", color="0.5", lw=0.8, connectionstyle="arc3,rad=-0.2"))
axI.annotate("relaxes between bursts\n(rate set by $R_{leak}$)", xy=(1.35, 0.53),
             xytext=(1.55, 0.18), fontsize=8.5, color="0.3", ha="left", va="center",
             arrowprops=dict(arrowstyle="->", color="0.5", lw=0.8, connectionstyle="arc3,rad=0.25"))
fig.suptitle(r"Short-term memory: $R_{leak}$ sets how long recent activity is retained",
             fontsize=11, y=0.96)
plt.show()

# ---- cell 4 --------------------------------------------------------
if QUICK:
    h = device.simulate_habituation()
    tg, I, f = h["tg"], h["I"], h["f"]
    print("checks:", {k: round(v, 3) for k, v in h["checks"].items()},
          "-> reproduced" if h["reproduced"] else "-> CHECK params")
else:
    d = np.load(DM / "habit_data.npz")
    tg, I, f = d["tg"], d["I"], d["f"]

tmin = tg / 60.0            # Ch4 Fig 4l is in minutes
I = I / I.max()

fig = plt.figure(figsize=(8.0, 5.0))
gs = fig.add_gridspec(2, 1, height_ratios=[0.8, 2.6], hspace=0.18)
axR, axI = fig.add_subplot(gs[0]), fig.add_subplot(gs[1])

def ticks(t0, t1, rate):
    n = max(2, int((t1 - t0) * rate / 18))
    return np.linspace(t0, t1, n, endpoint=False)
for t0, t1, rate in [(0, 108, 1.0), (108, 168, 10.0), (168, 276, 1.0)]:
    for tt in ticks(t0, t1, rate):
        axR.plot([tt / 60, tt / 60], [0, 1], color="0.6", lw=1.0)
axR.set_xlim(0, tmin[-1]); axR.set_ylim(0, 1.5); axR.set_yticks([])
axR.set_ylabel("pre-syn.\nspikes", fontsize=9)
for s in ("top", "right", "left"):
    axR.spines[s].set_visible(False)
axR.axvspan(108 / 60, 168 / 60, color=INDIGO, alpha=0.12, lw=0)
axR.text(54 / 60, 1.18, "1 Hz", ha="center", fontsize=9, color="0.3")
axR.text(138 / 60, 1.18, "10 Hz", ha="center", fontsize=9, color=INDIGO)
axR.text(222 / 60, 1.18, "1 Hz", ha="center", fontsize=9, color="0.3")
axR.tick_params(labelbottom=False)

axI.axvspan(108 / 60, 168 / 60, color=INDIGO, alpha=0.12, lw=0)
axI.plot(tmin, I, color=GREEN, lw=1.9)
axI.set_xlim(0, tmin[-1]); axI.set_ylim(0, 1.08)
axI.set_xlabel("Time (min)"); axI.set_ylabel(r"Device current, $I/I_{\max}$")
for s in ("top", "right"):
    axI.spines[s].set_visible(False)
axI.annotate("settles low under\nsustained high rate", xy=(150 / 60, I[np.argmin(np.abs(tg - 150))]),
             xytext=(2.0, 0.95), fontsize=8.5, color="0.3", ha="center", va="top",
             arrowprops=dict(arrowstyle="->", color="0.5", lw=0.8))
axI.annotate("recovers when\nrate returns to 1 Hz", xy=(255 / 60, I[np.argmin(np.abs(tg - 255))]),
             xytext=(3.25, 0.26), fontsize=8.5, color="0.3", ha="left", va="center",
             arrowprops=dict(arrowstyle="->", color="0.5", lw=0.8))
fig.suptitle("Extended model reproduces the measured rate habituation (cf. Chapter 4, Fig. 4l)",
             fontsize=11, y=0.96)
plt.show()

# ---- cell 5 --------------------------------------------------------
d = np.load(DM / "ito_decay_data.npz", allow_pickle=True)
betas_ito = np.asarray(d["beta"], float)
taus_ito = np.asarray(d["tau"], float)

# panel A (example traces + stretched-exp fits) needs the raw .xls sweeps; render it if
# present, otherwise skip gracefully (the fixture carries only the per-device fits).
import glob, os
RAW = os.path.join(str(paths.data_dir().parents[0]), "mnn-torch", "data", "ITO data")
raw_xls = sorted(glob.glob(os.path.join(RAW, "*voltage*bias*.xls")))
if raw_xls:
    print(f"raw ITO .xls present ({len(raw_xls)}) -- example-trace panel could be re-fit here")
else:
    print("raw ITO .xls sweeps absent -- example-trace panel skipped; replaying the "
          "committed per-device beta distribution instead")

# beta distribution vs the gold rise (Fig 2 of gen_fig_ito_decay.py)
fig, axB = plt.subplots(figsize=(6.6, 4.4))
axB.hist(betas_ito, bins=np.linspace(0, 2.2, 23), color=GREEN, alpha=0.85, edgecolor="white")
axB.axvline(np.median(betas_ito), color=INDIGO, lw=2.2,
            label=fr"ITO decay median $\beta$={np.median(betas_ito):.2f}")
axB.axvline(1.0, color="0.5", ls=":", lw=1.4, label=r"$\beta=1$ (single exponential)")
axB.axvline(2.0, color=RED, ls="--", lw=2.0, label=r"gold rise $\beta\approx2$ (compressed)")
axB.set_xlabel(r"Stretched-exponent $\beta$"); axB.set_ylabel("Number of ITO devices")
axB.legend(fontsize=8.0, loc="upper right", bbox_to_anchor=(0.995, 0.84),
           frameon=True, framealpha=0.95, edgecolor="none")
for s in ("top", "right"):
    axB.spines[s].set_visible(False)
_clean(axB)
fig.suptitle(fr"Dispersive ITO decay ($\beta\!<\!1$), n={len(betas_ito)} devices", fontsize=11, y=0.97)
fig.tight_layout(); plt.show()

print(f"ITO decay: n={len(betas_ito)}  median beta={np.median(betas_ito):.2f} "
      f"(mean {betas_ito.mean():.2f} +/- {betas_ito.std():.2f})  median tau={np.median(taus_ito):.2f} s")

# ---- cell 6 --------------------------------------------------------
# Uncomment to re-lock the device-model fixtures as a subprocess (fast -- seconds):
# import subprocess, sys
# subprocess.run([sys.executable, "-m", "mrl_trace.device", "--kww", "--habituation", "--full"])
print("see the markdown above for the full-scale command")
\end{lstlisting}

\begin{lstlisting}[caption={experiments/schematics.ipynb},label={lst:siox-nb-schematics}]
# ---- cell 1 --------------------------------------------------------
import os
import numpy as np
import matplotlib

import matplotlib.pyplot as plt
from matplotlib.patches import (Circle, Rectangle, FancyBboxPatch, Wedge,
                                FancyArrowPatch, Polygon)

# ---- palette -------------------------------------------------------------
INK    = "#1f2a37"
CHAMB  = "#eef2f6"
XRAY   = "#c0392b"     # X-ray beam (crimson, the single warm accent)
EBEAM  = "#2f4b8f"     # photoelectron path (navy)
ION    = "#3aa07a"     # Ar+ sputter beam (teal)
ANALY  = "#cfd8e0"
LENS   = "#9aa6b2"
DETBOX = "#3aa07a"
TXT    = "#1f2a37"
ITO_C  = "#cdd9ee"     # ITO layer (light navy tint)
SIOX_C = "#e3eaef"     # SiOx (very light)
MO_C   = "#aeb8c2"     # Mo (metal grey)
SI_C   = "#cdbfa6"     # Si substrate (tan, hatched)

FS_LABEL = 14
FS_SMALL = 12
FS_PM    = 16

plt.rcParams.update({"font.family": "DejaVu Sans", "font.size": 13})

fig, ax = plt.subplots(figsize=(9.0, 7.4))
ax.set_xlim(0, 12.4); ax.set_ylim(0, 10.6); ax.axis("off")
ax.set_aspect("equal")

def arrow(p0, p1, color=INK, lw=2.2, z=6, ls="-", ms=16):
    ax.add_patch(FancyArrowPatch(p0, p1, arrowstyle="-|>", mutation_scale=ms,
                 color=color, lw=lw, ls=ls, zorder=z, shrinkA=0, shrinkB=0))

def wire(pts, lw=2.1, color=INK, z=2, ls="-"):
    xs, ys = zip(*pts)
    ax.plot(xs, ys, color=color, lw=lw, ls=ls, solid_capstyle="round", zorder=z)

# =====================================================================
# UHV chamber
# =====================================================================
CX, CY, CR = 4.35, 4.85, 3.15
ax.add_patch(Circle((CX, CY), CR, facecolor=CHAMB, edgecolor=INK, lw=2.6, zorder=1))
# UHV chamber label (leader to lower-left rim; text sits in clear space below
# the circle so it does not sit on top of the leader line)
xrim, yrim = CX + CR*np.cos(np.deg2rad(212)), CY + CR*np.sin(np.deg2rad(212))
ax.annotate("UHV chamber", xy=(xrim, yrim), xytext=(0.10, 1.35),
            fontsize=FS_LABEL, color=TXT, va="top", ha="left",
            arrowprops=dict(arrowstyle="-", color=INK, lw=1.3,
                            shrinkB=2, shrinkA=4))

# pumps below chamber
ax.add_patch(Rectangle((CX-0.34, CY-CR-0.42), 0.68, 0.45,
                       facecolor="white", edgecolor=INK, lw=1.9, zorder=1))
ax.add_patch(Rectangle((CX-0.92, CY-CR-0.76), 1.84, 0.34,
                       facecolor=LENS, edgecolor=INK, lw=1.9, zorder=1))
ax.add_patch(Rectangle((CX-0.58, CY-CR-1.48), 1.16, 0.72,
                       facecolor="white", edgecolor=INK, lw=1.9, zorder=1))
ax.text(CX, CY-CR-1.66, "pumps", ha="center", va="top",
        fontsize=FS_LABEL, color=TXT)

# =====================================================================
# Sample = the layered stack, in place on the stage INSIDE the chamber
# =====================================================================
sx, shw = 3.70, 0.92            # stack centre x, half-width
sl, sr = sx-shw, sx+shw
layers = [  # (name, y0, y1, facecolor, hatch)
    ("ITO (50 nm)",     4.55, 5.20, ITO_C,  None),
    (r"SiO$_x$ (25 nm)", 4.18, 4.55, SIOX_C, None),
    ("Mo (150 nm)",     3.55, 4.18, MO_C,   None),
    ("Si substrate",    2.95, 3.55, SI_C,   "//"),
]
for name, y0, y1, fc, hatch in layers:
    ax.add_patch(Rectangle((sl, y0), 2*shw, y1-y0, facecolor=fc,
                 edgecolor=INK, lw=1.8, hatch=hatch, zorder=4))
    # short leader to a label that sits just right of the stack, well inside
    # the chamber wall (kept clear of the circle border)
    ax.annotate(name, xy=(sr, (y0+y1)/2), xytext=(sr+0.22, (y0+y1)/2),
                fontsize=FS_SMALL, color=TXT, va="center", ha="left",
                zorder=7,
                arrowprops=dict(arrowstyle="-", color=INK, lw=1.0))
ST_TOP = 5.20
# sample stage / holder beneath the stack
ax.add_patch(Rectangle((sx-0.55, 2.60), 1.10, 0.35,
                       facecolor=LENS, edgecolor=INK, lw=1.5, zorder=3))

# =====================================================================
# X-ray source: mounted OUTSIDE the chamber wall (upper left); the beam
# enters through the wall and strikes the stack surface. The source body sits
# beyond the circle so it does not lie inside the chamber.
# =====================================================================
ang = np.deg2rad(38)
hx, hy = sl+0.32, ST_TOP+0.12                # hit point at the stack surface
bx1, by1 = hx - 2.85*np.cos(ang), hy + 2.85*np.sin(ang)   # source mouth (outside)
# barrel, drawn pointing back away from the chamber
dpx, dpy = 0.30*np.sin(ang), 0.30*np.cos(ang)
ax.add_patch(Polygon([(bx1-dpx, by1-dpy), (bx1+dpx, by1+dpy),
                      (bx1+dpx + 0.60*np.cos(ang), by1+dpy + 0.60*np.sin(ang)),
                      (bx1-dpx + 0.60*np.cos(ang), by1-dpy + 0.60*np.sin(ang))],
             closed=True, facecolor=LENS, edgecolor=INK, lw=1.7, zorder=3))
ax.text(bx1 + 0.60*np.cos(ang) + 0.05, by1 + 0.60*np.sin(ang) + 0.12,
        "X-ray\nsource", ha="center", va="bottom", fontsize=FS_LABEL, color=TXT)
# beam: from source mouth to the surface
arrow((bx1, by1), (hx, hy), color=XRAY, lw=2.3, ls=(0, (5, 2)), z=6, ms=16)
# h-nu label offset above the beam line (clear of the dotted beam & source)
mx, my = (bx1+hx)/2, (by1+hy)/2
ax.text(mx + 0.46, my + 0.20, r"$h\nu$", ha="center", va="center",
        fontsize=FS_LABEL, color=XRAY, zorder=7)

# =====================================================================
# Ar+ sputter ion beam -> stack surface (steep, from the right side so it
# clears the lens column and its label sits in clear space at the rim)
# =====================================================================
arrow((sx+1.35, ST_TOP+0.78), (sx+0.50, ST_TOP+0.03), color=ION, lw=2.3,
      ls=(0, (2, 2)), z=6, ms=16)
ax.text(sx+1.42, ST_TOP+0.82, r"Ar$^{+}$ sputter", ha="left", va="bottom",
        fontsize=FS_SMALL, color=ION, zorder=7)

# small "increasing etch depth" axis just left of the stack
arrow((sl-0.42, ST_TOP-0.05), (sl-0.42, 3.05), color=INK, lw=1.7, z=5, ms=12)
ax.text(sl-0.68, (ST_TOP+3.05)/2, "etch depth", rotation=90, ha="center",
        va="center", fontsize=FS_SMALL, color=TXT)

# =====================================================================
# photoelectron -> lens -> analyser -> multiplier -> spectrum
# =====================================================================
LCX = sx
# photoelectron leaves the stack surface straight up into the lens column
arrow((sx, ST_TOP+0.05), (LCX, ST_TOP+0.78), color=EBEAM, lw=2.1, z=6, ms=15)
ax.text(sx+0.30, ST_TOP+0.45, r"e$^{-}$", ha="left", va="center",
        fontsize=FS_SMALL, color=EBEAM, zorder=7)
# lens plates (clearly above the stack)
for yy in (ST_TOP+0.95, ST_TOP+1.27, ST_TOP+1.59):
    ax.add_patch(Rectangle((LCX-0.40, yy-0.06), 0.80, 0.13,
                           facecolor=LENS, edgecolor=INK, lw=1.4, zorder=5))
ax.annotate("lens system", xy=(LCX+0.40, ST_TOP+1.40), xytext=(LCX+0.75, ST_TOP+1.70),
            fontsize=FS_LABEL, color=TXT, va="center", ha="left",
            arrowprops=dict(arrowstyle="-", color=INK, lw=1.3))

# hemispherical analyser (above the chamber)
HX, HY = LCX, 9.05
R_out, R_in = 1.12, 0.66
ax.add_patch(Wedge((HX, HY), R_out, 0, 180, width=0.34,
                   facecolor=ANALY, edgecolor=INK, lw=2.1, zorder=3))
ax.add_patch(Wedge((HX, HY), R_in, 0, 180, width=0.30,
                   facecolor=ANALY, edgecolor=INK, lw=2.1, zorder=3))
ax.text(HX, HY + R_out + 0.16, "electron energy analyser",
        ha="center", va="bottom", fontsize=FS_LABEL, color=TXT)
ax.text(HX, HY + (R_in+R_out)/2 + 0.02, "$-$", ha="center", va="center",
        fontsize=FS_PM, color=INK, zorder=6)
ax.text(HX, HY + R_in - 0.40, "$+$", ha="center", va="center",
        fontsize=FS_PM, color=INK, zorder=6)
th = np.linspace(np.pi, 0, 60); rr = (R_in+R_out)/2
ax.plot(HX + rr*np.cos(th), HY + rr*np.sin(th), color=EBEAM, lw=1.8,
        ls=(0, (4, 2)), zorder=4)
xen = HX - (R_in+R_out)/2; xex = HX + (R_in+R_out)/2
wire([(LCX, ST_TOP+1.65), (LCX, HY-0.55), (xen, HY-0.55), (xen, HY)],
     color=EBEAM, lw=2.1, z=4)

# electron multiplier (detector) to the upper right -- box sized to the text
DETX, DETY = 9.35, HY
DET_HW, DET_HH = 1.30, 0.62
arrow((xex+0.02, HY), (DETX-DET_HW-0.02, DETY), color=EBEAM, lw=2.1, z=5, ms=15)
ax.add_patch(FancyBboxPatch((DETX-DET_HW, DETY-DET_HH), 2*DET_HW, 2*DET_HH,
             boxstyle="round,pad=0.02,rounding_size=0.10",
             facecolor="white", edgecolor=DETBOX, lw=2.3, zorder=4))
ax.text(DETX, DETY, "electron\nmultiplier", ha="center", va="center",
        fontsize=FS_LABEL, color=TXT, zorder=5)

# XPS spectrum panel below the detector -- larger so the trace + label fit
SPW, SPH = 2.60, 1.65
SPX, SPY = DETX-SPW/2, 6.05
arrow((DETX, DETY-DET_HH), (DETX, SPY+SPH+0.02), color=EBEAM, lw=2.1, z=5, ms=14)
ax.add_patch(FancyBboxPatch((SPX, SPY), SPW, SPH,
             boxstyle="round,pad=0.02,rounding_size=0.08",
             facecolor="white", edgecolor=INK, lw=1.9, zorder=4))
gx = np.linspace(0, 1, 200)
def peak(x0, w, a): return a*np.exp(-((gx-x0)/w)**2)
gy = 0.12 + peak(0.34, 0.05, 0.58) + peak(0.64, 0.07, 0.95)
ax.plot(SPX + 0.26 + gx*(SPW-0.52), SPY + 0.30 + gy*(SPH-0.74),
        color=DETBOX, lw=1.8, zorder=5)
ax.text(SPX+SPW/2, SPY+SPH-0.06, "XPS spectrum", ha="center", va="top",
        fontsize=FS_SMALL, color=TXT, zorder=6)

plt.show()

# ---- cell 2 --------------------------------------------------------
import os
import numpy as np
import matplotlib

import matplotlib.pyplot as plt
from matplotlib.patches import Rectangle

INK   = "#2b2b2b"
MO    = "#9aa6b2"
OXIDE = "#eef1f4"
GOLD  = "#e0a93b"
ITO   = "#3aa07a"
ITOr  = "#6b4f2a"
TRAP  = "#3aa07a"
PROT  = "#c0392b"
GREEN = "#3aa07a"
NEU   = "#9aa6b2"

TRAP_PTS = None   # shared trap distribution, set in main()

def _stack(ax, reduced, top_color, top_label):
    """One MIM stack centred in the axes: Mo (grounded) / SiOx / top electrode (-V).

    Bias convention matches the device measurement: a negative bias is applied to
    the top electrode and the Mo bottom electrode is grounded, so the field drives
    the positive mobile species (a proton, with oxygen vacancies as the coupled
    species) UP toward the negatively-biased top electrode."""
    ax.set_xlim(0, 4); ax.set_ylim(0, 3.7); ax.axis("off")
    x0, w = 0.7, 2.6
    y_mo, h_mo = 0.25, 0.45
    y_ox, h_ox = 0.70, 1.7
    y_top, h_top = 2.40, 0.45
    # Mo bottom electrode (grounded)
    ax.add_patch(Rectangle((x0, y_mo), w, h_mo, fc=MO, ec=INK, lw=1.0, zorder=2))
    ax.text(x0+w/2, y_mo+h_mo/2, "Mo", ha="center", va="center", fontsize=8, color="white")
    # ground symbol on Mo
    gx = x0 + w + 0.18
    ax.plot([x0+w, gx], [y_mo+h_mo/2, y_mo+h_mo/2], color=INK, lw=1.0)
    for i, hw in enumerate([0.13, 0.08, 0.04]):
        yy = y_mo+h_mo/2 - 0.10 - 0.06*i
        ax.plot([gx-hw, gx+hw], [yy, yy], color=INK, lw=1.0)
    ax.plot([gx, gx], [y_mo+h_mo/2, y_mo+h_mo/2-0.10], color=INK, lw=1.0)
    # SiOx
    ax.add_patch(Rectangle((x0, y_ox), w, h_ox, fc=OXIDE, ec=INK, lw=1.0, zorder=2))
    ax.text(x0+0.08, y_ox+h_ox-0.12, r"SiO$_x$", ha="left", va="top", fontsize=8, color=INK)
    # top electrode + applied -V
    if reduced:
        ax.add_patch(Rectangle((x0, y_top), w, h_top, fc=ITOr, ec=INK, lw=1.0, zorder=2, hatch="////"))
    else:
        ax.add_patch(Rectangle((x0, y_top), w, h_top, fc=top_color, ec=INK, lw=1.0, zorder=2))
    ax.text(x0+w/2, y_top+h_top/2, top_label, ha="center", va="center", fontsize=8, color="white")
    ax.annotate(r"$-V$ applied", xy=(x0+w/2, y_top+h_top), xytext=(x0+w/2, y_top+h_top+0.55),
                ha="center", fontsize=8.0, color=INK,
                arrowprops=dict(arrowstyle="-|>", color=INK, lw=1.3))
    # distributed traps (same pattern in both)
    xs = x0 + 0.18 + (w-0.36)*TRAP_PTS[:, 0]
    ys = y_ox + 0.20 + (h_ox-0.40)*TRAP_PTS[:, 1]
    ax.scatter(xs, ys, s=15, facecolor="white", edgecolor=TRAP, lw=1.0, zorder=4)
    # mobile positive species drifts UP toward the negative top electrode
    for fx in np.linspace(0.28, 0.72, 4):
        x = x0 + fx*w
        ax.annotate("", xy=(x, y_top-0.04), xytext=(x, y_ox+0.35),
                    arrowprops=dict(arrowstyle="-|>", color=PROT, lw=1.3, alpha=0.9))
        ax.text(x+0.07, y_ox+0.62, r"H$^{+}$", color=PROT, fontsize=6.6, va="center", zorder=5)
    return dict(x0=x0, w=w, y_top=y_top, y_ox=y_ox, h_ox=h_ox)

def panel_a(ax):
    ax.text(-0.06, 1.12, "(a)", transform=ax.transAxes, fontsize=13, fontweight="bold")
    ax.text(0.5, 1.02, "Ti/Au: inert contact", transform=ax.transAxes,
            ha="center", fontsize=9.5, fontweight="bold")
    g = _stack(ax, reduced=False, top_color=GOLD, top_label="Au")
    # accumulation layer (+) just under the inert gold
    ax.add_patch(Rectangle((g["x0"], g["y_top"]-0.15), g["w"], 0.15, fc=PROT, ec="none", alpha=0.6, zorder=3))
    for fx in np.linspace(0.15, 0.85, 6):
        ax.text(g["x0"]+fx*g["w"], g["y_top"]-0.075, "+", color="white",
                ha="center", va="center", fontsize=8, fontweight="bold", zorder=5)
    ax.text(2.0, 0.00, r"H$^{+}$ accumulates at inert Au", ha="center", fontsize=8.0, color=PROT)
    ax.text(2.0, -0.30, "$\\rightarrow$ space charge screens the field", ha="center", fontsize=7.8, color=PROT)

def panel_b(ax):
    ax.text(-0.06, 1.12, "(b)", transform=ax.transAxes, fontsize=13, fontweight="bold")
    ax.text(0.5, 1.02, "ITO: reduced by protons", transform=ax.transAxes,
            ha="center", fontsize=9.5, fontweight="bold")
    g = _stack(ax, reduced=True, top_color=ITO, top_label="ITO")
    for fx in np.linspace(0.2, 0.8, 5):
        ax.text(g["x0"]+fx*g["w"], g["y_top"]-0.01, r"$\times$", color=PROT,
                ha="center", va="center", fontsize=9, fontweight="bold", zorder=5)
    ax.text(2.0, 0.00, r"H$^{+}$ reduces ITO, consumed", ha="center", fontsize=8.0, color=GREEN)
    ax.text(2.0, -0.30, "$\\rightarrow$ no space charge accumulates", ha="center", fontsize=7.8, color=GREEN)

def _transient_axes(ax):
    ax.set_xlim(0, 6); ax.set_ylim(0, 1.08)
    ax.set_xlabel("time (s)", fontsize=8.5); ax.set_ylabel(r"$|I|$", fontsize=9)
    ax.tick_params(labelsize=7.5)

def panel_c(ax):
    """Gold transient: rise then slow tau_d decay -> masks the discharge."""
    ax.text(-0.06, 1.12, "(c)", transform=ax.transAxes, fontsize=13, fontweight="bold")
    ax.set_title("Au: slow decay dominates", fontsize=9.5, fontweight="bold", pad=4)
    t = np.linspace(0, 6, 400)
    rise = 1 - np.exp(-(t/0.35)**2)
    I = rise*np.exp(-t/9.0); I = I/I.max()
    ax.plot(t, I, color=PROT, lw=2.4)
    _transient_axes(ax)
    # (tau_d arrow and "trap discharge hidden underneath" text removed -- the
    #  caption states the slow proton decay masks the faster trap discharge)

def panel_d(ax):
    """ITO transient: bare dispersive discharge exposed -> tau_leak."""
    ax.text(-0.06, 1.12, "(d)", transform=ax.transAxes, fontsize=13, fontweight="bold")
    ax.set_title("ITO: bare discharge exposed", fontsize=9.5, fontweight="bold", pad=4)
    t = np.linspace(0, 6, 400)
    rise = 1 - np.exp(-(t/0.30)**2)
    I = rise*np.exp(-(t/1.3)**0.54); I = I/I.max()
    ax.plot(t, I, color=TRAP, lw=2.4)
    _transient_axes(ax)
    # (tau_leak arrow/annotation removed -- the caption states ITO exposes the
    #  bare discharge tau_leak)

def panel_e(ax):
    """Why dispersive: distribution of rates -> stretched exponential."""
    ax.text(-0.06, 1.12, "(e)", transform=ax.transAxes, fontsize=13, fontweight="bold")
    ax.set_title("Why dispersive", fontsize=9.5, fontweight="bold", pad=4)
    t = np.linspace(0, 6, 400)
    for tt in [0.4, 0.8, 1.6, 3.2]:
        ax.plot(t, np.exp(-t/tt), color=NEU, lw=0.9, alpha=0.6)
    ax.plot(t, np.exp(-(t/1.3)**0.54), color=TRAP, lw=2.6, label=r"sum: $\beta\!=\!0.54$")
    ax.plot(t, np.exp(-t/1.3), color=PROT, lw=1.5, ls=(0, (4, 2)), label=r"single: $\beta\!=\!1$")
    ax.set_xlim(0, 6); ax.set_ylim(0, 1.02)
    ax.set_xlabel("time (arb.)", fontsize=8.5); ax.set_ylabel("norm. $I$", fontsize=8.5)
    ax.tick_params(labelsize=7.5)
    ax.legend(fontsize=7.4, loc="upper right", framealpha=0.92)
    # ("grey: individual release rates" floating text removed -- caption explains
    #  the spread of release rates summing to a stretched exponential)

def panel_f(ax):
    """One physics, two faces: compressed fill mirrors stretched discharge."""
    ax.text(-0.06, 1.12, "(f)", transform=ax.transAxes, fontsize=13, fontweight="bold")
    ax.set_title("One physics, two faces", fontsize=9.5, fontweight="bold", pad=4)
    t = np.linspace(0, 1, 200)
    ax.plot(t, 1-np.exp(-(t/0.18)**2), color=GOLD, lw=2.4, label=r"fill $\beta\!\approx\!2$")
    ax.plot(t, np.exp(-(t/0.30)**0.54), color=TRAP, lw=2.4, label=r"discharge $\beta\!\approx\!0.5$")
    ax.set_xlim(0, 1); ax.set_ylim(0, 1.05)
    ax.set_xlabel("time (arb.)", fontsize=8.5); ax.set_ylabel("norm. $I$", fontsize=8.5)
    ax.set_xticks([0, 0.5, 1.0]); ax.tick_params(labelsize=7.5)
    # floating "fill/discharge" labels and the italic subtitle removed; the two
    # curves are distinguished by a compact legend and explained in the caption
    ax.legend(fontsize=7.4, loc="center right", framealpha=0.92)

global TRAP_PTS
TRAP_PTS = np.random.RandomState(7).rand(11, 2)
fig, axes = plt.subplots(2, 3, figsize=(11.0, 6.4))
panel_a(axes[0, 0]); panel_b(axes[0, 1]); panel_e(axes[0, 2])
panel_c(axes[1, 0]); panel_d(axes[1, 1]); panel_f(axes[1, 2])
fig.subplots_adjust(left=0.06, right=0.97, top=0.92, bottom=0.12, wspace=0.32, hspace=0.42)

plt.show()

# ---- cell 3 --------------------------------------------------------
import os
import numpy as np
import matplotlib

import matplotlib.pyplot as plt
from matplotlib.patches import FancyArrowPatch, Circle, Rectangle, FancyBboxPatch

INK   = "#2b2b2b"
CB    = "#c8d0d8"
TRAP  = "#3aa07a"     # volatile tag / trapped electrons
ELEC  = "#3aa07a"
FIL   = "#6b4f2a"     # filament / structural non-volatile weight
CAP   = "#c0392b"     # reward / capture-write
GREEN = "#3aa07a"     # coincidence
ASSUM = "#b07cc6"     # ASSUMED coupling (purple, to flag prediction)
NEU   = "#9aa6b2"

def _trapped_electrons(ax, x0, y0, n, filled):
    """Row of trap sites; 'filled' of n carry an electron (the surviving tag)."""
    xs = np.linspace(x0, x0+2.0, n)
    for i, x in enumerate(xs):
        ax.plot([x-0.12, x+0.12], [y0, y0], color=TRAP, lw=2.0, zorder=3)
        if i < filled:
            ax.add_patch(Circle((x, y0+0.13), 0.10, fc=ELEC, ec=INK, lw=0.7, zorder=5))
    return xs

def _filament(ax, cx, y_bot, y_top, strength):
    """A vertical filament whose width/height encodes the non-volatile weight G."""
    w = 0.10 + 0.28*strength
    ax.add_patch(FancyBboxPatch((cx-w/2, y_bot), w, (y_top-y_bot)*(0.4+0.6*strength),
                 boxstyle="round,pad=0.01", fc=FIL, ec=INK, lw=1.0, zorder=4))

def panel_a(ax):
    ax.text(-0.06, 1.10, "(a)", transform=ax.transAxes, fontsize=13, fontweight="bold")
    ax.set_title("Two state variables in one cell", fontsize=9.5, fontweight="bold", pad=4)
    ax.set_xlim(0, 10); ax.set_ylim(0, 10); ax.axis("off")
    # electrodes
    ax.add_patch(Rectangle((1.0, 8.4), 8.0, 0.5, fc=NEU, ec=INK, lw=1.0)); ax.text(9.1, 8.65, "top", fontsize=7.2, va="center")
    ax.add_patch(Rectangle((1.0, 0.7), 8.0, 0.5, fc=NEU, ec=INK, lw=1.0)); ax.text(9.1, 0.95, "Mo", fontsize=7.2, va="center")
    ax.text(0.9, 4.8, r"SiO$_x$", fontsize=8, ha="right", color=INK)
    # volatile tag: trapped electrons (left)
    _trapped_electrons(ax, 1.7, 5.2, 4, filled=3)
    ax.text(2.5, 6.1, "volatile tag", color=TRAP, fontsize=8.4, ha="center", fontweight="bold")
    ax.text(2.5, 3.9, "trapped electrons\n$e(t)$, decays over $\\tau_{\\mathrm{leak}}$", color=TRAP,
            fontsize=7.4, ha="center")
    # non-volatile weight: filament drawn at the far right, labels to its LEFT (clear of it)
    _filament(ax, 8.5, 1.2, 8.4, strength=0.45)
    ax.text(6.6, 6.1, "non-volatile weight", color=FIL, fontsize=8.4, ha="center", fontweight="bold")
    ax.text(6.6, 3.9, "filament conductance $G$\n(measured ladder)", color=FIL, fontsize=7.4, ha="center")
    ax.annotate("", xy=(8.2, 4.8), xytext=(7.4, 4.8),
                arrowprops=dict(arrowstyle="-", color=FIL, lw=0.8))
    # divider note
    ax.text(5.0, 2.2, "distinct states: the tag alone does not move the weight",
            fontsize=7.0, ha="center", color=NEU, style="italic")

def _cell(ax, e_level, committed, title, ok, note):
    """Shared cell drawing for (b)/(c): tag occupancy e_level, weight grows if committed."""
    ax.set_xlim(0, 10); ax.set_ylim(0, 10); ax.axis("off")
    ax.set_title(title, fontsize=9.5, fontweight="bold", pad=4)
    ax.add_patch(Rectangle((1.0, 8.4), 8.0, 0.5, fc=NEU, ec=INK, lw=1.0))
    ax.add_patch(Rectangle((1.0, 0.7), 8.0, 0.5, fc=NEU, ec=INK, lw=1.0))
    nfill = int(round(4*e_level))
    _trapped_electrons(ax, 1.5, 5.0, 5, filled=nfill)
    ax.text(2.5, 5.9, f"tag $e(t_R)\\!\\approx\\!{e_level:.1f}$", color=TRAP, fontsize=8.0, ha="center")
    # the ASSUMED coupling: trap-assisted switching, drawn as a labelled box that names
    # the physical mechanism (the box IS the trapped-charge gating of the filament write)
    ax.add_patch(FancyBboxPatch((3.3, 6.45), 3.6, 1.35, boxstyle="round,pad=0.04",
                 fc="white", ec=ASSUM, lw=1.6, zorder=5))
    ax.text(5.1, 7.45, "trap-assisted switching", color=ASSUM, fontsize=7.6, ha="center", va="center", fontweight="bold", zorder=7)
    ax.text(5.1, 7.06, r"gate $g(e)$", color=ASSUM, fontsize=7.2, ha="center", va="center", zorder=7)
    ax.text(5.1, 6.70, "(assumed coupling)", color=ASSUM, fontsize=6.6, ha="center", va="center", style="italic", zorder=7)
    # reward pulse arriving from the top electrode, stopping at the box top
    ax.add_patch(FancyArrowPatch((5.1, 9.55), (5.1, 7.86), arrowstyle="-|>",
                 mutation_scale=13, color=CAP, lw=2.0, zorder=6))
    ax.text(5.1, 9.8, "reward pulse", color=CAP, fontsize=7.8, ha="center")
    # tag modulates the gate (left -> gate)
    ax.add_patch(FancyArrowPatch((3.2, 5.3), (3.7, 6.5), arrowstyle="-|>",
                 mutation_scale=9, color=TRAP, lw=1.2, ls=(0,(2,1.3)),
                 connectionstyle="arc3,rad=0.2", zorder=4))
    # gated write: gate -> filament (thickness/alpha scale with how open the gate is)
    a_alpha = 1.0 if ok else 0.22
    ax.add_patch(FancyArrowPatch((6.7, 6.9), (8.3, 4.6), arrowstyle="-|>",
                 mutation_scale=12, color=CAP, lw=1.2+1.8*e_level, alpha=a_alpha,
                 connectionstyle="arc3,rad=-0.25", zorder=5))
    # filament: grows only if committed
    _filament(ax, 8.5, 1.2, 8.4, strength=0.45 + (0.4 if committed else 0.0))
    if committed:
        ax.text(7.0, 3.0, r"$\Delta G$ committed", color=CAP, fontsize=7.8, ha="center", fontweight="bold")
    else:
        ax.text(7.0, 3.0, r"$\Delta G\!\approx\!0$", color=NEU, fontsize=7.8, ha="center")
    ax.text(5.0, 0.12, note, fontsize=7.2, ha="center", color=INK if ok else NEU, style="italic")

def panel_b(ax):
    ax.text(-0.06, 1.10, "(b)", transform=ax.transAxes, fontsize=13, fontweight="bold")
    _cell(ax, e_level=0.8, committed=True, title="Reward within the window", ok=True,
          note="electrons still trapped $\\Rightarrow$ gate open $\\Rightarrow$ structural write")

def panel_c(ax):
    ax.text(-0.06, 1.10, "(c)", transform=ax.transAxes, fontsize=13, fontweight="bold")
    _cell(ax, e_level=0.1, committed=False, title="Reward too late", ok=False,
          note="traps emptied $\\Rightarrow$ gate shut $\\Rightarrow$ no write")

def panel_d(ax):
    ax.text(-0.06, 1.10, "(d)", transform=ax.transAxes, fontsize=13, fontweight="bold")
    ax.set_title("The assumed coupling $g(e)$", fontsize=9.5, fontweight="bold", pad=4)
    e = np.linspace(0, 1, 300)
    for p, col, lab in [(1.0, NEU, "$p=1$"), (3.0, ASSUM, "$p=3$ (soft threshold)")]:
        g = e**p/(e**p + 0.25**p)
        ax.plot(e, g, color=col, lw=2.4 if p > 1 else 1.6,
                ls="-" if p > 1 else (0,(4,2)), label=lab)
    ax.set_xlim(0, 1); ax.set_ylim(0, 1.05)
    ax.set_xlabel(r"surviving tag $e(t_R)$", fontsize=8.5)
    ax.set_ylabel(r"capture fraction $g(e)$", fontsize=8.5)
    ax.tick_params(labelsize=7.5)
    ax.legend(fontsize=7.4, loc="lower right", framealpha=0.92)
    ax.text(0.04, 0.86, r"$\Delta G=\kappa(R\!-\!b)\,g(e)\,\Delta G_0$", transform=ax.transAxes,
            fontsize=8.2, color=CAP)
    ax.text(0.04, 0.70, "single ASSUMED element;\ntarget of the $\\Delta G$-delay\nmeasurement",
            transform=ax.transAxes, fontsize=7.2, color=ASSUM, style="italic")

fig, axes = plt.subplots(2, 2, figsize=(10.0, 7.2))
panel_a(axes[0, 0]); panel_b(axes[0, 1])
panel_c(axes[1, 0]); panel_d(axes[1, 1])
fig.suptitle("Prediction: electron-level capture by a second-order cell [Eq. (10)]",
             fontsize=11, fontweight="bold", y=0.98)
fig.subplots_adjust(left=0.06, right=0.97, top=0.90, bottom=0.09, wspace=0.20, hspace=0.40)

plt.show()

# ---- cell 4 --------------------------------------------------------
import os
import numpy as np
import matplotlib

import matplotlib.pyplot as plt
from matplotlib.patches import FancyArrowPatch, Circle, Rectangle, FancyBboxPatch

INK   = "#2b2b2b"
CB    = "#cdd6df"     # conduction band shade
TRAP  = "#3aa07a"     # trap level / trapped electrons / trace
ELEC  = "#3aa07a"
CAP   = "#c0392b"     # reward / capture-write
GREEN = "#c75c2e"     # coincidence (spike accent, distinct from teal trace)
NEU   = "#9aa6b2"
PAL   = [plt.cm.viridis(x) for x in (0.85, 0.5, 0.15)]   # tau_leak curves: sequential viridis

def panel_a(ax):
    ax.text(-0.06, 1.10, "(a)", transform=ax.transAxes, fontsize=13, fontweight="bold")
    ax.set_title("Capture, then emission at a trap", fontsize=9.6, fontweight="bold", pad=4)
    ax.set_xlim(0, 10); ax.set_ylim(0, 10); ax.axis("off")

    # conduction band
    ax.fill_between([0.5, 9.5], 8.2, 10, color=CB, zorder=1)
    ax.plot([0.5, 9.5], [8.2, 8.2], color=INK, lw=1.4, zorder=2)
    ax.text(9.4, 9.1, "conduction band", ha="right", fontsize=8.2, color=INK)

    # A short SEQUENTIAL cascade of trap sites at energetically-distributed depths
    # (faithful to Ch5: carriers hop through a chain v_n1 -> v_n2 -> ... -> v_nk;
    # the sites are spatially and energetically distributed, giving a SPREAD of
    # release times on discharge). Capture fills the chain; emission empties it.
    traps = [(2.2, 6.7), (4.4, 5.7), (6.6, 4.4)]   # distributed depths along the chain
    for j, (x, y) in enumerate(traps):
        ax.plot([x-0.5, x+0.5], [y, y], color=TRAP, lw=2.2, zorder=3)
        ax.add_patch(Circle((x, y+0.17), 0.15, fc=ELEC, ec=INK, lw=0.8, zorder=5))
        ax.text(x, y-0.5, rf"$v_{{n{j+1}}}$", ha="center", fontsize=7.4, color=TRAP)
        # sequential hop to the next site (the cascade rise)
        if j < len(traps)-1:
            xn, yn = traps[j+1]
            ax.add_patch(FancyArrowPatch((x+0.5, y), (xn-0.5, yn), arrowstyle="-|>",
                         mutation_scale=10, color=GREEN, lw=1.4,
                         connectionstyle="arc3,rad=-0.15", zorder=4))
    # capture from band into the chain head (during coincidence)
    ax.add_patch(FancyArrowPatch((2.2, 8.1), (2.2, 6.95), arrowstyle="-|>",
                 mutation_scale=12, color=GREEN, lw=1.7, zorder=4))
    # emission back to band, once the drive stops (dispersive: a spread of rates)
    for (x, y) in traps:
        ax.add_patch(FancyArrowPatch((x+0.30, y+0.30), (x+0.30, 8.1), arrowstyle="-|>",
                     mutation_scale=10, color=CAP, lw=1.2, ls=(0,(2,1.3)), zorder=4))
    # depth bracket: distribution of trap energies
    ax.annotate("", xy=(8.6, 8.2), xytext=(8.6, 4.4),
                arrowprops=dict(arrowstyle="<->", color=INK, lw=1.0))
    ax.text(8.8, 6.3, r"spread of$\ E_T$", fontsize=8.0, color=INK, rotation=90, va="center")
    # labels
    ax.text(2.0, 1.45, "capture (coincidence):", color=GREEN, fontsize=8.0, ha="center")
    ax.text(2.0, 0.95, "fills the sequential chain", color=GREEN, fontsize=7.2, ha="center")
    ax.text(6.8, 1.45, r"emission (drive off): rate $\propto e^{\gamma\sqrt{V}}$", color=CAP, fontsize=8.0, ha="center")
    ax.text(6.8, 0.95, "distributed $E_T$ $\\Rightarrow$ spread of release times", color=CAP, fontsize=7.2, ha="center")

def panel_b(ax):
    ax.text(-0.06, 1.10, "(b)", transform=ax.transAxes, fontsize=13, fontweight="bold")
    ax.set_title("Trapped population $=$ eligibility trace", fontsize=9.6, fontweight="bold", pad=4)
    dt = 0.01
    t = np.arange(0, 10, dt)
    tau, beta = 2.2, 0.6
    # build during a brief coincidence, then dispersive decay (clamp t-1>=0 so the
    # fractional power never sees a negative base)
    peak = 1 - np.exp(-(1.0/0.18)**2)
    td = np.maximum(t - 1.0, 0.0)
    e = np.where(t < 1.0, 1 - np.exp(-(t/0.18)**2),
                 peak*np.exp(-(td/tau)**beta))
    e = e / e.max()
    ax.fill_between(t, 0, e, color=TRAP, alpha=0.15)
    ax.plot(t, e, color=TRAP, lw=2.4)
    ax.axvspan(0, 1.0, color=GREEN, alpha=0.25)
    ax.text(0.5, 1.06, "coincidence", color=GREEN, fontsize=7.8, ha="center")
    # tau_leak marker at 1/e of peak
    ax.annotate("", xy=(1.0+tau, 0.05), xytext=(1.0, 0.05),
                arrowprops=dict(arrowstyle="<->", color=INK, lw=1.0))
    ax.text(1.0+tau/2, 0.12, r"$\tau_{\mathrm{leak}}=R_{\mathrm{leak}}C$", ha="center", fontsize=8.2, color=INK)
    # (floating "occupancy of trapped electrons = e(t)" annotation removed)
    ax.set_xlim(0, 10); ax.set_ylim(0, 1.12)
    ax.set_xlabel("time after coincidence (s)", fontsize=8.5)
    ax.set_ylabel("trapped population $e(t)$", fontsize=8.5)
    ax.tick_params(labelsize=7.5)

def panel_c(ax):
    ax.text(-0.06, 1.10, "(c)", transform=ax.transAxes, fontsize=13, fontweight="bold")
    ax.set_title("Reward captures the surviving trace", fontsize=9.6, fontweight="bold", pad=4)
    dt = 0.01
    t = np.arange(0, 10, dt)
    tau, beta = 2.2, 0.6
    peak = 1 - np.exp(-(1.0/0.18)**2)
    td = np.maximum(t - 1.0, 0.0)
    e = np.where(t < 1.0, 1 - np.exp(-(t/0.18)**2),
                 peak*np.exp(-(td/tau)**beta))
    e = e / e.max()
    ax.fill_between(t, 0, e, color=TRAP, alpha=0.15)
    ax.plot(t, e, color=TRAP, lw=2.2)
    # two reward times
    for tR, ok in [(3.0, True), (8.0, False)]:
        eR = e[int(tR/dt)]
        ax.axvline(tR, color=CAP, lw=1.6, ls=(0, (4, 2)), alpha=1.0 if ok else 0.55)
        ax.plot([tR], [eR], "o", color=CAP, ms=6, zorder=5)
        ax.text(tR, 1.04, "reward", color=CAP, fontsize=7.6, ha="center", alpha=1.0 if ok else 0.6)
        # committed weight change bar (proportional to e(t_R))
        ax.add_patch(Rectangle((tR+0.12, 0.0), 0.5, eR*0.8, fc=CAP, ec="none",
                               alpha=0.8 if ok else 0.3, zorder=4))
    # (floating "within window..." and "too late..." annotations removed; the
    #  reward markers and the Delta-w equation below convey the same point)
    ax.text(5.0, -0.30, r"$\Delta w_{ij}=\eta\,(R-b)\,e_{ij}(t_R)$  [Eq.~(3)]",
            transform=ax.transAxes, ha="center", fontsize=8.4, color=CAP)
    ax.set_xlim(0, 10); ax.set_ylim(0, 1.12)
    ax.set_xlabel("action$\\rightarrow$reward delay (s)", fontsize=8.5)
    ax.set_ylabel("eligibility $e(t)$", fontsize=8.5)
    ax.tick_params(labelsize=7.5)

def panel_d(ax):
    ax.text(-0.06, 1.10, "(d)", transform=ax.transAxes, fontsize=13, fontweight="bold")
    ax.set_title(r"$\tau_{\mathrm{leak}}$ sets the credit window", fontsize=9.6, fontweight="bold", pad=4)
    dt = 0.01
    t = np.arange(0, 10, dt)
    peak = 1 - np.exp(-(1.0/0.18)**2)
    for tau, col, lab in [(0.8, PAL[2], r"$\tau_{\mathrm{leak}}$ short (shallow traps)"),
                          (2.2, PAL[1], r"$\tau_{\mathrm{leak}}$ mid"),
                          (5.0, PAL[0], r"$\tau_{\mathrm{leak}}$ long (deep traps)")]:
        td = np.maximum(t - 1.0, 0.0)
        e = np.where(t < 1.0, 1 - np.exp(-(t/0.18)**2),
                     peak*np.exp(-(td/tau)**0.6))
        e = e/e.max()
        ax.plot(t, e, color=col, lw=2.2, label=lab)
    ax.axvspan(0, 1.0, color=GREEN, alpha=0.18)
    ax.set_xlim(0, 10); ax.set_ylim(0, 1.08)
    ax.set_xlabel("time after coincidence (s)", fontsize=8.5)
    ax.set_ylabel("eligibility $e(t)$", fontsize=8.5)
    ax.tick_params(labelsize=7.5)
    ax.legend(fontsize=7.2, loc="upper right", framealpha=0.92)
    # (floating "fast-forgetting <-> integrating" annotation removed; the legend
    #  of short/mid/long tau_leak curves already conveys this)

fig, axes = plt.subplots(2, 2, figsize=(10.0, 7.2))
panel_a(axes[0, 0]); panel_b(axes[0, 1])
panel_c(axes[1, 0]); panel_d(axes[1, 1])
fig.subplots_adjust(left=0.07, right=0.97, top=0.92, bottom=0.10, wspace=0.26, hspace=0.42)

plt.show()

# ---- cell 5 --------------------------------------------------------
import os
import numpy as np
import matplotlib

import matplotlib.pyplot as plt
from matplotlib.patches import (FancyBboxPatch, FancyArrowPatch, Circle, Rectangle,
                                Patch, Ellipse, Polygon, Wedge)

INK    = "#2b2b2b"
BIO    = "#b07cc6"
AXON   = "#e7b27a"     # axon terminal
SPINE  = "#b07cc6"     # dendritic spine
TAGB   = "#3aa07a"     # chemical tag (green)
VES    = "#8c5a3c"     # vesicles
DOPA   = "#c0392b"     # dopamine / reward
LTP    = "#b07cc6"     # LTP growth
TRAP   = "#2f4b8f"     # trap / electron
LEAK   = "#c75c2e"     # dispersive leak
REW    = "#c0392b"     # global scalar (R-b)
ELEC   = "#9aa6b2"     # electrode
OX     = "#eef1f4"     # oxide
FIL    = "#6b4f2a"     # filament
PRED   = "#b07cc6"     # predicted
NEU    = "#9aa6b2"

# row y-centres (top -> bottom) and their times -- generously spaced
ROWS = [(13.4, "0 s",  "action"),
        (9.9,  "3 s",  "wait"),
        (6.4,  "5 s",  "reward"),
        (2.6,  "10 s", "capture")]
XL = 3.3     # biology column centre
XR = 10.7    # silicon column centre
XC = 7.0     # timeline x
SY = 0.9     # synapse glyph scale
SC = 1.05    # cell glyph scale

# ----------------------------------------------------------------- synapse glyph
def synapse(ax, cx, cy, grow=0.0, lit=False, tag=None, faded=False):
    """A conventional chemical synapse: a presynaptic axon terminal (bouton) with
    vesicles on top, a synaptic cleft, and a postsynaptic mushroom dendritic spine
    (head on a neck into the dendrite shaft) below, receptors in the head. Signal
    flows top->bottom. Late-LTP enlarges the spine head and adds receptors; the tag
    is a mark in the spine head."""
    s = SY
    # ---- presynaptic axon + bouton (top) ----
    ax.add_patch(Rectangle((cx-0.10*s, cy+1.05*s), 0.20*s, 0.55*s, fc=AXON, ec=INK, lw=1.2, zorder=3))  # axon stalk
    bout_w, bout_h = 0.95*s, 0.62*s
    if lit:
        ax.add_patch(Ellipse((cx, cy+0.62*s), bout_w+0.14*s, bout_h+0.14*s, fc="none", ec="#e0a93b", lw=2.6, zorder=2))
    ax.add_patch(Ellipse((cx, cy+0.62*s), bout_w, bout_h, fc=AXON, ec=INK, lw=1.6, zorder=4))            # bouton
    # vesicles clustered toward the active zone (bottom of the bouton)
    rng = np.random.RandomState(2)
    for i in range(5):
        vx = cx-0.30*s + 0.60*s*rng.rand(); vy = cy+0.45*s + 0.34*s*rng.rand()
        ax.add_patch(Circle((vx, vy), 0.075*s, fc="white", ec=VES, lw=1.0, zorder=5))
    # ---- synaptic cleft ----
    ax.add_patch(Rectangle((cx-0.50*s, cy+0.20*s), 1.0*s, 0.06*s, fc="none", ec="none", zorder=3))
    if lit:  # neurotransmitter released into the cleft
        for k in range(4):
            ax.add_patch(Circle((cx-0.28*s+0.56*s*rng.rand(), cy+0.20*s+0.06*s*rng.rand()),
                         0.05*s, fc=TAGB, ec="none", alpha=0.85, zorder=6))
    # ---- postsynaptic mushroom spine (head + neck into dendrite) ----
    head_r = (0.42 + 0.26*grow)*s
    head_y = cy - 0.10*s - head_r
    ax.add_patch(Rectangle((cx-0.22*s, head_y-0.55*s), 0.44*s, 0.5*s, fc=SPINE, ec=INK, lw=1.2, zorder=2))  # dendrite shaft
    ax.add_patch(Rectangle((cx-0.09*s, head_y-0.10*s), 0.18*s, 0.32*s,
                 fc=(LTP if grow > 0 else SPINE), ec=INK, lw=1.1, zorder=3))                               # neck
    ax.add_patch(Circle((cx, head_y), head_r, fc=(LTP if grow > 0 else SPINE), ec=INK, lw=1.6,
                 alpha=0.95 if grow > 0 else 1.0, zorder=4))                                               # head
    # receptors on the head facing the cleft (top arc; more when grown)
    nrec = 3 + int(round(3*grow))
    for a in np.linspace(-0.7, 0.7, nrec):
        rx = cx + head_r*np.sin(a); ry = head_y + head_r*np.cos(a)
        ax.add_patch(Rectangle((rx-0.045*s, ry-0.05*s), 0.09*s, 0.12*s, fc=INK, ec="none", zorder=5))
    # the chemical tag inside the spine head
    if tag is not None:
        tc = "#a9d9c5" if faded else TAGB
        ax.add_patch(Circle((cx, head_y), 0.15*s, fc=tc, ec=INK, lw=0.9, zorder=6,
                     alpha=0.5 if faded else 1.0))

# ----------------------------------------------------------------- SiOx cell glyph
def cell(ax, cx, cy, fil=0.0, trapped=0):
    s = SC
    ax.add_patch(Rectangle((cx-0.95*s, cy+0.50*s), 1.9*s, 0.20*s, fc=ELEC, ec=INK, lw=0.9, zorder=4))   # top electrode
    ax.add_patch(Rectangle((cx-0.95*s, cy-0.70*s), 1.9*s, 0.20*s, fc=ELEC, ec=INK, lw=0.9, zorder=4))   # Mo bottom
    ax.add_patch(Rectangle((cx-0.95*s, cy-0.50*s), 1.9*s, 1.00*s, fc=OX, ec=INK, lw=0.9, zorder=3))     # SiOx
    ax.text(cx-0.80*s, cy+0.40*s, r"SiO$_x$", fontsize=7.0, color=INK, va="top", zorder=6)
    # filament (thickness = non-volatile weight)
    fw = (0.10 + 0.42*fil)*s
    if fil > 0:
        ax.add_patch(Rectangle((cx+0.55*s-fw/2, cy-0.50*s), fw, 1.00*s, fc=FIL, ec=INK, lw=0.7, zorder=5))
    # trap sites (filled = trapped electrons), left of the filament
    rng = np.random.RandomState(7)
    for i in range(6):
        x = cx-0.62*s + 0.55*s*rng.rand(); y = cy-0.34*s + 0.68*s*rng.rand()
        filled = i < trapped
        ax.add_patch(Circle((x, y), 0.075*s, fc=(TRAP if filled else "white"), ec=TRAP, lw=1.0, zorder=6))

# ----------------------------------------------------------------- main
fig, ax = plt.subplots(figsize=(11.0, 13.4))
ax.set_xlim(0, 14); ax.set_ylim(0, 15.8); ax.axis("off")

# column headers (in the top margin, clear of the first row band which tops at 14.85)
ax.text(XL, 15.5, "BIOLOGY", ha="center", fontsize=15, fontweight="bold", color=BIO)
ax.text(XL, 15.12, "synaptic tagging and capture", ha="center", fontsize=9.5, color=BIO, style="italic")
ax.text(XR, 15.5, "SILICON", ha="center", fontsize=15, fontweight="bold", color=TRAP)
ax.text(XR, 15.12, "subthreshold SiO$_x$ memristor", ha="center", fontsize=9.5, color=TRAP, style="italic")

# segmented timeline: a separate arrow into each node, with the event name
# sitting in the GAP just above the node (no overlap with the arrow or label)
node_r = 0.56
seg_tops = [14.85] + [ROWS[i][0]-node_r for i in range(len(ROWS)-1)]
for (y, t, lab), top in zip(ROWS, seg_tops):
    # faint full-width row band
    ax.add_patch(Rectangle((0.4, y-1.72), 13.2, 3.36, fc="#f4f5f7", ec="none", zorder=0))
    # event-name label sits in the gap above the node, beside the segment
    ax.text(XC, (top+y+node_r)/2, lab, ha="center", va="center", fontsize=11,
            color=INK, fontweight="bold", style="italic", zorder=7,
            bbox=dict(boxstyle="round,pad=0.18", fc="white", ec="none"))
    # arrow for this segment, stopping at the node top
    ax.add_patch(FancyArrowPatch((XC, top), (XC, y+node_r), arrowstyle="-|>",
                 mutation_scale=20, color=INK, lw=2.2, zorder=2))
    # the node circle with the timestamp
    ax.add_patch(Circle((XC, y), node_r, fc="white", ec=INK, lw=2.0, zorder=6))
    ax.text(XC, y, t, ha="center", va="center", fontsize=11.5, fontweight="bold", color=INK, zorder=7)
# tail arrow below the last node
ax.add_patch(FancyArrowPatch((XC, ROWS[-1][0]-node_r), (XC, 0.95), arrowstyle="-|>",
             mutation_scale=20, color=INK, lw=2.2, zorder=2))
ax.text(XC, 0.6, "time", ha="center", fontsize=10.5, color=INK, fontweight="bold")

# ---------------- BIOLOGY ----------------
y = ROWS[0][0]; synapse(ax, XL, y, grow=0.0, lit=True, tag=TAGB)
ax.text(XL, y-1.18, "terminal fires; sets a transient tag", ha="center", fontsize=10.0, color=TAGB)
y = ROWS[1][0]; synapse(ax, XL, y, grow=0.0, tag=TAGB, faded=True)
ax.text(XL, y-1.18, "the chemical tag fades", ha="center", fontsize=10.0, color=TAGB)
y = ROWS[2][0]; synapse(ax, XL, y, grow=0.0, tag=TAGB)
for dx, dy in [(-1.7, 1.0), (-1.1, 1.25), (1.3, 1.05), (1.9, 0.65), (0.5, 1.35), (-0.3, 1.1)]:
    ax.add_patch(Circle((XL+dx, y+dy), 0.12, fc=DOPA, ec="none", alpha=0.85, zorder=5))
ax.add_patch(FancyArrowPatch((XL+0.9, y+1.0), (XL+0.1, y+0.2), arrowstyle="-|>",
             mutation_scale=12, color=DOPA, lw=1.5, zorder=6))
ax.text(XL, y-1.18, "dopamine binds only the surviving tag", ha="center", fontsize=10.0, color=DOPA)
y = ROWS[3][0]; synapse(ax, XL, y, grow=1.0, tag=None)
ax.text(XL, y-1.58, "late-LTP: the spine enlarges", ha="center", fontsize=10.0, color=LTP)

# ---------------- SILICON ----------------
y = ROWS[0][0]; cell(ax, XR, y, fil=0.28, trapped=6)
ax.add_patch(FancyArrowPatch((XR-1.8, y), (XR-1.0, y), arrowstyle="-|>",
             mutation_scale=13, color=TRAP, lw=1.7, zorder=6))
ax.text(XR, y-1.30, "LIF fires; electrons into deep traps", ha="center", fontsize=10.0, color=TRAP)
y = ROWS[1][0]; cell(ax, XR, y, fil=0.28, trapped=2)
for dx in (-0.45, -0.20, 0.05):
    ax.add_patch(FancyArrowPatch((XR+dx, y+0.40), (XR+dx, y+0.82), arrowstyle="-|>",
                 mutation_scale=9, color=LEAK, lw=1.2, ls=(0,(2,1)), zorder=6))
# dispersive-leak inset, tucked into the right margin beside this cell (not floating)
ix, iy, iw, ih = XR+1.25, y-0.45, 1.05, 0.95
tt = np.linspace(0, 1, 60); kww = np.exp(-(tt/0.35)**0.54)
ax.plot(ix+tt*iw, iy+kww*ih, color=LEAK, lw=2.0, zorder=6)
ax.plot([ix, ix, ix+iw], [iy+ih, iy, iy], color=NEU, lw=0.8, zorder=5)   # tiny axes
ax.text(ix+iw*0.55, iy+ih*0.78, r"$\beta\!<\!1$", fontsize=8.6, color=LEAK, ha="left")
ax.text(ix+iw*0.5, iy-0.22, "leak", fontsize=7.2, color=LEAK, ha="center", style="italic")
ax.text(XR, y-1.30, "electrons detrap: dispersive leak", ha="center", fontsize=10.0, color=LEAK)
y = ROWS[2][0]; cell(ax, XR, y, fil=0.28, trapped=2)
# the (R-b) pulse is gated by the surviving tag: it must meet the trapped
# electrons (left of the cell), not act on the filament directly -- the
# filament change is the consequence, shown in the next (capture) row
ax.add_patch(FancyArrowPatch((XR+1.8, y+0.95), (XR-0.25, y+0.18), arrowstyle="-|>",
             mutation_scale=14, color=REW, lw=2.3, zorder=7))
ax.text(XR+1.95, y+1.05, r"$(R\!-\!b)$", ha="center", fontsize=11, color=REW, fontweight="bold")
ax.text(XR, y-1.30, "global scalar meets trapped electrons", ha="center", fontsize=10.0, color=REW)
y = ROWS[3][0]; cell(ax, XR, y, fil=1.0, trapped=1)
ax.text(XR, y-1.30, "non-volatile weight committed", ha="center", fontsize=10.0, color=FIL)
ax.text(XR, y-1.72, r"three-factor write: $\Delta w=\eta(R-b)\,e(t_R)$", ha="center", fontsize=8.8, color=FIL, style="italic")

fig.suptitle("From a synaptic tag to a silicon tag: a biological-to-device translation map",
             fontsize=13.5, fontweight="bold", y=0.975)
fig.subplots_adjust(left=0.01, right=0.99, top=0.93, bottom=0.015)

plt.show()

# ---- cell 6 --------------------------------------------------------
import os
import numpy as np
import matplotlib

import matplotlib.pyplot as plt
from matplotlib.patches import (Circle, FancyBboxPatch, FancyArrowPatch,
                                Rectangle)

INK = "#2b2b2b"; CUE = "#e0a93b"; DIST = "#9aa6b2"; REW = "#c0392b"
TRACE = "#3aa07a"; DEV = "#eaf0fb"; OUT = "#cfd8e0"

plt.rcParams.update({"font.family": "DejaVu Sans"})
fig = plt.figure(figsize=(10.0, 4.4))
gsL = fig.add_axes([0.015, 0.04, 0.52, 0.92]); gsL.axis("off")
gsR = fig.add_axes([0.60, 0.04, 0.39, 0.92]); gsR.axis("off")
gsL.set_xlim(0, 10.6); gsL.set_ylim(0, 10)
gsR.set_xlim(0, 10); gsR.set_ylim(0, 10)

def arrow(ax, p0, p1, color=INK, lw=1.8, ms=12, ls="-", z=5, rad=0.0):
    cs = f"arc3,rad={rad}" if rad else None
    ax.add_patch(FancyArrowPatch(p0, p1, arrowstyle="-|>", mutation_scale=ms,
                 color=color, lw=lw, ls=ls, zorder=z, shrinkA=0, shrinkB=0,
                 connectionstyle=cs))

# ============================ PANEL A ============================
gsL.text(-0.6, 9.7, "(a)", fontsize=15, fontweight="bold", va="top")

# --- network sits in the upper band (y ~ 4.5-9) ---
ins_y = [8.4, 7.3, 6.2, 5.1]
labels = ["cue", "distractor", "distractor", r"$\vdots$"]
incol = [CUE, DIST, DIST, DIST]
ix = 1.5
out = (6.4, 6.75)
synx = (ix + out[0]) / 2          # synapse-box column
bus_x = 8.3                        # reward bus column (right of everything)

for y, lab, c in zip(ins_y, labels, incol):
    if lab == r"$\vdots$":
        gsL.text(ix, y, r"$\vdots$", fontsize=22.4, ha="center", va="center", color=INK)
        continue
    gsL.add_patch(Circle((ix, y), 0.40, facecolor="white", edgecolor=c, lw=2.2, zorder=4))
    gsL.text(ix - 0.70, y, lab, fontsize=14, ha="right", va="center", color=c)
    # synapse box on the wire to the output
    gsL.add_patch(FancyBboxPatch((synx-0.34, y-0.21), 0.68, 0.42,
                  boxstyle="round,pad=0.02,rounding_size=0.06",
                  facecolor=DEV, edgecolor=INK, lw=1.2, zorder=5))
    gsL.plot([ix+0.40, synx-0.34], [y, y], color=c, lw=1.5, zorder=3)
    arrow(gsL, (synx+0.34, y), (out[0]-0.55, out[1]), color=c, lw=1.5, ms=10, z=3)
gsL.text(synx, ins_y[0]+0.62, "device synapses", fontsize=12.6,
         ha="center", color=INK, style="italic")

# LIF output neuron
gsL.add_patch(Circle(out, 0.55, facecolor=OUT, edgecolor=INK, lw=2.4, zorder=4))
gsL.text(out[0], out[1], "LIF", fontsize=14, ha="center", va="center", color=INK, fontweight="bold")
gsL.text(out[0], out[1]-0.85, "output", fontsize=13.3, ha="center", va="top", color=INK)

# --- reward delivered on a single clean vertical bus (no crossing lines) ---
# bus runs down the right side; a short horizontal stub taps each synapse box.
bus_top = ins_y[0] + 0.05
bus_bot = 3.7
gsL.plot([bus_x, bus_x], [bus_bot, bus_top], color=REW, lw=2.0, zorder=2)
for y in ins_y[:3]:
    arrow(gsL, (bus_x, y), (synx+0.36, y), color=REW, lw=1.2, ms=9,
          ls=(0, (3, 2)), z=2)
    gsL.add_patch(Circle((bus_x, y), 0.06, facecolor=REW, edgecolor=REW, zorder=3))
# reward source box feeding the bus from below (centred on the bus, fully inside axes)
rb_w, rb_h = 2.8, 1.0
rb_cx, rb_cy = bus_x, 3.15
gsL.add_patch(FancyBboxPatch((rb_cx-rb_w/2, rb_cy-rb_h/2), rb_w, rb_h,
              boxstyle="round,pad=0.03,rounding_size=0.1",
              facecolor="white", edgecolor=REW, lw=2.0, zorder=5))
gsL.text(rb_cx, rb_cy+0.18, r"reward $R(t)$", fontsize=14, ha="center", va="center", color=REW, zorder=6)
gsL.text(rb_cx, rb_cy-0.22, "global third factor", fontsize=11.2, ha="center", va="center", color=REW, style="italic", zorder=6)
gsL.plot([bus_x, bus_x], [rb_cy+rb_h/2, bus_bot], color=REW, lw=2.0, zorder=2)
gsL.text(bus_x+0.18, (bus_bot+bus_top)/2, "broadcast", fontsize=10.5, ha="left",
         va="center", color=REW, rotation=90, style="italic")

# --- timeline in its own band at the bottom (y ~ 0.4-2.0), well separated ---
ty = 1.1
gsL.annotate("", xy=(9.8, ty), xytext=(0.4, ty),
             arrowprops=dict(arrowstyle="-|>", color=INK, lw=1.3))
gsL.text(9.7, ty-0.45, "time", fontsize=11.9, ha="right", color=INK)
# cue epoch block (below the axis), label below it
gsL.add_patch(Rectangle((1.0, ty-0.30), 1.8, 0.30, facecolor=CUE, alpha=0.45, edgecolor="none"))
gsL.text(1.9, ty-0.55, "cue epoch", fontsize=11.9, ha="center", va="top", color=CUE)
# delay span (above the axis), label above it
gsL.annotate("", xy=(7.0, ty+0.30), xytext=(2.8, ty+0.30),
             arrowprops=dict(arrowstyle="<->", color=INK, lw=1.1))
gsL.text(4.9, ty+0.42, r"action$\to$reward delay $D$", fontsize=11.9, ha="center", va="bottom", color=INK)
# reward marker (on the axis), label above the tick, clear of the delay span
gsL.plot([7.0, 7.0], [ty-0.18, ty+0.18], color=REW, lw=2.4)
gsL.text(7.0, ty-0.30, "reward", fontsize=11.9, ha="center", va="top", color=REW)

# ============================ PANEL B ============================
gsR.text(-0.6, 9.7, "(b)", fontsize=15, fontweight="bold", va="top")
# 1. coincidence
gsR.add_patch(FancyBboxPatch((0.5, 8.0), 9.0, 1.05, boxstyle="round,pad=0.03,rounding_size=0.1",
              facecolor="white", edgecolor=INK, lw=1.6))
gsR.text(5.0, 8.72, "pre--post coincidence (signed, leak-dominant)", fontsize=10.8, ha="center", color=INK)
gsR.text(2.6, 8.25, r"causal: $+1$", fontsize=11.9, ha="center", color=CUE)
gsR.text(7.2, 8.25, r"acausal: $-\lambda$", fontsize=11.9, ha="center", color=REW)
arrow(gsR, (5.0, 8.0), (5.0, 7.35), lw=1.6)
# 2. device gate box (widened to match the coincidence box so the subtitle fits)
gsR.add_patch(FancyBboxPatch((0.5, 5.1), 9.0, 2.1, boxstyle="round,pad=0.03,rounding_size=0.1",
              facecolor=DEV, edgecolor=INK, lw=1.8))
gsR.text(5.0, 6.85, "device transient gate", fontsize=13.3, ha="center", color=INK, fontweight="bold")
gsR.text(5.0, 6.40, r"trap cascade ($k\!\approx\!3$, rise) $+$ $R_{\mathrm{leak}}$ (relax)", fontsize=10.8, ha="center", color=INK)
# mini trace inside
tx = np.linspace(0, 1, 120)
e = (1 - np.exp(-(tx/0.18)**2)) * np.exp(-tx/0.45); e = e/e.max()
gsR.plot(2.2 + tx*5.0, 5.35 + e*0.75, color=TRACE, lw=1.8)
gsR.text(7.6, 5.7, r"$e(t)$", fontsize=12.6, color=TRACE, ha="left")
arrow(gsR, (5.0, 5.1), (5.0, 4.45), lw=1.6)
# 3. eligibility trace label
gsR.text(5.0, 4.12, r"eligibility trace $e_{ij}(t)$,  retention $\tau_{\mathrm{leak}}\!=\!R_{\mathrm{leak}}C$",
         fontsize=12.3, ha="center", color=TRACE)
arrow(gsR, (5.0, 3.75), (5.0, 3.05), lw=1.6)
# 4. update box
gsR.add_patch(FancyBboxPatch((1.0, 1.6), 8.0, 1.35, boxstyle="round,pad=0.03,rounding_size=0.1",
              facecolor="white", edgecolor=REW, lw=2.0))
gsR.text(5.0, 2.50, "three-factor update at reward", fontsize=13.3, ha="center", color=REW, fontweight="bold")
gsR.text(5.0, 1.92, r"$\Delta w_{ij} = \eta\,(R-b)\,e_{ij}(t_R)$", fontsize=16.8, ha="center", color=INK)
plt.show()

# ---- cell 7 --------------------------------------------------------
import os
import numpy as np
import matplotlib

import matplotlib.pyplot as plt
from matplotlib.patches import (Circle, FancyBboxPatch, FancyArrowPatch,
                                Rectangle)

INK = "#2b2b2b"; STATE = "#e0a93b"; ACT = "#2f4b8f"; REW = "#c0392b"
TRACE = "#3aa07a"; DEV = "#eaf0fb"; OUT = "#cfd8e0"; GRID = "#9aa6b2"

fig = plt.figure(figsize=(9.8, 4.3))
A = fig.add_axes([0.015, 0.04, 0.575, 0.92]); A.axis("off")
B = fig.add_axes([0.625, 0.04, 0.365, 0.92]); B.axis("off")
A.set_xlim(0, 11); A.set_ylim(0, 9.6)
B.set_xlim(0, 10); B.set_ylim(0, 9)

def arrow(ax, p0, p1, color=INK, lw=1.7, ms=12, ls="-", z=5):
    ax.add_patch(FancyArrowPatch(p0, p1, arrowstyle="-|>", mutation_scale=ms,
                 color=color, lw=lw, ls=ls, zorder=z, shrinkA=0, shrinkB=0))

# ============================ PANEL A: crossbar ============================
A.text(-0.2, 8.8, "(a)", fontsize=15, fontweight="bold", va="top")

rows_y = [7.1, 6.0]                       # two state wordlines
cols_x = [4.2, 5.7]                       # two action bitlines
row_x0, row_x1 = 1.7, 6.6
col_y0, col_y1 = 4.6, 7.6

# state input wordlines (rows)
state_lab = [r"$s_1$", r"$s_2$"]
for y, lab in zip(rows_y, state_lab):
    A.plot([row_x0, row_x1], [y, y], color=STATE, lw=2.0, zorder=2)
    A.add_patch(Circle((row_x0 - 0.45, y), 0.34, facecolor="white",
                       edgecolor=STATE, lw=2.0, zorder=4))
    A.text(row_x0 - 0.45, y, lab, fontsize=9.5, ha="center", va="center", color=STATE)
A.text(row_x0 - 0.95, (rows_y[0] + rows_y[1]) / 2, "state\ninputs", fontsize=8.4,
       ha="right", va="center", color=STATE, style="italic")

# action bitlines (cols) + compound-cell synapses at crosspoints
for x in cols_x:
    A.plot([x, x], [col_y0, col_y1], color=ACT, lw=2.0, zorder=2)

def compound_cell(ax, cx, cy):
    """One synapse drawn as its true compound cell: a non-volatile WEIGHT device (left,
    dark) + a subthreshold TRACE device (right, light) + a small active-gating element
    (the access transistor / selector) marked beneath. Compact so the 2x2 array stays
    legible while honestly showing two devices per synapse rather than one glyph."""
    # weight device (left square, filled)
    ax.add_patch(FancyBboxPatch((cx - 0.30, cy - 0.15), 0.27, 0.30,
                 boxstyle="round,pad=0.01,rounding_size=0.04",
                 facecolor=DEV, edgecolor=INK, lw=1.0, zorder=5))
    # trace device (right square, light fill, trace colour edge)
    ax.add_patch(FancyBboxPatch((cx + 0.03, cy - 0.15), 0.27, 0.30,
                 boxstyle="round,pad=0.01,rounding_size=0.04",
                 facecolor="white", edgecolor=TRACE, lw=1.3, zorder=5))
    # active-gating element: a small filled circle (transistor/selector) under the pair,
    # tapping the third-factor line
    ax.add_patch(Circle((cx, cy - 0.32), 0.07, facecolor=REW, edgecolor=INK,
                        lw=0.8, zorder=6))

for y in rows_y:
    for x in cols_x:
        compound_cell(A, x, y)
# compact inline legend to the RIGHT of the array (clear of the left-margin labels),
# stacked above the I=G^T V annotation
leg_x = cols_x[1] + 1.15
leg_y = col_y1 + 0.55
A.add_patch(FancyBboxPatch((leg_x, leg_y - 0.10), 0.20, 0.22, boxstyle="round,pad=0.01",
            facecolor=DEV, edgecolor=INK, lw=0.8, zorder=5))
A.text(leg_x + 0.30, leg_y, "weight device", fontsize=6.8, ha="left", va="center", color=INK)
A.add_patch(FancyBboxPatch((leg_x, leg_y - 0.50), 0.20, 0.22, boxstyle="round,pad=0.01",
            facecolor="white", edgecolor=TRACE, lw=1.0, zorder=5))
A.text(leg_x + 0.30, leg_y - 0.40, "trace device", fontsize=6.8, ha="left", va="center", color=TRACE)
A.add_patch(Circle((leg_x + 0.10, leg_y - 0.80), 0.07, facecolor=REW, edgecolor=INK, lw=0.7, zorder=6))
A.text(leg_x + 0.30, leg_y - 0.80, "active gating", fontsize=6.8, ha="left", va="center", color=REW)
A.text(leg_x, leg_y + 0.40, "each synapse:", fontsize=7.0, ha="left", va="center",
       color=INK, style="italic")

# bitline current accumulation label
A.text(cols_x[1] + 1.15, (col_y0 + col_y1) / 2, r"$\mathbf{I}=\mathbf{G}^{\!\top}\mathbf{V}$",
       fontsize=10.5, ha="left", va="center", color=ACT)
A.text(cols_x[1] + 1.15, (col_y0 + col_y1) / 2 - 0.55, "(VMM along\nbitlines)",
       fontsize=7.6, ha="left", va="center", color=ACT, style="italic")

# action LIF neurons below the columns
act_y = 3.2
act_lab = [r"$a_1$", r"$a_2$"]
for x, lab in zip(cols_x, act_lab):
    arrow(A, (x, col_y0), (x, act_y + 0.5), color=ACT, lw=1.6, ms=10, z=3)
    A.add_patch(Circle((x, act_y), 0.42, facecolor=OUT, edgecolor=INK, lw=2.0, zorder=4))
    A.text(x, act_y, "LIF", fontsize=8.2, ha="center", va="center", color=INK, fontweight="bold", zorder=6)
    # label beside the neuron (not below) so it clears the WTA arrow on the bitline
    A.text(x + 0.55, act_y, lab, fontsize=9.0, ha="left", va="center", color=ACT, zorder=6)
A.text(cols_x[1] + 1.15, act_y - 0.5, "action\nneurons", fontsize=8.2, ha="left",
       va="center", color=ACT, style="italic")

# winner-take-all action selection
wta_y = 1.55
A.add_patch(FancyBboxPatch((cols_x[0] - 0.9, wta_y - 0.34), 3.3, 0.68,
            boxstyle="round,pad=0.03,rounding_size=0.1",
            facecolor="white", edgecolor=INK, lw=1.5, zorder=5))
A.text((cols_x[0] + cols_x[1]) / 2, wta_y, r"WTA: $a=\arg\max_j\!\sum_t s^{\mathrm{post}}_j$",
       fontsize=8.6, ha="center", va="center", color=INK, zorder=6)
for x in cols_x:
    arrow(A, (x, act_y - 0.45), (x, wta_y + 0.36), color=INK, lw=1.3, ms=8, z=3)

# environment / reward (contingent, delayed)
env_x = 9.55
A.add_patch(FancyBboxPatch((env_x - 1.15, wta_y - 0.42), 2.3, 1.5,
            boxstyle="round,pad=0.03,rounding_size=0.1",
            facecolor="white", edgecolor=REW, lw=1.8, zorder=5))
A.text(env_x, wta_y + 0.72, "environment", fontsize=8.6, ha="center", va="center", color=REW, zorder=6)
A.text(env_x, wta_y + 0.28, r"$R=\mathbf{1}(a{=}a^\star_s)$", fontsize=8.6, ha="center", va="center", color=INK, zorder=6)
A.text(env_x, wta_y - 0.16, r"delay $D$", fontsize=8.0, ha="center", va="center", color=REW, style="italic", zorder=6)
arrow(A, (cols_x[1] + 1.75, wta_y), (env_x - 1.2, wta_y), color=INK, lw=1.4, ms=10, z=3)

# global third-factor line: R-b broadcast back to every crosspoint
gf_y = 8.95
A.plot([env_x, env_x], [wta_y + 1.08, gf_y], color=REW, lw=1.6, ls=(0, (4, 2)), zorder=2)
A.plot([row_x0 - 0.1, env_x], [gf_y, gf_y], color=REW, lw=1.6, ls=(0, (4, 2)), zorder=2)
A.text((row_x0 + cols_x[1]) / 2 - 0.3, gf_y + 0.26,
       r"global third factor $(R-b)$ broadcast to all synapses",
       fontsize=8.2, ha="center", va="bottom", color=REW)
for x in cols_x:
    arrow(A, (x, gf_y), (x, rows_y[0] + 0.22), color=REW, lw=1.0, ms=7, ls=(0, (2, 2)), z=2)

# (generalization annotation removed -- the caption states the array is S x A)

# ============================ PANEL B: one crosspoint ============================
B.text(-0.2, 8.8, "(b)", fontsize=15, fontweight="bold", va="top")
B.text(5.0, 8.5, "single-layer synapse: one-term update", fontsize=9.0, ha="center", va="top",
       color=INK, fontweight="bold")

# weight device (non-volatile, switching regime) -- differential pair
B.add_patch(FancyBboxPatch((0.7, 5.7), 8.6, 1.9,
            boxstyle="round,pad=0.03,rounding_size=0.1",
            facecolor=DEV, edgecolor=INK, lw=1.5))
B.text(5.0, 7.25, "weight device (non-volatile, switching regime)", fontsize=8.2, ha="center", color=INK)
B.text(5.0, 6.45, r"$w_{ij}\;\propto\;G^{+}_{ij}-G^{-}_{ij}$", fontsize=11, ha="center", color=INK)
arrow(B, (5.0, 5.7), (5.0, 5.05), lw=1.5)

# trace device (subthreshold regime) -> local eligibility
B.add_patch(FancyBboxPatch((0.7, 3.0), 8.6, 2.0,
            boxstyle="round,pad=0.03,rounding_size=0.1",
            facecolor="white", edgecolor=TRACE, lw=1.6))
B.text(5.0, 4.6, "trace device (subthreshold regime)", fontsize=8.2, ha="center", color=TRACE)
tx = np.linspace(0, 1, 120)
e = (1 - np.exp(-(tx / 0.18) ** 2)) * np.exp(-tx / 0.45); e = e / e.max()
B.plot(1.5 + tx * 3.0, 3.35 + e * 0.85, color=TRACE, lw=1.8)
B.text(5.2, 3.75, r"local eligibility $e_{ij}(t)$", fontsize=8.6, ha="left", va="center", color=TRACE)
B.text(5.2, 3.30, r"retention $\tau_{\mathrm{leak}}$", fontsize=7.8, ha="left", va="center", color=TRACE, style="italic")
arrow(B, (5.0, 3.0), (5.0, 2.35), lw=1.5)

# access transistor performs the reward-gated multiply-write
B.add_patch(FancyBboxPatch((0.4, 0.7), 9.2, 1.6,
            boxstyle="round,pad=0.03,rounding_size=0.1",
            facecolor="white", edgecolor=REW, lw=1.9))
B.text(5.0, 1.92, "reward-gated multiply-write (active element)", fontsize=8.2, ha="center", color=REW)
B.text(5.0, 1.38, r"$\Delta w_{ij}=\eta\,(R-b)\;e_{ij}(t_R)$",
       fontsize=11.5, ha="center", color=INK)
B.text(2.95, 0.88, "global $(R-b)$", fontsize=7.4, ha="center", va="center", color=REW, style="italic")
B.text(6.55, 0.88, "local trace", fontsize=7.4, ha="center", va="center", color=TRACE, style="italic")
# ("no weight transport, no per-synapse gradient" removed -- stated in the text)
plt.show()

# ---- cell 8 --------------------------------------------------------
import os
import numpy as np
import matplotlib

import matplotlib.pyplot as plt
from matplotlib.patches import FancyBboxPatch, Circle, FancyArrowPatch, Rectangle

# palette matched to crossbar.png
GREEN = "#3aa07a"      # state-input lines
BLUE = "#2f4b8f"       # device synapses / LIF / forward path
LBLUE = "#eaf0fb"      # device-gate fill
RED = "#c0392b"        # global reward broadcast
PURPLE = "#b07cc6"     # fixed-random feedback (DFA), distinct from forward blue
ORANGE = "#e0a93b"     # homeostasis (local activity regulation)
GREY = "#9aa6b2"
INK = "#2b2b2b"

def _box(ax, xy, w, h, fc, ec, **kw):
    ax.add_patch(FancyBboxPatch(xy, w, h, boxstyle="round,pad=0.012,rounding_size=0.02",
                                fc=fc, ec=ec, lw=kw.pop("lw", 1.4), zorder=kw.pop("z", 3), **kw))

def _compound_node(ax, cx, cy, s=0.085):
    """A crossbar synapse drawn as a compact two-tone split square: left half = the
    non-volatile WEIGHT device, right half = the subthreshold TRACE device. Signals the
    two-device compound cell at array scale without three sub-glyphs per tiny node."""
    ax.add_patch(Rectangle((cx - s, cy - s), s, 2 * s, fc="#cbd6e0", ec=INK, lw=0.9, zorder=4))
    ax.add_patch(Rectangle((cx, cy - s), s, 2 * s, fc="white", ec=BLUE, lw=1.0, zorder=4))

def panel_a(ax, tag=True):
    ax.set_xlim(0, 10); ax.set_ylim(-0.3, 7.6); ax.axis("off")
    if tag:
        ax.text(0.1, 7.3, "(a)", fontsize=15, fontweight="bold")

    # ---- global reward broadcast (top, red dashed) -------------------------------
    ax.plot([0.7, 9.4], [7.05, 7.05], color=RED, ls="--", lw=1.6, zorder=2)
    ax.text(5.05, 7.32, r"global third factor $(R-b)$ broadcast to every crosspoint",
            color=RED, fontsize=9.5, ha="center")

    # ---- layer 1 array: inputs (rows) x hidden (cols) ----------------------------
    in_y = [6.15, 5.5]                      # F input lines (2 shown)
    h_x = [2.3, 3.1, 3.9]                   # H hidden columns (3 shown)
    h_lif_y = 4.05
    ax.text(0.15, 5.82, "state\ninputs", color=GREEN, fontsize=9, va="center")
    for x in h_x:                           # hidden bitlines drop to the LIF row
        ax.plot([x, x], [6.45, h_lif_y + 0.2], color=BLUE, lw=1.4, zorder=1)
    for y in in_y:
        ax.add_patch(Circle((1.25, y), 0.11, fc="white", ec=GREEN, lw=1.6, zorder=4))
        ax.plot([1.36, 4.25], [y, y], color=GREEN, lw=2.0, zorder=2)
    for y in in_y:                          # W1 synapses: compound (weight|trace) cells
        for x in h_x:
            _compound_node(ax, x, y, s=0.085)
    ax.text(3.1, 6.62, r"$W_1$ (input$\rightarrow$hidden array)",
            color=BLUE, fontsize=8.5, ha="center")

    # ---- hidden LIF layer (with homeostasis rings) -------------------------------
    for x in h_x:
        ax.add_patch(Circle((x, h_lif_y), 0.32, fc="none", ec=ORANGE, lw=1.3,
                            ls=(0, (2, 1.5)), zorder=4))      # homeostasis ring
        ax.add_patch(Circle((x, h_lif_y), 0.20, fc="#eaf0fb", ec=BLUE, lw=1.5, zorder=5))
        ax.text(x, h_lif_y, "LIF", color=BLUE, fontsize=6.5, ha="center", va="center")
    ax.text(4.35, h_lif_y, "hidden\nneurons", color=BLUE, fontsize=8.5, va="center")
    ax.annotate("", xy=(1.85, h_lif_y + 0.05), xytext=(1.1, h_lif_y + 0.05),
                arrowprops=dict(arrowstyle="-|>", color=ORANGE, lw=1.5))
    ax.text(0.95, h_lif_y + 0.05, "local\nhomeostasis", color=ORANGE, fontsize=8.0,
            ha="center", va="center")

    # ---- layer 2 array: hidden (rows) x action (cols) ----------------------------
    a_x = [6.4, 7.2]                        # A action columns (2 shown)
    row2_y = [2.45, 2.05, 1.65]             # H hidden rows feeding layer 2
    bus_x = 5.55
    for i, x in enumerate(h_x):             # hidden activity routed to layer-2 rows
        ax.plot([x, bus_x], [h_lif_y - 0.32, h_lif_y - 0.32], color=BLUE, lw=1.0, zorder=1)
    ax.plot([bus_x, bus_x], [h_lif_y - 0.32, row2_y[-1]], color=BLUE, lw=1.0, zorder=1)
    for y in row2_y:
        ax.plot([bus_x, 7.6], [y, y], color=BLUE, lw=1.4, zorder=1)
    for x in a_x:
        ax.plot([x, x], [2.7, 0.55], color=BLUE, lw=1.4, zorder=1)
        for y in row2_y:
            _compound_node(ax, x, y, s=0.080)
    ax.text(6.8, 2.78, r"$W_2$ (hidden$\rightarrow$action array)",
            color=BLUE, fontsize=8.5, ha="center")
    # legend: each array node is a compound cell (weight half + trace half)
    ax.add_patch(Rectangle((8.15, 6.05), 0.14, 0.28, fc="#cbd6e0", ec=INK, lw=0.8, zorder=5))
    ax.add_patch(Rectangle((8.29, 6.05), 0.14, 0.28, fc="white", ec=BLUE, lw=0.9, zorder=5))
    ax.text(8.55, 6.19, "compound cell:\nweight | trace device", color=INK, fontsize=6.6, va="center")

    # ---- action LIF + WTA --------------------------------------------------------
    for x in a_x:
        ax.add_patch(Circle((x, 0.35), 0.20, fc="#eaf0fb", ec=BLUE, lw=1.5, zorder=5))
        ax.text(x, 0.35, "LIF", color=BLUE, fontsize=6.5, ha="center", va="center")
    ax.text(7.65, 0.35, "action\nneurons", color=BLUE, fontsize=8.5, va="center")

    # ---- DFA fixed-random feedback array (purple) --------------------------------
    # output error projected to the hidden layer through a FIXED RANDOM matrix (no W2^T)
    fb = FancyArrowPatch((6.55, 0.15), (3.1, 3.55), connectionstyle="arc3,rad=0.30",
                         arrowstyle="-|>", mutation_scale=13, color=PURPLE, lw=1.9,
                         ls=(0, (5, 2)), zorder=2)
    ax.add_patch(fb)
    ax.text(3.05, 1.15, r"fixed random feedback $B$:  $L_h = B\,L_a$" "\n"
            r"per-neuron credit, no $W_2^{\top}$ transport",
            color=PURPLE, fontsize=8.0, ha="center", va="center")

    # (array-dimension annotation removed -- the caption describes the two
    #  stacked W_1, W_2 arrays flanking the hidden layer)

def panel_b(ax, tag=True):
    ax.set_xlim(0, 10); ax.set_ylim(0, 7.4); ax.axis("off")
    if tag:
        ax.text(0.1, 7.05, "(b)", fontsize=15, fontweight="bold")
    ax.text(5.0, 6.95, "hidden-layer synapse: two-term update", fontsize=10.5, fontweight="bold", ha="center")

    # weight device (non-volatile)
    _box(ax, (1.4, 5.45), 7.2, 0.95, "#f1f3f5", "#444")
    ax.text(5.0, 6.12, r"weight device (non-volatile, switching)", fontsize=9.0, ha="center")
    ax.text(5.0, 5.72, r"$w_{ij}\;\propto\;G^{+}_{ij}-G^{-}_{ij}$", fontsize=11, ha="center")
    # arrows flow top->bottom (weight -> gate -> update): arrowhead (xy) at the LOWER box,
    # tail (xytext) at the upper box, so the head points DOWN into the next stage.
    ax.annotate("", xy=(5.0, 4.95), xytext=(5.0, 5.45),
                arrowprops=dict(arrowstyle="-|>", color=INK, lw=1.4))

    # trace device (subthreshold) / eligibility
    _box(ax, (1.4, 3.85), 7.2, 1.05, LBLUE, BLUE)
    ax.text(5.0, 4.62, "trace device (subthreshold regime)", color=BLUE, fontsize=9.0, ha="center")
    ax.text(5.0, 4.18, r"local eligibility $e_{ij}(t)$,  retention $\tau_{\mathrm{leak}}$",
            color=BLUE, fontsize=8.5, ha="center", style="italic")
    # tiny eligibility glyph
    tt = np.linspace(0, 1, 60)
    gl = (1 - np.exp(-(tt / 0.18) ** 2)) * np.exp(-tt / 0.5)
    ax.plot(1.7 + tt * 0.9, 4.0 + gl * 0.55, color=BLUE, lw=1.3)
    ax.annotate("", xy=(5.0, 3.4), xytext=(5.0, 3.85),
                arrowprops=dict(arrowstyle="-|>", color=INK, lw=1.4))

    # local update box: two contributions; reward-gated multiply-write by an active
    # element (per-cell selector/transistor or shared peripheral CMOS), not the memristors
    _box(ax, (1.0, 0.5), 8.0, 2.75, "white", RED, lw=1.6)
    ax.text(5.0, 2.92, r"reward-gated update $\Delta w_{ij}$ (active element)", color=RED, fontsize=9.0, ha="center")
    # term 1: feedback-aligned credit
    ax.text(5.0, 2.42, r"$\eta\,(R-b)\,[B\,L]_{j}\;\,e_{ij}(t_R)$",
            fontsize=11, ha="center", color=PURPLE)
    ax.text(5.0, 2.08, "per-neuron credit (DFA)", fontsize=7.6, ha="center",
            color=PURPLE, style="italic")
    # plus
    ax.text(5.0, 1.62, r"$+$", fontsize=12, ha="center", color=INK)
    # term 2: homeostatic regulation
    ax.text(5.0, 1.18, r"$\eta_h\,(\rho^{\ast}-\bar\rho_j)\;\,w_{ij}$",
            fontsize=11, ha="center", color=ORANGE)
    ax.text(5.0, 0.84, "local homeostasis", fontsize=7.6, ha="center",
            color=ORANGE, style="italic")
    # (grey "rho_j: neuron's own activity / rho*: shared target" gloss removed
    #  -- both symbols are defined in the surrounding text, Eq. 7.14 paragraph)

    # local/global annotations
    ax.text(2.0, 0.12, "global scalar", color=RED, fontsize=7.4, ha="center")
    ax.text(5.0, 0.12, "fixed-random, local", color=PURPLE, fontsize=7.4, ha="center")
    ax.text(8.0, 0.12, "self-activity, local", color=ORANGE, fontsize=7.4, ha="center")

# combined two-panel version (kept for reference)
fig, (axA, axB) = plt.subplots(1, 2, figsize=(11.0, 4.0),
                               gridspec_kw={"width_ratios": [1.32, 1.0]})
panel_a(axA)
panel_b(axB)
fig.subplots_adjust(left=0.01, right=0.99, top=0.99, bottom=0.04, wspace=0.04)

plt.show()

# split panels: each rendered standalone at full width so the schematic text is
# legible when the figure is placed at ~0.7-0.8\textwidth in the thesis.
figA = plt.figure(figsize=(7.2, 5.2))
axA1 = figA.add_axes([0.01, 0.02, 0.98, 0.96])
panel_a(axA1, tag=False)
outA = os.path.join(OUT, "fig_deep_crossbar_a.png")
plt.show()
figB = plt.figure(figsize=(6.2, 4.6))
axB1 = figB.add_axes([0.01, 0.02, 0.98, 0.96])
panel_b(axB1, tag=False)
outB = os.path.join(OUT, "fig_deep_crossbar_b.png")
plt.show()

# ---- cell 9 --------------------------------------------------------
import os
import numpy as np
import matplotlib

import matplotlib.pyplot as plt
from matplotlib.patches import FancyArrowPatch, Circle, FancyBboxPatch, Patch

INK    = "#222222"
POS    = "#c0392b"     # positive feedback / runaway (red)
HOMEO  = "#e0a93b"     # homeostasis (orange)
EEG    = "#b07cc6"     # EEG / noisy reward (indigo)
NET    = "#2f4b8f"
GREEN  = "#3aa07a"
NEU    = "#9aa6b2"
SET    = "#3aa07a"     # set-point

def feedback_ring(ax, cx, cy, r, color, sign, ccw=True):
    """A circular feedback arrow with a +/- sign in the middle."""
    th0, th1 = (200, -110) if ccw else (-20, 250)
    arc = np.linspace(np.deg2rad(th0), np.deg2rad(th1), 100)
    ax.plot(cx+r*np.cos(arc), cy+r*np.sin(arc), color=color, lw=2.6, zorder=4)
    # arrowhead at the end
    a = arc[-1]; da = (arc[-1]-arc[-2])
    ax.add_patch(FancyArrowPatch((cx+r*np.cos(a-da), cy+r*np.sin(a-da)),
                 (cx+r*np.cos(a), cy+r*np.sin(a)), arrowstyle="-|>",
                 mutation_scale=16, color=color, lw=2.6, zorder=5))
    ax.text(cx, cy, sign, ha="center", va="center", fontsize=20, fontweight="bold",
            color=color, zorder=6)

def panel_a(ax):
    ax.text(-0.04, 1.08, "(a)", transform=ax.transAxes, fontsize=13, fontweight="bold")
    ax.set_title("Associative plasticity is positive feedback", fontsize=9.4, fontweight="bold", pad=6)
    ax.set_xlim(0, 10); ax.set_ylim(0, 10); ax.axis("off")
    # the loop: fires more -> weights grow -> fires more
    feedback_ring(ax, 5.0, 6.3, 1.7, POS, "$+$", ccw=True)
    ax.text(5.0, 8.5, "neuron fires more", ha="center", fontsize=8.0, color=POS)
    ax.text(8.0, 6.3, "its weights\ngrow", ha="center", fontsize=8.0, color=POS)
    ax.text(2.0, 6.3, "drive\nrises", ha="center", fontsize=8.0, color=POS)
    # runaway trace
    t = np.linspace(0, 1, 100)
    run = 1/(1+np.exp(-(t-0.45)*12))           # logistic blow-up to the rail
    ax.plot(1.2+t*7.6, 0.8+run*2.6, color=POS, lw=2.2)
    ax.axhline(0.8+2.6, xmin=0.12, xmax=0.88, color=NEU, lw=0.8, ls=(0,(3,2)))
    ax.text(8.9, 0.8+2.6, "rail", fontsize=6.8, color=NEU, va="center")
    ax.text(5.0, 0.35, "left alone: a hidden unit locks high $\\to$ policy collapses",
            ha="center", fontsize=7.4, color=POS, style="italic")

def panel_b(ax):
    ax.text(-0.04, 1.08, "(b)", transform=ax.transAxes, fontsize=13, fontweight="bold")
    ax.set_title("Homeostasis is the negative-feedback counterweight", fontsize=9.4, fontweight="bold", pad=6)
    ax.set_xlim(0, 10); ax.set_ylim(0, 10); ax.axis("off")
    feedback_ring(ax, 5.0, 6.3, 1.7, HOMEO, "$-$", ccw=False)
    ax.text(5.0, 8.5, "rate above $\\rho^{\\!*}$", ha="center", fontsize=8.0, color=HOMEO)
    ax.text(8.1, 6.3, "weights\npulled down", ha="center", fontsize=8.0, color=HOMEO)
    ax.text(1.9, 6.3, "rate\nreturns", ha="center", fontsize=8.0, color=HOMEO)
    # rate settling to set-point
    t = np.linspace(0, 1, 200)
    settle = 0.5 + 0.5*np.exp(-t*5)*np.cos(t*14)   # damped to set-point
    ax.plot(1.2+t*7.6, 0.8+settle*2.6, color=HOMEO, lw=2.2)
    ax.axhline(0.8+0.5*2.6, xmin=0.12, xmax=0.88, color=SET, lw=1.0, ls=(0,(4,2)))
    ax.text(8.9, 0.8+0.5*2.6, r"$\rho^{\!*}$", fontsize=8.5, color=SET, va="center")
    ax.text(5.0, 0.35, "each neuron holds its own rate near a shared set-point",
            ha="center", fontsize=7.4, color=HOMEO, style="italic")

def panel_c(ax):
    rng = np.random.RandomState(3)
    n = 200; trial = np.arange(n)
    # noisy reward strip at the top
    rew = 0.5 + 0.5*np.sign(np.sin(trial*0.5)) * (rng.rand(n) < 0.85)  # noisy +/- reward
    ax.plot(trial, 9.3 + 0.001*trial, color="none")
    for i in range(0, n, 2):
        c = EEG if rew[i] > 0.5 else NEU
        ax.plot([trial[i], trial[i]], [9.0, 9.0+0.35*(0.4+0.6*(rew[i]>0.5))], color=c, lw=0.8, alpha=0.7)
    ax.text(n/2, 9.7, "noisy single-trial reward (EEG-decoded)", ha="center", fontsize=8.0, color=EEG)

    # without homeostasis: firing rate diverges to a rail (positive feedback unchecked)
    rho = 0.5; x = rho; traj_no = []
    for i in range(n):
        x += 0.018*(rew[i]-0.5)*4 + 0.03*(x-rho) + 0.01*rng.randn()  # pos feedback (x-rho) destabilises
        x = np.clip(x, 0, 1); traj_no.append(x)
    ax.plot(trial, np.array(traj_no)*6.8+0.8, color=POS, lw=2.2,
            label="no homeostasis (positive feedback runs away)")

    # with homeostasis: bounded near set-point (negative-feedback counterweight)
    x = rho; traj_h = []
    for i in range(n):
        x += 0.018*(rew[i]-0.5)*4 + 0.03*(x-rho) - 0.16*(x-rho) + 0.01*rng.randn()  # neg feedback dominates
        x = np.clip(x, 0, 1); traj_h.append(x)
    ax.plot(trial, np.array(traj_h)*6.8+0.8, color=HOMEO, lw=2.4,
            label="with homeostasis (held near the set-point)")

    ax.axhline(rho*6.8+0.8, color=SET, lw=1.0, ls=(0,(4,2)))
    ax.text(n+2, rho*6.8+0.8, r"set-point $\rho^{\!*}$", fontsize=8.5, color=SET, va="center")
    ax.text(n+2, 1.0*6.8+0.8-0.18, "saturation rail", fontsize=7.4, color=NEU, va="center")
    # name the two control regimes directly on the curves
    ax.annotate("runaway: unit locks high,\npolicy collapses", xy=(150, traj_no[150]*6.8+0.8),
                xytext=(96, 8.1), fontsize=7.6, color=POS, va="center",
                arrowprops=dict(arrowstyle="->", color=POS, lw=1.0))
    ax.annotate("bounded: extra reward noise\nabsorbed", xy=(150, traj_h[150]*6.8+0.8),
                xytext=(70, 2.0), fontsize=7.6, color=HOMEO, va="center",
                arrowprops=dict(arrowstyle="->", color=HOMEO, lw=1.0))

    ax.set_xlim(0, n+30); ax.set_ylim(0, 10.2)
    ax.set_xlabel("trial", fontsize=9.5)
    ax.set_ylabel("hidden-unit firing rate", fontsize=9.5)
    ax.set_yticks([]); ax.tick_params(labelsize=8.5)
    ax.legend(loc="lower left", fontsize=8.2, framealpha=0.95, bbox_to_anchor=(0.01, 0.02))

fig, ax = plt.subplots(figsize=(7.2, 4.4))
panel_c(ax)
fig.suptitle("Why homeostasis stabilises learning under a noisy reward",
             fontsize=11.5, fontweight="bold", y=0.98)
fig.subplots_adjust(left=0.07, right=0.99, top=0.90, bottom=0.12)

plt.show()

# ---- cell 10 --------------------------------------------------------
import os
import numpy as np
import matplotlib

import matplotlib.pyplot as plt
from matplotlib.patches import Patch
from matplotlib.lines import Line2D

# ---- palette (matches the manuscript green-blue scheme) ---------------------
TEAL = "#3aa07a"   # device measured bands / markers
INDIGO = "#2f4b8f"  # tuned RL retention settings
BRICK = "#c0392b"  # biological window accent
INK = "#2b2b2b"
GRID = "#c8d0d8"

plt.rcParams.update({
    "font.family": "DejaVu Sans",
    "font.size": 9,
    "axes.edgecolor": INK,
    "text.color": INK,
    "axes.labelcolor": INK,
    "xtick.color": INK,
    "ytick.color": INK,
})

# ---- data (all real; see module docstring) ----------------------------------
BIO = (0.3, 10.0)           # biological eligibility window (s)
TAU_R = (1.9, 14.5)         # MEASURED cascade rise constant (s)
TAU_D = (36.0, 537.0)       # MEASURED space-charge decay constant (s)
TAU_LEAK_MEAS = 1.3         # MEASURED held-bias trap-discharge retention (ITO median, s)
TAU_LEAK_FF = 3.6           # field-free retention (Poole-Frenkel extrapolation, s)
TAU_LEAK = [5, 10, 20]      # MODEL-SET retention settings swept in RL (s)

# the genuine overlap stripe = biological window  AND  device measured dynamics
# i.e. [max(0.3, 1.3), min(10, 14.5)] = [1.3, 10]
OVERLAP = (max(BIO[0], TAU_LEAK_MEAS), min(BIO[1], TAU_R[1]))

# y-rows (top to bottom)
Y_BIO = 3
Y_R = 2
Y_D = 1
Y_LEAK = 0
BAR_H = 0.30  # half-height of bands

fig, ax = plt.subplots(figsize=(7.2, 3.0))

# ---- key-message overlap stripe (the true intersection ~1.3-10 s) -----------
ax.axvspan(OVERLAP[0], OVERLAP[1], color=TEAL, alpha=0.08, zorder=0)
ax.text(np.sqrt(OVERLAP[0] * OVERLAP[1]), 4.05,
        "device dynamics meet\nthe biological window",
        ha="center", va="center", fontsize=7.0, style="italic", color="#5a6b62",
        linespacing=1.05, zorder=5)

# ---- row 1: biological eligibility window (hatched band) --------------------
ax.fill_between(BIO, Y_BIO - BAR_H, Y_BIO + BAR_H, facecolor="none",
                edgecolor=BRICK, hatch="////", linewidth=1.1, zorder=3)
ax.text(BIO[1] * 1.35, Y_BIO,
        "biological eligibility window (~0.3-10 s)",
        va="center", ha="left", fontsize=8.5, color=BRICK)

# ---- row 2: MEASURED device rise tau_r (solid band) -------------------------
ax.fill_between(TAU_R, Y_R - BAR_H, Y_R + BAR_H, color=TEAL, alpha=0.85,
                zorder=3)
ax.text(TAU_R[1] * 1.35, Y_R,
        r"device rise  $\tau_r$  (measured, 1.9-14.5 s)",
        va="center", ha="left", fontsize=8.5, color=INK)

# ---- row 3: MEASURED device decay tau_d (solid band, lighter shade) ---------
ax.fill_between(TAU_D, Y_D - BAR_H, Y_D + BAR_H, color=TEAL, alpha=0.45,
                zorder=3)
ax.text(TAU_D[0] * 0.62, Y_D,
        r"device decay  $\tau_d$  (measured, 36-537 s)",
        va="center", ha="right", fontsize=8.5, color=INK)

# ---- row 4: tau_leak -- MEASURED retention (filled) + RL sweep (open) --------
# the held-bias ITO retention (filled) and its field-free Poole-Frenkel
# extrapolation (filled, lighter): the credit window the paper rests on, both
# sitting inside the biological window. An arrow marks the field correction.
ax.annotate("", xy=(TAU_LEAK_FF, Y_LEAK), xytext=(TAU_LEAK_MEAS, Y_LEAK),
            arrowprops=dict(arrowstyle="-|>", color="#5a6b62", lw=1.1,
                            shrinkA=5, shrinkB=5), zorder=4)
ax.plot([TAU_LEAK_MEAS], [Y_LEAK], marker="o", ms=9.5, mfc=TEAL,
        mec=INK, mew=1.0, zorder=5)
ax.annotate(r"$1.3$ (held bias)", (TAU_LEAK_MEAS, Y_LEAK),
            textcoords="offset points", xytext=(0, 10), ha="center",
            fontsize=6.8, color=INK)
ax.plot([TAU_LEAK_FF], [Y_LEAK], marker="o", ms=9.5, mfc=TEAL, alpha=0.6,
        mec=INK, mew=1.0, zorder=5)
ax.annotate(r"$3.6$ (field-free)", (TAU_LEAK_FF, Y_LEAK),
            textcoords="offset points", xytext=(0, -14), ha="center",
            fontsize=6.8, color=INK)
# the tunable retention settings swept in the RL experiments (open markers)
for t in TAU_LEAK:
    ax.plot([t], [Y_LEAK], marker="o", ms=8.5, mfc="white",
            mec=INDIGO, mew=1.6, zorder=4)
    ax.annotate(f"{t}", (t, Y_LEAK), textcoords="offset points",
                xytext=(0, 9), ha="center", fontsize=7, color=INDIGO)
ax.text(TAU_LEAK[-1] * 1.6, Y_LEAK,
        r"retention  $\tau_{\mathrm{leak}}$  (measured $+$ tuned)",
        va="center", ha="left", fontsize=8.5, color=INDIGO)

# ---- axes cosmetics ---------------------------------------------------------
ax.set_xscale("log")
ax.set_xlim(0.1, 1000)
ax.set_ylim(-0.7, 4.5)
ax.set_xlabel("time constant / window (s)")
ax.set_yticks([])
ax.spines["top"].set_visible(False)
ax.spines["right"].set_visible(False)
ax.spines["left"].set_visible(False)
ax.tick_params(axis="x", which="both", length=3)
ax.grid(True, axis="x", which="major", color=GRID, lw=0.6, zorder=0)
ax.set_axisbelow(True)

# ---- legend: MEASURED vs MODEL-SET vs BIOLOGICAL ----------------------------
legend_handles = [
    Patch(facecolor=TEAL, alpha=0.85, label="measured device dynamics"),
    Line2D([0], [0], marker="o", color="none", mfc=TEAL, mec=INK,
           mew=1.0, ms=8, label=r"measured $\tau_{\mathrm{leak}}$ (credit window)"),
    Line2D([0], [0], marker="o", color="none", mfc="white", mec=INDIGO,
           mew=1.6, ms=8, label="tuned retention (RL)"),
    Patch(facecolor="none", edgecolor=BRICK, hatch="////",
          label="biological (literature)"),
]
ax.legend(handles=legend_handles, loc="upper right",
          bbox_to_anchor=(1.0, 0.99), frameon=False, fontsize=7.5,
          handlelength=1.6, borderaxespad=0.2, labelspacing=0.35)

fig.tight_layout()

plt.show()
plt.close(fig)

# ---- cell 11 --------------------------------------------------------
import os
import numpy as np
import matplotlib

import matplotlib.pyplot as plt
from matplotlib.patches import Patch
from matplotlib.lines import Line2D

# ---- palette (matches the thesis green-indigo scheme) -----------------------
TEAL = "#3aa07a"    # device measured bands / markers
INDIGO = "#2f4b8f"  # tuned RL retention settings
BRICK = "#c0392b"   # biological window accent
INK = "#2b2b2b"
GRID = "#c8d0d8"

plt.rcParams.update({
    "font.family": "DejaVu Sans",
    "font.size": 9,
    "axes.edgecolor": INK,
    "text.color": INK,
    "axes.labelcolor": INK,
    "xtick.color": INK,
    "ytick.color": INK,
})

# ---- data (all real; per the Chapter 7 caption) -----------------------------
BIO = (0.3, 10.0)           # biological eligibility window (s)
TAU_R = (1.9, 14.5)         # MEASURED cascade rise constant (s)
TAU_D = (36.0, 537.0)       # MEASURED space-charge decay constant (s)
TAU_LEAK_MEAS = 1.3         # MEASURED held-bias trap-discharge retention (ITO median, s)
TAU_LEAK_FF = 3.6           # field-free retention (Poole-Frenkel extrapolation, s)
TAU_LEAK = [5, 10, 20]      # MODEL-SET retention settings swept in RL (s)

OVERLAP = (max(BIO[0], TAU_LEAK_MEAS), min(BIO[1], TAU_R[1]))  # true intersection ~1.3-10 s

Y_BIO = 3
Y_R = 2
Y_D = 1
Y_LEAK = 0
BAR_H = 0.30

fig, ax = plt.subplots(figsize=(7.2, 3.0))

# ---- key-message overlap stripe (the true intersection ~1.3-10 s), no text ----
ax.axvspan(OVERLAP[0], OVERLAP[1], color=TEAL, alpha=0.08, zorder=0)

# ---- row 1: biological eligibility window (hatched band) --------------------
ax.fill_between(BIO, Y_BIO - BAR_H, Y_BIO + BAR_H, facecolor="none",
                edgecolor=BRICK, hatch="////", linewidth=1.1, zorder=3)

# ---- row 2: MEASURED device rise tau_r (solid band) -------------------------
ax.fill_between(TAU_R, Y_R - BAR_H, Y_R + BAR_H, color=TEAL, alpha=0.85, zorder=3)

# ---- row 3: MEASURED device decay tau_d (solid band, lighter shade) ---------
ax.fill_between(TAU_D, Y_D - BAR_H, Y_D + BAR_H, color=TEAL, alpha=0.45, zorder=3)

# ---- row 4: tau_leak -- MEASURED retention (filled) + RL sweep (open) --------
# held-bias ITO retention (filled) and its field-free Poole-Frenkel extrapolation
# (filled, lighter); an arrow marks the field correction. No inline numbers.
ax.annotate("", xy=(TAU_LEAK_FF, Y_LEAK), xytext=(TAU_LEAK_MEAS, Y_LEAK),
            arrowprops=dict(arrowstyle="-|>", color="#5a6b62", lw=1.1,
                            shrinkA=5, shrinkB=5), zorder=4)
ax.plot([TAU_LEAK_MEAS], [Y_LEAK], marker="o", ms=9.5, mfc=TEAL,
        mec=INK, mew=1.0, zorder=5)
ax.plot([TAU_LEAK_FF], [Y_LEAK], marker="o", ms=9.5, mfc=TEAL, alpha=0.6,
        mec=INK, mew=1.0, zorder=5)
# the tunable retention settings swept in the RL experiments (open markers), no numbers
for t in TAU_LEAK:
    ax.plot([t], [Y_LEAK], marker="o", ms=8.5, mfc="white",
            mec=INDIGO, mew=1.6, zorder=4)

# ---- axes cosmetics ---------------------------------------------------------
ax.set_xscale("log")
ax.set_xlim(0.1, 1000)
ax.set_ylim(-0.7, 4.5)
ax.set_xlabel("time constant / window (s)")
ax.set_yticks([])
ax.spines["top"].set_visible(False)
ax.spines["right"].set_visible(False)
ax.spines["left"].set_visible(False)
ax.tick_params(axis="x", which="both", length=3)
ax.grid(True, axis="x", which="major", color=GRID, lw=0.6, zorder=0)
ax.set_axisbelow(True)

# ---- legend: everything the removed inline labels used to carry --------------
legend_handles = [
    Patch(facecolor="none", edgecolor=BRICK, hatch="////",
          label=r"biological eligibility window ($\sim$0.3--10 s)"),
    Patch(facecolor=TEAL, alpha=0.85, label=r"device rise $\tau_r$ (measured)"),
    Patch(facecolor=TEAL, alpha=0.45, label=r"device decay $\tau_d$ (measured)"),
    Line2D([0], [0], marker="o", color="none", mfc=TEAL, mec=INK,
           mew=1.0, ms=8, label=r"$\tau_{\mathrm{leak}}$, held-bias (measured)"),
    Line2D([0], [0], marker="o", color="none", mfc=TEAL, alpha=0.6, mec=INK,
           mew=1.0, ms=8, label=r"$\tau_{\mathrm{leak}}$, field-free (extrapolated)"),
    Line2D([0], [0], marker="o", color="none", mfc="white", mec=INDIGO,
           mew=1.6, ms=8, label=r"$\tau_{\mathrm{leak}}$, tuned (RL settings)"),
]
ax.legend(handles=legend_handles, loc="upper right",
          bbox_to_anchor=(1.0, 1.02), frameon=False, fontsize=7.2,
          handlelength=1.6, borderaxespad=0.2, labelspacing=0.4)

fig.tight_layout()
plt.show()
plt.close(fig)
\end{lstlisting}

\begin{lstlisting}[caption={experiments/REPRODUCE\_ALL.ipynb},label={lst:siox-nb-REPRODUCE-ALL}]
# ---- cell 1 --------------------------------------------------------
import mrl_trace, mrl_trace.paths as paths
print("mrl_trace OK -- data/results:", paths.results_dir())

# ---- cell 2 --------------------------------------------------------
import glob, os
import nbformat
from nbclient import NotebookClient
from nbclient.exceptions import CellExecutionError

HERE = os.path.abspath(".")
notebooks = sorted(glob.glob(os.path.join(HERE, "0?_*.ipynb"))) + [
    os.path.join(HERE, "device_model.ipynb"),
    os.path.join(HERE, "schematics.ipynb"),
]
ok = 0
for path in notebooks:
    name = os.path.basename(path)
    nb = nbformat.read(path, as_version=4)
    client = NotebookClient(nb, timeout=1800, kernel_name="siox",
                            resources={"metadata": {"path": HERE}})
    try:
        client.execute()
        n_fig = sum(1 for c in nb.cells if c.cell_type == "code"
                    for o in c.get("outputs", [])
                    if "image/png" in o.get("data", {}))
        print(f"OK   {name:38s} ({n_fig} figures)")
        ok += 1
    except CellExecutionError as e:
        print(f"FAIL {name:38s} {str(e).splitlines()[-1][:80]}")
print(f"\n{ok}/{len(notebooks)} topic notebooks executed cleanly (replay mode).")
\end{lstlisting}

